% Generated mechanically from frozen input p04/ja/coherent.tex.
% Do not edit this generated file; edit the bound locale source instead.

% OK, start here.
%

\setcounter{chapter}{29}
\chapter{スキームのコホモロジー}



\phantomsection
\label{coherent-section-phantom}


\section{はじめに}
\label{coherent-section-introduction}

\noindent
本章では、まず準連接層のコホモロジーに関するいくつかの結果を証明する。
基本的な参考文献は \cite{EGA} である。
その後、Noether 的な設定における連接層のコホモロジーを詳しく論じる。
\cite{FAC} を参照せよ。







\section{準連接層の {\v C}ech コホモロジー}
\label{coherent-section-cech-quasi-coherent}

\noindent
$X$ をスキームとする。
$U \subset X$ をアフィン開集合とする。
$U$ の {\it 標準開被覆}とは、
$\mathcal{U} : U = \bigcup_{i = 1}^n D(f_i)$
という形の被覆であったことを思い出そう。ここで
$f_1, \ldots, f_n \in \Gamma(U, \mathcal{O}_X)$ は単位イデアルを生成する。
スキーム、定義
\ref{schemes-definition-standard-covering} を参照せよ。

\begin{lemma}
\label{coherent-lemma-cech-cohomology-quasi-coherent-trivial}
$X$ をスキームとする。
$\mathcal{F}$ を準連接 $\mathcal{O}_X$-加群とする。
$\mathcal{U} : U = \bigcup_{i = 1}^n D(f_i)$ を、$X$ のあるアフィン開集合の
標準開被覆とする。
このとき、$\check{H}^p(\mathcal{U}, \mathcal{F}) = 0$ が
すべての $p > 0$ に対して成り立つ。
\end{lemma}

\begin{proof}
$U = \Spec(A)$ と書く。ここで $A$ はある環である。
言い換えると、$f_1, \ldots, f_n$ は $A$ の単位イデアルを生成する
$A$ の元である。
$\mathcal{F}|_U = \widetilde{M}$ と書く。ここで、ある $A$-加群を $M$ と表した。
明らかに {\v C}ech 複体
$\check{\mathcal{C}}^\bullet(\mathcal{U}, \mathcal{F})$
は複体
$$
\prod\nolimits_{i_0} M_{f_{i_0}} \to
\prod\nolimits_{i_0i_1} M_{f_{i_0}f_{i_1}} \to
\prod\nolimits_{i_0i_1i_2} M_{f_{i_0}f_{i_1}f_{i_2}} \to
\ldots
$$
と同一視される。示すべきことは、拡大複体
\begin{equation}
\label{coherent-equation-extended}
0 \to
M \to
\prod\nolimits_{i_0} M_{f_{i_0}} \to
\prod\nolimits_{i_0i_1} M_{f_{i_0}f_{i_1}} \to
\prod\nolimits_{i_0i_1i_2} M_{f_{i_0}f_{i_1}f_{i_2}} \to
\ldots
\end{equation}
が完全であることである（その切り詰めは代数、補題
\ref{algebra-lemma-cover-module} で調べた）。
複体 (\ref{coherent-equation-extended}) が各素イデアル $\mathfrak p$ で局所化した後に
完全であることを示せば十分である。代数、補題
\ref{algebra-lemma-characterize-zero-local} を参照せよ。
実際には、$\mathfrak p$ で局所化した拡大複体がゼロにホモトピックであることを示す。

\medskip\noindent
添字 $i$ で $f_i \not \in \mathfrak p$ となるものが存在する。
そのような元を一つ選んで $i_{\text{fix}}$ と固定する。すると
$M_{f_{i_{\text{fix}}}, \mathfrak p} = M_{\mathfrak p}$ である。また同様に、
積 $f_{i_0} \ldots f_{i_p}$ で局所化し、さらに $\mathfrak p$ で局所化するとき、
任意の $f_{i_j}$ は、$i_j = i_{\text{fix}}$ ならば省くことができる。
ホモトピー
$$
h :
\prod\nolimits_{i_0 \ldots i_{p + 1}}
M_{f_{i_0} \ldots f_{i_{p + 1}}, \mathfrak p}
\longrightarrow
\prod\nolimits_{i_0 \ldots i_p}
M_{f_{i_0} \ldots f_{i_p}, \mathfrak p}
$$
を規則
$$
h(s)_{i_0 \ldots i_p} = s_{i_{\text{fix}} i_0 \ldots i_p}
$$
によって定める。
（これはコホモロジー、補題 \ref{cohomology-lemma-homology-complex} の
証明におけるホモトピーと「双対的」である。）
言い換えると、
$h : \prod_{i_0} M_{f_{i_0}, \mathfrak p} \to M_\mathfrak p$
は因子
$M_{f_{i_{\text{fix}}}, \mathfrak p} = M_{\mathfrak p}$ への射影であり、
一般に写像 $h$ は因子
$M_{f_{i_{\text{fix}}} f_{i_1} \ldots f_{i_{p + 1}}, \mathfrak p}
= M_{f_{i_1} \ldots f_{i_{p + 1}}, \mathfrak p}$ への射影である。計算すると
\begin{align*}
(dh + hd)(s)_{i_0 \ldots i_p}
& =
\sum\nolimits_{j = 0}^p
(-1)^j
h(s)_{i_0 \ldots \hat i_j \ldots i_p}
+
d(s)_{i_{\text{fix}} i_0 \ldots i_p}\\
& =
\sum\nolimits_{j = 0}^p
(-1)^j
s_{i_{\text{fix}} i_0 \ldots \hat i_j \ldots i_p}
+
s_{i_0 \ldots i_p}
+
\sum\nolimits_{j = 0}^p
(-1)^{j + 1}
s_{i_{\text{fix}} i_0 \ldots \hat i_j \ldots i_p} \\
& =
s_{i_0 \ldots i_p}
\end{align*}
となる。従って、望んだ通り恒等写像はゼロにホモトピックである。
\end{proof}

\noindent
次の補題は特に、任意のアフィンスキーム $X$ と任意の準連接層
$\mathcal{F}$（$X$ 上のもの）に対して
$$
H^p(X, \mathcal{F}) = 0
$$
がすべての $p > 0$ について成り立つことを述べている。

\begin{lemma}
\label{coherent-lemma-quasi-coherent-affine-cohomology-zero}
\begin{slogan}
Serre 消滅：アフィンスキーム上では準連接加群の高次コホモロジーは消滅する。
\end{slogan}
$X$ をスキームとする。
$\mathcal{F}$ を準連接 $\mathcal{O}_X$-加群とする。
任意のアフィン開集合 $U \subset X$ に対して
$H^p(U, \mathcal{F}) = 0$ がすべての $p > 0$ について成り立つ。
\end{lemma}

\begin{proof}
コホモロジー、補題
\ref{cohomology-lemma-cech-vanish-basis} を適用する。
位相の基 $\mathcal{B}$ として $X$ のすべてのアフィン開集合を用いる。
ここで扱う位相は $X$ のものである。
開被覆の集合 $\text{Cov}$ として、$X$ のアフィン開集合の標準開被覆を用いる。
次に、コホモロジー、補題
\ref{cohomology-lemma-cech-vanish-basis} の条件 (1)、(2)、(3) が
成り立つことを確認する。アフィンスキーム内の標準開集合どうしの交わりは
再び標準開集合である。従って性質 (1) が成り立つ。
これらの被覆は $\mathcal{B}$ の任意の元の開被覆の共終系をなす。スキーム、補題
\ref{schemes-lemma-standard-open} を参照せよ。
従って (2) が成り立つ。
最後に、条件 (3) は補題
\ref{coherent-lemma-cech-cohomology-quasi-coherent-trivial} から従う。
\end{proof}

\noindent
ここで、アフィンスキーム上の準連接層のコホモロジー消滅の相対版を与える。

\begin{lemma}
\label{coherent-lemma-relative-affine-vanishing}
$f : X \to S$ をスキームの射とする。
$\mathcal{F}$ を準連接 $\mathcal{O}_X$-加群とする。
$f$ がアフィンならば、$R^if_*\mathcal{F} = 0$ が
すべての $i > 0$ に対して成り立つ。
\end{lemma}

\begin{proof}
コホモロジー、補題
\ref{cohomology-lemma-describe-higher-direct-images} によれば、層
$R^if_*\mathcal{F}$ は、前層
$V \mapsto H^i(f^{-1}(V), \mathcal{F}|_{f^{-1}(V)})$
に伴う層である。
仮定により、$V$ がアフィンならば $f^{-1}(V)$ はアフィンである。
射、定義 \ref{morphisms-definition-affine} を参照せよ。
補題 \ref{coherent-lemma-quasi-coherent-affine-cohomology-zero} により、
$H^i(f^{-1}(V), \mathcal{F}|_{f^{-1}(V)}) = 0$
が $V$ がアフィンならば成り立つ。$S$ はアフィン開集合からなる基をもつので、結論を得る。
\end{proof}

\begin{lemma}
\label{coherent-lemma-relative-affine-cohomology}
$f : X \to S$ をスキームのアフィン射とする。
$\mathcal{F}$ を準連接 $\mathcal{O}_X$-加群とする。
このとき、$H^i(X, \mathcal{F}) = H^i(S, f_*\mathcal{F})$ が
すべての $i \geq 0$ に対して成り立つ。
\end{lemma}

\begin{proof}
補題 \ref{coherent-lemma-relative-affine-vanishing} と Leray スペクトル系列から従う。
コホモロジー、補題 \ref{cohomology-lemma-apply-Leray} を参照せよ。
\end{proof}

\noindent
次の二つの補題は、いつ {\v C}ech コホモロジーを用いて
準連接加群のコホモロジーを計算できるかを説明する。

\begin{lemma}
\label{coherent-lemma-affine-diagonal}
$X$ をスキームとする。次の条件は同値である。
\begin{enumerate}
\item $X$ はアフィンな対角射 $\Delta : X \to X \times X$ をもつ。
\item アフィン開集合 $U, V \subset X$ に対して、交わり
$U \cap V$ はアフィンである。
\item 開被覆 $\mathcal{U} : X = \bigcup_{i \in I} U_i$ で、
$U_{i_0 \ldots i_p}$ がすべての $p \ge 0$ およびすべての
$i_0, \ldots, i_p \in I$ に対してアフィン開集合となるものが存在する。
\end{enumerate}
特に、$X$ が分離的ならばこれらの条件が成り立つ。
\end{lemma}

\begin{proof}
$X$ がアフィンな対角射をもつと仮定する。
$U, V \subset X$ をアフィン開集合とする。
このとき $U \cap V = \Delta^{-1}(U \times V)$ はアフィンである。
従って (2) が成り立つ。(2) が (3) を含意することは直ちに分かる。
逆に、(3) のような $X$ の被覆が存在するならば、
$X \times X = \bigcup U_i \times U_{i'}$ はアフィン開被覆であり、
$\Delta^{-1}(U_i \times U_{i'}) = U_i \cap U_{i'}$
がアフィンであることが分かる。このとき、射、補題
\ref{morphisms-lemma-characterize-affine} により $\Delta$ はアフィン射である。
最後の主張はスキーム、補題
\ref{schemes-lemma-characterize-separated} から従う。
\end{proof}

\begin{lemma}
\label{coherent-lemma-cech-cohomology-quasi-coherent}
$X$ をスキームとする。
$\mathcal{U} : X = \bigcup_{i \in I} U_i$ を、
$U_{i_0 \ldots i_p}$ がすべての $p \ge 0$ とすべての
$i_0, \ldots, i_p \in I$ に対してアフィン開集合となるような開被覆とする。
この場合、任意の準連接層 $\mathcal{F}$ に対し、
$$
\check{H}^p(\mathcal{U}, \mathcal{F}) = H^p(X, \mathcal{F})
$$
が $\Gamma(X, \mathcal{O}_X)$-加群としてすべての $p$ について成り立つ。
\end{lemma}

\begin{proof}
補題 \ref{coherent-lemma-quasi-coherent-affine-cohomology-zero} によれば、これは
コホモロジー、補題
\ref{cohomology-lemma-cech-spectral-sequence-application}
の特殊な場合である。
\end{proof}









\section{コホモロジーの消滅}
\label{coherent-section-vanishing}

\noindent
アフィンスキーム上では、任意の準連接層の高次コホモロジー群が消滅することを見た
（補題 \ref{coherent-lemma-quasi-coherent-affine-cohomology-zero}）。
実は、このことはアフィンスキームを特徴づける。二つの形を与える。

\begin{lemma}
\label{coherent-lemma-quasi-compact-h1-zero-covering}
\begin{reference}
\cite{Serre-criterion}, \cite[II, Theorem 5.2.1 (d') and IV (1.7.17)]{EGA}
\end{reference}
\begin{slogan}
アフィン性の Serre 判定法。
\end{slogan}
$X$ をスキームとする。次を仮定する。
\begin{enumerate}
\item $X$ は準コンパクトである。
\item 任意の準連接イデアル層
$\mathcal{I} \subset \mathcal{O}_X$ に対して $H^1(X, \mathcal{I}) = 0$ である。
\end{enumerate}
このとき $X$ はアフィンである。
\end{lemma}

\begin{proof}
$x \in X$ を閉点とする。$U \subset X$ を $x$ のアフィン開近傍とする。
$U = \Spec(A)$ と書き、$\mathfrak m \subset A$ を $x$ に対応する極大イデアルとする。
$Z = X \setminus U$ および $Z' = Z \cup \{x\}$ と置く。
スキーム、補題 \ref{schemes-lemma-reduced-closed-subscheme} により、
準連接イデアル層 $\mathcal{I}$ および $\mathcal{I}'$ で、
被約閉部分スキーム $Z$ および $Z'$ をそれぞれ切り出すものが存在する。
短完全列
$$
0 \to \mathcal{I}' \to \mathcal{I} \to \mathcal{I}/\mathcal{I}' \to 0.
$$
を考える。$x$ は $X$ の閉点であり、$x \not \in Z$ であるから、
$\mathcal{I}/\mathcal{I}'$ は $x$ に台をもつ。実際、
$\mathcal{I}/\mathcal{I'}$ の $U$ への制限は $A$-加群
$A/\mathfrak m$ に対応する。従って
$\Gamma(X, \mathcal{I}/\mathcal{I'})
= A/\mathfrak m$ である。
仮定により $H^1(X, \mathcal{I}') = 0$ なので、大域切断
$f \in \Gamma(X, \mathcal{I})$ で、元 $1 \in A/\mathfrak m$ に写るものが存在する。
ここで像は $\mathcal{I}/\mathcal{I'}$ の切断として考える。明らかに
$x \in X_f \subset U$ である。従って
$X_f = D(f_A)$ であり、ここで $f_A$ は $f$ の
$A = \Gamma(U, \mathcal{O}_X)$ における像である。特に $X_f$ はアフィンである。

\medskip\noindent
$W = \bigcup X_f$ を考える。ここで合併はすべての
$f \in \Gamma(X, \mathcal{O}_X)$ で $X_f$ がアフィンとなるものにわたる。
明らかに $W$ は $X$ の開集合である。
上の議論により、$X$ のすべての閉点は $W$ に含まれる。
閉部分集合 $X \setminus W$（$X$ のもの）も準コンパクトである
（位相、補題 \ref{topology-lemma-closed-in-quasi-compact} を参照）。
従って、空でなければ閉点をもつ
（位相、補題 \ref{topology-lemma-quasi-compact-closed-point} を参照）。
これはすべての閉点が $W$ にあることに矛盾する。
従って $X = W$ である。

\medskip\noindent
$f_1, \ldots, f_n \in \Gamma(X, \mathcal{O}_X)$ を有限個選び、
$X = X_{f_1} \cup \ldots \cup X_{f_n}$ かつ各
$X_{f_i}$ がアフィンとなるようにする。上で見たように、これは可能である。
性質、補題 \ref{properties-lemma-characterize-affine} により、証明を終えるには
$f_1, \ldots, f_n$ が
$\Gamma(X, \mathcal{O}_X)$ の単位イデアルを生成することを示せば十分である。
短完全列
$$
\xymatrix{
0 \ar[r] &
\mathcal{F} \ar[r] &
\mathcal{O}_X^{\oplus n} \ar[rr]^{f_1, \ldots, f_n} & &
\mathcal{O}_X \ar[r] &
0
}
$$
を考える。$f_1, \ldots, f_n$ によって定まる矢印は、
開集合 $X_{f_i}$ が $X$ を被覆するため全射である。
この全射の核を $\mathcal{F}$ とする。
$\mathcal{F}$ はフィルトレーション
$$
0 = \mathcal{F}_0 \subset \mathcal{F}_1 \subset
\ldots \subset \mathcal{F}_n = \mathcal{F}
$$
をもち、各劣商 $\mathcal{F}_i/\mathcal{F}_{i - 1}$ は
準連接イデアル層と同型である。
実際、$\mathcal{F}_i$ を、$\mathcal{F}$ と、最初の $i$ 個の
$\mathcal{O}_X^{\oplus n}$ の直和因子との交わりとして取ればよい。
補題の仮定により、
$H^1(X, \mathcal{F}_i/\mathcal{F}_{i - 1}) = 0$ がすべての $i$ について成り立つ。
このことから $H^1(X, \mathcal{F}_2) = 0$ が従う。実際、それは
$H^1(X, \mathcal{F}_1)$ と $H^1(X, \mathcal{F}_2/\mathcal{F}_1)$ の間に挟まれている。
同様に続けて、$H^1(X, \mathcal{F}) = 0$ を得る。
従って写像
$$
\xymatrix{
\bigoplus\nolimits_{i = 1, \ldots, n} \Gamma(X, \mathcal{O}_X)
\ar[rr]^{f_1, \ldots, f_n} & &
\Gamma(X, \mathcal{O}_X)
}
$$
は望んだ通り全射である。
\end{proof}

\noindent
$X$ が Noether スキームならば、任意の準連接イデアル層は自動的に
連接イデアル層であり、有限型の準連接イデアル層でもあることに注意せよ。
従って、この場合には前の補題と次の補題がともに適用できる。

\begin{lemma}
\label{coherent-lemma-quasi-separated-h1-zero-covering}
\begin{reference}
\cite{Serre-criterion}, \cite[II, Theorem 5.2.1]{EGA}
\end{reference}
\begin{slogan}
アフィン性の Serre 判定法。
\end{slogan}
$X$ をスキームとする。次を仮定する。
\begin{enumerate}
\item $X$ は準コンパクトである。
\item $X$ は準分離的である。
\item $H^1(X, \mathcal{I}) = 0$ が、有限型の任意の準連接イデアル層
$\mathcal{I}$ に対して成り立つ。
\end{enumerate}
このとき $X$ はアフィンである。
\end{lemma}

\begin{proof}
性質、補題
\ref{properties-lemma-quasi-coherent-colimit-finite-type} により、
任意の準連接イデアル層は有限型の準連接イデアル層の有向余極限である。
コホモロジー、補題
\ref{cohomology-lemma-quasi-separated-cohomology-colimit} により、
$X$ 上でコホモロジーを取ることは有向余極限と可換である。
従って、$H^1(X, \mathcal{I}) = 0$ が任意の準連接イデアル層
（$X$ 上のもの）に対して成り立つ。
言い換えると、補題 \ref{coherent-lemma-quasi-compact-h1-zero-covering} を適用できる。
\end{proof}

\noindent
上で与えた議論を用いると、可逆層が豊富となるための十分条件を得られる。
ただし、この条件は必要条件ではないことを読者に注意しておく。

\begin{lemma}
\label{coherent-lemma-quasi-compact-h1-zero-invertible}
$X$ をスキームとする。$\mathcal{L}$ を可逆 $\mathcal{O}_X$-加群とする。
次を仮定する。
\begin{enumerate}
\item $X$ は準コンパクトである。
\item 任意の準連接イデアル層
$\mathcal{I} \subset \mathcal{O}_X$ に対して、ある $n \geq 1$ が存在して
$H^1(X, \mathcal{I} \otimes_{\mathcal{O}_X} \mathcal{L}^{\otimes n}) = 0$
となる。
\end{enumerate}
このとき $\mathcal{L}$ は豊富である。
\end{lemma}

\begin{proof}
補題 \ref{coherent-lemma-quasi-compact-h1-zero-covering} とまったく同じ方法で証明する。
$x \in X$ を閉点とする。$U \subset X$ を $x$ のアフィン開近傍で、
$\mathcal{L}|_U \cong \mathcal{O}_U$ となるものとする。
$U = \Spec(A)$ と書き、$\mathfrak m \subset A$ を $x$ に対応する極大イデアルとする。
$Z = X \setminus U$ および $Z' = Z \cup \{x\}$ と置く。
スキーム、補題 \ref{schemes-lemma-reduced-closed-subscheme} により、
準連接イデアル層 $\mathcal{I}$ および $\mathcal{I}'$ で、
被約閉部分スキーム $Z$ および $Z'$ をそれぞれ切り出すものが存在する。
短完全列
$$
0 \to \mathcal{I}' \to \mathcal{I} \to \mathcal{I}/\mathcal{I}' \to 0.
$$
を考える。すべての $n \geq 1$ に対して短完全列
$$
0 \to \mathcal{I}' \otimes_{\mathcal{O}_X} \mathcal{L}^{\otimes n}
\to \mathcal{I} \otimes_{\mathcal{O}_X} \mathcal{L}^{\otimes n} \to
\mathcal{I}/\mathcal{I}' \otimes_{\mathcal{O}_X} \mathcal{L}^{\otimes n} \to 0.
$$
を得る。仮定により、$n$ を
$H^1(X, \mathcal{I}' \otimes_{\mathcal{O}_X} \mathcal{L}^{\otimes n}) = 0$
となるように選べる。
$x$ は $X$ の閉点であり、$x \not \in Z$ なので、
$\mathcal{I}/\mathcal{I}'$ は $x$ に台をもつ。実際、
$\mathcal{I}/\mathcal{I'}$ の $U$ への制限は $A$-加群
$A/\mathfrak m$ に対応する。$\mathcal{L}$ は $U$ 上自明であるから、
$\mathcal{I}/\mathcal{I}' \otimes_{\mathcal{O}_X} \mathcal{L}^{\otimes n}$
の $U$ への制限も $A$-加群 $A/\mathfrak m$ に対応する。
従って
$\Gamma(X, \mathcal{I}/\mathcal{I'} \otimes_{\mathcal{O}_X}
\mathcal{L}^{\otimes n}) = A/\mathfrak m$
である。$n$ の選び方により、大域切断
$s \in \Gamma(X, \mathcal{I} \otimes_{\mathcal{O}_X} \mathcal{L}^{\otimes n})$
で元 $1 \in A/\mathfrak m$ に写るものが存在する。
明らかに $x \in X_s \subset U$ である。実際、$s$ は $Z$ の点で消える。
従って $X_s = D(f)$ である。ここで
$f \in A$ は $s$ の $A \cong \Gamma(U, \mathcal{L}^{\otimes n})$ における
像である。特に $X_s$ はアフィンである。

\medskip\noindent
$W = \bigcup X_s$ を考える。ここで合併は、すべての
$s \in \Gamma(X, \mathcal{L}^{\otimes n})$ と $n \geq 1$ で
$X_s$ がアフィンとなるものにわたる。明らかに $W$ は $X$ の開集合である。
上の議論により、$X$ のすべての閉点は $W$ に含まれる。
閉部分集合 $X \setminus W$（$X$ のもの）も準コンパクトである
（位相、補題 \ref{topology-lemma-closed-in-quasi-compact} を参照）。
従って、空でなければ閉点をもつ
（位相、補題 \ref{topology-lemma-quasi-compact-closed-point} を参照）。
これはすべての閉点が $W$ にあることに矛盾する。
従って $X = W$ である。このことは、性質、定義
\ref{properties-definition-ample} により $\mathcal{L}$ が豊富であることを意味する。
\end{proof}

\noindent
補題 \ref{coherent-lemma-quasi-compact-h1-zero-invertible} には、有限型イデアル層を用いる
変種がある。必要になれば、ここで定式化して証明する。

\begin{lemma}
\label{coherent-lemma-criterion-affine-morphism}
$f : X \to Y$ を準コンパクトな射とし、$X$ と $Y$ は準分離的であるとする。
$R^1f_*\mathcal{I} = 0$ が任意の準連接イデアル層
$\mathcal{I}$（$X$ 上のもの）に対して成り立つならば、$f$ はアフィンである。
\end{lemma}

\begin{proof}
$V \subset Y$ をアフィン開部分スキームとする。
$U = f^{-1}(V)$ がアフィンであることを示さなければならない。
包含射 $V \to Y$ はスキーム、補題
\ref{schemes-lemma-quasi-compact-permanence} により準コンパクトである。
従って、基底変換 $U \to X$ は準コンパクトである。スキーム、補題
\ref{schemes-lemma-quasi-compact-preserved-base-change} を参照せよ。
ゆえに、任意の準連接イデアル層 $\mathcal{I}$（$U$ 上のもの）は
$X$ 上の準連接イデアル層へ延長できる。性質、補題
\ref{properties-lemma-extend-trivial} を参照せよ。
$R^1f_*$ の構成は $Y$ 上局所的である
（コホモロジー、節 \ref{cohomology-section-locality}）から、
補題の仮定により $R^1(U \to V)_*\mathcal{I} = 0$ である。
従って、Leray スペクトル系列
（コホモロジー、補題 \ref{cohomology-lemma-Leray}）により、
$H^1(U, \mathcal{I}) = H^1(V, (U \to V)_*\mathcal{I})$ である。
スキーム、補題 \ref{schemes-lemma-push-forward-quasi-coherent} により
$(U \to V)_*\mathcal{I}$ は準連接であるから、補題
\ref{coherent-lemma-quasi-coherent-affine-cohomology-zero} により
$H^1(V, (U \to V)_*\mathcal{I}) = 0$ である。
従って、補題 \ref{coherent-lemma-quasi-compact-h1-zero-covering} により
$U$ はアフィンである。
\end{proof}














\section{高次順像の準連接性}
\label{coherent-section-quasi-coherence}

\noindent
アフィン・スキーム上では、準連接加群の高次コホモロジー群が零になることを見た。
（準）分離的な準コンパクト・スキーム $X$ については、このことから、ある次数を
超える準連接層のコホモロジー群が消滅することが従う。しかし、$X$ の
コホモロジー次元が有限であるとは限らない。というのも、それは
{\it すべての} $\mathcal{O}_X$-加群のコホモロジーの消滅によって定義されるからである。

\begin{lemma}[帰納原理]
\label{coherent-lemma-induction-principle}
\begin{reference}
\cite[Proposition 3.3.1]{BvdB}
\end{reference}
$X$ を準コンパクトかつ準分離的なスキームとする。$P$ を $X$ の準コンパクトな
開集合の性質とする。次を仮定する。
\begin{enumerate}
\item $P$ は $X$ のすべてのアフィン開集合について成り立つ。
\item $U$ が準コンパクトな開集合、$V$ がアフィン開集合であり、
$P$ が $U$, $V$, および $U \cap V$ について成り立つならば、
$P$ は $U \cup V$ についても成り立つ。
\end{enumerate}
このとき $P$ は $X$ のすべての準コンパクトな開集合について成り立ち、
特に $X$ について成り立つ。
\end{lemma}

\begin{proof}
まず、$P$ が {\it 分離的な} 準コンパクト開集合 $W \subset X$ について
成り立つことを帰納法で示す。実際、このような開集合は
$W = U_1 \cup \ldots \cup U_n$ と書け、$n$ に関する帰納法を行える。
その際、性質 (2) を $U = U_1 \cup \ldots \cup U_{n - 1}$ および
$V = U_n$ に適用する。これは、
$U \cap V = (U_1 \cap U_n) \cup \ldots \cup (U_{n - 1} \cap U_n)$
も $n - 1$ 個のアフィン開部分スキームの合併だから許される。実際、これは
アフィン開集合 $U_i$ と $U_n$（いずれも $W$ のもの）にスキーム、補題
\ref{schemes-lemma-characterize-separated} を適用すれば従う。
以上を踏まえると、任意の準コンパクト開集合 $W \subset X$ に対し、
$W$ を被覆するのに必要なアフィン開集合の個数について、先ほどと同じ仕方で
帰納法を行える。ここでは、準コンパクト開集合 $U_i \cap U_n$ が
アフィン・スキーム $U_n$ の開部分スキームとして分離的であることを用いる。
\end{proof}

\begin{lemma}
\label{coherent-lemma-vanishing-nr-affines}
\begin{slogan}
対角射がアフィンなスキームでは、準連接加群のコホモロジーは、被覆に必要な
アフィン開集合の個数以上の次数で消滅する。
\end{slogan}
$X$ を対角射がアフィンな準コンパクト・スキームとする
（例えば $X$ が分離的ならばよい）。
$t = t(X)$ を $X$ を被覆するのに必要なアフィン開集合の最小個数とする。
このとき $H^n(X, \mathcal{F}) = 0$ が、すべての $n \geq t$ と
すべての準連接層 $\mathcal{F}$ に対して成り立つ。
\end{lemma}

\begin{proof}
第一の証明。
$t$ に関する帰納法による。
$t = 1$ ならば、結果は補題
\ref{coherent-lemma-quasi-coherent-affine-cohomology-zero} から従う。
$t > 1$ ならば、$X = U \cup V$ と書く。ここで $V$ はアフィン開集合であり、
$U = U_1 \cup \ldots \cup U_{t - 1}$ は $t - 1$ 個のアフィン開集合の合併である。
この場合、
$U \cap V =  (U_1 \cap V) \cup \ldots (U_{t - 1} \cap V)$
も $t - 1$ 個のアフィン開部分スキームの合併であることに注意せよ。
実際、対角射がアフィンなので、二つのアフィン開集合の共通部分は
アフィンである。補題 \ref{coherent-lemma-affine-diagonal} を参照せよ。
Mayer--Vietoris 長完全系列
$$
0 \to
H^0(X, \mathcal{F}) \to
H^0(U, \mathcal{F}) \oplus H^0(V, \mathcal{F}) \to
H^0(U \cap V, \mathcal{F}) \to
H^1(X, \mathcal{F}) \to \ldots
$$
を適用する。コホモロジー、補題
\ref{cohomology-lemma-mayer-vietoris} を参照せよ。
帰納法により、群 $H^i(U, \mathcal{F})$,
$H^i(V, \mathcal{F})$, $H^i(U \cap V, \mathcal{F})$ は
$i \geq t - 1$ に対して零である。直ちに、$H^i(X, \mathcal{F})$ は
$i \geq t$ に対して零であることが従う。

\medskip\noindent
第二の証明。
$\mathcal{U} : X = \bigcup_{i = 1}^t U_i$ を有限アフィン開被覆とする。
$X$ は対角射がアフィンなので、多重共通部分
$U_{i_0 \ldots i_p}$ はすべてアフィンである。補題
\ref{coherent-lemma-affine-diagonal} を参照せよ。
補題 \ref{coherent-lemma-cech-cohomology-quasi-coherent} により、{\v C}ech
コホモロジー群 $\check{H}^p(\mathcal{U}, \mathcal{F})$ は
コホモロジー群と一致する。コホモロジー、補題
\ref{cohomology-lemma-alternating-usual} により、{\v C}ech コホモロジー群は
交代 {\v C}ech 複体 $\check{\mathcal{C}}_{alt}^\bullet(\mathcal{U}, \mathcal{F})$
を用いて計算できる。被覆は $t$ 個の元からなるので、直ちに
$\check{\mathcal{C}}_{alt}^p(\mathcal{U}, \mathcal{F}) = 0$ が
すべての $p \geq t$ に対して成り立つ。従って結果が従う。
\end{proof}

\begin{lemma}
\label{coherent-lemma-affine-diagonal-universal-delta-functor}
$X$ を対角射がアフィンな準コンパクト・スキームとする
（例えば $X$ が分離的ならばよい）。このとき次が成り立つ。
\begin{enumerate}
\item 準連接 $\mathcal{O}_X$-加群 $\mathcal{F}$ が与えられると、埋め込み
$\mathcal{F} \to \mathcal{F}'$（準連接 $\mathcal{O}_X$-加群のもの）で、
$H^p(X, \mathcal{F}') = 0$ がすべての $p \geq 1$ に対して成り立つものが存在する。
\item $\{H^n(X, -)\}_{n \geq 0}$ は普遍 $\delta$-関手であり、
$\QCoh(\mathcal{O}_X)$ から $\textit{Ab}$ への関手である。
\end{enumerate}
\end{lemma}

\begin{proof}
$X = \bigcup U_i$ を有限アフィン開被覆とする。
$U = \coprod U_i$ とおき、$j : U \to X$ を、与えられた開埋め込み
$U_i \to X$ を誘導する射とする。$U$ はアフィン・スキームであり、$X$ の対角射は
アフィンなので、射 $j$ はアフィンである。射、補題
\ref{morphisms-lemma-affine-permanence} を参照せよ。
任意の $\mathcal{O}_X$-加群 $\mathcal{F}$ に対し、標準射
$\mathcal{F} \to j_*j^*\mathcal{F}$ がある。
この射が単射であることは茎で調べれば分かる。すなわち、$x \in U_i$ ならば
分解
$$
\mathcal{F}_x \to (j_*j^*\mathcal{F})_x
\to (j^*\mathcal{F})_{x'} = \mathcal{F}_x
$$
がある。ここで $x' \in U$ は、点 $x$ を $U_i \subset U$ の点として
見たものである。
さて、$\mathcal{F}$ が準連接ならば、$j^*\mathcal{F}$ はアフィン・スキーム
$U$ 上準連接であり、従って補題
\ref{coherent-lemma-quasi-coherent-affine-cohomology-zero} により高次コホモロジーが消滅する。
このとき $H^p(X, j_*j^*\mathcal{F}) = 0$ が $p > 0$ に対して、補題
\ref{coherent-lemma-relative-affine-cohomology} により成り立つ。実際 $j$ はアフィンである。
これで (1) が証明された。
最後に、写像
$H^p(X, \mathcal{F}) \to H^p(X, j_*j^*\mathcal{F})$
は零であり、(2) はホモロジー、補題
\ref{homology-lemma-efface-implies-universal} から従う。
\end{proof}

\begin{lemma}
\label{coherent-lemma-vanishing-nr-affines-quasi-separated}
$X$ を準コンパクトかつ準分離的なスキームとする。
$X = U_1 \cup \ldots \cup U_n$ を、各 $U_i$ が準コンパクトかつ分離的
（例えばアフィン）であるような開被覆とする。次のようにおく。
$$
d = \max\nolimits_{I \subset \{1, \ldots, n\}}
\left(|I| + t(\bigcap\nolimits_{i \in I} U_i) - 1\right)
$$
ここで $t(U)$ はスキーム $U$ を被覆するのに必要なアフィン開集合の
最小個数である。このとき $H^p(X, \mathcal{F}) = 0$ が、すべての
$p \geq d$ とすべての準連接層 $\mathcal{F}$ に対して成り立つ。
\end{lemma}

\begin{proof}
$X$ は準分離的であり $U_i$ は準コンパクトなので、数
$t(\bigcap_{i \in I} U_i)$ は有限である。$n$ に関する帰納法で証明する。
$n = 1$ ならば、結果は補題 \ref{coherent-lemma-vanishing-nr-affines} から従う。
$n > 1$ ならば、$X = U \cup V$ と書く。ここで
$U = U_1 \cup \ldots \cup U_{n - 1}$ および $V = U_n$ である。
Mayer--Vietoris 長完全系列
$$
0 \to
H^0(X, \mathcal{F}) \to
H^0(U, \mathcal{F}) \oplus H^0(V, \mathcal{F}) \to
H^0(U \cap V, \mathcal{F}) \to
H^1(X, \mathcal{F}) \to \ldots
$$
を適用する。コホモロジー、補題
\ref{cohomology-lemma-mayer-vietoris} を参照せよ。
$q \geq d$ に対して証明を終えるため、
$H^q(V, \mathcal{F})$, $H^q(U, \mathcal{F})$, および
$H^{q - 1}(U \cap V, \mathcal{F})$ が消滅することを示す。
$n = 1$ の場合から、$H^q(V, \mathcal{F}) = 0$ が
$q \geq t(V) = t(U_n)$ に対して成り立つ。$t(V) \leq d$ なので、これは所望の
主張を与える。帰納法の仮定により $H^q(U, \mathcal{F}) = 0$ が
$$
q \geq \max\nolimits_{I \subset \{1, \ldots, n - 1\}}
\left(|I| + t(\bigcap\nolimits_{i \in I} U_i) - 1\right)
$$
に対して成り立つ。右辺の整数は $d$ 以下なので、これも所望の主張を与える。
最後に、開集合
$U \cap V = (U_1 \cap U_n) \cup \ldots \cup (U_{n - 1} \cap U_n)$
に帰納法の仮定を適用すると、$H^q(U \cap V, \mathcal{F}) = 0$ が
$$
q \geq \max\nolimits_{I \subset \{1, \ldots, n - 1\}}
\left(|I| + t(U_n \cap \bigcap\nolimits_{i \in I} U_i) - 1\right)
$$
に対して成り立つ。右辺の整数は $d$ より真に小さいので、補題が従う。
\end{proof}

\begin{lemma}
\label{coherent-lemma-quasi-coherence-higher-direct-images}
$f : X \to S$ をスキームの射とする。
$f$ が準分離的かつ準コンパクトであると仮定する。
\begin{enumerate}
\item 任意の準連接 $\mathcal{O}_X$-加群 $\mathcal{F}$ に対し、高次順像
$R^pf_*\mathcal{F}$ は $S$ 上準連接である。
\item $S$ が準コンパクトならば、整数 $n = n(X, S, f)$ が存在し、
$R^pf_*\mathcal{F} = 0$ が、すべての $p \geq n$ と任意の準連接層
$\mathcal{F}$（$X$ 上のもの）に対して成り立つ。
\item 実際、$S$ が準コンパクトならば、$n = n(X, S, f)$ を次のように
選べる。スキームの任意の射 $S' \to S$ に対し、
$R^p(f')_*\mathcal{F}' = 0$ が、$p \geq n$ と任意の準連接層
$\mathcal{F}'$（$X'$ 上のもの）について成り立つ。ここで
$f' : X' = S' \times_S X \to S'$ は $f$ の基底変換である。
\end{enumerate}
\end{lemma}

\begin{proof}
まず (1) を証明する。補題の仮定のもとで層
$R^0f_*\mathcal{F} = f_*\mathcal{F}$ は、スキーム、補題
\ref{schemes-lemma-push-forward-quasi-coherent} により準連接であることに注意する。
コホモロジー、補題
\ref{cohomology-lemma-localize-higher-direct-images} を用いると、
高次順像の形成が開部分スキームへの制限と可換であることが分かる。
準連接性は $S$ 上局所的なので、次の段落で扱う場合に帰着する。

\medskip\noindent
$S$ がアフィンの場合の (1) の証明。帰納原理を用いる。
$f$ は準コンパクトかつ準分離的なので、$X$ は準コンパクトかつ
準分離的である。準コンパクト開集合 $U \subset X$ と $a = f|_U$ に対し、
$P(U)$ を、$R^pa_*\mathcal{F}$ が $S$ 上準連接である、という性質とする。
ここで $\mathcal{F}$ は $U$ 上の任意の準連接加群であり、$p \geq 0$ である。
$P(X)$ は (1) そのものなので、
補題 \ref{coherent-lemma-induction-principle} の条件 (1) と (2) が成り立つことを
示せば十分である。$U$ がアフィンならば、$P(U)$ は成り立つ。実際、
$R^pa_*\mathcal{F} = 0$ が $p \geq 1$ に対して成り立つ
（補題 \ref{coherent-lemma-relative-affine-vanishing} および射、補題
\ref{morphisms-lemma-morphism-affines-affine} による）。また、最初の段落で
$p = 0$ の場合には結果が成り立つことをすでに述べた。
次に、$U \subset X$ を準コンパクト開集合、$V \subset X$ をアフィン開集合とし、
$P(U)$, $P(V)$, $P(U \cap V)$ が成り立つと仮定する。
$a = f|_U$, $b = f|_V$, $c = f|_{U \cap V}$, $g = f|_{U \cup V}$ とおく。
すると、任意の準連接 $\mathcal{O}_{U \cup V}$-加群 $\mathcal{F}$ に対して、
相対 Mayer--Vietoris 系列
$$
0 \to
g_*\mathcal{F} \to
a_*(\mathcal{F}|_U) \oplus b_*(\mathcal{F}|_V) \to
c_*(\mathcal{F}|_{U \cap V}) \to
R^1g_*\mathcal{F} \to \ldots
$$
がある。コホモロジー、補題
\ref{cohomology-lemma-relative-mayer-vietoris} を参照せよ。
$P(U)$, $P(V)$, $P(U \cap V)$ により、
$R^pa_*(\mathcal{F}|_U)$, $R^pb_*(\mathcal{F}|_V)$ および
$R^pc_*(\mathcal{F}|_{U \cap V})$ はすべて準連接である。
スキーム、節 \ref{schemes-section-quasi-coherent} の準連接層に関する結果を
用いると、各層 $R^pg_*\mathcal{F}$ が準連接であることが従う。実際、それは
左側にある準連接層間の写像の余核と、右側にある準連接層間の写像の核とからなる
短完全系列の中央に位置する。従って $P(U \cup V)$ が成り立ち、(1) の証明が完了する。

\medskip\noindent
次に (3)、従って特に (2) を証明する。有限アフィン開被覆
$S = \bigcup_{j = 1, \ldots m} S_j$ を選ぶ。各 $j$ に対して有限アフィン開被覆
$f^{-1}(S_j) = \bigcup_{i = 1, \ldots t_j} U_{ji} $ を選ぶ。次のようにおく。
$$
d_j = \max\nolimits_{I \subset \{1, \ldots, t_j\}}
\left(|I| + t(\bigcap\nolimits_{i \in I} U_{ji})\right)
$$
これは補題 \ref{coherent-lemma-vanishing-nr-affines-quasi-separated} で得られる整数である。
$n(X, S, f) = \max d_j$ が機能すると主張する。

\medskip\noindent
実際、$S' \to S$ をスキームの射とし、$\mathcal{F}'$ を
$X' = S' \times_S X$ 上の準連接層とする。
$R^pf'_*\mathcal{F}' = 0$ が $p \geq n(X, S, f)$ に対して成り立つことを示したい。
この問題は $S'$ 上局所的なので、$S'$ はアフィンで、ある $S_j$ へ写ると
仮定してよい（その添字を $j$ とする）。このとき
$X' = S' \times_{S_j} f^{-1}(S_j)$ はアフィン開集合
$S' \times_{S_j} U_{ji}$（$i = 1, \ldots t_j$）で被覆され、共通部分
$$
\bigcap\nolimits_{i \in I} S' \times_{S_j} U_{ji} =
S' \times_{S_j} \bigcap\nolimits_{i \in I} U_{ji}
$$
は基底変換前と同じ個数のアフィン開集合で被覆される。
補題 \ref{coherent-lemma-vanishing-nr-affines-quasi-separated} を適用すると
$H^p(X', \mathcal{F}') = 0$ を得る。証明の第一部により、各
$R^qf'_*\mathcal{F}'$ が準連接であることはすでに分かっており、従って
アフィン・スキーム $S'$ 上でその高次コホモロジー群は消滅する。
ゆえに、コホモロジー、補題 \ref{cohomology-lemma-apply-Leray} により
$H^0(S', R^pf'_*\mathcal{F}') = H^p(X', \mathcal{F}') = 0$ である。
$R^pf'_*\mathcal{F}'$ は準連接なので、$R^pf'_*\mathcal{F}' = 0$ と結論する。
\end{proof}

\begin{lemma}
\label{coherent-lemma-quasi-coherence-higher-direct-images-application}
$f : X \to S$ をスキームの射とする。
$f$ が準分離的かつ準コンパクトであると仮定する。
$S$ がアフィンであると仮定する。
任意の準連接 $\mathcal{O}_X$-加群 $\mathcal{F}$ に対し、
$$
H^q(X, \mathcal{F}) = H^0(S, R^qf_*\mathcal{F})
$$
がすべての $q \in \mathbf{Z}$ に対して成り立つ。
\end{lemma}

\begin{proof}
$E_2^{p, q} = H^p(S, R^qf_*\mathcal{F})$ であり
$H^{p + q}(X, \mathcal{F})$ に収束する Leray スペクトル系列を考える。
コホモロジー、補題 \ref{cohomology-lemma-Leray} を参照せよ。
補題 \ref{coherent-lemma-quasi-coherence-higher-direct-images} により、層
$R^qf_*\mathcal{F}$ は準連接である。
補題 \ref{coherent-lemma-quasi-coherent-affine-cohomology-zero} により、
$E_2^{p, q} = 0$ が $p > 0$ のとき成り立つ。
従ってスペクトル系列は $E_2$ で退化し、主張が従う。
一般原理については、コホモロジー、補題
\ref{cohomology-lemma-apply-Leray} (2) も参照せよ。
\end{proof}

\section{コホモロジーと基底変換 I}
\label{coherent-section-cohomology-and-base-change}

\noindent
$f : X \to S$ をスキームの射とする。
$\mathcal{F}$ を $X$ 上の準連接層とする。
さらに、$g : S' \to S$ をスキームの任意の射とする。
$X' = X_{S'} = S' \times_S X$ を $X$ の基底変換と書き、
$f' : X' \to S'$ を $f$ の基底変換と書く。
また $g' : X' \to X$ を射影と書き、
$\mathcal{F}' = (g')^*\mathcal{F}$ とおく。
この状況を表す図式は次の通りである。
\begin{equation}
\label{coherent-equation-base-change-diagram}
\vcenter{
\xymatrix{
\mathcal{F}' = (g')^*\mathcal{F} &
X' \ar[r]_{g'} \ar[d]_{f'} &
X \ar[d]^f &
\mathcal{F} \\
Rf'_*\mathcal{F}' &
S' \ar[r]^g &
S &
Rf_*\mathcal{F}
}
}
\end{equation}
念頭にある基底変換性質の最も単純な場合を述べる。

\begin{lemma}
\label{coherent-lemma-affine-base-change}
$f : X \to S$ をスキームの射とする。
$\mathcal{F}$ を準連接 $\mathcal{O}_X$-加群とする。
$f$ がアフィンであると仮定する。この場合、
$f_*\mathcal{F} \cong Rf_*\mathcal{F}$ は準連接層であり、すべての基底変換図式
（\ref{coherent-equation-base-change-diagram}）に対して
$$
g^*f_*\mathcal{F} = f'_*(g')^*\mathcal{F}.
$$
が成り立つ。
\end{lemma}

\begin{proof}
高次順像の消滅は補題 \ref{coherent-lemma-relative-affine-vanishing} である。
この主張は $S$ と $S'$ 上局所的である。従って
$X = \Spec(A)$, $S = \Spec(R)$,
$S' = \Spec(R')$ かつ $\mathcal{F} = \widetilde{M}$ と仮定してよい。
ここで、ある $A$-加群 $M$ を用いた。
$\mathcal{F}$ の引戻しと順像を記述するため、スキーム、補題
\ref{schemes-lemma-widetilde-pullback} を用いる。すなわち、
$X' = \Spec(R' \otimes_R A)$ であり、$\mathcal{F}'$ は
$(R' \otimes_R A) \otimes_A M$ に付随する準連接層である。
従って補題は、$R'$-加群としての等式
$$
(R' \otimes_R A) \otimes_A M = R' \otimes_R M
$$
に帰着する。
\end{proof}

\noindent
多くの状況では、次に述べるコホモロジーと基底変換の特別な場合を知れば十分である。
これは節 \ref{coherent-section-cohomology-and-base-change-derived} のより強い結果から
直ちに従うが、非常に重要なので独立した証明を与えるに値する。

\begin{lemma}[平坦基底変換]
\label{coherent-lemma-flat-base-change-cohomology}
スキームのデカルト図式
$$
\xymatrix{
X' \ar[d]_{f'} \ar[r]_{g'} & X \ar[d]^f \\
S' \ar[r]^g & S
}
$$
を考える。$\mathcal{F}$ を準連接 $\mathcal{O}_X$-加群とし、その引戻しを
$\mathcal{F}' = (g')^*\mathcal{F}$ とする。
$g$ が平坦で、$f$ が準コンパクトかつ準分離的であると仮定する。
任意の $i \geq 0$ に対して次が成り立つ。
\begin{enumerate}
\item コホモロジー、補題
\ref{cohomology-lemma-base-change-map-flat-case} の基底変換写像は同型
$$
g^*R^if_*\mathcal{F} \longrightarrow R^if'_*\mathcal{F}',
$$
である。
\item $S = \Spec(A)$ かつ $S' = \Spec(B)$ ならば、
$H^i(X, \mathcal{F}) \otimes_A B = H^i(X', \mathcal{F}')$ である。
\end{enumerate}
\end{lemma}

\begin{proof}
(1) でコホモロジー、補題
\ref{cohomology-lemma-base-change-map-flat-case} を用いてよい。実際、射、補題
\ref{morphisms-lemma-base-change-flat} により $g'$ は平坦である。
これを踏まえると、(1) は (2) から従う。実際、(1) は $S'$ 上局所的なので、
$S$ と $S'$ はアフィンであると仮定してよい。言い換えると、(2) と同様に
$S = \Spec(A)$ および $S' = \Spec(B)$ である。
このとき $R^if_*\mathcal{F}$ は準連接なので
（補題 \ref{coherent-lemma-quasi-coherence-higher-direct-images}）、それは準連接
$\mathcal{O}_S$-加群であり、$A$-加群
$H^0(S, R^if_*\mathcal{F}) = H^i(X, \mathcal{F})$ に付随する
（等式は補題 \ref{coherent-lemma-quasi-coherence-higher-direct-images-application} による）。
同様に、$R^if'_*\mathcal{F}'$ は準連接 $\mathcal{O}_{S'}$-加群であり、
$B$-加群 $H^i(X', \mathcal{F}')$ に付随する。
$g$ による引戻しは加群上の $- \otimes_A B$ に対応するので
（スキーム、補題 \ref{schemes-lemma-widetilde-pullback}）、
(2) を証明すれば十分である。

\medskip\noindent
$A \to B$ を平坦な環準同型とする。
$X$ を $A$ 上の準コンパクトかつ準分離的なスキームとする。
$\mathcal{F}$ を準連接 $\mathcal{O}_X$-加群とする。
$X_B = X \times_{\Spec(A)} \Spec(B)$ とおき、$\mathcal{F}_B$ を
$\mathcal{F}$ の引戻しとする。写像
$$
H^i(X, \mathcal{F}) \otimes_A B \longrightarrow H^i(X_B, \mathcal{F}_B)
$$
（補題の主張で参照された構成による）が同型であることを示したい。

\medskip\noindent
$X$ が分離的な場合、アフィン開被覆
$\mathcal{U} : X = U_1 \cup \ldots \cup U_t$ を選び、
$$
\check{H}^p(\mathcal{U}, \mathcal{F}) = H^p(X, \mathcal{F}),
$$
であることを思い出す。補題
\ref{coherent-lemma-cech-cohomology-quasi-coherent} を参照せよ。
基底変換により得られる
$\mathcal{U}_B : X_B = (U_1)_B \cup \ldots \cup (U_t)_B$ についても、
各 $(U_i)_B$ はアフィンであり、$X_B$ は分離的である。従って
$$
\check{H}^p(\mathcal{U}_B, \mathcal{F}_B) = H^p(X_B, \mathcal{F}_B).
$$
さらに、{\v C}ech 複体の間には関係
$$
\check{\mathcal{C}}^\bullet(\mathcal{U}_B, \mathcal{F}_B) =
\check{\mathcal{C}}^\bullet(\mathcal{U}, \mathcal{F}) \otimes_A B
$$
がある。これは補題 \ref{coherent-lemma-affine-base-change} から従う。
$A \to B$ は平坦なので、コホモロジーを取っても同じことが成り立つ。

\medskip\noindent
$X$ が準分離的な場合、アフィン開被覆
$\mathcal{U} : X = U_1 \cup \ldots \cup U_t$ を選ぶ。
コホモロジー、補題 \ref{cohomology-lemma-cech-spectral-sequence} の
{\v C}ech-コホモロジー・スペクトル系列を用いる。
このスペクトル系列を避けたい読者は、補題
\ref{coherent-lemma-quasi-coherence-higher-direct-images} の証明と同様に、
Mayer--Vietoris 系列と $t$ に関する帰納法を用いればよい。
スペクトル系列の $E_2$-項は
$E_2^{p, q} = \check{H}^p(\mathcal{U}, \underline{H}^q(\mathcal{F}))$
であり、$H^{p + q}(X, \mathcal{F})$ に収束する。
同様に、$E_2$-項が
$E_2^{p, q} = \check{H}^p(\mathcal{U}_B, \underline{H}^q(\mathcal{F}_B))$
であり、$H^{p + q}(X_B, \mathcal{F}_B)$ に収束するスペクトル系列がある。
共通部分 $U_{i_0 \ldots i_p}$ は準コンパクトかつ分離的なので、
証明の第二段落の結果は
$\check{H}^p(\mathcal{U}_B, \underline{H}^q(\mathcal{F}_B)) =
\check{H}^p(\mathcal{U}, \underline{H}^q(\mathcal{F})) \otimes_A B$
を与える。$A \to B$ が平坦であることを用いると、
$H^i(X, \mathcal{F}) \otimes_A B \to H^i(X_B, \mathcal{F}_B)$
はすべての $i$ について同型であると結論でき、主張が従う。
\end{proof}

\begin{lemma}[有限局所自由基底変換]
\label{coherent-lemma-finite-locally-free-base-change-cohomology}
スキームのデカルト図式
$$
\xymatrix{
Y \ar[d]_{g} \ar[r]_h & X \ar[d]^f \\
\Spec(B) \ar[r] & \Spec(A)
}
$$
を考える。$\mathcal{F}$ を準連接 $\mathcal{O}_X$-加群とし、その引戻しを
$\mathcal{G} = h^*\mathcal{F}$ とする。
$B$ が有限局所自由 $A$-加群ならば、
$H^i(X, \mathcal{F}) \otimes_A B = H^i(Y, \mathcal{G})$ である。
\end{lemma}

\noindent
{\bf 警告}：この補題と補題 \ref{coherent-lemma-flat-base-change-cohomology} の違いを
理解していない限り、この補題を用いてはならない。

\begin{proof}
$X$ が分離的な場合、アフィン開被覆
$\mathcal{U} : X = \bigcup_{i \in I} U_i$ を選び、
$$
\check{H}^p(\mathcal{U}, \mathcal{F}) = H^p(X, \mathcal{F}),
$$
であることを思い出す。補題
\ref{coherent-lemma-cech-cohomology-quasi-coherent} を参照せよ。
$\mathcal{V} : Y = \bigcup_{i \in I} g^{-1}(U_i)$ を、対応する $Y$ の
アフィン開被覆とする。開集合
$V_i = g^{-1}(U_i) = U_i \times_{\Spec(A)} \Spec(B)$ はアフィンであり、
$Y$ は分離的である。従って
$$
\check{H}^p(\mathcal{V}, \mathcal{G}) = H^p(Y, \mathcal{G}).
$$
{\v C}ech 複体の写像
$$
\check{\mathcal{C}}^\bullet(\mathcal{U}, \mathcal{F}) \otimes_A B
\longrightarrow
\check{\mathcal{C}}^\bullet(\mathcal{V}, \mathcal{G})
$$
は同型であると主張する。実際、$B$ は $A$-加群として有限表示なので、
$B$ とのテンソル積を $A$ 上で取る操作は積と可換する。代数、命題
\ref{algebra-proposition-fp-tensor} を参照せよ。
従って、写像
$\Gamma(U_{i_0 \ldots i_p}, \mathcal{F}) \otimes_A B \to
\Gamma(V_{i_0 \ldots i_p}, \mathcal{G})$
が同型であることを示せば十分であり、これは補題
\ref{coherent-lemma-affine-base-change} から従う。
$A \to B$ は平坦なので、コホモロジーを取っても同じことが成り立つ。

\medskip\noindent
一般の場合も、アフィン開被覆
$\mathcal{U} : X = \bigcup_{i \in I} U_i$ と、対応する被覆
$\mathcal{V} : Y = \bigcup_{i \in I} V_i$（ここで
$V_i = g^{-1}(U_i)$）を用いて、まったく同様に論じる。
コホモロジー、補題 \ref{cohomology-lemma-cech-spectral-sequence} の
{\v C}ech-コホモロジー・スペクトル系列を用いる。
このスペクトル系列の $E_2$-項は
$E_2^{p, q} = \check{H}^p(\mathcal{U}, \underline{H}^q(\mathcal{F}))$
であり、$H^{p + q}(X, \mathcal{F})$ に収束する。
同様に、$E_2$-項が
$E_2^{p, q} = \check{H}^p(\mathcal{V}, \underline{H}^q(\mathcal{G}))$
であり、$H^{p + q}(Y, \mathcal{G})$ に収束するスペクトル系列がある。
共通部分 $U_{i_0 \ldots i_p}$ は分離的なので、前段落の結果は同型
$\Gamma(U_{i_0 \ldots i_p}, \underline{H}^q(\mathcal{F})) \otimes_A B
\to \Gamma(V_{i_0 \ldots i_p}, \underline{H}^q(\mathcal{G}))$
を与える。$- \otimes_A B$ が積と可換し、かつ完全であることを用いると、
$\check{H}^p(\mathcal{U}, \underline{H}^q(\mathcal{F})) \otimes_A B
\to \check{H}^p(\mathcal{V}, \underline{H}^q(\mathcal{G}))$
は同型である。$A \to B$ が平坦であることを用いると、
$H^i(X, \mathcal{F}) \otimes_A B \to H^i(Y, \mathcal{G})$
はすべての $i$ について同型であると結論でき、主張が従う。
\end{proof}

\section{余極限と高次順像}
\label{coherent-section-colimits}

\noindent
この種の一般的な結果は、コホモロジー、節
\ref{cohomology-section-limits}、層、補題
\ref{sheaves-lemma-directed-colimits-sections}、および加群、補題
\ref{modules-lemma-finite-presentation-quasi-compact-colimit} に見いだせる。

\begin{lemma}
\label{coherent-lemma-colimit-cohomology}
$f : X \to S$ をスキームの準コンパクトかつ準分離的な射とする。
$\mathcal{F} = \colim \mathcal{F}_i$ を $X$ 上のアーベル層の
フィルター余極限とする。このとき、任意の $p \geq 0$ に対して
$$
R^pf_*\mathcal{F} = \colim R^pf_*\mathcal{F}_i.
$$
が成り立つ。
\end{lemma}

\begin{proof}
コホモロジー、補題
\ref{cohomology-lemma-higher-direct-image-colimit} を適用する。
アフィン開集合は $S$ の位相の基底をなすので、アフィン開集合
$U \subset S$ に対して
$H^p(f^{-1}U, \mathcal{F}) = \colim H^p(f^{-1}U, \mathcal{F}_i)$
が成り立つことを示せば十分である。
$f^{-1}U$ は準コンパクトかつ準分離的なので、コホモロジー、補題
\ref{cohomology-lemma-quasi-separated-cohomology-colimit} を用いて結論を得る。
（$f^{-1}U$ のアフィン開集合からなる基底が同補題の仮定を満たすからである。）
\end{proof}

\section{コホモロジーと基底変換 II}
\label{coherent-section-cohomology-and-base-change-derived}

\noindent
$f : X \to S$ をスキームの射とし、$\mathcal{F}$ を準連接
$\mathcal{O}_X$-加群とする。$f$ が準コンパクトかつ準分離的ならば、
$Rf_*\mathcal{F}$ を $S$ 上の準連接層の複体で表したい。
これは、層 $R^if_*\mathcal{F}$ が準連接であること
（$S$ が準コンパクトで対角射がアフィンならば成り立つ）、および
$D_\QCoh(S)$ が $D(\QCoh(\mathcal{O}_S))$ と同値であることから従う。
スキームの導来圏、命題
\ref{perfect-proposition-quasi-compact-affine-diagonal} を参照せよ。

\medskip\noindent
本節では、良い基底変換性質をもつ明示的な複体を生み出す別の方法を用いる。
$f$ と $S$ が分離的である場合、より一般には対角射がアフィンである場合、
この構成は特に容易である。これは圧倒的に最もよく用いられる場合なので、
別に扱う。

\begin{lemma}
\label{coherent-lemma-separated-case-relative-cech}
$f : X \to S$ をスキームの射とする。
$\mathcal{F}$ を準連接 $\mathcal{O}_X$-加群とする。
$X$ は準コンパクトであり、$X$ と $S$ の対角射はアフィンであると仮定する
（例えば $X$ と $S$ が分離的ならばよい）。
この場合、$Rf_*\mathcal{F}$ は次のように計算できる。
\begin{enumerate}
\item 有限アフィン開被覆
$\mathcal{U} : X = \bigcup_{i = 1, \ldots, n} U_i$ を選ぶ。
\item $i_0, \ldots, i_p \in \{1, \ldots, n\}$ に対し、
$f_{i_0 \ldots i_p} : U_{i_0 \ldots i_p} \to S$ を、$f$ の共通部分
$U_{i_0 \ldots i_p} = U_{i_0} \cap \ldots \cap U_{i_p}$ への制限とする。
\item $\mathcal{F}_{i_0 \ldots i_p}$ を、$\mathcal{F}$ の
$U_{i_0 \ldots i_p}$ への制限に等しいものとする。
\item 次のようにおく。
$$
\check{\mathcal{C}}^p(\mathcal{U}, f, \mathcal{F}) =
\bigoplus\nolimits_{i_0 \ldots i_p}
f_{i_0 \ldots i_p *} \mathcal{F}_{i_0 \ldots i_p}
$$
また、微分
$d : \check{\mathcal{C}}^p(\mathcal{U}, f, \mathcal{F})
\to \check{\mathcal{C}}^{p + 1}(\mathcal{U}, f, \mathcal{F})$
をコホモロジー、式 (\ref{cohomology-equation-d-cech}) と同様に定義する。
\end{enumerate}
このとき、複体 $\check{\mathcal{C}}^\bullet(\mathcal{U}, f, \mathcal{F})$ は
$S$ 上の準連接層の複体であり、同型
$$
\check{\mathcal{C}}^\bullet(\mathcal{U}, f, \mathcal{F})
\longrightarrow
Rf_*\mathcal{F}
$$
を $D^{+}(S)$ において備える。この同型は準連接層
$\mathcal{F}$ に関して関手的である。
\end{lemma}

\begin{proof}
コホモロジー、補題 \ref{cohomology-lemma-covering-resolution} の分解
$\mathcal{F} \to {\mathfrak C}^\bullet(\mathcal{U}, \mathcal{F})$ を考える。
等式
$\check{\mathcal{C}}^\bullet(\mathcal{U}, f, \mathcal{F}) =
f_*{\mathfrak C}^\bullet(\mathcal{U}, \mathcal{F})$
が準連接 $\mathcal{O}_S$-加群の複体として成り立つ。射
$j_{i_0 \ldots i_p} : U_{i_0 \ldots i_p} \to X$ と射
$f_{i_0 \ldots i_p} : U_{i_0 \ldots i_p} \to S$ は、射、補題
\ref{morphisms-lemma-affine-permanence} および補題
\ref{coherent-lemma-affine-diagonal} によりアフィンである。従って
$R^qj_{i_0 \ldots i_p *}\mathcal{F}_{i_0 \ldots i_p}$ と
$R^qf_{i_0 \ldots i_p *}\mathcal{F}_{i_0 \ldots i_p}$ はともに
$q > 0$ に対して零である
（補題 \ref{coherent-lemma-relative-affine-vanishing}）。
$f \circ j_{i_0 \ldots i_p} = f_{i_0 \ldots i_p}$ と、コホモロジー、補題
\ref{cohomology-lemma-relative-Leray} のスペクトル系列を用いると、
$R^qf_*(j_{i_0 \ldots i_p *}\mathcal{F}_{i_0 \ldots i_p}) = 0$
が $q > 0$ に対して成り立つと結論できる。
複体 ${\mathfrak C}^\bullet(\mathcal{U}, \mathcal{F})$ の各項は層
$j_{i_0 \ldots i_p *}\mathcal{F}_{i_0 \ldots i_p}$ の有限直和なので、
Leray の非輪状性補題
（導来圏、補題 \ref{derived-lemma-leray-acyclicity}）を用いて
$$
Rf_* \mathcal{F} = f_*{\mathfrak C}^\bullet(\mathcal{U}, \mathcal{F}) =
\check{\mathcal{C}}^\bullet(\mathcal{U}, f, \mathcal{F})
$$
を得る。これは所望の等式である。
\end{proof}

\noindent
次に、基底変換を行うと何が起こるかを考える。

\begin{lemma}
\label{coherent-lemma-base-change-complex}
図式 (\ref{coherent-equation-base-change-diagram}) の記法を用いる。
$f : X \to S$ と $\mathcal{F}$ が補題
\ref{coherent-lemma-separated-case-relative-cech} の仮定を満たすとする。
有限アフィン開被覆 $\mathcal{U} : X = \bigcup U_i$ を選ぶ。これは $X$ の被覆である。
標準同型
$$
g^*\check{\mathcal{C}}^\bullet(\mathcal{U}, f, \mathcal{F})
\longrightarrow
Rf'_*\mathcal{F}'
$$
が $D^{+}(S')$ において存在する。さらに、$S' \to S$ がアフィンならば、実際
$$
g^*\check{\mathcal{C}}^\bullet(\mathcal{U}, f, \mathcal{F})
=
\check{\mathcal{C}}^\bullet(\mathcal{U}', f', \mathcal{F}')
$$
である。ここで $\mathcal{U}' : X' = \bigcup U_i'$ であり、
$U_i' = (g')^{-1}(U_i) = U_{i, S'}$ もアフィンである。
\end{lemma}

\begin{proof}
$U_i' = (g')^{-1}(U_i) = U_{i, S'}$ と定義できる。これは
$S'$ が $S$ 上アフィンであるかどうかにかかわらない。
$\mathcal{U}' : X' = \bigcup U_i'$ を $X'$ の誘導された被覆とする。
この場合、
$$
g^*\check{\mathcal{C}}^\bullet(\mathcal{U}, f, \mathcal{F})
=
\check{\mathcal{C}}^\bullet(\mathcal{U}', f', \mathcal{F}')
$$
であると主張する。ここで
$\check{\mathcal{C}}^\bullet(\mathcal{U}', f', \mathcal{F}')$ は補題
\ref{coherent-lemma-separated-case-relative-cech} とまったく同じ仕方で定義される。
これは $S'$ 上局所的に考え、アフィン射の場合
（補題 \ref{coherent-lemma-affine-base-change}）から明らかである。
さらに、補題 \ref{coherent-lemma-separated-case-relative-cech} の証明とまったく同様に、
同型
$$
\check{\mathcal{C}}^\bullet(\mathcal{U}', f', \mathcal{F}')
\longrightarrow
Rf'_*\mathcal{F}'
$$
が $D^{+}(S')$ において存在することが分かる。実際、射
$U_i' \to X'$ と $U_i' \to S'$ は依然としてアフィンである
（アフィン射の基底変換だからである）。詳細は省略する。
\end{proof}

\noindent
上の補題は、複体
$$
\mathcal{K}^\bullet = \check{\mathcal{C}}^\bullet(\mathcal{U}, f, \mathcal{F})
$$
が $S$ 上の準連接層からなる下に有界な複体であり、
$f : X \to S$ の高次順像を {\it 普遍的に} 計算することを述べている。
これはこの特定の複体に関する性質であって、一般には
$\check{\mathcal{C}}^\bullet(\mathcal{U}, f, \mathcal{F})$ を
擬同型な複体で置き換えても保たれない！
言い換えると、これは導来圏で意味をもつ主張ではない。その理由は、引戻し
$g^*\mathcal{K}^\bullet$ は一般に導来引戻し $Lg^*\mathcal{K}^\bullet$
（$\mathcal{K}^\bullet$ のもの）と {\it 等しくない} からである！

\medskip\noindent
この主張を証明できる、より一般的な場合を述べる。
$S$ が分離的であるという条件は、ほとんどの応用で害をなさない。
なぜなら、通常は全導来像のある局所的性質を証明するために用いるからである。
証明はかなり複雑になり、超被覆を用いる。これは、超被覆をどのように使えるかを
示す好例でもある。

\begin{lemma}
\label{coherent-lemma-hypercoverings}
$f : X \to S$ をスキームの射とする。
$\mathcal{F}$ を $X$ 上の準連接層とする。
$f$ は準コンパクトかつ準分離的であり、$S$ は準コンパクトかつ分離的であると
仮定する。下に有界な複体 $\mathcal{K}^\bullet$（準連接
$\mathcal{O}_S$-加群からなるもの）で、次の性質をもつものが存在する。任意の射
$g : S' \to S$ に対し、複体 $g^*\mathcal{K}^\bullet$ は
$Rf'_*\mathcal{F}'$ の代表である。記法は図式
（\ref{coherent-equation-base-change-diagram}）の通りである。
\end{lemma}

\begin{proof}
（$f$ も分離的ならば、補題 \ref{coherent-lemma-base-change-complex} を参照せよ。）
仮定から、特に $X$ はスキームとして準コンパクトかつ準分離的である。
$\mathcal{B}$ を $X$ のアフィン開集合全体の集合とする。超被覆、補題
\ref{hypercovering-lemma-quasi-separated-quasi-compact-hypercovering} により、
超被覆 $K = (I, \{U_i\})$ で、各 $I_n$ が有限であり各 $U_i$ が $X$ の
アフィン開集合であるものを選べる。超被覆、補題
\ref{hypercovering-lemma-cech-spectral-sequence} により、$E_2$-項が
$$
E_2^{p, q} = \check{H}^p(K, \underline{H}^q(\mathcal{F}))
$$
であり、$H^{p + q}(X, \mathcal{F})$ に収束するスペクトル系列がある。
$\check{H}^p(K, \underline{H}^q(\mathcal{F}))$ は、複体
$$
\prod\nolimits_{i \in I_0} H^q(U_i, \mathcal{F})
\to
\prod\nolimits_{i \in I_1} H^q(U_i, \mathcal{F})
\to
\prod\nolimits_{i \in I_2} H^q(U_i, \mathcal{F})
\to \ldots
$$
の第 $p$ コホモロジー群であることに注意する。
各 $U_i$ はアフィンなので、これは $q = 0$ でない限り零であり、その場合には
$$
\prod\nolimits_{i \in I_0} \mathcal{F}(U_i)
\to
\prod\nolimits_{i \in I_1} \mathcal{F}(U_i)
\to
\prod\nolimits_{i \in I_2} \mathcal{F}(U_i)
\to \ldots
$$
を得る。従って、$R\Gamma(X, \mathcal{F})$ はこの複体によって計算される。

\medskip\noindent
任意の $n$ と $i \in I_n$ に対し、$f_i : U_i \to S$ を $f$ の
$U_i$ への制限とする。$S$ は分離的で $U_i$ はアフィンなので、この射は
アフィンである。準連接層の複体
$$
\mathcal{K}^\bullet = (
\prod\nolimits_{i \in I_0} f_{i, *}\mathcal{F}|_{U_i}
\to
\prod\nolimits_{i \in I_1} f_{i, *}\mathcal{F}|_{U_i}
\to
\prod\nolimits_{i \in I_2} f_{i, *}\mathcal{F}|_{U_i}
\to \ldots )
$$
を $S$ 上で考える。超被覆、補題
\ref{hypercovering-lemma-cech-spectral-sequence} と同様に、写像
$\mathcal{K}^\bullet \to Rf_*\mathcal{F}$ を $D(\mathcal{O}_S)$ 内に得る。
そのために $\mathcal{F}$ の入射分解を選ぶ（詳細は省略する）。
任意のアフィン・スキーム $V$ と射
$g : V \to S$ を考える。このとき、基底変換 $X_V$ は、基底変換により得られる
超被覆 $K_V = (I, \{U_{i, V}\})$ をもつ。さらに、
$g^*f_{i, *}\mathcal{F} = f_{i, V, *}(g')^*\mathcal{F}|_{U_{i, V}}$ である。
従って、上の議論は $\Gamma(V, g^*\mathcal{K}^\bullet)$ が
$R\Gamma(X_V, (g')^*\mathcal{F})$ を計算することを示す。
複体の等式は $S'$ 上 Zariski 局所的に証明すれば十分なので、
これで補題の証明が完了する。
\end{proof}

\noindent
次の補題は平坦基底変換の変種である。

\begin{lemma}
\label{coherent-lemma-base-change-tensor-with-flat}
スキームのデカルト図式
$$
\xymatrix{
X' \ar[d]_{f'} \ar[r]_{g'} & X \ar[d]^f \\
S' \ar[r]^g & S
}
$$
を考える。$\mathcal{F}$ を準連接 $\mathcal{O}_X$-加群とする。
$\mathcal{G}$ を準連接 $\mathcal{O}_{S'}$-加群とし、これは $S$ 上平坦であるとする。
$f$ は準コンパクトかつ準分離的であると仮定する。
任意の $i \geq 0$ に対し、同一視
$$
\mathcal{G} \otimes_{\mathcal{O}_{S'}} g^*R^if_*\mathcal{F} =
R^if'_*\left((f')^*\mathcal{G}
\otimes_{\mathcal{O}_{X'}} (g')^*\mathcal{F}\right)
$$
がある。
\end{lemma}

\begin{proof}
左辺から右辺への写像を構成しよう。まず、基底変換写像
$Lg^*Rf_*\mathcal{F} \to Rf'_*L(g')^*\mathcal{F}$ がある。また、随伴写像
$\mathcal{G} \to Rf'_*L(f')^*\mathcal{G}$ もある。相対カップ積を用いると、
次を得る。
\begin{align*}
\mathcal{G}
\otimes_{\mathcal{O}_{S'}}^\mathbf{L}
Lg^*Rf_*\mathcal{F}
& \to
Rf'_*L(f')^*\mathcal{G}
\otimes_{\mathcal{O}_{S'}}^\mathbf{L}
Rf'_*L(g')^*\mathcal{F} \\
& \to
Rf'_*\left(L(f')^*\mathcal{G}
\otimes_{\mathcal{O}_{X'}}^\mathbf{L}
L(g')^*\mathcal{F}\right) \\
& \to
Rf'_*\left((f')^*\mathcal{G}
\otimes_{\mathcal{O}_{X'}} (g')^*\mathcal{F}\right)
\end{align*}
中央の矢印では相対カップ積を用いた。コホモロジー、注意
\ref{cohomology-remark-cup-product} を参照せよ。
合成の始域は、コホモロジー、補題
\ref{cohomology-lemma-variant-derived-pullback} により
$$
\mathcal{G} \otimes_{\mathcal{O}_{S'}}^\mathbf{L} Lg^*Rf_*\mathcal{F} =
\mathcal{G} \otimes_{g^{-1}\mathcal{O}_S}^\mathbf{L} g^{-1}Rf_*\mathcal{F}
$$
である。$\mathcal{G}$ は $g^{-1}\mathcal{O}_S$-加群の層として平坦であり、
$g^{-1}$ は完全関手なので、これは第 $i$ コホモロジー層が
$\mathcal{G} \otimes_{g^{-1}\mathcal{O}_S} g^{-1}R^if_*\mathcal{F} =
\mathcal{G} \otimes_{\mathcal{O}_{S'}} g^*R^if_*\mathcal{F}$
である複体である。このようにして、補題の等式における左辺から右辺への
大域的な写像を得る。この写像が同型であることを示すには、$S'$ 上局所的に
考えてよい。従って $S$ と $S'$ はアフィン・スキームであると仮定する。

\medskip\noindent
$S$ と $S'$ がアフィンの場合の証明。
$S = \Spec(A)$, $S' = \Spec(B)$ とし、$\mathcal{G}$ は
$B$-加群 $N$ に対応し、これは $A$-平坦であると仮定する。
$S$ はアフィンなので、$X$ は準コンパクトかつ準分離的である。
超被覆の議論を用いて証明を終える。$X$ が分離的であるか対角射が
アフィンであるなら、{\v C}ech 被覆を用いることもできる。
$\mathcal{B}$ を $X$ のアフィン開集合全体の集合とする。超被覆、補題
\ref{hypercovering-lemma-quasi-separated-quasi-compact-hypercovering} により、
超被覆 $K = (I, \{U_i\})$ を $X$ に対して、各 $I_n$ が有限であり、各 $U_i$ が
$X$ のアフィン開集合となるように選べる。超被覆、補題
\ref{hypercovering-lemma-cech-spectral-sequence} により、$E_2$-項が
$$
E_2^{p, q} = \check{H}^p(K, \underline{H}^q(\mathcal{F}))
$$
であり、$H^{p + q}(X, \mathcal{F})$ に収束するスペクトル系列がある。
各 $U_i$ はアフィンで $\mathcal{F}$ は準連接なので、
$\underline{H}^q(\mathcal{F})$ の $U_i$ 上の値は $q > 0$ に対して零である。
従ってスペクトル系列は退化し、コホモロジー加群
$H^q(X, \mathcal{F})$ は
$$
\prod\nolimits_{i \in I_0} \mathcal{F}(U_i)
\to
\prod\nolimits_{i \in I_1} \mathcal{F}(U_i)
\to
\prod\nolimits_{i \in I_2} \mathcal{F}(U_i)
\to \ldots
$$
によって計算されると結論する。
次に、超被覆の $S'$ への基底変換は
$X' = S' \times_S X$ の超被覆であることに注意する。
スキーム $S' \times_S U_i$ もアフィンであり、
$$
\left((f')^*\mathcal{G} \otimes_{\mathcal{O}_{S'}} (g')^*\mathcal{F}\right)
(S' \times_S U_i) =
N \otimes_A \mathcal{F}(U_i)
$$
である。このようにして、コホモロジー加群
$H^q(X', (f')^*\mathcal{G} \otimes_{\mathcal{O}_{S'}} (g')^*\mathcal{F})$
は
$$
N \otimes_A \left(
\prod\nolimits_{i \in I_0} \mathcal{F}(U_i)
\to
\prod\nolimits_{i \in I_1} \mathcal{F}(U_i)
\to
\prod\nolimits_{i \in I_2} \mathcal{F}(U_i)
\to \ldots
\right)
$$
によって計算されると結論する。$N$ は $A$ 上平坦なので、
$$
H^q(X', (f')^*\mathcal{G} \otimes_{\mathcal{O}_{S'}} (g')^*\mathcal{F})
=
N \otimes_A H^q(X, \mathcal{F})
$$
と結論する。これは示すべき主張を代数に移したものなので、証明は完了する。
\end{proof}

\section{射影空間のコホモロジー}
\label{coherent-section-cohomology-projective-space}

\noindent
本節では、$\mathbf{P}^n_S$ における構造層の捻りのコホモロジーを、
スキーム $S$ 上で計算する。$\mathbf{P}^n_S$ は、構成、定義
\ref{constructions-definition-projective-space} においてファイバー積
$
\mathbf{P}^n_S = S \times_{\Spec(\mathbf{Z})} \mathbf{P}^n_{\mathbf{Z}}
$
として定義されたことを思い出そう。構成、補題
\ref{constructions-lemma-projective-space-bundle} において、これは
$$
\mathbf{P}^n_S = \underline{\text{Proj}}_S(\mathcal{O}_S[T_0, \ldots, T_n])
$$
に等しいことが示された。特に、射影空間は射影束の特別な場合である。
$S = \Spec(R)$ がアフィンならば
$$
\mathbf{P}^n_S = \mathbf{P}^n_R = \text{Proj}(R[T_0, \ldots, T_n]).
$$
である。これらの同一視はすべて互いに整合し、捻った構造層
$\mathcal{O}_{\mathbf{P}^n_S}(d)$ の構成とも整合する。

\medskip\noindent
結果を述べる前に、いくつか記法を導入する。
$R$ を環とする。$R[T_0, \ldots, T_n]$ は次数付き
$R$-代数であり、各 $T_i$ は次数 $1$ の斉次元であることを思い出そう。
$(R[T_0, \ldots, T_n])_d$ で次数 $d$ の斉次成分を表す。
これは有限自由 $R$-加群であり、階数は $\binom{n + d}{d}$（$d \geq 0$
の場合）で、それ以外は零である。その基底は単項式
$T_0^{e_0} \ldots T_n^{e_n}$ で、$\sum e_i = d$ を満たすものからなる。
次の記法も用いる。$R[\frac{1}{T_0}, \ldots, \frac{1}{T_n}]$ は
$\mathbf{Z}$-次数付き環を表し、$\frac{1}{T_i}$ の次数は $-1$ である。
特に、以下の主張に現れる $\mathbf{Z}$-次数付き
$R[\frac{1}{T_0}, \ldots, \frac{1}{T_n}]$-加群
$$
\frac{1}{T_0 \ldots T_n} R[\frac{1}{T_0}, \ldots, \frac{1}{T_n}]
$$
は次数 $\geq -n$ で零であり、生成元 $\frac{1}{T_0 \ldots T_n}$ 上、
次数 $-n - 1$ では自由であり、階数 $(-1)^n\binom{n + d}{d}$ の
自由加群となるのは $d \leq -n - 1$ の場合である。

\begin{lemma}
\label{coherent-lemma-cohomology-projective-space-over-ring}
\begin{reference}
\cite[III Proposition 2.1.12]{EGA}
\end{reference}
$R$ を環とし、$n \geq 0$ を整数とする。このとき
$$
H^q(\mathbf{P}^n, \mathcal{O}_{\mathbf{P}^n_R}(d)) =
\left\{
\begin{matrix}
(R[T_0, \ldots, T_n])_d & \text{if} & q = 0,\ d \geq 0 \\
\left(\frac{1}{T_0 \ldots T_n} R[\frac{1}{T_0}, \ldots, \frac{1}{T_n}]\right)_d
& \text{if} & q = n,\ d < 0 \\
0 & \text{if} & q, n, d\text{ not as above}
\end{matrix}
\right.
$$
が $R$-加群として成り立つ。
\end{lemma}

\begin{proof}
標準アフィン開被覆
$$
\mathcal{U} : \mathbf{P}^n_R = \bigcup\nolimits_{i = 0}^n D_{+}(T_i)
$$
を用い、{\v C}ech 複体によってコホモロジーを計算する。
補題 \ref{coherent-lemma-cech-cohomology-quasi-coherent} によりこれは許される。
実際、有限個のアフィン $D_{+}(T_i)$ の任意の共通部分も標準アフィン
開集合である（構成、節 \ref{constructions-section-proj} を参照）。
実際、コホモロジー、補題 \ref{cohomology-lemma-ordered-alternating} および
\ref{cohomology-lemma-alternating-usual} に従い、交代 {\v C}ech 複体または
順序付き {\v C}ech 複体を用いることができる。

\medskip\noindent
$\{0, \ldots, n\}$ には通常の順序を用いる。従って対象となる複体の項は
$$
\check{\mathcal{C}}_{ord}^p(\mathcal{U}, \mathcal{O}_{\mathbf{P}_R}(d))
=
\bigoplus\nolimits_{i_0 < \ldots < i_p}
(R[T_0, \ldots, T_n, \frac{1}{T_{i_0} \ldots T_{i_p}}])_d
$$
である。また、写像は通常の公式
$$
d(s)_{i_0 \ldots i_{p + 1}} =
\sum\nolimits_{j = 0}^{p + 1} (-1)^j s_{i_0 \ldots \hat i_j \ldots i_{p + 1}}
$$
で与えられる。コホモロジー、節
\ref{cohomology-section-alternating-cech} を参照せよ。
この複体の各項は自然な $\mathbf{Z}^{n + 1}$-次数付けをもつ。
すなわち、単項式 $T_0^{e_0} \ldots T_n^{e_n}$ の重みを
$(e_0, \ldots, e_n) \in \mathbf{Z}^{n + 1}$ と宣言して得られる。
上の微分がこの次数付けを保つことは明らかである。公式では
$$
\check{\mathcal{C}}_{ord}^\bullet(\mathcal{U}, \mathcal{O}_{\mathbf{P}_R}(d))
=
\bigoplus\nolimits_{\vec{e} \in \mathbf{Z}^{n + 1}}
\check{\mathcal{C}}^\bullet(\vec{e})
$$
である。ただし右辺のすべての直和因子が現れるわけではない（後述）。
従って複体のコホモロジー加群を計算するには、次数付き成分の
コホモロジーを計算し、最後に直和を取れば十分である。

\medskip\noindent
$\vec{e} = (e_0, \ldots, e_n) \in \mathbf{Z}^{n + 1}$ を固定する。
この重みが上の複体に現れるためには $e_0 + \ldots + e_n = d$ と
仮定する必要がある（そうでなければ、もちろん構造層の別の捻りに現れる）。
これを仮定して
$$
NEG(\vec{e}) = \{i \in \{0, \ldots, n\} \mid e_i < 0\}.
$$
とおく。この記法のもとで、上の {\v C}ech 複体の重み $\vec{e}$ の直和因子
$\check{\mathcal{C}}^\bullet(\vec{e})$ は次の項をもつ。
$$
\check{\mathcal{C}}^p(\vec{e})
=
\bigoplus\nolimits_{i_0 < \ldots < i_p,
\ NEG(\vec{e}) \subset \{i_0, \ldots, i_p\}}
R \cdot T_0^{e_0} \ldots T_n^{e_n}
$$
言い換えると、$i_0 < \ldots < i_p$ であって $NEG(\vec{e})$ が
$\{i_0 \ldots i_p\}$ に含まれないものに対応する項は零である。
複体 $\check{\mathcal{C}}^\bullet(\vec{e})$ の微分は、依然として
上とまったく同じ公式で与えられる。

\medskip\noindent
$NEG(\vec{e}) = \{0, \ldots, n\}$、すなわちすべての指数 $e_i$ が
負であると仮定する。この場合、複体
$\check{\mathcal{C}}^\bullet(\vec{e})$ はただ一つの項、すなわち
$\check{\mathcal{C}}^n(\vec{e}) =
R \cdot \frac{1}{T_0^{-e_0} \ldots T_n^{-e_n}}$ のみをもつ。従って
$$
H^q(\check{\mathcal{C}}^\bullet(\vec{e})) =
\left\{
\begin{matrix}
R \cdot \frac{1}{T_0^{-e_0} \ldots T_n^{-e_n}} & \text{if} & q = n \\
0 & \text{if} & \text{else}
\end{matrix}
\right.
$$
である。これらすべての項の直和は、補題で述べたとおり次数 $n$ における値
$$
\left(\frac{1}{T_0 \ldots T_n} R[\frac{1}{T_0}, \ldots, \frac{1}{T_n}]\right)_d
$$
を明らかに与える。また、これらの項は他の次数のコホモロジーには寄与しない
（これも補題の主張と一致する）。

\medskip\noindent
$NEG(\vec{e}) = \emptyset$ と仮定する。この場合、複体
$\check{\mathcal{C}}^\bullet(\vec{e})$ は直和因子 $R$ を、すべての
$i_0 < \ldots < i_p$ に対してもつ。
$\check{\mathcal{C}}^\bullet(\vec{e})$ を別の複体と比較しよう。
すなわち、アフィン開被覆
$$
\mathcal{V} : \Spec(R) = \bigcup\nolimits_{i \in \{0, \ldots, n\}} V_i
$$
を考える。ここで、$V_i = \Spec(R)$ とすべての $i$ に対しておく。
交代 {\v C}ech 複体
$$
\check{\mathcal{C}}_{ord}^\bullet(\mathcal{V}, \mathcal{O}_{\Spec(R)})
$$
を考える。上と同じ理由により、これは $\Spec(R)$ 上の構造層の
コホモロジーを計算する。従って
$H^p(
\check{\mathcal{C}}_{ord}^\bullet(\mathcal{V}, \mathcal{O}_{\Spec(R)})
) = R$ は $p = 0$ のとき成り立ち、$0$ は $p > 0$ のとき成り立つ。
これらの事実については、補題
\ref{coherent-lemma-cech-cohomology-quasi-coherent-trivial} とその証明を参照せよ。
また、$\check{\mathcal{C}}_{ord}^\bullet(\mathcal{V}, \mathcal{O}_{\Spec(R)})$
も直和因子 $R$ を各 $i_0 < \ldots < i_p$ に対してもち、
$\check{\mathcal{C}}^\bullet(\vec{e})$ とまったく同じ微分をもつ。
言い換えると、これらは同型な複体であり、従って同じコホモロジーをもつ。
ゆえに
$$
H^q(\check{\mathcal{C}}^\bullet(\vec{e})) =
\left\{
\begin{matrix}
R \cdot T_0^{e_0} \ldots T_n^{e_n} & \text{if} & q = 0 \\
0 & \text{if} & \text{else}
\end{matrix}
\right.
$$
と結論する。

$NEG(\vec{e}) = \emptyset$ の場合である。これらすべての項の直和は、
補題で述べたとおり次数 $0$ における値
$$
(R[T_0, \ldots, T_n])_d
$$
を明らかに与える。また、これらの項は他の次数のコホモロジーには寄与しない
（これも補題の主張と一致する）。

\medskip\noindent
補題の証明を終えるには、複体 $\check{\mathcal{C}}^\bullet(\vec{e})$ が、
$NEG(\vec{e})$ が空でも $\{0, \ldots, n\}$ と等しくもないとき
非輪状であることを示さなければならない。
添字 $i_{\text{fix}} \not \in NEG(\vec{e})$ を一つ選ぶ
（そのような添字は存在する）。写像
$$
h :
\check{\mathcal{C}}^{p + 1}(\vec{e})
\to
\check{\mathcal{C}}^p(\vec{e})
$$
を、$i_0 < \ldots < i_p$ に対して次の規則で定める。
$$
h(s)_{i_0 \ldots i_p} =
\left\{
\begin{matrix}
0 & \text{if} & p \not \in \{0, \ldots, n - 1\} \\
0 & \text{if} & i_{\text{fix}} \in \{i_0, \ldots, i_p\} \\
s_{i_{\text{fix}} i_0 \ldots i_p} & \text{if} & i_{\text{fix}} < i_0 \\
(-1)^a s_{i_0 \ldots i_{a - 1} i_{\text{fix}} i_a \ldots i_p} &
\text{if} & i_{a - 1} < i_{\text{fix}} < i_a \\
(-1)^{p + 1} s_{i_0 \ldots i_p i_{\text{fix}}} &
\text{if} & i_p < i_{\text{fix}}
\end{matrix}
\right.
$$
補題 \ref{coherent-lemma-cech-cohomology-quasi-coherent-trivial} の証明と比較されたい。
これは
$$
NEG(\vec{e}) \subset \{i_0, \ldots, i_p\}
\Leftrightarrow
NEG(\vec{e}) \subset \{i_{\text{fix}}, i_0, \ldots, i_p\}
$$
が成り立つため、意味をもつ。補題
\ref{coherent-lemma-cech-cohomology-quasi-coherent-trivial} の証明とまったく同じ
組合せ論的計算\footnote{
例えば、$i_0 < \ldots < i_p$ が
$i_{\text{fix}} \not \in \{i_0, \ldots, i_p\}$ を満たし、
$i_{a - 1} < i_{\text{fix}} < i_a$ が、ある $1 \leq a \leq p$ に対して
成り立つと仮定する。このとき
\begin{align*}
&
(dh + hd)(s)_{i_0 \ldots i_p} \\
& =
\sum\nolimits_{j = 0}^p
(-1)^j
h(s)_{i_0 \ldots \hat i_j \ldots i_p}
+
(-1)^a d(s)_{i_0 \ldots i_{a - 1} i_{\text{fix}} i_a \ldots i_p}\\
& =
\sum\nolimits_{j = 0}^{a - 1}
(-1)^{j + a - 1}
s_{i_0 \ldots \hat i_j \ldots i_{a - 1} i_{\text{fix}} i_a \ldots i_p} +
\sum\nolimits_{j = a}^p
(-1)^{j + a}
s_{i_0 \ldots i_{a - 1} i_{\text{fix}} i_a \ldots \hat i_j \ldots i_p}
+ \\
&
\sum\nolimits_{j = 0}^{a - 1}
(-1)^{a + j}
s_{i_0 \ldots \hat i_j \ldots i_{a - 1} i_{\text{fix}} i_a \ldots i_p} +
(-1)^{2a} s_{i_0 \ldots i_p} +
\sum\nolimits_{j = a}^p
(-1)^{a + j + 1}
s_{i_0 \ldots i_{a - 1} i_{\text{fix}} i_a \ldots \hat i_j \ldots i_p}
\\
& =
s_{i_0 \ldots i_p}
\end{align*}
となり、所望の結果を得る。他の場合も同様である。}
により
$$
(hd + dh)(s)_{i_0 \ldots i_p}
=
s_{i_0 \ldots i_p}
$$
を得る。従って複体 $\check{\mathcal{C}}^\bullet(\vec{e})$ の恒等写像は
零写像と鎖ホモトピックであり、この複体は非輪状である。
\end{proof}

\noindent
次の補題では、自由 $R$-加群の対
$$
R[T_0, \ldots, T_n]
\times
\frac{1}{T_0 \ldots T_n} R[\frac{1}{T_0}, \ldots, \frac{1}{T_n}]
\longrightarrow
R
$$
を用いる。これは規則
$$
(f, g)
\longmapsto
\text{coefficient of }
\frac{1}{T_0 \ldots T_n}
\text{ in }fg.
$$
で定義される。言い換えると、基底元 $T_0^{e_0} \ldots T_n^{e_n}$ と
基底元 $T_0^{d_0} \ldots T_n^{d_n}$ の対が $1$ となるのは、
$e_i + d_i = -1$ がすべての $i$ に対して成り立つとき、かつそのときに限り、
それ以外の場合は零となる。この対を用いると同一視
$$
\left(\frac{1}{T_0 \ldots T_n} R[\frac{1}{T_0}, \ldots, \frac{1}{T_n}]\right)_d
=
\Hom_R((R[T_0, \ldots, T_n])_{-n - 1 - d}, R)
$$
を得る。従って補題 \ref{coherent-lemma-cohomology-projective-space-over-ring} の結果は
次のように言い換えられる。
\begin{equation}
\label{coherent-equation-identify}
H^q(\mathbf{P}^n, \mathcal{O}_{\mathbf{P}^n_R}(d)) =
\left\{
\begin{matrix}
(R[T_0, \ldots, T_n])_d & \text{if} & q = 0 \\
0 & \text{if} & q \not = 0, n \\
\Hom_R((R[T_0, \ldots, T_n])_{-n - 1 - d}, R)
& \text{if} & q = n
\end{matrix}
\right.
\end{equation}

\begin{lemma}
\label{coherent-lemma-identify-functorially}
式 (\ref{coherent-equation-identify}) の同一視は、環準同型 $R \to R'$ に関する
基底変換と整合する。さらに、任意の斉次元
$f \in R[T_0, \ldots, T_n]$ を次数 $m$ とすると、$f$ 倍写像
$$
\mathcal{O}_{\mathbf{P}^n_R}(d)
\longrightarrow
\mathcal{O}_{\mathbf{P}^n_R}(d + m)
$$
がコホモロジー群上に誘導する写像は、式 (\ref{coherent-equation-identify}) の
同一視を通じて、$f$ 倍写像を $H^0$ 上に誘導し、$f$ 倍写像
$$
(R[T_0, \ldots, T_n])_{-n - 1 - (d + m)}
\longrightarrow
(R[T_0, \ldots, T_n])_{-n - 1 - d}
$$
の双対写像を $H^n$ 上に誘導する。
\end{lemma}

\begin{proof}
$R \to R'$ を環準同型とする。$\mathcal{U}$ を
$\mathbf{P}^n_R$ の標準アフィン開被覆、$\mathcal{U}'$ を
$\mathbf{P}^n_{R'}$ の標準アフィン開被覆とする。$\mathcal{U}'$ は被覆
$\mathcal{U}$ の、標準射 $\mathbf{P}^n_{R'} \to \mathbf{P}^n_R$ による
引戻しである。従って {\v C}ech 複体の写像
$$
\gamma :
\check{\mathcal{C}}_{ord}^\bullet(\mathcal{U},
\mathcal{O}_{\mathbf{P}_R}(d))
\longrightarrow
\check{\mathcal{C}}_{ord}^\bullet(\mathcal{U}',
\mathcal{O}_{\mathbf{P}_{R'}}(d))
$$
があり、コホモロジー、補題
\ref{cohomology-lemma-functoriality-cech} により、これはコホモロジー上の
写像と整合する。補題
\ref{coherent-lemma-cohomology-projective-space-over-ring} の証明での計算から、
この {\v C}ech 複体の写像が問題のコホモロジー群の同一視と整合することは
明らかである。（実際、$R$ 上の {\v C}ech 複体の基底元は、単に
$R'$ 上の {\v C}ech 複体の対応する基底元へ写る。）これで補題の第一の主張を得る。

\medskip\noindent
次に環 $R$ を固定し、二つの斉次多項式
$f, g \in R[T_0, \ldots, T_n]$ を考え、両者の次数を $m$ とする。
コホモロジーは加法的関手なので、
$f + g$ 倍写像が誘導する写像は、$f$ 倍写像が誘導する写像と
$g$ 倍写像が誘導する写像との和である。またコホモロジーは関手なので、
積 $fg$ による乗法についても、$f, g$ がいずれも斉次であるとき
（同じ次数である必要はない）、同様の結果が成り立つ。
従って補題の第二の主張を確かめるには、$f = x \in R$ の場合と
$f = T_i$ の場合を証明すれば十分である。元 $x \in R$ による
乗法の場合、ここに現れる各コホモロジー群または複体は
$R$-加群または $R$-加群の複体の構造をもつので、結果は従う。
最後に、$T_i$ 倍写像を $\mathcal{O}_{\mathbf{P}^n_R}$-線形写像
$$
\mathcal{O}_{\mathbf{P}^n_R}(d)
\longrightarrow
\mathcal{O}_{\mathbf{P}^n_R}(d + 1)
$$
として考える。

$H^0$ に関する主張は明らかである。$H^n$ に関する主張については、
$T_i$ 倍写像を {\v C}ech 複体の写像
$$
\check{\mathcal{C}}_{ord}^\bullet(\mathcal{U},
\mathcal{O}_{\mathbf{P}_R}(d))
\longrightarrow
\check{\mathcal{C}}_{ord}^\bullet(\mathcal{U},
\mathcal{O}_{\mathbf{P}_R}(d + 1))
$$
として考える。補題 \ref{coherent-lemma-cohomology-projective-space-over-ring} の証明で
導入した記法を用いる。分解
$$
\check{\mathcal{C}}_{ord}^\bullet(\mathcal{U}, \mathcal{O}_{\mathbf{P}_R}(d))
=
\bigoplus\nolimits_{\vec{e} \in \mathbf{Z}^{n + 1}, \ \sum e_i = d}
\check{\mathcal{C}}^\bullet(\vec{e})
$$
および
$$
\check{\mathcal{C}}_{ord}^\bullet(\mathcal{U},
\mathcal{O}_{\mathbf{P}_R}(d + 1))
=
\bigoplus\nolimits_{\vec{e} \in \mathbf{Z}^{n + 1}, \ \sum e_i = d + 1}
\check{\mathcal{C}}^\bullet(\vec{e})
$$
に関して、$T_i$ 倍写像の作用を考える。これは部分複体
$\check{\mathcal{C}}^\bullet(\vec{e})$ を部分複体
$\check{\mathcal{C}}^\bullet(\vec{e} + \vec{b}_i)$ へ写すことが明らかである。
ここで $\vec{b}_i = (0, \ldots, 0, 1, 0, \ldots, 0))$ は第 $i$ 基底ベクトルである。
言い換えると、$H^n$ の $\vec{e}$ に対応する直和因子で、
$e_i < 0$ かつ $\sum e_i = d$ を満たすものを、$H^n$ の
$\vec{e} + \vec{b}_i$ に対応する直和因子へ写す
（後者は $e_i + b_i \geq 0$ ならば零である）。これが、$T_i$ 倍写像
$$
(R[T_0, \ldots, T_n])_{-n - 1 - (d + 1)}
\longrightarrow
(R[T_0, \ldots, T_n])_{-n - 1 - d}
$$
の双対写像の作用に正確に対応することは容易に分かる。これで補題が証明された。
\end{proof}

\noindent
相対版を述べる前に、いくつか記法が必要である。すなわち、
$\mathcal{O}_S[T_0, \ldots, T_n]$ は次数付き $\mathcal{O}_S$-加群であり、
各 $T_i$ は次数 $1$ の斉次元であることを思い出そう。
$(\mathcal{O}_S[T_0, \ldots, T_n])_d$ で次数 $d$ の斉次成分を表す。
これは階数 $\binom{n + d}{d}$ の有限局所自由層（$S$ 上）である。

\begin{lemma}
\label{coherent-lemma-cohomology-projective-space-over-base}
$S$ をスキームとし、$n \geq 0$ を整数とする。構造射
$$
f : \mathbf{P}^n_S \longrightarrow S.
$$
を考える。このとき
$$
R^qf_*(\mathcal{O}_{\mathbf{P}^n_S}(d)) =
\left\{
\begin{matrix}
(\mathcal{O}_S[T_0, \ldots, T_n])_d & \text{if} & q = 0 \\
0 & \text{if} & q \not = 0, n \\
\SheafHom_{\mathcal{O}_S}(
(\mathcal{O}_S[T_0, \ldots, T_n])_{- n - 1 - d}, \mathcal{O}_S)
& \text{if} & q = n
\end{matrix}
\right.
$$
が成り立つ。
\end{lemma}

\begin{proof}
省略する。ヒント：これは、(\ref{coherent-equation-identify}) の同一視が、補題
\ref{coherent-lemma-identify-functorially} によりアフィン基底変換と整合することから従う。
\end{proof}

\noindent
次に、有限局所自由層に付随する射影束についての版を述べる。
$S$ をスキームとし、$\mathcal{E}$ を有限局所自由
$\mathcal{O}_S$-加群で定数階数 $n + 1$ のものとする。加群、節
\ref{modules-section-locally-free} を参照せよ。この場合、
$\text{Sym}(\mathcal{E})$ を次数付き $\mathcal{O}_S$-加群とみなし、
$\mathcal{E}$ が次数 $1$ の次数付き部分をなすものとする。また
$\text{Sym}^d(\mathcal{E})$ は次数 $d$ の斉次成分である。
これは階数 $\binom{n + d}{d}$ の有限局所自由層（$S$ 上）である。
われわれの正規化では
$$
\pi :
\mathbf{P}(\mathcal{E})
=
\underline{\text{Proj}}_S(\text{Sym}(\mathcal{E}))
\longrightarrow
S
$$
であり、自然な写像
$\text{Sym}^d(\mathcal{E}) \to \pi_*\mathcal{O}_{\mathbf{P}(\mathcal{E})}(d)$
があることを思い出そう。

\begin{lemma}
\label{coherent-lemma-cohomology-projective-bundle}
$S$ をスキームとし、$n \geq 1$ とする。$\mathcal{E}$ を有限局所自由
$\mathcal{O}_S$-加群で定数階数 $n + 1$ のものとする。構造射
$$
\pi : \mathbf{P}(\mathcal{E}) \longrightarrow S.
$$
を考える。このとき
$$
R^q\pi_*(\mathcal{O}_{\mathbf{P}(\mathcal{E})}(d)) =
\left\{
\begin{matrix}
\text{Sym}^d(\mathcal{E}) & \text{if} & q = 0 \\
0 & \text{if} & q \not = 0, n \\
\SheafHom_{\mathcal{O}_S}(
\text{Sym}^{- n - 1 - d}(\mathcal{E})
\otimes_{\mathcal{O}_S}
\wedge^{n + 1}\mathcal{E},
\mathcal{O}_S)
& \text{if} & q = n
\end{matrix}
\right.
$$
が成り立つ。これらの同一視は、基底変換および局所自由層の同型と整合する。
\end{lemma}

\begin{proof}
標準写像
$$
\pi^*\mathcal{E} \longrightarrow \mathcal{O}_{\mathbf{P}(\mathcal{E})}(1)
$$
を考え、$1$ だけ下に捻って
$$
\pi^*(\mathcal{E})(-1) \longrightarrow \mathcal{O}_{\mathbf{P}(\mathcal{E})}
$$
を得る。これは階数 $n + 1$ の局所自由層から構造層への全射である。
従って対応する Koszul 複体は完全である（代数続論、補題
\ref{more-algebra-lemma-homotopy-koszul-abstract}）。言い換えると、
完全複体
$$
0 \to
\pi^*(\wedge^{n + 1}\mathcal{E})(-n - 1) \to
\ldots \to
\pi^*(\wedge^i\mathcal{E})(-i) \to
\ldots \to
\pi^*\mathcal{E}(-1) \to
\mathcal{O}_{\mathbf{P}(\mathcal{E})} \to 0
$$
がある。項 $\pi^*(\wedge^i\mathcal{E})(-i)$ は次数 $-i$ にあるとみなす。
導来圏、補題 \ref{derived-lemma-two-ss-complex-functor} の第一スペクトル系列を
用いて、この非輪状複体の高次順像を計算する。すなわち、項
$$
E_1^{p, q} = R^q\pi_*\left(\pi^*(\wedge^{-p}\mathcal{E})(p)\right)
$$
をもち零へ収束するスペクトル系列がある。射影公式
（コホモロジー、補題 \ref{cohomology-lemma-projection-formula}）により
$$
E_1^{p, q} = \wedge^{-p} \mathcal{E} \otimes_{\mathcal{O}_S}
R^q\pi_*\left(\mathcal{O}_{\mathbf{P}(\mathcal{E})}(p)\right).
$$
である。$S$ 上局所的に層 $\mathcal{E}$ は自明、すなわち
$\mathcal{O}_S^{\oplus n + 1}$ と同型なので、$S$ 上局所的に射
$\mathbf{P}(\mathcal{E}) \to S$ を $\mathbf{P}^n_S \to S$ と同一視できる。
従って $S$ 上局所的に、補題
\ref{coherent-lemma-cohomology-projective-space-over-ring},
\ref{coherent-lemma-identify-functorially}, または
\ref{coherent-lemma-cohomology-projective-space-over-base} の結果を用いることができる。
このとき $E_1^{p, q} = 0$ であるのは、$(p, q)$ が $(0, 0)$ または
$(-n - 1, n)$ でない場合である。非零の項は
\begin{align*}
E_1^{0, 0} & = \pi_*\mathcal{O}_{\mathbf{P}(\mathcal{E})} = \mathcal{O}_S \\
E_1^{-n - 1, n} & =
R^n\pi_*\left(\pi^*(\wedge^{n + 1}\mathcal{E})(-n - 1)\right) =
\wedge^{n + 1}\mathcal{E} \otimes_{\mathcal{O}_S}
R^n\pi_*\left(\mathcal{O}_{\mathbf{P}(\mathcal{E})}(-n - 1)\right)
\end{align*}
である。従ってスペクトル系列には、写像
$d_{n + 1}^{-n - 1, n} : E_{n + 1}^{-n - 1, n} \to E_{n + 1}^{0, 0}$
という非零となりうる微分がただ一つあり、これは同型でなければならない
（スペクトル系列が $0$ 層へ収束するからである）。従って
$E_1^{p, q} = E_{n + 1}^{p, q}$ であり、標準同型
$$
\wedge^{n + 1}\mathcal{E} \otimes_{\mathcal{O}_S}
R^n\pi_*\left(\mathcal{O}_{\mathbf{P}(\mathcal{E})}(-n - 1)\right) =
R^n\pi_*\left(\pi^*(\wedge^{n + 1}\mathcal{E})(-n - 1)\right)
\xrightarrow{d_{n + 1}^{-n - 1, n}}
\mathcal{O}_S
$$
を得る。$\wedge^{n + 1}\mathcal{E}$ は可逆層なので、これは
$R^n\pi_*\mathcal{O}_{\mathbf{P}(\mathcal{E})}(-n - 1)$ も可逆であり、
$\wedge^{n + 1}\mathcal{E}$ の逆と標準的に同型であることを含意する。
言い換えると、補題の $d = - n - 1$ の場合を証明した。

\medskip\noindent
$S$ 上局所的に考えれば、補題
\ref{coherent-lemma-cohomology-projective-space-over-ring},
\ref{coherent-lemma-identify-functorially}, または
\ref{coherent-lemma-cohomology-projective-space-over-base} のコホモロジー計算から、
補題の消滅に関する主張が直ちに分かる。また、$R^0\pi_*$ に関する結果も
明らかである。実際、標準写像
$\text{Sym}^d(\mathcal{E}) \to \pi_* \mathcal{O}_{\mathbf{P}(\mathcal{E})}(d)$
がすべての $d$ に対して存在する。残るのは、
$R^n\pi_*\mathcal{O}_{\mathbf{P}(\mathcal{E})}(d)$ の記述が
$d < -n - 1$ に対して正しいことを示す
ことである。このため、写像
$$
\pi^*(\text{Sym}^{-d - n - 1}(\mathcal{E}))
\otimes_{\mathcal{O}_{\mathbf{P}(\mathcal{E})}}
\mathcal{O}_{\mathbf{P}(\mathcal{E})}(d)
\longrightarrow
\mathcal{O}_{\mathbf{P}(\mathcal{E})}(-n - 1)
$$
を考える。$R^n\pi_*$ を適用し、射影公式（上記参照）を用いると、写像
$$
\text{Sym}^{-d - n - 1}(\mathcal{E})
\otimes_{\mathcal{O}_S}
R^n\pi_*(\mathcal{O}_{\mathbf{P}(\mathcal{E})}(d))
\longrightarrow
R^n\pi_*\mathcal{O}_{\mathbf{P}(\mathcal{E})}(-n - 1) =
(\wedge^{n + 1}\mathcal{E})^{\otimes -1}
$$
を得る（最後の等式は上で示した）。再び補題
\ref{coherent-lemma-cohomology-projective-space-over-ring},
\ref{coherent-lemma-identify-functorially}, または
\ref{coherent-lemma-cohomology-projective-space-over-base} の局所計算により、
この写像は $R^n\pi_*(\mathcal{O}_{\mathbf{P}(\mathcal{E})}(d))$ と
$\text{Sym}^{-d - n - 1}(\mathcal{E}) \otimes \wedge^{n + 1}(\mathcal{E})$
との間に所望の完全対を誘導する。
\end{proof}

\section{局所 Noether スキーム上の連接層}
\label{coherent-section-coherent-sheaves}

\noindent
任意の環付き空間上の連接加群の概念を、加群、節
\ref{modules-section-coherent} で定義した。非 Noether スキーム上の連接層を
考えることもできるが、連接層を考えるときは常に基礎スキームが局所 Noether
であると仮定する。局所 Noether スキーム上の連接層は次のように特徴づけられる。

\begin{lemma}
\label{coherent-lemma-coherent-Noetherian}
$X$ を局所 Noether スキームとし、$\mathcal{F}$ を
$\mathcal{O}_X$-加群とする。次は同値である。
\begin{enumerate}
\item $\mathcal{F}$ は連接である。
\item $\mathcal{F}$ は準連接かつ有限型の $\mathcal{O}_X$-加群である。
\item $\mathcal{F}$ は有限表示の $\mathcal{O}_X$-加群である。
\item 任意のアフィン開集合 $\Spec(A) = U \subset X$ に対して
$\mathcal{F}|_U = \widetilde M$ であり、$M$ は有限生成 $A$-加群である。
\item アフィン開被覆 $X = \bigcup U_i$, $U_i = \Spec(A_i)$ が存在し、
各 $\mathcal{F}|_{U_i} = \widetilde M_i$ について $M_i$ は有限生成
$A_i$-加群である。
\end{enumerate}
特に $\mathcal{O}_X$ は連接であり、任意の可逆 $\mathcal{O}_X$-加群は
連接であり、さらに一般に任意の有限局所自由 $\mathcal{O}_X$-加群は連接である。
\end{lemma}

\begin{proof}
含意 (1) $\Rightarrow$ (2) および (1) $\Rightarrow$ (3) は一般に成り立つ。
加群、補題 \ref{modules-lemma-coherent-finite-presentation} を参照せよ。
$\mathcal{F}$ が有限表示されていれば $\mathcal{F}$ は準連接である。
加群、補題 \ref{modules-lemma-finite-presentation-quasi-coherent} を参照せよ。
従って (3) $\Rightarrow$ (2) も成り立つ。

\medskip\noindent
$\mathcal{F}$ を準連接かつ有限型の $\mathcal{O}_X$-加群と仮定する。
性質、補題 \ref{properties-lemma-finite-type-module} により、任意の
アフィン開集合 $\Spec(A) = U \subset X$ 上で
$\mathcal{F}|_U = \widetilde M$ であり、$M$ は有限生成 $A$-加群である。
$A$ は Noether なので、$M$ は有限自由加群による表示
$$
A^{\oplus m} \to A^{\oplus n} \to M \to 0.
$$
をもつ。従って性質、補題
\ref{properties-lemma-finite-presentation-module} により
$\mathcal{F}$ は有限表示である。言い換えると (2) $\Rightarrow$ (3) である。

\medskip\noindent
加群、補題 \ref{modules-lemma-coherent-structure-sheaf} により、(3) が (1) を
含意することを示すには $\mathcal{O}_X$ が連接であることを示せば十分である。
従って示すべきことは次である。任意の開集合 $U \subset X$ と、有限個の切断
$f_i \in \mathcal{O}_X(U)$（$i = 1, \ldots, n$）が与えられたとき、写像
$\bigoplus_{i = 1, \ldots, n} \mathcal{O}_U \to \mathcal{O}_U$ の核は
有限型である。有限型であることは局所的性質なので、任意の $x \in U$ の
近傍で確かめれば十分である。従って $U = \Spec(A)$ はアフィンであると
仮定してよい。この場合、$f_1, \ldots, f_n \in A$ は $A$ の元である。
$A$ は Noether なので（性質、補題
\ref{properties-lemma-locally-Noetherian} を参照）、核 $K$ は、写像
$\bigoplus_{i = 1, \ldots, n} A \to A$ の核として有限生成 $A$-加群である。
例えば代数、補題 \ref{algebra-lemma-Noetherian-basic} を参照せよ。
関手\ $\widetilde{ }$\ は完全なので（スキーム、補題
\ref{schemes-lemma-spec-sheaves} を参照）、完全列
$$
\widetilde K \to
\bigoplus\nolimits_{i = 1, \ldots, n} \mathcal{O}_U \to
\mathcal{O}_U
$$
を得る。再び性質、補題
\ref{properties-lemma-finite-type-module} により $\widetilde K$ は有限型である。
従って (1), (2), (3) はすべて同値である。

\medskip\noindent
性質、補題 \ref{properties-lemma-finite-type-module} から (2) が (4) を
含意する。(4) が (5) を含意することは自明である。スキーム、節
\ref{schemes-section-quasi-coherent} の議論から (5) は
$\mathcal{F}$ が準連接であることを含意し、また (5) が
$\mathcal{F}$ の有限型性を含意することも明らかである。従って (5) は
(2) を含意し、証明が完了する。
\end{proof}

\begin{lemma}
\label{coherent-lemma-coherent-abelian-Noetherian}
$X$ を局所 Noether スキームとする。連接 $\mathcal{O}_X$-加群の圏は
アーベル圏である。より正確には、連接 $\mathcal{O}_X$-加群の射の核と余核は
連接である。連接層の任意の拡大は連接である。
\end{lemma}

\begin{proof}
これは加群、補題 \ref{modules-lemma-coherent-abelian} を特別な場合に
言い換えたものである。
\end{proof}

\noindent
次の補題は、連接 $\mathcal{O}_X$-加群の圏について、一般の環付き空間
$X$ 上では常に成り立つとは限らない。

\begin{lemma}
\label{coherent-lemma-coherent-Noetherian-quasi-coherent-sub-quotient}
$X$ を局所 Noether スキームとし、$\mathcal{F}$ を連接
$\mathcal{O}_X$-加群とする。$\mathcal{F}$ の任意の準連接部分加群は
連接である。$\mathcal{F}$ の任意の準連接商加群も連接である。
\end{lemma}

\begin{proof}
$X$ はアフィン、$X = \Spec(A)$ であると仮定してよい。性質、補題
\ref{properties-lemma-locally-Noetherian} により $A$ は Noether である。
補題 \ref{coherent-lemma-coherent-Noetherian} により、これは代数の主張に帰着する。
対応する代数的主張は、有限生成 $A$-加群の商または部分加群が有限生成
$A$-加群であるというものである。例えば代数、補題
\ref{algebra-lemma-Noetherian-basic} を参照せよ。
\end{proof}

\begin{lemma}
\label{coherent-lemma-tensor-hom-coherent}
$X$ を局所 Noether スキームとし、$\mathcal{F}$, $\mathcal{G}$ を連接
$\mathcal{O}_X$-加群とする。このとき $\mathcal{O}_X$-加群
$\mathcal{F} \otimes_{\mathcal{O}_X} \mathcal{G}$ および
$\SheafHom_{\mathcal{O}_X}(\mathcal{F}, \mathcal{G})$ は連接である。
\end{lemma}

\begin{proof}
加群、補題
\ref{modules-lemma-internal-hom-locally-kernel-direct-sum} において
$\SheafHom_{\mathcal{O}_X}(\mathcal{F}, \mathcal{G})$ が連接であることが
示されている。テンソル積に関する結果は、加群、補題
\ref{modules-lemma-tensor-product-permanence} である。
\end{proof}

\begin{lemma}
\label{coherent-lemma-local-isomorphism}
$X$ を局所 Noether スキームとし、$\mathcal{F}$, $\mathcal{G}$ を連接
$\mathcal{O}_X$-加群とする。$\varphi : \mathcal{G} \to \mathcal{F}$ を
$\mathcal{O}_X$-加群の準同型とし、$x \in X$ とする。
\begin{enumerate}
\item $\mathcal{F}_x = 0$ ならば、開近傍 $U \subset X$ で、$x$ を含み
$\mathcal{F}|_U = 0$ を満たすものが存在する。
\item $\varphi_x : \mathcal{G}_x \to \mathcal{F}_x$ が単射ならば、
開近傍 $U \subset X$ で、$x$ を含み $\varphi|_U$ が単射となるものが存在する。
\item $\varphi_x : \mathcal{G}_x \to \mathcal{F}_x$ が全射ならば、
開近傍 $U \subset X$ で、$x$ を含み $\varphi|_U$ が全射となるものが存在する。
\item $\varphi_x : \mathcal{G}_x \to \mathcal{F}_x$ が全単射ならば、
開近傍 $U \subset X$ で、$x$ を含み $\varphi|_U$ が同型となるものが存在する。
\end{enumerate}
\end{lemma}

\begin{proof}
加群、補題
\ref{modules-lemma-finite-type-surjective-on-stalk},
\ref{modules-lemma-finite-type-stalk-zero}, および
\ref{modules-lemma-finite-type-to-coherent-injective-on-stalk}
を参照せよ。
\end{proof}

\begin{lemma}
\label{coherent-lemma-map-stalks-local-map}
$X$ を局所 Noether スキームとし、$\mathcal{F}$, $\mathcal{G}$ を連接
$\mathcal{O}_X$-加群とする。$x \in X$ とし、
$\psi : \mathcal{G}_x \to \mathcal{F}_x$ を
$\mathcal{O}_{X, x}$-加群の写像と仮定する。このとき開近傍
$U \subset X$ で、$x$ を含むものと、写像
$\varphi : \mathcal{G}|_U \to \mathcal{F}|_U$ で
$\varphi_x = \psi$ を満たすものが存在する。
\end{lemma}

\begin{proof}
補題 \ref{coherent-lemma-coherent-Noetherian} により、これは加群、補題
\ref{modules-lemma-stalk-internal-hom} の言い換えである。
\end{proof}

\begin{lemma}
\label{coherent-lemma-coherent-support-closed}
$X$ を局所 Noether スキームとし、$\mathcal{F}$ を連接
$\mathcal{O}_X$-加群とする。このとき $\text{Supp}(\mathcal{F})$ は閉であり、
$\mathcal{F}$ は $\mathcal{F}$ のスキーム論的台の上の連接層から来る。
射、定義 \ref{morphisms-definition-scheme-theoretic-support} を参照せよ。
\end{lemma}

\begin{proof}
$i : Z \to X$ を $\mathcal{F}$ のスキーム論的台とし、$\mathcal{G}$ を
$Z$ 上の有限型準連接層で $i_*\mathcal{G} \cong \mathcal{F}$ を満たすものとする。
$Z = \text{Supp}(\mathcal{F})$ なので台は閉である。スキーム $Z$ は、射、補題
\ref{morphisms-lemma-immersion-locally-finite-type} および
\ref{morphisms-lemma-finite-type-noetherian} により局所 Noether である。
最後に、補題 \ref{coherent-lemma-coherent-Noetherian} により $\mathcal{G}$ は連接
$\mathcal{O}_Z$-加群である。
\end{proof}

\begin{lemma}
\label{coherent-lemma-i-star-equivalence}
$i : Z \to X$ を局所 Noether スキームの閉埋め込みとし、
$\mathcal{I} \subset \mathcal{O}_X$ を $Z$ を切り出す準連接イデアル層とする。
関手 $i_*$ は、連接 $\mathcal{O}_X$-加群で $\mathcal{I}$ により
零化されるものの圏と、
連接 $\mathcal{O}_Z$-加群の圏との間の圏同値を誘導する。
\end{lemma}

\begin{proof}
射、補題 \ref{morphisms-lemma-i-star-equivalence} により、この関手は充満忠実である。
$\mathcal{F}$ を連接 $\mathcal{O}_X$-加群で $\mathcal{I}$ により
零化されるものとする。
射、補題 \ref{morphisms-lemma-i-star-equivalence} により、
$\mathcal{F} = i_*\mathcal{G}$ と書ける。ここで $\mathcal{G}$ は
$Z$ 上のある準連接層である。加群、補題
\ref{modules-lemma-i-star-reflects-finite-type} により、
$\mathcal{G}$ は有限型である。従って補題
\ref{coherent-lemma-coherent-Noetherian} により $\mathcal{G}$ は連接である。
ゆえに、この関手は本質的全射でもあり、所望の結論を得る。
\end{proof}

\begin{lemma}
\label{coherent-lemma-finite-pushforward-coherent}
$f : X \to Y$ をスキームの射とし、$\mathcal{F}$ を準連接
$\mathcal{O}_X$-加群とする。$f$ は有限で $Y$ は局所 Noether であると仮定する。
このとき $R^pf_*\mathcal{F} = 0$ は $p > 0$ に対して成り立ち、
$f_*\mathcal{F}$ は $\mathcal{F}$ が連接ならば連接である。
\end{lemma}

\begin{proof}
高次順像は補題 \ref{coherent-lemma-relative-affine-vanishing} により、また有限射が
（定義により）アフィンであることから消滅する。仮定から $X$ も局所 Noether
であることに注意せよ（射、補題
\ref{morphisms-lemma-finite-type-noetherian} を参照）。従って主張は意味をもつ。
$\Spec(A) = V \subset Y$ をアフィン開部分集合とする。射、定義
\ref{morphisms-definition-integral} により
$f^{-1}(V) = \Spec(B)$ であり、$A \to B$ は有限である。
補題 \ref{coherent-lemma-coherent-Noetherian} により、補題の主張は次の代数的事実に
帰着する。$M$ が有限生成 $B$-加群ならば、$M$ は $A$-加群とみても有限生成である。
代数、補題 \ref{algebra-lemma-finite-module-over-finite-extension} を参照せよ。
\end{proof}

\noindent
補題の状況では、高次順像も消滅するので連接である。これが局所 Noether
スキーム間の任意の固有射について成り立つことを、命題
\ref{coherent-proposition-proper-pushforward-coherent} で示す。

\begin{lemma}
\label{coherent-lemma-coherent-support-dimension-0}
$X$ を局所 Noether スキームとし、$\mathcal{F}$ を
$\dim(\text{Supp}(\mathcal{F})) \leq 0$ を満たす連接層とする。
このとき $\mathcal{F}$ は大域切断で生成され、
$H^i(X, \mathcal{F}) = 0$ は $i > 0$ に対して成り立つ。
\end{lemma}

\begin{proof}
補題 \ref{coherent-lemma-coherent-support-closed} により、
$\mathcal{F} = i_*\mathcal{G}$ と書ける。ここで
$i : Z \to X$ は $\mathcal{F}$ のスキーム論的台の包含であり、
$\mathcal{G}$ は連接 $\mathcal{O}_Z$-加群である。
$Z$ の次元は $0$ なので、$Z$ はアフィンの非交和である
（性質、補題 \ref{properties-lemma-locally-Noetherian-dimension-0}）。
従って $\mathcal{G}$ は大域生成され、$\mathcal{G}$ の高次コホモロジー群は零である
（補題 \ref{coherent-lemma-quasi-coherent-affine-cohomology-zero}）。
ゆえに $\mathcal{F} = i_*\mathcal{G}$ は大域生成される。
$\mathcal{F}$ と $\mathcal{G}$ のコホモロジーは一致する
（閉埋め込みはアフィンなので補題
\ref{coherent-lemma-relative-affine-cohomology} が適用される）から、
$\mathcal{F}$ の高次コホモロジー群は零である。
\end{proof}

\begin{lemma}
\label{coherent-lemma-pushforward-coherent-on-open}
$X$ をスキームとし、$j : U \to X$ を開集合の包含とする。
$T \subset X$ を $U$ に含まれる閉部分集合とする。
$\mathcal{F}$ が連接 $\mathcal{O}_U$-加群で
$\text{Supp}(\mathcal{F}) \subset T$ を満たすならば、$j_*\mathcal{F}$ は連接
$\mathcal{O}_X$-加群である。
\end{lemma}

\begin{proof}
開被覆 $X = U \cup (X \setminus T)$ を考える。
$j_*\mathcal{F}|_U = \mathcal{F}$ は連接であり、
$j_*\mathcal{F}|_{X \setminus T} = 0$ も連接である。
従って $j_*\mathcal{F}$ は連接である。
\end{proof}

\section{Noether スキーム上の連接層}
\label{coherent-section-coherent-quasi-compact}

\noindent
本節では、Noether スキーム上の連接層のいくつかの性質を述べる。

\begin{lemma}
\label{coherent-lemma-acc-coherent}
$X$ を Noether スキームとし、$\mathcal{F}$ を連接
$\mathcal{O}_X$-加群とする。$\mathcal{F}$ の準連接部分加群について
昇鎖条件が成り立つ。言い換えると、準連接部分加群の任意の列
$$
\mathcal{F}_1 \subset \mathcal{F}_2 \subset \ldots \subset \mathcal{F}
$$
が与えられれば、$\mathcal{F}_n = \mathcal{F}_{n + 1} = \ldots $ となる
ある $n \geq 0$ が存在する。
\end{lemma}

\begin{proof}
有限アフィン開被覆を選ぶ。この被覆の各元上では、代数、補題
\ref{algebra-lemma-Noetherian-basic} により列は停止する。
従って補題が従う。
\end{proof}

\begin{lemma}
\label{coherent-lemma-power-ideal-kills-sheaf}
$X$ を Noether スキームとし、$\mathcal{F}$ を $X$ 上の連接層とする。
$\mathcal{I} \subset \mathcal{O}_X$ を閉部分スキーム
$Z \subset X$ に対応する準連接イデアル層とする。このとき、ある
$n \geq 0$ に対して $\mathcal{I}^n\mathcal{F} = 0$ となることと、
$\text{Supp}(\mathcal{F}) \subset Z$ が集合論的に成り立つこととは同値である。
\end{lemma}

\begin{proof}
$X$ は Noether 環のスペクトルによる有限被覆をもつので、これは代数、補題
\ref{algebra-lemma-Noetherian-power-ideal-kills-module}
から直ちに従う。
\end{proof}

\begin{lemma}[Artin-Rees]
\label{coherent-lemma-Artin-Rees}
$X$ を Noether スキームとし、$\mathcal{F}$ を $X$ 上の連接層とする。
$\mathcal{G} \subset \mathcal{F}$ を準連接部分層とし、
$\mathcal{I} \subset \mathcal{O}_X$ を準連接イデアル層とする。
このとき、ある $c \geq 0$ が存在して、すべての $n \geq c$ に対し
$$
\mathcal{I}^{n - c}(\mathcal{I}^c\mathcal{F} \cap \mathcal{G})
=
\mathcal{I}^n\mathcal{F} \cap \mathcal{G}
$$
が成り立つ。
\end{lemma}

\begin{proof}
$X$ は Noether 環のスペクトルによる有限被覆をもつので、これは代数、補題
\ref{algebra-lemma-Artin-Rees} から直ちに従う。
\end{proof}

\begin{lemma}
\label{coherent-lemma-directed-colimit-coherent}
$X$ を Noether スキームとする。任意の準連接 $\mathcal{O}_X$-加群は、
その連接部分加群のフィルター余極限である。
\end{lemma}

\begin{proof}
これは性質、補題
\ref{properties-lemma-quasi-coherent-colimit-finite-type} の言い換えである。
実際、補題 \ref{coherent-lemma-coherent-Noetherian} により、有限型準連接
$\mathcal{O}_X$-加群は連接である。
\end{proof}

\begin{lemma}
\label{coherent-lemma-homs-over-open}
$X$ を Noether スキームとする。$\mathcal{F}$ を準連接
$\mathcal{O}_X$-加群とし、$\mathcal{G}$ を連接
$\mathcal{O}_X$-加群とする。$\mathcal{I} \subset \mathcal{O}_X$ を
準連接イデアル層とする。対応する閉部分スキームを
$Z \subset X$ と表し、$U = X \setminus Z$ と置く。このとき標準同型
$$
\colim_n \Hom_{\mathcal{O}_X}(\mathcal{I}^n\mathcal{G}, \mathcal{F})
\longrightarrow
\Hom_{\mathcal{O}_U}(\mathcal{G}|_U, \mathcal{F}|_U).
$$
が存在する。特に同型
$$
\colim_n \Hom_{\mathcal{O}_X}(
\mathcal{I}^n, \mathcal{F})
\longrightarrow
\Gamma(U, \mathcal{F}).
$$
が存在する。
\end{lemma}

\begin{proof}
まず第二の写像が同型であることを示す。性質、補題
\ref{properties-lemma-sections-over-quasi-compact-open} により単射である。
$\mathcal{F}$ はその連接部分加群の合併なので、性質、補題
\ref{properties-lemma-quasi-coherent-colimit-finite-type}
（および補題 \ref{coherent-lemma-coherent-Noetherian}）を参照すると、全射性を
示す際には $\mathcal{F}$ が連接であると仮定してよいし、そうする。
$\mathcal{F}_n$ を、$\mathcal{F}$ の切断のうち $\mathcal{I}^n$ により
零化されるものからなる準連接部分層とする。性質、補題
\ref{properties-lemma-sections-over-quasi-compact-open} を参照せよ。
$\mathcal{F}_1 \subset \mathcal{F}_2 \subset \ldots$ なので、補題
\ref{coherent-lemma-acc-coherent} により、
$\mathcal{F}_n = \mathcal{F}_{n + 1} = \ldots $ となる
ある $n \geq 0$ が存在する。$\mathcal{H} = \mathcal{F}_n$ と置く。
ここで $n$ はこの整数である。Artin--Rees（補題
\ref{coherent-lemma-Artin-Rees}）により、ある $c \geq 0$ が存在して
$\mathcal{I}^m\mathcal{F} \cap \mathcal{H}
\subset \mathcal{I}^{m - c}\mathcal{H}$ となる。$m = n + c$ と選べば
$\mathcal{I}^m\mathcal{F} \cap \mathcal{H} \subset \mathcal{I}^n\mathcal{H}
= 0$ を得る。従って $\mathcal{F}' = \mathcal{I}^m\mathcal{F}$ と置くと、
$\mathcal{F}' \cap \mathcal{F}_n = 0$ かつ
$\mathcal{F}'|_U = \mathcal{F}|_U$ である。特に、部分層
$(\mathcal{F}')_N$、すなわち $\mathcal{I}^N$ により零化される切断の層は、
すべての $N \geq 0$ に対して
零である。従って性質、補題
\ref{properties-lemma-sections-over-quasi-compact-open} により、次の可換図式の
上の水平矢印は全単射である。
$$
\xymatrix{
\colim_n \Hom_{\mathcal{O}_X}(
\mathcal{I}^n, \mathcal{F}')
\ar[r] \ar[d] &
\Gamma(U, \mathcal{F}') \ar[d] \\
\colim_n \Hom_{\mathcal{O}_X}(
\mathcal{I}^n, \mathcal{F})
\ar[r] &
\Gamma(U, \mathcal{F})
}
$$
右の垂直矢印も全単射なので、下の水平矢印は所望の通り全射である。

\medskip\noindent
次に、補題の第一の矢印が全単射であることを示す。補題
\ref{coherent-lemma-coherent-Noetherian} により層 $\mathcal{G}$ は有限表示であり、
従って層
$\mathcal{H} = \SheafHom_{\mathcal{O}_X}(\mathcal{G}, \mathcal{F})$
は準連接である。スキーム、節
\ref{schemes-section-quasi-coherent} を参照せよ。定義により
$$
\mathcal{H}(U)
=
\Hom_{\mathcal{O}_U}(\mathcal{G}|_U, \mathcal{F}|_U)
$$
である。補題の第一の矢印の右辺に属する $\psi$ を選ぶ、すなわち
$\psi \in \mathcal{H}(U)$ とする。先ほど示した結果を
$\mathcal{H}$ に適用すると、ある $n \geq 0$ と写像
$\varphi : \mathcal{I}^n \to \mathcal{H}$ が存在し、
$\psi$ を $U$ への制限として与える。加群、補題
\ref{modules-lemma-internal-hom} により
$\varphi$ は写像
$$
\varphi :
\mathcal{I}^n \otimes_{\mathcal{O}_X} \mathcal{G}
\longrightarrow
\mathcal{F}.
$$
に対応する。これはほぼ欲しい写像であるが、その始域は
$\mathcal{I}^n$ と $\mathcal{G}$ の積ではなくテンソル積である。
$n$ を増やすことで、この差が無関係になることを示す。短完全列
$$
0 \to \mathcal{K} \to
\mathcal{I}^n \otimes_{\mathcal{O}_X} \mathcal{G} \to
\mathcal{I}^n\mathcal{G} \to 0
$$
を考える。ここで $\mathcal{K}$ は核として定める。
$\mathcal{I}^n\mathcal{K} = 0$ であることに注意せよ（証明は省略する）。
再び Artin--Rees により、十分大きなある $m$ に対して
$$
\mathcal{K}
\cap
\mathcal{I}^m(\mathcal{I}^n \otimes_{\mathcal{O}_X} \mathcal{G})
=
0
$$
となる。言い換えると、写像
$$
\mathcal{I}^m(\mathcal{I}^n \otimes_{\mathcal{O}_X} \mathcal{G})
\longrightarrow
\mathcal{I}^{n + m}\mathcal{G}
$$
は同型である。この部分加群への $\varphi'$ を $\varphi$ の制限とし、
これを写像
$\mathcal{I}^{m + n}\mathcal{G} \to \mathcal{F}$ とみなす。
すると $\varphi'$ は、補題の第一の矢印の左辺の元であって、その矢印により
$\psi$ に写るものを与える。言い換えると、矢印の全射性が示された。
単射性の証明は省略する。
\end{proof}

\begin{lemma}
\label{coherent-lemma-extend-coherent}
$X$ を局所 Noether スキームとし、$\mathcal{F}$, $\mathcal{G}$ を連接
$\mathcal{O}_X$-加群とする。$U \subset X$ を開集合とし、
$\varphi : \mathcal{F}|_U \to \mathcal{G}|_U$ を
$\mathcal{O}_U$-加群写像とする。このとき、連接部分加群
$\mathcal{F}' \subset \mathcal{F}$ で、$\mathcal{F}$ と $U$ 上で一致し、
$\varphi$ が $\varphi' : \mathcal{F}' \to \mathcal{G}$ に延長されるものが
存在する。
\end{lemma}

\begin{proof}
$\mathcal{I} \subset \mathcal{O}_X$ を、$X \setminus U$ 上の被約誘導
スキーム構造を切り出す連接イデアル層とする。$X$ が Noether ならば、補題
\ref{coherent-lemma-homs-over-open} により、
$\mathcal{F}' = \mathcal{I}^n\mathcal{F}$ とある $n$ に対して取ることができる。
一般の場合は、Zorn の補題を用いてこの場合から従う。

\medskip\noindent
三つ組 $(U', \mathcal{F}', \varphi')$ 全体の集合を考える。ここで
$U \subset U' \subset X$ は開集合であり、
$\mathcal{F}' \subset \mathcal{F}|_{U'}$ は $\mathcal{F}$ と $U$ 上で
一致する連接部分層であり、
$\varphi' : \mathcal{F}' \to \mathcal{G}|_{U'}$ は
$\varphi$ を $U$ への制限として与える。三つ組の順序を
$(U'', \mathcal{F}'', \varphi'') \geq (U', \mathcal{F}', \varphi')$
が成り立つことと、$U'' \supset U'$,
$\mathcal{F}''|_{U'} = \mathcal{F}'$, および
$\varphi''|_{U'} = \varphi'$ がすべて成り立つこととが同値となるように定める。
三つ組 $(U_i, \mathcal{F}_i, \varphi_i)$ の全順序部分集合が与えられれば、
$\mathcal{F}_i$ を部分層 $\mathcal{F}'$ に貼り合わせることができる。
これは $\mathcal{F}$ の $U' = \bigcup U_i$ 上の部分層であり、
$\varphi$ を写像
$\varphi' : \mathcal{F}' \to \mathcal{G}|_{U'}$ に延長できることは明らかである。
従って、三つ組の任意の全順序部分集合は上界をもつ。最後に、
$(U', \mathcal{F}', \varphi')$ を任意の三つ組とし、
$U' \not = X$ と仮定する。このとき、アフィン開集合
$W \subset X$ で $U'$ に含まれないものを選べる。第一段落の結果により、部分層
$\mathcal{F}'|_{W \cap U'}$ と制限 $\varphi'|_{W \cap U'}$ を、ある部分層
$\mathcal{F}'' \subset \mathcal{F}|_W$ と写像
$\varphi'' : \mathcal{F}'' \to \mathcal{G}|_W$ に延長できる。
$(\mathcal{F}', \varphi')$ と $(\mathcal{F}'', \varphi'')$ の
$W \cap U'$ 上での一致とは、まさにこれらを三つ組
$(U' \cup W, \mathcal{F}''', \varphi''')$ に延長できることを意味する。
従って、任意の極大三つ組 $(U', \mathcal{F}', \varphi')$ は
（Zorn の補題により存在し）$U' = X$ を満たす。これで証明が完了する。
\end{proof}

\section{深さ}
\label{coherent-section-depth}

\noindent
本節では、局所 Noether スキーム上の連接加群の深さと性質
$(S_k)$ について少し述べる。この概念は局所 Noether スキームについて
既に性質、節 \ref{properties-section-Rk} で論じたことに注意せよ。

\begin{definition}
\label{coherent-definition-depth}
$X$ を局所 Noether スキームとし、$\mathcal{F}$ を連接
$\mathcal{O}_X$-加群とする。$k \geq 0$ を整数とする。
\begin{enumerate}
\item $\mathcal{F}$ が {\it ある点で深さ $k$ をもつ}とは、その点を
$x$ と書き、それが $X$ に属するとき、
$\text{depth}_{\mathcal{O}_{X, x}}(\mathcal{F}_x) = k$ であることをいう。
\item $X$ が {\it ある点で深さ $k$ をもつ}とは、その点を $x$ と書き、
それが $X$ に属するとき、$\text{depth}(\mathcal{O}_{X, x}) = k$
であることをいう。
\item $\mathcal{F}$ が性質 {\it $(S_k)$} をもつとは、
$$
\text{depth}_{\mathcal{O}_{X, x}}(\mathcal{F}_x)
\geq \min(k, \dim(\text{Supp}(\mathcal{F}_x)))
$$
がすべての $x \in X$ に対して成り立つことをいう。
\item $X$ が性質 {\it $(S_k)$} をもつとは、$\mathcal{O}_X$ が性質
$(S_k)$ をもつことをいう。
\end{enumerate}
\end{definition}

\noindent
任意の連接層は条件 $(S_0)$ を満たす。条件 $(S_1)$ は埋め込まれた随伴点を
もたないことと同値である。因子、補題
\ref{divisors-lemma-S1-no-embedded} を参照せよ。

\begin{lemma}
\label{coherent-lemma-hom-into-depth}
$X$ を局所 Noether スキームとする。$\mathcal{F}$, $\mathcal{G}$ を連接
$\mathcal{O}_X$-加群とし、$x \in X$ とする。
\begin{enumerate}
\item $\mathcal{G}_x$ の深さが $\geq 1$ ならば、
$\SheafHom_{\mathcal{O}_X}(\mathcal{F}, \mathcal{G})_x$
の深さも $\geq 1$ である。
\item $\mathcal{G}_x$ の深さが $\geq 2$ ならば、
$\Hom_{\mathcal{O}_X}(\mathcal{F}, \mathcal{G})_x$ の深さも
$\geq 2$ である。
\end{enumerate}
\end{lemma}

\begin{proof}
補題 \ref{coherent-lemma-tensor-hom-coherent} により
$\SheafHom_{\mathcal{O}_X}(\mathcal{F}, \mathcal{G})$ は連接
$\mathcal{O}_X$-加群である。連接加群は有限表示なので
（補題 \ref{coherent-lemma-coherent-Noetherian}）、茎を取ることは
$\SheafHom$ および $\Hom$ を取ることと可換である。加群、補題
\ref{modules-lemma-stalk-internal-hom} を参照せよ。従って局所環上の有限生成加群の
場合に帰着し、これは代数詳論、補題
\ref{more-algebra-lemma-hom-into-depth} である。
\end{proof}

\begin{lemma}
\label{coherent-lemma-hom-into-S2}
$X$ を局所 Noether スキームとする。$\mathcal{F}$, $\mathcal{G}$ を連接
$\mathcal{O}_X$-加群とする。
\begin{enumerate}
\item $\mathcal{G}$ が性質 $(S_1)$ をもつならば、
$\SheafHom_{\mathcal{O}_X}(\mathcal{F}, \mathcal{G})$ も性質 $(S_1)$ をもつ。
\item $\mathcal{G}$ が性質 $(S_2)$ をもつならば、
$\SheafHom_{\mathcal{O}_X}(\mathcal{F}, \mathcal{G})$ も性質 $(S_2)$ をもつ。
\end{enumerate}
\end{lemma}

\begin{proof}
補題 \ref{coherent-lemma-hom-into-depth} と定義から直ちに従う。
\end{proof}

\noindent
性質、補題 \ref{properties-lemma-scheme-CM-iff-all-Sk} において、局所
Noether スキームが Cohen--Macaulay であることと、$(S_k)$ がすべての
$k$ に対して成り立つこととは同値であることを見た。従って次の定義を導入するのは
自然である。この定義は、すべての茎が Cohen--Macaulay 加群であるという条件と
同値である。

\begin{definition}
\label{coherent-definition-Cohen-Macaulay}
$X$ を局所 Noether スキームとし、$\mathcal{F}$ を連接
$\mathcal{O}_X$-加群とする。$\mathcal{F}$ が {\it Cohen--Macaulay}
であるとは、$(S_k)$ がすべての $k \geq 0$ に対して成り立つことをいう。
\end{definition}

\begin{lemma}
\label{coherent-lemma-Cohen-Macaulay-over-regular}
$X$ を正則スキームとし、$\mathcal{F}$ を連接
$\mathcal{O}_X$-加群とする。次は同値である。
\begin{enumerate}
\item $\mathcal{F}$ は Cohen--Macaulay であり、
$\text{Supp}(\mathcal{F}) = X$ である。
\item $\mathcal{F}$ は階数 $> 0$ の有限局所自由加群である。
\end{enumerate}
\end{lemma}

\begin{proof}
$x \in X$ とする。(2) が成り立つならば、$\mathcal{F}_x$ は自由
$\mathcal{O}_{X, x}$-加群で、階数は $> 0$ である。従って正則局所環は
Cohen--Macaulay なので（代数、補題
\ref{algebra-lemma-regular-ring-CM}）、
$\text{depth}(\mathcal{F}_x) = \dim(\mathcal{O}_{X, x})$ である。
逆に (1) が成り立つならば、$\mathcal{F}_x$ は
$\mathcal{O}_{X, x}$ 上の極大 Cohen--Macaulay 加群である
（代数、定義 \ref{algebra-definition-maximal-CM}）。従って代数、補題
\ref{algebra-lemma-regular-mcm-free} により $\mathcal{F}_x$ は自由である。
\end{proof}

\section{連接層のデヴィサージュ}
\label{coherent-section-devissage}

\noindent
$X$ を Noether スキームとする。整な閉部分スキーム
$i : Z \to X$ を考える。$i_*\mathcal{G}$ の形の連接層を考えると
便利なことが多い。ここで $\mathcal{G}$ は $Z$ 上の連接層である。
特に、$\mathcal{G}$ が階数 $1$ の捩れなし層である場合に、このような層に
関心がある。例えば $\mathcal{G}$ は $Z$ 上の非零イデアル層でもよく、
さらに特別には $\mathcal{G} = \mathcal{O}_Z$ でもよい。

\medskip\noindent
本節を通じて、連接層は有限型準連接層と同じものであり、連接層の
準連接な部分商は連接であるという事実を用いる。節
\ref{coherent-section-coherent-sheaves} を参照せよ。連接層の台は閉である。
加群、補題 \ref{modules-lemma-support-finite-type-closed} を参照せよ。

\begin{lemma}
\label{coherent-lemma-prepare-filter-support}
$X$ を Noether スキームとする。
$\mathcal{F}$ を $X$ 上の連接層とする。
$\text{Supp}(\mathcal{F}) = Z \cup Z'$ であり、$Z$, $Z'$ は閉であると仮定する。
このとき、連接層の短完全列
$$
0 \to \mathcal{G}' \to \mathcal{F} \to \mathcal{G} \to 0
$$
で、$\text{Supp}(\mathcal{G}') \subset Z'$ および
$\text{Supp}(\mathcal{G}) \subset Z$ を満たすものが存在する。
\end{lemma}

\begin{proof}
$\mathcal{I} \subset \mathcal{O}_X$ を、$Z$ 上の被約誘導閉部分スキーム構造を
定めるイデアル層とする。スキーム、補題
\ref{schemes-lemma-reduced-closed-subscheme} を参照せよ。
部分層 $\mathcal{G}'_n = \mathcal{I}^n\mathcal{F}$ と商
$\mathcal{G}_n = \mathcal{F}/\mathcal{I}^n\mathcal{F}$ を考える。
各 $n$ に対して短完全列
$$
0 \to \mathcal{G}'_n \to \mathcal{F} \to \mathcal{G}_n \to 0
$$
がある。$x$ を $Z' \setminus Z$ の任意の点とすると、
$\mathcal{I}_x = \mathcal{O}_{X, x}$ であり、従って
$\mathcal{G}_{n, x} = 0$ である。ゆえに
$\text{Supp}(\mathcal{G}_n) \subset Z$ である。$X \setminus Z'$ は
Noether スキームであることに注意せよ。従って補題
\ref{coherent-lemma-power-ideal-kills-sheaf} により、ある $n$ が存在して
$\mathcal{G}'_n|_{X \setminus Z'} =
\mathcal{I}^n\mathcal{F}|_{X \setminus Z'} = 0$ となる。
そのような $n$ に対して $\text{Supp}(\mathcal{G}'_n) \subset Z'$ である。
従って $\mathcal{G}' = \mathcal{G}'_n$ および
$\mathcal{G} = \mathcal{G}_n$ とおけばよい。
\end{proof}

\begin{lemma}
\label{coherent-lemma-prepare-filter-irreducible}
$X$ を Noether スキームとする。
$i : Z \to X$ を整な閉部分スキームとする。
$\xi \in Z$ を生成点とする。
$\mathcal{F}$ を $X$ 上の連接層とする。
$\mathcal{F}_\xi$ が $\mathfrak m_\xi$ で零化されると仮定する。
このとき、整数 $r \geq 0$、連接イデアル層
$\mathcal{I} \subset \mathcal{O}_Z$、および連接層の単射
$$
i_*\left(\mathcal{I}^{\oplus r}\right) \to \mathcal{F}
$$
で、$\xi$ のある近傍上で同型となるものが存在する。
\end{lemma}

\begin{proof}
$\mathcal{J} \subset \mathcal{O}_X$ を $Z$ のイデアル層とする。
$\mathcal{F}' \subset \mathcal{F}$ を、$\mathcal{F}$ の局所切断のうち
$\mathcal{J}$ で零化されるものからなる部分層とする。これは性質、補題
\ref{properties-lemma-sections-annihilated-by-ideal} により準連接層である。
さらに、$\mathcal{F}'_\xi = \mathcal{F}_\xi$ である。これは
$\mathcal{J}_\xi = \mathfrak m_\xi$ と性質、補題
\ref{properties-lemma-sections-annihilated-by-ideal} の (3) による。補題
\ref{coherent-lemma-local-isomorphism} により、$\mathcal{F}' \to \mathcal{F}$ は
$\xi$ のある近傍上で同型を誘導する。従って $\mathcal{F}$ を
$\mathcal{F}'$ で置き換え、$\mathcal{F}$ が $\mathcal{J}$ で零化されると
仮定してよい。

\medskip\noindent
$\mathcal{J}\mathcal{F} = 0$ と仮定する。補題
\ref{coherent-lemma-i-star-equivalence} により、$\mathcal{F} = i_*\mathcal{G}$ と書ける。
ここで $\mathcal{G}$ は $Z$ 上のある連接層である。射
$\mathcal{I}^{\oplus r} \to \mathcal{G}$ で、生成点 $\xi$ のある近傍上で
同型となるものを見つけられると仮定する。ここでこの点は $Z$ の生成点である。
このとき $i_*$ を適用すれば（これは左完全である）、補題の結論を得る。
従って $X = Z$ の場合に帰着した。

\medskip\noindent
$Z = X$ が生成点 $\xi$ をもつ整 Noether スキームであると仮定する。
この場合、$\mathcal{O}_{X, \xi} = \kappa(\xi)$ は $X$ の関数体であることに
注意せよ。$\mathcal{F}_\xi$ は有限 $\mathcal{O}_\xi$-加群なので、
$r = \dim_{\kappa(\xi)} \mathcal{F}_\xi$ は有限である。従って層
$\mathcal{O}_X^{\oplus r}$ と $\mathcal{F}$ の $\xi$ における茎は同型である。
補題 \ref{coherent-lemma-map-stalks-local-map} により、空でない開集合
$U \subset X$ と射
$\psi : \mathcal{O}_X^{\oplus r}|_U \to \mathcal{F}|_U$ で、
$\xi$ において同型となるものが存在する。従って補題
\ref{coherent-lemma-local-isomorphism} により、これは $\xi$ のある近傍上で同型である。
スキーム、補題 \ref{schemes-lemma-reduced-closed-subscheme} により、
準連接イデアル層 $\mathcal{I} \subset \mathcal{O}_X$ が存在し、それに付随する
閉部分スキーム $Y \subset X$ は $U$ の補集合である。補題
\ref{coherent-lemma-homs-over-open} により、ある $n \geq 0$ と射
$\mathcal{I}^n(\mathcal{O}_X^{\oplus r}) \to \mathcal{F}$ が存在し、
元の $\psi$ を $U$ 上で復元する。
$\mathcal{I}^n(\mathcal{O}_X^{\oplus r}) = (\mathcal{I}^n)^{\oplus r}$
なので、補題にあるような写像を得る。$X$ は整であり、この写像は
$X$ の生成点で単射なので、これは単射である（容易な証明は省略する）。
\end{proof}

\begin{lemma}
\label{coherent-lemma-coherent-filter}
$X$ を Noether スキームとする。
$\mathcal{F}$ を $X$ 上の連接層とする。
連接部分層によるフィルトレーション
$$
0 = \mathcal{F}_0 \subset \mathcal{F}_1 \subset
\ldots \subset \mathcal{F}_m = \mathcal{F}
$$
で、各 $j = 1, \ldots, m$ に対して、整な閉部分スキーム
$Z_j \subset X$ と非零連接イデアル層
$\mathcal{I}_j \subset \mathcal{O}_{Z_j}$ が存在し、
$$
\mathcal{F}_j/\mathcal{F}_{j - 1}
\cong (Z_j \to X)_* \mathcal{I}_j
$$
となるものが存在する。
\end{lemma}

\begin{proof}
集合
$$
\mathcal{T} =
\left\{
\begin{matrix}
Z \subset X
\text{ closed such that there exists a coherent sheaf }
\mathcal{F} \\
\text{ with }
\text{Supp}(\mathcal{F}) = Z
\text{ for which the lemma is wrong}
\end{matrix}
\right\}
$$
を考える。$\mathcal{T}$ が空であることを示したい。そうでなければ、
$X$ は Noether なので、極小元 $Z \in \mathcal{T}$ を選べる。これは、
ある連接層 $\mathcal{F}$（$X$ 上のもの）が存在し、その台が $Z$ であり、
補題が成り立たないことを意味する。台が空の層は零層だけであり、それに対して
補題は成り立つので、明らかに $Z \not = \emptyset$ である
（$m = 0$ とすればよい）。

\medskip\noindent
$Z$ が既約でなければ、$Z = Z_1 \cup Z_2$ と書ける。ここで
$Z_1, Z_2$ は閉で、$Z$ より真に小さい。このとき補題
\ref{coherent-lemma-prepare-filter-support} を適用して、連接層の短完全列
$$
0 \to
\mathcal{G}_1 \to
\mathcal{F} \to
\mathcal{G}_2 \to 0
$$
で、$\text{Supp}(\mathcal{G}_i) \subset Z_i$ を満たすものを得る。
$Z$ の極小性により、各 $\mathcal{G}_i$ は補題の主張にあるような
フィルトレーションをもつ。$\mathcal{F}$ 上に誘導されるフィルトレーションを
考えると矛盾に至る。従って $Z$ は既約である。

\medskip\noindent
$Z$ が既約であると仮定する。$\mathcal{J}$ を、$Z$ の被約誘導閉部分スキーム
構造を切り出すイデアル層とする。スキーム、補題
\ref{schemes-lemma-reduced-closed-subscheme} を参照せよ。補題
\ref{coherent-lemma-power-ideal-kills-sheaf} により、ある $n \geq 0$ に対して
$\mathcal{J}^n\mathcal{F} = 0$ である。従ってフィルトレーション
$$
0 = \mathcal{J}^n\mathcal{F} \subset \mathcal{J}^{n - 1}\mathcal{F}
\subset \ldots \subset \mathcal{J}\mathcal{F} \subset \mathcal{F}
$$
を得て、その各逐次部分商は $\mathcal{J}$ で零化される。従って、これらの
部分商がそれぞれ補題の主張にあるようなフィルトレーションをもてば、
$\mathcal{F}$ もそうである。言い換えれば、$\mathcal{J}$ が
$\mathcal{F}$ を零化すると仮定してよい。

\medskip\noindent
$Z$ が既約で $\mathcal{J}\mathcal{F} = 0$ である場合、補題
\ref{coherent-lemma-prepare-filter-irreducible} を適用できる。これにより短完全列
$$
0 \to
i_*(\mathcal{I}^{\oplus r}) \to
\mathcal{F} \to
\mathcal{Q} \to 0
$$
を得る。ここで $\mathcal{Q}$ は商として定義される。先ほど引用した補題により
$\mathcal{Q}$ は $\xi$ のある近傍で零なので、$\mathcal{Q}$ の台は
$Z$ より真に小さい。従って $\mathcal{Q}$ は、$Z$ の極小性により所望の型の
フィルトレーションをもつ。しかし、そうすれば明らかに
$\mathcal{F}$ も同じ型のフィルトレーションをもち、最後の矛盾を得る。
\end{proof}

\begin{lemma}
\label{coherent-lemma-property-initial}
$X$ を Noether スキームとする。
$\mathcal{P}$ を $X$ 上の連接層の性質とする。次を仮定する。
\begin{enumerate}
\item 連接層の任意の短完全列
$$
0 \to \mathcal{F}_1 \to \mathcal{F} \to \mathcal{F}_2 \to 0
$$
に対し、$\mathcal{F}_i$, $i = 1, 2$ が性質 $\mathcal{P}$ をもてば、
$\mathcal{F}$ もそうである。
\item 任意の整な閉部分スキーム $Z \subset X$ と任意の準連接
イデアル層 $\mathcal{I} \subset \mathcal{O}_Z$ に対し、性質
$\mathcal{P}$ が $i_*\mathcal{I}$ に対して成り立つ。
\end{enumerate}
このとき、性質 $\mathcal{P}$ は $X$ 上のすべての連接層に対して成り立つ。
\end{lemma}

\begin{proof}
まず、$\mathcal{F}$ が連接層で、連接部分層によるフィルトレーション
$$
0 = \mathcal{F}_0 \subset \mathcal{F}_1 \subset
\ldots \subset \mathcal{F}_m = \mathcal{F}
$$
をもち、各 $\mathcal{F}_i/\mathcal{F}_{i - 1}$ が性質
$\mathcal{P}$ をもつならば、$\mathcal{F}$ もそうである。
これは $\mathcal{P}$ の性質 (1) から従う。一方、補題
\ref{coherent-lemma-coherent-filter} により、任意の $\mathcal{F}$ は (2) にある形の
逐次部分商をもつフィルトレーションを備える。従って補題が従う。
\end{proof}

\begin{lemma}
\label{coherent-lemma-property-irreducible}
$X$ を Noether スキームとする。$Z_0 \subset X$ を既約閉部分集合とし、
$\xi$ をその生成点とする。$\mathcal{P}$ を $X$ 上の連接層で台が
$Z_0$ に含まれるものの性質とし、次を仮定する。
\begin{enumerate}
\item 連接層の任意の短完全列において、三つの層のうち二つが性質
$\mathcal{P}$ をもてば、残りの層も同じ性質をもつ。
\item 任意の整な閉部分スキーム $Z \subset Z_0 \subset X$、
$Z \not = Z_0$、および任意の準連接イデアル層
$\mathcal{I} \subset \mathcal{O}_Z$ に対し、性質
$\mathcal{P}$ が $(Z \to X)_*\mathcal{I}$ に対して成り立つ。
\item 次を満たす連接層 $\mathcal{G}$（$X$ 上のもの）が存在する。
\begin{enumerate}
\item $\text{Supp}(\mathcal{G}) = Z_0$,
\item $\mathcal{G}_\xi$ は $\mathfrak m_\xi$ で零化される、
\item $\dim_{\kappa(\xi)} \mathcal{G}_\xi = 1$、かつ
\item 性質 $\mathcal{P}$ が $\mathcal{G}$ に対して成り立つ。
\end{enumerate}
\end{enumerate}
このとき、性質 $\mathcal{P}$ は、すべての連接層 $\mathcal{F}$（$X$ 上で
台が $Z_0$ に含まれるもの）に対して成り立つ。
\end{lemma}

\begin{proof}
まず、連接層 $\mathcal{F}$ で台が $Z_0$ に含まれるものが、連接部分層による
フィルトレーション
$$
0 = \mathcal{F}_0 \subset \mathcal{F}_1 \subset
\ldots \subset \mathcal{F}_m = \mathcal{F}
$$
をもち、各 $\mathcal{F}_i/\mathcal{F}_{i - 1}$ が性質
$\mathcal{P}$ をもてば、$\mathcal{F}$ もそうである。また、$\mathcal{F}$ が
性質 $\mathcal{P}$ をもち、
$\mathcal{F}_i/\mathcal{F}_{i - 1}$ のうち一つを除くすべてが性質
$\mathcal{P}$ をもてば、最後の一つもそうである。
これは仮定 (1) から従う。

\medskip\noindent
最初の応用として、台が $Z_0$ に真に含まれる任意の連接層が性質
$\mathcal{P}$ をもつと結論できる。実際、そのような層はフィルトレーション
（補題 \ref{coherent-lemma-coherent-filter} を参照）をもち、その部分商は (2) により
性質 $\mathcal{P}$ をもつ。

\medskip\noindent
$\mathcal{G}$ を (3) にあるものとする。補題
\ref{coherent-lemma-prepare-filter-irreducible} により、イデアル層 $\mathcal{I}$（$Z_0$
上のもの）、整数 $r \geq 1$、および短完全列
$$
0 \to
\left((Z_0 \to X)_*\mathcal{I}\right)^{\oplus r} \to
\mathcal{G} \to
\mathcal{Q} \to 0
$$
が存在し、$\mathcal{Q}$ の台は $Z_0$ に真に含まれる。(3)(c) により
$r = 1$ である。$\mathcal{Q}$ も性質 $\mathcal{P}$ をもつので、
$(Z_0 \to X)_*\mathcal{I}$ が性質 $\mathcal{P}$ をもつと結論できる。

\medskip\noindent
次に、$\mathcal{I}' \not = 0$ を $Z_0$ 上の別の準連接イデアル層とする。
このとき交わり $\mathcal{I}'' = \mathcal{I}' \cap \mathcal{I}$ を考えられ、
二つの短完全列
$$
0 \to
(Z_0 \to X)_*\mathcal{I}'' \to
(Z_0 \to X)_*\mathcal{I} \to
\mathcal{Q} \to 0
$$
および
$$
0 \to
(Z_0 \to X)_*\mathcal{I}'' \to
(Z_0 \to X)_*\mathcal{I}' \to
\mathcal{Q}' \to 0.
$$
を得る。連接層 $\mathcal{Q}$ と $\mathcal{Q}'$ の台は $Z_0$ に
真に含まれることに注意せよ。従って $\mathcal{Q}$ と
$\mathcal{Q}'$ は性質 $\mathcal{P}$ をもつ（上記参照）。ゆえに (1) を用いて、
$(Z_0 \to X)_*\mathcal{I}''$ と $(Z_0 \to X)_*\mathcal{I}'$ の双方も
性質 $\mathcal{P}$ をもつと結論できる。

\medskip\noindent
証明の最後の段階では、任意の連接層 $\mathcal{F}$（$X$ 上で台が
$Z_0$ に含まれるもの）がフィルトレーション（再び補題
\ref{coherent-lemma-coherent-filter} を参照）をもち、その部分商はすべて、今示したことに
より性質 $\mathcal{P}$ をもつことに注意すればよい。
\end{proof}

\begin{lemma}
\label{coherent-lemma-property}
$X$ を Noether スキームとする。
$\mathcal{P}$ を $X$ 上の連接層の性質とし、次を仮定する。
\begin{enumerate}
\item 連接層の任意の短完全列において、三つの層のうち二つが性質
$\mathcal{P}$ をもてば、残りの層も同じ性質をもつ。
\item 任意の整な閉部分スキーム $Z \subset X$ で生成点 $\xi$ をもつものに
対して、
次を満たす連接層 $\mathcal{G}$ が存在する。
\begin{enumerate}
\item $\text{Supp}(\mathcal{G}) = Z$,
\item $\mathcal{G}_\xi$ は $\mathfrak m_\xi$ で零化される、
\item $\dim_{\kappa(\xi)} \mathcal{G}_\xi = 1$、かつ
\item 性質 $\mathcal{P}$ が $\mathcal{G}$ に対して成り立つ。
\end{enumerate}
\end{enumerate}
このとき、性質 $\mathcal{P}$ は $X$ 上のすべての連接層に対して成り立つ。
\end{lemma}

\begin{proof}
補題 \ref{coherent-lemma-property-initial} により、任意の整な閉部分スキーム
$Z \subset X$ と任意の準連接イデアル層
$\mathcal{I} \subset \mathcal{O}_Z$ に対して、性質 $\mathcal{P}$ が
$(Z \to X)_*\mathcal{I}$ に対して成り立つことを示せば十分である。
これが成り立たなければ、$X$ は Noether なので、極小の整な閉部分スキーム
$Z_0 \subset X$ が存在し、ある準連接イデアル層について性質 $\mathcal{P}$ が
$(Z_0 \to X)_*\mathcal{I}_0$ に対して成り立たない。そのイデアル層は
$\mathcal{I}_0 \subset \mathcal{O}_{Z_0}$ である。一方、性質 $\mathcal{P}$ は
$(Z \to X)_*\mathcal{I}$ に対して、すべての整な閉部分スキーム
$Z \subset Z_0$、$Z \not = Z_0$ と準連接イデアル層
$\mathcal{I} \subset \mathcal{O}_Z$ について成り立つ。(2) により
$\mathcal{G}$ が $Z_0$ に対して存在するので、補題
\ref{coherent-lemma-property-irreducible} によれば、
これは起こりえない。
\end{proof}

\begin{lemma}
\label{coherent-lemma-property-irreducible-higher-rank-cohomological}
$X$ を Noether スキームとする。$Z_0 \subset X$ を生成点 $\xi$ をもつ
既約閉部分集合とする。$\mathcal{P}$ を $X$ 上の連接層の性質とし、
次を仮定する。
\begin{enumerate}
\item 連接層の任意の短完全列
$$
0 \to \mathcal{F}_1 \to \mathcal{F} \to \mathcal{F}_2 \to 0
$$
に対し、$\mathcal{F}_i$, $i = 1, 2$ が性質 $\mathcal{P}$ をもてば、
$\mathcal{F}$ もそうである。
\item 性質 $\mathcal{P}$ が $\mathcal{F}^{\oplus r}$ に対して、ある
$r \geq 1$ について成り立つならば、$\mathcal{F}$ に対しても成り立つ。
\item 任意の整な閉部分スキーム $Z \subset Z_0 \subset X$、
$Z \not = Z_0$、および任意の準連接イデアル層
$\mathcal{I} \subset \mathcal{O}_Z$ に対して、性質
$\mathcal{P}$ が $(Z \to X)_*\mathcal{I}$ に対して成り立つ。
\item 次を満たす連接層 $\mathcal{G}$ が存在する。
\begin{enumerate}
\item $\text{Supp}(\mathcal{G}) = Z_0$,
\item $\mathcal{G}_\xi$ は $\mathfrak m_\xi$ で零化される、かつ
\item 任意の準連接イデアル層 $\mathcal{J} \subset \mathcal{O}_X$ で
$\mathcal{J}_\xi = \mathcal{O}_{X, \xi}$ を満たすものに対し、準連接部分層
$\mathcal{G}' \subset \mathcal{J}\mathcal{G}$ で
$\mathcal{G}'_\xi = \mathcal{G}_\xi$ を満たし、性質 $\mathcal{P}$ が
$\mathcal{G}'$ に対して成り立つものが存在する。
\end{enumerate}
\end{enumerate}
このとき、性質 $\mathcal{P}$ は、すべての連接層 $\mathcal{F}$（$X$ 上で
台が $Z_0$ に含まれるもの）に対して成り立つ。
\end{lemma}

\begin{proof}
$\mathcal{F}$ が連接層で、連接部分層によるフィルトレーション
$$
0 = \mathcal{F}_0 \subset \mathcal{F}_1 \subset
\ldots \subset \mathcal{F}_m = \mathcal{F}
$$
をもち、各 $\mathcal{F}_i/\mathcal{F}_{i - 1}$ が性質
$\mathcal{P}$ をもてば、$\mathcal{F}$ もそうであることに
注意せよ。これは仮定 (1) から従う。

\medskip\noindent
最初の応用として、台が $Z_0$ に真に含まれる任意の連接層は性質
$\mathcal{P}$ をもつと結論できる。実際、そのような層はフィルトレーション
（補題 \ref{coherent-lemma-coherent-filter} を参照）をもち、その部分商は (3) により
性質 $\mathcal{P}$ をもつ。

\medskip\noindent
$i : Z_0 \to X$ で閉埋め込みを表す。(4) にある連接層
$\mathcal{G}$ を考える。補題 \ref{coherent-lemma-prepare-filter-irreducible} により、
イデアル層 $\mathcal{I}$（$Z_0$ 上のもの）と短完全列
$$
0 \to
i_*\mathcal{I}^{\oplus r} \to
\mathcal{G} \to
\mathcal{Q} \to 0
$$
が存在し、$\mathcal{Q}$ の台は $Z_0$ に真に含まれる。特に、
$r > 0$ かつ $\mathcal{I}$ は非零である。これは $\mathcal{G}$ の台が
$Z_0$ に等しいからである。$\mathcal{I}' \subset \mathcal{I}$ を、
$Z_0$ 上の、$\mathcal{I}$ に含まれる任意の非零準連接イデアル層とする。
すると短完全列
$$
0 \to
i_*(\mathcal{I}')^{\oplus r} \to
\mathcal{G} \to
\mathcal{Q}' \to 0
$$
も得られ、$\mathcal{Q}'$ の台は $Z_0$ に真に含まれる。
$\mathcal{J} \subset \mathcal{O}_X$ を、$\mathcal{Q}'$ の台を切り出す
準連接イデアル層とする（例えば、$\mathcal{Q}'$ の台上の被約誘導閉部分スキーム
構造に対応するイデアル）。このとき
$\mathcal{J}_\xi = \mathcal{O}_{X, \xi}$ である。補題
\ref{coherent-lemma-power-ideal-kills-sheaf} により、
$\mathcal{J}^n\mathcal{Q}' = 0$ となる $n$ が存在する。従って
$\mathcal{J}^n\mathcal{G} \subset i_*(\mathcal{I}')^{\oplus r}$ である。
補題の仮定 (4)(c) により、準連接部分層
$\mathcal{G}' \subset \mathcal{J}^n\mathcal{G}$ で、
$\mathcal{G}'_\xi = \mathcal{G}_\xi$ かつ性質 $\mathcal{P}$ をもつものが
存在する。従って短完全列
$$
0 \to \mathcal{G}' \to
i_*(\mathcal{I}')^{\oplus r} \to
\mathcal{Q}'' \to 0
$$
を得て、$\mathcal{Q}''$ の台は $Z_0$ に真に含まれる。ゆえに冒頭の注意と
補題の性質 (1) により、$i_*(\mathcal{I}')^{\oplus r}$ は
$\mathcal{P}$ を満たす。従って (2) により $i_*\mathcal{I}'$ は
$\mathcal{P}$ を満たす。最後に、任意の準連接イデアル層
$\mathcal{I}'' \subset \mathcal{O}_{Z_0}$ に対して
$\mathcal{I}' = \mathcal{I}'' \cap \mathcal{I}$ とおけば、短完全列
$$
0 \to
i_*(\mathcal{I}') \to
i_*(\mathcal{I}'') \to
\mathcal{Q}''' \to 0
$$
を得て、$\mathcal{Q}'''$ の台は $Z_0$ に真に含まれる。従って性質
$\mathcal{P}$ が $i_*\mathcal{I}''$ に対しても成り立つと結論できる。

\medskip\noindent
証明の最後の段階では、任意の連接層 $\mathcal{F}$（$X$ 上で台が
$Z_0$ に含まれるもの）がフィルトレーション（再び補題
\ref{coherent-lemma-coherent-filter} を参照）をもち、その部分商はすべて今示したことにより
性質 $\mathcal{P}$ をもつことに注意すればよい。
\end{proof}

\begin{lemma}
\label{coherent-lemma-property-higher-rank-cohomological}
$X$ を Noether スキームとする。
$\mathcal{P}$ を $X$ 上の連接層の性質とし、次を仮定する。
\begin{enumerate}
\item 連接層の任意の短完全列
$$
0 \to \mathcal{F}_1 \to \mathcal{F} \to \mathcal{F}_2 \to 0
$$
に対し、$\mathcal{F}_i$, $i = 1, 2$ が性質 $\mathcal{P}$ をもてば、
$\mathcal{F}$ もそうである。
\item 性質 $\mathcal{P}$ が $\mathcal{F}^{\oplus r}$ に対して、ある
$r \geq 1$ について成り立つならば、$\mathcal{F}$ に対しても成り立つ。
\item 任意の整な閉部分スキーム $Z \subset X$ で生成点 $\xi$ をもつものに
対して、
次を満たす連接層 $\mathcal{G}$ が存在する。
\begin{enumerate}
\item $\text{Supp}(\mathcal{G}) = Z$,
\item $\mathcal{G}_\xi$ は $\mathfrak m_\xi$ で零化される、かつ
\item 任意の準連接イデアル層 $\mathcal{J} \subset \mathcal{O}_X$ で
$\mathcal{J}_\xi = \mathcal{O}_{X, \xi}$ を満たすものに対し、準連接部分層
$\mathcal{G}' \subset \mathcal{J}\mathcal{G}$ で
$\mathcal{G}'_\xi = \mathcal{G}_\xi$ を満たし、性質 $\mathcal{P}$ が
$\mathcal{G}'$ に対して成り立つものが存在する。
\end{enumerate}
\end{enumerate}
このとき、性質 $\mathcal{P}$ は $X$ 上のすべての連接層に対して成り立つ。
\end{lemma}

\begin{proof}
補題 \ref{coherent-lemma-property-irreducible-higher-rank-cohomological} と
補題 \ref{coherent-lemma-property} が補題 \ref{coherent-lemma-property-irreducible} から従うのと
まったく同じ仕方で従う。
\end{proof}

\section{有限射とアフィン性}
\label{coherent-section-finite-affine}

\noindent
本節では、前節までの結果を用いて、Noether アフィンスキームの有限射による像が
アフィンであることを示す。後で、この結果がより一般に成り立つことを見る
（Limits, Lemma \ref{limits-lemma-affine} および
Proposition \ref{limits-proposition-affine} を参照）。

\begin{lemma}
\label{coherent-lemma-finite-morphism-Noetherian}
$f : Y \to X$ をスキームの射とする。
$f$ が有限かつ全射であり、$X$ が局所 Noether であると仮定する。
$Z \subset X$ を生成点 $\xi$ をもつ整な閉部分スキームとする。このとき、
ある連接層 $\mathcal{F}$ が $Y$ 上に存在し、$f_*\mathcal{F}$ の台が
$Z$ に等しく、$(f_*\mathcal{F})_\xi$ が $\mathfrak m_\xi$ で
零化される。
\end{lemma}

\begin{proof}
Morphisms, Lemma \ref{morphisms-lemma-finite-type-noetherian} により
$Y$ は局所 Noether である。$f$ は全射なのでファイバー $Y_\xi$ は空でない。
$\xi' \in Y$ を、その像が $\xi$ となるように選び、
$Z' = \overline{\{\xi'\}}$ とおく。
Schemes, Lemma \ref{schemes-lemma-reduced-closed-subscheme} により、
$Z' \subset Y$ を被約閉部分スキームとみなしてよい。
従って層 $\mathcal{F} = (Z' \to Y)_*\mathcal{O}_{Z'}$ は
$Y$ 上の連接層である（補題 \ref{coherent-lemma-finite-pushforward-coherent} を参照）。
可換図式
$$
\xymatrix{
Z' \ar[r]_{i'} \ar[d]_{f'} &
Y \ar[d]^f \\
Z \ar[r]^i &
X
}
$$
を考える。$f_*\mathcal{F} = i_*f'_*\mathcal{O}_{Z'}$ であることが分かる。
従って $f_*\mathcal{F}$ の $\xi$ における茎は
$f'_*\mathcal{O}_{Z'}$ の $\xi$ における茎である。$Z'$ は生成点
$\xi'$ をもつ整スキームなので、$\xi'$ は $Z'$ の点のうち
$\xi$ の上にある唯一の点であることに注意する。Algebra, Lemmas
\ref{algebra-lemma-finite-is-integral} および
\ref{algebra-lemma-integral-no-inclusion} を参照。
従って $f'_*\mathcal{O}_{Z'}$ の $\xi$ における茎は
$\mathcal{O}_{Z', \xi'} = \kappa(\xi')$ に等しい。特に
$f_*\mathcal{F}$ の $\xi$ における茎は零でない。このことと、
$f_*\mathcal{F}$ が $i_*f'_*(\text{something})$ の形であることから補題が従う。
\end{proof}

\begin{lemma}
\label{coherent-lemma-affine-morphism-projection-ideal}
$f : Y \to X$ をスキームの射とする。
$\mathcal{F}$ を $Y$ 上の準連接層とする。
$\mathcal{I}$ を $X$ 上の準連接イデアル層とする。
射 $f$ がアフィンならば
$\mathcal{I}f_*\mathcal{F} = f_*(f^{-1}\mathcal{I}\mathcal{F})$ である。
\end{lemma}

\begin{proof}
記法の意味は次のとおりである。$f^{-1}$ は完全関手なので、
$f^{-1}\mathcal{I}$ は $f^{-1}\mathcal{O}_X$ のイデアル層である。
写像 $f^\sharp : f^{-1}\mathcal{O}_X \to \mathcal{O}_Y$ を介して、
これは $\mathcal{F}$ に作用する。このとき
$f^{-1}\mathcal{I}\mathcal{F}$ は、$as$ の形の局所切断の和で生成される
部分層である。ここで $a$ は $f^{-1}\mathcal{I}$ の局所切断であり、
$s$ は $\mathcal{F}$ の局所切断である。これは準連接
$\mathcal{O}_Y$-部分加群かつ $\mathcal{F}$ の部分層である。実際、自然な写像
$f^*\mathcal{I} \otimes_{\mathcal{O}_Y} \mathcal{F} \to \mathcal{F}$ の像でもある。

\medskip\noindent
以上を踏まえれば証明は直接的である。実際、問題は局所的なので
$X$ はアフィンであると仮定してよい。$f$ はアフィンだから
$Y$ もアフィンである。従って
$Y = \Spec(B)$, $X = \Spec(A)$, $\mathcal{F} = \widetilde{M}$,
および $\mathcal{I} = \widetilde{I}$ と書ける。この場合、補題の主張は
$$
I(M_A) = ((IB)M)_A
$$
という等式に帰着する。ここで $M_A$ は、$A$-加群としてみた
$B$-加群 $M$ を表す。
\end{proof}

\begin{lemma}
\label{coherent-lemma-image-affine-finite-morphism-affine-Noetherian}
$f : Y \to X$ をスキームの射とする。次を仮定する。
\begin{enumerate}
\item $f$ は有限である、
\item $f$ は全射である、
\item $Y$ はアフィンである、かつ
\item $X$ は Noether である。
\end{enumerate}
このとき $X$ はアフィンである。
\end{lemma}

\begin{proof}
補題の仮定のもとで、任意の連接 $\mathcal{O}_X$-加群 $\mathcal{F}$ に対して
$H^1(X, \mathcal{F}) = 0$ であることを示す。特に、これにより
$H^1(X, \mathcal{I}) = 0$ が $\mathcal{O}_X$ の任意の準連接イデアル層
について成り立つ。このとき補題
\ref{coherent-lemma-quasi-compact-h1-zero-covering} または補題
\ref{coherent-lemma-quasi-separated-h1-zero-covering} のいずれからも、
$X$ がアフィンであることが従う。

\medskip\noindent
連接層の性質 $\mathcal{P}$ を次のように定める。$\mathcal{F}$ が
$X$ 上の連接層ならば、
$$
\mathcal{P}(\mathcal{F}) \Leftrightarrow H^1(X, \mathcal{F}) = 0.
$$
という規則で定める。補題 \ref{coherent-lemma-property-higher-rank-cohomological} を
適用する。そのため、この補題の (1), (2), (3) を $\mathcal{P}$ について
確かめなければならない。性質 (1) は層の短完全列に付随する
コホモロジー長完全列から従う。性質 (2) は $H^1(X, -)$ が加法的関手で
あることから従う。(3) を示すため、$Z \subset X$ を生成点 $\xi$ をもつ
整な閉部分スキームとする。補題 \ref{coherent-lemma-finite-morphism-Noetherian} により、
連接層 $\mathcal{F}$ が $Y$ 上に存在し、その $f_*\mathcal{F}$ の台が
$Z$ に等しく、$(f_*\mathcal{F})_\xi$ が $\mathfrak m_\xi$ で
零化される。$\mathcal{G} = f_*\mathcal{F}$ ととればよいと
主張する。補題 \ref{coherent-lemma-property-higher-rank-cohomological} の
(3)(c) だけを確かめればよい。そこで
$\mathcal{J} \subset \mathcal{O}_X$ を
$\mathcal{J}_\xi = \mathcal{O}_{X, \xi}$ を満たす準連接イデアル層とする。
有限射はアフィンなので、補題 \ref{coherent-lemma-affine-morphism-projection-ideal} により
$\mathcal{J}\mathcal{G} = f_*(f^{-1}\mathcal{J}\mathcal{F})$ である。また、
補題 \ref{coherent-lemma-affine-morphism-projection-ideal} の証明で指摘したように、
$f^{-1}\mathcal{J}\mathcal{F}$ は準連接 $\mathcal{O}_Y$-加群である。
$Y$ はアフィンなので
$H^1(Y, f^{-1}\mathcal{J}\mathcal{F}) = 0$ である。補題
\ref{coherent-lemma-quasi-coherent-affine-cohomology-zero} を参照。さらに $f$ は有限、
従ってアフィンなので、
$$
H^1(X, \mathcal{J}\mathcal{G}) =
H^1(X, f_*(f^{-1}\mathcal{J}\mathcal{F})) =
H^1(Y, f^{-1}\mathcal{J}\mathcal{F}) = 0
$$
が補題 \ref{coherent-lemma-relative-affine-cohomology} から従う。従って準連接部分層
$\mathcal{G}' = \mathcal{J}\mathcal{G}$ は $\mathcal{P}$ を満たす。
これは補題 \ref{coherent-lemma-property-higher-rank-cohomological} の性質 (3)(c) を
望みどおり確かめている。
\end{proof}

\section{Proj 上の連接層 I}
\label{coherent-section-coherent-proj}

\noindent
本節では $\text{Proj}(A)$ 上の連接層を論じる。ここで $A$ は
$A_1$ により $A_0$ 上で生成される Noether 次数付き環である。次節では
$A$ が次数 $1$ の元で生成されない場合に何が起こるかを論じる。
まず、Noether 環上の射影空間について、諸結果を一つにまとめた形で述べる。

\begin{lemma}
\label{coherent-lemma-coherent-projective}
$R$ を Noether 環とする。
$n \geq 0$ を整数とする。
任意の連接層 $\mathcal{F}$ が $\mathbf{P}^n_R$ 上にあるとき、次が成り立つ。
\begin{enumerate}
\item ある $r \geq 0$ と
$d_1, \ldots, d_r \in \mathbf{Z}$、および全射
$$
\bigoplus\nolimits_{j = 1, \ldots, r}
\mathcal{O}_{\mathbf{P}^n_R}(d_j)
\longrightarrow
\mathcal{F}.
$$
が存在する。
\item $H^i(\mathbf{P}^n_R, \mathcal{F}) = 0$ である。ただし
$0 \leq i \leq n$ の場合を除く。
\item 任意の $i$ に対して、コホモロジー群
$H^i(\mathbf{P}^n_R, \mathcal{F})$ は有限 $R$-加群である。
\item $i > 0$ ならば、$H^i(\mathbf{P}^n_R, \mathcal{F}(d)) = 0$ が
十分大きいすべての $d$ に対して成り立つ。
\item 任意の $k \in \mathbf{Z}$ に対して、次数付き
$R[T_0, \ldots, T_n]$-加群
$$
\bigoplus\nolimits_{d \geq k} H^0(\mathbf{P}^n_R, \mathcal{F}(d))
$$
は有限 $R[T_0, \ldots, T_n]$-加群である。
\end{enumerate}
\end{lemma}

\begin{proof}
$\mathcal{O}_{\mathbf{P}^n_R}(1)$ がスキーム
$\mathbf{P}^n_R$ 上の豊富な可逆層であることを用いる。これは定義から
直ちに従う。実際、$\mathbf{P}^n_R$ は標準アフィン開集合 $D_{+}(T_i)$ で
被覆される。従って Properties, Proposition
\ref{properties-proposition-characterize-ample} により、任意の有限型準連接
$\mathcal{O}_{\mathbf{P}^n_R}$-加群は
$\mathcal{O}_{\mathbf{P}^n_R}(1)$ のテンソル冪の有限直和の商である。
一方、補題 \ref{coherent-lemma-coherent-Noetherian} により、$R$ 上の射影空間では
連接層と有限型準連接層は同じものである。従って (1) を得る。

\medskip\noindent
射影 $n$-空間 $\mathbf{P}^n_R$ は $n + 1$ 個のアフィン開集合、すなわち
標準開集合 $D_{+}(T_i)$, $i = 0, \ldots, n$ で被覆される。Constructions,
Lemma \ref{constructions-lemma-standard-covering-projective-space} を参照。
従って任意の準連接層 $\mathcal{F}$ が $\mathbf{P}^n_R$ 上にあるとき、
$H^i(\mathbf{P}^n_R, \mathcal{F}) = 0$ である。ただし $i \geq n + 1$ とする。
補題 \ref{coherent-lemma-vanishing-nr-affines} を参照。ゆえに (2) が成り立つ。

\medskip\noindent
$\mathbf{P}^n_R$ 上のすべての連接層について、$i$ に関する降下帰納法により
(3) と (4) を同時に証明する。(2) により、$i \geq n + 1$ について結果が
成り立つことは明らかである。$i + 1$ について結果が既知であり、
$i$ について示したいと仮定する。（$i = 0$ ならば (4) は空虚である。）
$\mathcal{F}$ を $\mathbf{P}^n_R$ 上の連接層とする。(1) のような全射を選び、
その核を $\mathcal{G}$ と書くと、短完全列
$$
0 \to \mathcal{G} \to
\bigoplus\nolimits_{j = 1, \ldots, r}
\mathcal{O}_{\mathbf{P}^n_R}(d_j)
\to
\mathcal{F} \to 0
$$
を得る。補題 \ref{coherent-lemma-coherent-abelian-Noetherian} により
$\mathcal{G}$ は連接である。コホモロジー長完全列から完全列
$$
H^i(\mathbf{P}^n_R, \bigoplus\nolimits_{j = 1, \ldots, r}
\mathcal{O}_{\mathbf{P}^n_R}(d_j))
\to
H^i(\mathbf{P}^n_R, \mathcal{F})
\to
H^{i + 1}(\mathbf{P}^n_R, \mathcal{G}).
$$
を得る。帰納法の仮定により右辺の $R$-加群は有限であり、補題
\ref{coherent-lemma-cohomology-projective-space-over-ring} により左辺の
$R$-加群も有限である。$R$ は Noether なので、直ちに
$H^i(\mathbf{P}^n_R, \mathcal{F})$ は有限 $R$-加群であることが従う。
これで主張 (3) の帰納段階が証明された。
$\mathcal{O}_{\mathbf{P}^n_R}(d)$ は可逆であるから、
$\mathbf{P}^n_R$ 上のツイストは完全関手である（可逆層とのテンソル積で
得られるためである。Constructions, Definition
\ref{constructions-definition-twist} を参照）。これは、すべての
$d \in \mathbf{Z}$ に対して列
$$
0 \to \mathcal{G}(d) \to
\bigoplus\nolimits_{j = 1, \ldots, r}
\mathcal{O}_{\mathbf{P}^n_R}(d_j + d)
\to
\mathcal{F}(d) \to 0
$$
が短完全であることを意味する。これから得られるコホモロジー列は
$$
H^i(\mathbf{P}^n_R, \bigoplus\nolimits_{j = 1, \ldots, r}
\mathcal{O}_{\mathbf{P}^n_R}(d_j + d))
\to
H^i(\mathbf{P}^n_R, \mathcal{F}(d))
\to
H^{i + 1}(\mathbf{P}^n_R, \mathcal{G}(d)).
$$
である。帰納法の仮定により、右辺の加群は $d \gg 0$ に対して零である。
また、補題 \ref{coherent-lemma-cohomology-projective-space-over-ring} の計算により、
左辺の加群は $d \geq -\min\{d_j\}$ かつ $i \geq 1$ ならば零である。
従って主張 (4) の帰納段階も成り立つ。これで (3) と (4) の証明が完了する。

\medskip\noindent
(5) を証明するため、十分大きいすべての $d$ に対して写像
$$
H^0(\mathbf{P}^n_R, \bigoplus\nolimits_{j = 1, \ldots, r}
\mathcal{O}_{\mathbf{P}^n_R}(d_j + d))
\to
H^0(\mathbf{P}^n_R, \mathcal{F}(d))
$$
が全射であることに注意する。これは今証明した
$H^1(\mathbf{P}^n_R, \mathcal{G}(d))$ の消滅による。言い換えると、加群
$$
M_k
=
\bigoplus\nolimits_{d \geq k} H^0(\mathbf{P}^n_R, \mathcal{F}(d))
$$
は、$k$ が十分大きければ、対応する加群
$$
N_k
=
\bigoplus\nolimits_{d \geq k} H^0(\mathbf{P}^n_R,
\bigoplus\nolimits_{j = 1, \ldots, r}
\mathcal{O}_{\mathbf{P}^n_R}(d_j + d)
)
$$
の商である。$k$ が十分小さいとき（例えば、$k < -d_j$ がすべての
$j$ について成り立つとき）、
$$
N_k = \bigoplus\nolimits_{j = 1, \ldots, r}
R[T_0, \ldots, T_n](d_j)
$$
である。これは Section \ref{coherent-section-cohomology-projective-space} での計算による。
特に有限生成である。
任意の $k \in \mathbf{Z}$ をとる。
$k_{-} \ll k \ll k_{+}$ を選ぶ。図式
$$
\xymatrix{
N_{k_{-}} & N_{k_{+}} \ar[d] \ar[l] \\
M_k & M_{k_{+}} \ar[l]
}
$$
を考える。ここで縦の矢印は上の全射であり、横の矢印は明らかな包含写像である。
上で述べたことにより、$N_{k_{-}}$ は有限生成
$R[T_0, \ldots, T_n]$-加群である。従って $N_{k_{+}}$ は有限生成
$R[T_0, \ldots, T_n]$-加群の部分加群であり、環
$R[T_0, \ldots, T_n]$ が Noether なので、これも有限生成である。
縦の矢印は全射だから、$M_{k_{+}}$ は有限生成
$R[T_0, \ldots, T_n]$-加群である。商 $M_k/M_{k_{+}}$ は
$R$-加群として有限である。実際、これは有限個の有限 $R$-加群
$H^0(\mathbf{P}^n_R, \mathcal{F}(d))$, $k \leq d < k_{+}$ の直和である。
ここで (3) の $i = 0$ の場合を用いたことに注意する。
従って $M_k/M_{k_{+}}$ はなおさら有限
$R[T_0, \ldots, T_n]$-加群である。言い換えれば、$M_k$ を二つの有限
$R[T_0, \ldots, T_n]$-加群の間に挟んだことになり、証明が完了する。
\end{proof}

\begin{lemma}
\label{coherent-lemma-coherent-on-proj}
$A$ を次数付き環とし、$A_0$ は Noether であり、$A$ は
$A_1$ の有限個の元により $A_0$ 上で生成されるとする。
$X = \text{Proj}(A)$ とおく。このとき $X$ は Noether スキームである。
$\mathcal{F}$ を連接 $\mathcal{O}_X$-加群とする。
\begin{enumerate}
\item ある $r \geq 0$ と
$d_1, \ldots, d_r \in \mathbf{Z}$、および全射
$$
\bigoplus\nolimits_{j = 1, \ldots, r} \mathcal{O}_X(d_j)
\longrightarrow \mathcal{F}.
$$
が存在する。
\item 任意の $p$ に対して、コホモロジー群 $H^p(X, \mathcal{F})$ は有限
$A_0$-加群である。
\item $p > 0$ ならば、$H^p(X, \mathcal{F}(d)) = 0$ が
十分大きいすべての $d$ に対して成り立つ。
\item 任意の $k \in \mathbf{Z}$ に対して、次数付き $A$-加群
$$
\bigoplus\nolimits_{d \geq k} H^0(X, \mathcal{F}(d))
$$
は有限 $A$-加群である。
\end{enumerate}
\end{lemma}

\begin{proof}
仮定により、次数付き $A_0$-代数の全射
$$
A_0[T_0, \ldots, T_n] \longrightarrow A
$$
が存在する。ここで $\deg(T_j) = 1$, $j = 0, \ldots, n$ である。
Constructions, Lemma
\ref{constructions-lemma-surjective-graded-rings-generated-degree-1-map-proj}
により、これは閉埋め込み $i : X \to \mathbf{P}^n_{A_0}$ で
$i^*\mathcal{O}_{\mathbf{P}^n_{A_0}}(1) = \mathcal{O}_X(1)$ を
満たすものを定める。特に $X$ は Noether スキーム
$\mathbf{P}^n_{A_0}$ の閉部分スキームとして Noether である。
$\mathcal{F}$ に対する補題の諸結果は、補題 \ref{coherent-lemma-coherent-projective} の
対応する結果を、連接層 $i_*\mathcal{F}$（補題
\ref{coherent-lemma-i-star-equivalence}）に $\mathbf{P}^n_{A_0}$ 上で適用すれば
従うと主張する。
例えば、この補題により全射
$$
\bigoplus\nolimits_{j = 1, \ldots, r} \mathcal{O}_{\mathbf{P}^n_{A_0}}(d_j)
\longrightarrow i_*\mathcal{F}.
$$
が存在する。随伴により、これは写像
$\bigoplus_{j = 1, \ldots, r} \mathcal{O}_X(d_j) \longrightarrow \mathcal{F}$
に対応し、この写像も全射である。コホモロジーに関する主張は、補題
\ref{coherent-lemma-relative-affine-cohomology} による等式
$H^p(X, \mathcal{F}(d)) = H^p(\mathbf{P}^n_{A_0}, i_*\mathcal{F}(d))$
から従う。
\end{proof}

\begin{lemma}
\label{coherent-lemma-recover-tail-graded-module}
$A$ を次数付き環とし、$A_0$ は Noether であり、$A$ は
$A_1$ の有限個の元により $A_0$ 上で生成されるとする。$M$ を有限次数付き
$A$-加群とする。$X = \text{Proj}(A)$ とおき、$\widetilde{M}$ を
準連接 $\mathcal{O}_X$-加群、すなわち $X$ 上で $M$ に付随するものとする。
Constructions,
Lemma \ref{constructions-lemma-apply-modules} の写像
$$
M_n \longrightarrow \Gamma(X, \widetilde{M}(n))
$$
は、十分大きいすべての $n$ に対して同型である。
\end{lemma}

\begin{proof}
$M$ は有限 $A$-加群なので、$\widetilde{M}$ は有限型
$\mathcal{O}_X$-加群、すなわち連接 $\mathcal{O}_X$-加群である。
$N = \bigoplus_{n \in \mathbf{Z}} \Gamma(X, \widetilde{M}(n))$ とおく。
写像 $M \to N$ が次数付き $A$-加群の写像として、十分大きいすべての次数で
同型であることを示さなければならない。
Properties, Lemma \ref{properties-lemma-proj-quasi-coherent} により、
標準同型 $\widetilde{N} \to \widetilde{M}$ があり、それが誘導する写像
$N_n \to N_n = \Gamma(X, \widetilde{M}(n))$ は恒等写像である。
従って写像
$\widetilde{M} \to \widetilde{N} \to \widetilde{M}$
があり、すべての $n$ に対して図式
$$
\xymatrix{
M_n \ar[d] \ar[r] & N_n \ar[d] \ar@{=}[rd] \\
\Gamma(X, \widetilde{M}(n)) \ar[r] &
\Gamma(X, \widetilde{N}(n)) \ar[r]^{\cong} &
\Gamma(X, \widetilde{M}(n))
}
$$
は可換である。これは合成
$$
M_n \to \Gamma(X, \widetilde{M}(n))
\to \Gamma(X, \widetilde{N}(n)) \to \Gamma(X, \widetilde{M}(n))
$$
が標準写像 $M_n \to \Gamma(X, \widetilde{M}(n))$ に等しいことを意味する。
明らかに、このことから合成
$\widetilde{M} \to \widetilde{N} \to \widetilde{M}$ が恒等写像であると
分かる。従って $\widetilde{M} \to \widetilde{N}$ は同型である。
$K = \Ker(M \to N)$ および $Q = \Coker(M \to N)$ とおく。
関手 $M \mapsto \widetilde{M}$ は完全であった。Constructions,
Lemma \ref{constructions-lemma-proj-sheaves} を参照。
従って $\widetilde{K} = 0$ かつ $\widetilde{Q} = 0$ である。
$A$ は Noether 環であり、$M$ は有限生成 $A$-加群であり、$N$ は
次数付き $A$-加群であって、補題 \ref{coherent-lemma-coherent-on-proj} の最後の部分により
$N' = \bigoplus_{n \geq 0} N_n$ が有限生成であることを思い出す。
従って $K' = \bigoplus_{n \geq 0} K_n$ および
$Q' = \bigoplus_{n \geq 0} Q_n$ は有限 $A$-加群である。
$\widetilde{Q} = \widetilde{Q'}$ および
$\widetilde{K} = \widetilde{K'}$ であることに注意する。
従って証明を終えるには、有限 $A$-加群 $K$ で
$\widetilde{K} = 0$ を満たすものについて、非零斉次成分 $K_d$ で
$d \geq 0$ を満たすものが有限個しかないことを示せば十分である。
このため、$x_1, \ldots, x_r \in K$ を斉次生成元とし、その次数を
$d_1, \ldots, d_r$ とする。
$f_1, \ldots, f_n \in A_1$ を $A$ を $A_0$ 上生成する元とする。
各 $i$ と $j$ に対して、ある $n_{ij} \geq 0$ が存在して
$f_i^{n_{ij}} x_j = 0$ が $K_{d_j + n_{ij}}$ で成り立つ。そうでなければ
$x_i/f_i^{d_i} \in K_{(f_i)}$ は零でなく、すなわち $\widetilde{K}$ は
零でないからである。
このとき、$K_d$ は $d > \max_j(d_j + \sum_i n_{ij})$ に対して零である。
実際、$K_d$ の任意の元は項の和であり、各項は $f_i$ の単項式と
$x_j$ の一つとの積で、その全次数は $d$ である。
\end{proof}

\noindent
$A$ を次数付き環とし、$A_0$ は Noether であり、$A$ は
$A_1$ の有限個の元により $A_0$ 上で生成されるとする。
$A_+ = \bigoplus_{n > 0} A_n$ は無関係イデアルであった。
$M$ を次数付き $A$-加群とする。$M$ が $A_+$-冪捩れ加群であるとは、
すべての $x \in M$ に対して、ある $n \geq 1$ が存在し
$(A_+)^n x = 0$ となることをいう。More on Algebra, Definition
\ref{more-algebra-definition-f-power-torsion} を参照。$M$ が有限生成ならば、
これは $M_n = 0$ が $n \gg 0$ に対して成り立つことと同値である。
$A_+$-冪捩れ加群を単に捩れ加群と呼ぶこともある。次数付き
$A$-加群 $M$ について、$x \in M$ と $(A_+)^n x = 0$, $n > 0$ から
$x = 0$ が従うとき、捩れなしと呼ぶこともある。
$\text{Mod}_A$ で次数付き $A$-加群の圏を、
$\text{Mod}^{fg}_A$ で有限生成なものの充満部分圏を、
$\text{Mod}^{fg}_{A, torsion}$ で $M$ が $M_n = 0$ を
$n \gg 0$ に対して満たす加群の充満部分圏を表す。

\begin{proposition}
\label{coherent-proposition-coherent-modules-on-proj}
$A$ を次数付き環とし、$A_0$ は Noether であり、$A$ は
$A_1$ の有限個の元により $A_0$ 上で生成されるとする。
$X = \text{Proj}(A)$ とおく。関手 $M \mapsto \widetilde M$ は同値
$$
\text{Mod}^{fg}_A/\text{Mod}^{fg}_{A, torsion}
\longrightarrow
\textit{Coh}(\mathcal{O}_X)
$$
を誘導し、その擬逆関手は
$\mathcal{F} \longmapsto \bigoplus_{n \geq 0} \Gamma(X, \mathcal{F}(n))$
で与えられる。
\end{proposition}

\begin{proof}
部分圏 $\text{Mod}^{fg}_{A, torsion}$ は
$\text{Mod}^{fg}_A$ の Serre 部分圏である。Homology, Definition
\ref{homology-definition-serre-subcategory} を参照。これは上で与えた対象の
記述から明らかだが、More on Algebra, Lemma
\ref{more-algebra-lemma-I-power-torsion} からも従う。従って矢印の左側の
商圏は Homology, Lemma \ref{homology-lemma-serre-subcategory-is-kernel} で
定義されている。命題の関手を定義するには、関手
$M \mapsto \widetilde M$ が捩れ加群を $0$ に送ることを示せば十分である。
これは明らかである。実際、斉次な任意の $f \in A_+$ に対して加群
$M_f$ は零であり、従って $M_{(f)}$、すなわち $\widetilde M$ の
$D_+(f)$ 上での値も零である。

\medskip\noindent
補題 \ref{coherent-lemma-coherent-on-proj} により、提示した擬逆関手は意味をもつ。
実際、この補題は
$\mathcal{F} \longmapsto \bigoplus_{n \geq 0} \Gamma(X, \mathcal{F}(n))$
が関手 $\textit{Coh}(\mathcal{O}_X) \to \text{Mod}^{fg}_A$ であることを示す。
これを商関手
$\text{Mod}^{fg}_A \to \text{Mod}^{fg}_A/\text{Mod}^{fg}_{A, torsion}$
と合成できる。

\medskip\noindent
補題 \ref{coherent-lemma-recover-tail-graded-module} により、左から右へ進み、さらに左へ
戻る合成は恒等関手と同型である。

\medskip\noindent
最後に、$\mathcal{F}$ を連接 $\mathcal{O}_X$-加群とする。
$M = \bigoplus_{n \in \mathbf{Z}} \Gamma(X, \mathcal{F}(n))$ とおき、
これを次数付き $A$-加群とみなす。このとき本関手は $\mathcal{F}$ を
$M_{\geq 0} = \bigoplus_{n \geq 0} M_n$ に送る。
Properties, Lemma \ref{properties-lemma-proj-quasi-coherent} により、
標準写像 $\widetilde M \to \mathcal{F}$ は同型である。包含写像
$M_{\geq 0} \to M$ は同型
$\widetilde{M_{\geq 0}} \to \widetilde M$ を定めるので、右から左へ進み、
さらに右へ戻る合成も恒等関手と同型であると結論する。
\end{proof}

\section{Proj 上の連接層 II}
\label{coherent-section-coherent-proj-general}

\noindent
本節では $\text{Proj}(A)$ 上の連接層を論じる。ここで $A$ は Noether
次数付き環である。結果の大部分は、$A$ が次数 $1$ の元で生成される場合を
論じた前節の対応する結果から、少し工夫して導く。

\begin{lemma}
\label{coherent-lemma-coherent-on-proj-general}
$A$ を Noether 次数付き環とする。$X = \text{Proj}(A)$ とおく。このとき
$X$ は Noether スキームである。$\mathcal{F}$ を連接
$\mathcal{O}_X$-加群とする。
\begin{enumerate}
\item ある $r \geq 0$ と
$d_1, \ldots, d_r \in \mathbf{Z}$、および全射
$$
\bigoplus\nolimits_{j = 1, \ldots, r} \mathcal{O}_X(d_j)
\longrightarrow \mathcal{F}.
$$
が存在する。
\item 任意の $p$ に対して、コホモロジー群 $H^p(X, \mathcal{F})$ は有限
$A_0$-加群である。
\item $p > 0$ ならば、$H^p(X, \mathcal{F}(d)) = 0$ が
十分大きいすべての $d$ に対して成り立つ。
\item 任意の $k \in \mathbf{Z}$ に対して、次数付き $A$-加群
$$
\bigoplus\nolimits_{d \geq k} H^0(X, \mathcal{F}(d))
$$
は有限 $A$-加群である。
\end{enumerate}
\end{lemma}

\begin{proof}
主張を補題 \ref{coherent-lemma-coherent-on-proj} に帰着させて証明する。
Algebra, Lemmas \ref{algebra-lemma-graded-Noetherian} および
\ref{algebra-lemma-S-plus-generated} により、環 $A_0$ は Noether であり、
$A$ は $A_0$ 上で、正の次数をもつ斉次元
$f_1, \ldots, f_r$ の有限個により生成される。
$d$ を十分割り切れる整数とする。Algebra, Section
\ref{algebra-section-graded} の記法で $A' = A^{(d)}$ とおく。
このとき $A'$ は $A'_0 = A_0$ 上で次数 $1$ の元により生成される。
Algebra, Lemma \ref{algebra-lemma-uple-generated-degree-1} を参照。
従って補題 \ref{coherent-lemma-coherent-on-proj} を $X' = \text{Proj}(A')$ に
適用できる。

\medskip\noindent
Constructions, Lemma \ref{constructions-lemma-d-uple} により、
スキームの同型 $i : X \to X'$ と同型
$\mathcal{O}_X(nd) \to i^*\mathcal{O}_{X'}(n)$ が存在し、これらは写像
$A' \to A$、および写像
$A_n \to H^0(X, \mathcal{O}_X(n))$ と
$A'_n \to H^0(X', \mathcal{O}_{X'}(n))$ に整合する。
従って補題 \ref{coherent-lemma-coherent-on-proj} から $X$ が Noether であり、
(1) と (2) が成り立つことが従う。(3) と (4) を見るため、固定した任意の
$k$, $p$, $q$ に対して、上の整合性により
$$
\bigoplus\nolimits_{dn + q \geq k} H^p(X, \mathcal{F}(dn + q)) =
\bigoplus\nolimits_{dn + q \geq k} H^p(X', (i_*\mathcal{F}(q))(n)
$$
であることを用いる。$p > 0$ ならば、補題
\ref{coherent-lemma-coherent-on-proj} により、右辺は $k$ が $q$ に依存して
十分大きいとき消滅する。整数の $d$ を法とする合同類は有限個しかないので、
$\mathcal{F}$ が $X$ 上にある場合について (3) が成り立つ。$p = 0$ ならば、
補題 \ref{coherent-lemma-coherent-on-proj} により右辺は有限 $A'$-加群である。
合同類が有限個であることをもう一度用いると、
$\bigoplus_{n \geq k} H^0(X, \mathcal{F}(n))$ も有限 $A'$-加群である。
$A'$-加群構造は（上で述べた整合性により）$A$-加群構造から来るので、
これは $A$-加群としても有限である。
\end{proof}

\begin{lemma}
\label{coherent-lemma-recover-tail-graded-module-general}
$A$ を Noether 次数付き環とし、$d$ を $A$ の $A_0$ 上の生成元の
最小公倍数とする。$M$ を有限次数付き $A$-加群とする。
$X = \text{Proj}(A)$ とおき、$\widetilde{M}$ を準連接
$\mathcal{O}_X$-加群、すなわち $X$ 上で $M$ に付随するものとする。
$k \in \mathbf{Z}$ とする。
\begin{enumerate}
\item $N' = \bigoplus_{n \geq k} H^0(X, \widetilde{M(n)})$
は有限 $A$-加群である、
\item $N = \bigoplus_{n \geq k} H^0(X, \widetilde{M}(n))$
は有限 $A$-加群である、
\item 標準写像 $N \to N'$ が存在する、
\item $k$ が十分小さければ標準写像 $M \to N'$ が存在する、
\item 写像 $M_n \to N'_n$ は $n \gg 0$ に対して同型である、
\item $N_n \to N'_n$ は $d | n$ に対して同型である。
\end{enumerate}
\end{lemma}

\begin{proof}
(3) の写像 $N \to N'$ は、Constructions, Equation
(\ref{constructions-equation-multiply-more-generally}) から大域切断を
とることにより得られる。

\medskip\noindent
Constructions, Equation
(\ref{constructions-equation-global-sections-more-generally}) により、
次数付き $A$-加群の写像
$M \to \bigoplus_{n \in \mathbf{Z}} H^0(X, \widetilde{M(n)})$
が存在する。$M$ の生成元の次数が $\geq k$ ならば、像は部分加群
$N' \subset \bigoplus_{n \in \mathbf{Z}} H^0(X, \widetilde{M(n)})$
に含まれ、(4) の写像を得る。

\medskip\noindent
Algebra, Lemmas \ref{algebra-lemma-graded-Noetherian} および
\ref{algebra-lemma-S-plus-generated} により、環 $A_0$ は Noether であり、
$A$ は $A_0$ 上で、正の次数をもつ斉次元
$f_1, \ldots, f_r$ の有限個により生成される。
$d = \text{lcm}(\deg(f_i))$ とおく。すると、例えば Constructions,
Lemma \ref{constructions-lemma-where-invertible} により (6) が成り立つ。

\medskip\noindent
$M$ は有限 $A$-加群なので、$\widetilde{M}$ は有限型
$\mathcal{O}_X$-加群、すなわち連接 $\mathcal{O}_X$-加群である。
従って (2) は補題 \ref{coherent-lemma-coherent-on-proj-general} から従う。

\medskip\noindent
一つの工夫により (2) から (1) を導く。
$q \in \{0, \ldots, d - 1\}$ に対して
$$
{}^qN = \bigoplus\nolimits_{n + q \geq k} H^0(X, \widetilde{M(q)}(n))
$$
と書く。(2) により、これらは有限 $A$-加群である。Noether 環 $A$ は
$A^{(d)} = \bigoplus_{n \geq 0} A_{dn}$ 上有限である。実際、
もとの環は $f_i$ により $A^{(d)}$ 上生成され、$f_i^d \in A^{(d)}$ だからである。
従って ${}^qN$ は有限 $A^{(d)}$-加群である。
さらに $A^{(d)}$ は Noether である（Algebra, Lemma
\ref{algebra-lemma-dehomogenize-finite-type} から従う）。
従って $A^{(d)}$-部分加群
${}^qN^{(d)} = \bigoplus_{n \in \mathbf{Z}} {}^qN_{dn}$
は有限 $A^{(d)}$-加群である。同型
$\widetilde{M(dn + q)} = \widetilde{M(q)}(dn)$ を用いると
$$
N' =
\bigoplus\nolimits_{q \in \{0, \ldots, d - 1\}}
\bigoplus\nolimits_{dn + q \geq k}
H^0(X, \widetilde{M(q)}(dn)) =
\bigoplus\nolimits_{q \in \{0, \ldots, d - 1\}} {}^qN^{(d)}
$$
と書ける。従って $N'$ は $A^{(d)}$ 上有限であり、なおさら $A$ 上
有限である。ゆえに (1) は真である。

\medskip\noindent
$K$ を有限 $A$-加群とし、$\widetilde{K} = 0$ を満たすものとする。
$K_n = 0$ が $d|n$ かつ $n \gg 0$ に対して成り立つと主張する。
上と同様に論じると、$K^{(d)}$ は有限 $A^{(d)}$-加群である。
$x_1, \ldots, x_m \in K$ を $K^{(d)}$ の $A^{(d)}$ 上の斉次生成元とし、
その次数を $d_1, \ldots, d_m$ とする。ただし $d | d_j$ である。
各 $i$ と $j$ に対して、ある $n_{ij} \geq 0$ が存在して
$f_i^{n_{ij}} x_j = 0$ が $K_{d_j + n_{ij}}$ で成り立つ。そうでなければ
$x_j/f_i^{d_i/\deg(f_i)} \in K_{(f_i)}$ は零でなく、すなわち
$\widetilde{K}$ は零でないからである。ここでは
$\deg(f_i) | d | d_j$ がすべての $i, j$ に対して成り立つことを用いる。
$K_n$ は、$n$ が $d | n$ および
$n > \max_j (d_j + \sum_i n_{ij} \deg(f_i))$ を満たすとき零であると
結論する。実際、$K_n$ の任意の元は項の和であり、各項は
$f_i$ の単項式と $x_j$ の一つとの積で、その全次数は $n$ である。

\medskip\noindent
証明を終えるには、$M \to N'$ が十分大きいすべての次数で同型であることを
示さなければならない。写像 $N \to N'$ は同型
$\widetilde{N} \to \widetilde{N'}$ を誘導する。実際、アフィン開集合
$D_+(f_i) = D_+(f_i^d)$ 上では、対応する加群は
$N_{(f_i)} \cong N_{(f_i^d)} \cong N'_{(f_i^d)} \cong N'_{(f_i)}$
により同型である。これは性質 (6) による。
Properties, Lemma \ref{properties-lemma-proj-quasi-coherent} により、
標準同型 $\widetilde{N} \to \widetilde{M}$ がある。合成
$\widetilde{N} \to \widetilde{M} \to \widetilde{N'}$ は上の同型である
（証明は省略する。ヒント：標準アフィン開集合上で調べよ）。
従って写像 $M \to N'$ は同型 $\widetilde{M} \to \widetilde{N'}$ を誘導する。
$K = \Ker(M \to N')$ および $Q = \Coker(M \to N')$ とおく。
関手 $M \mapsto \widetilde{M}$ は完全であった。Constructions,
Lemma \ref{constructions-lemma-proj-sheaves} を参照。
従って $\widetilde{K} = 0$ かつ $\widetilde{Q} = 0$ である。
前段落の結果により、$K_n = 0$ および $Q_n = 0$ が
$d | n$ かつ $n \gg 0$ に対して成り立つ。ここでついに、$N'$ を
$N$ より優先して用いる利点が分かる。関手 $M \leadsto N'$ はシフトと整合する
（構成から直ちに分かる）。従って $M$ を $M(q)$ で置き換えて議論全体を
繰り返すと、$K_n = 0$ および $Q_n = 0$ が
$n \equiv q \bmod d$ かつ $n \gg 0$ に対して成り立つ。
$n$ を法とする合同類は有限個しかないので、証明が完了する。
\end{proof}

\noindent
$A$ を Noether 次数付き環とする。
$A_+ = \bigoplus_{n > 0} A_n$ は無関係イデアルであった。
Algebra, Lemmas \ref{algebra-lemma-graded-Noetherian} および
\ref{algebra-lemma-S-plus-generated} により、環 $A_0$ は Noether であり、
$A$ は $A_0$ 上で、正の次数をもつ斉次元
$f_1, \ldots, f_r$ の有限個により生成される。
$d = \text{lcm}(\deg(f_i))$ とおく。$M$ を次数付き $A$-加群とする。
この状況で、斉次元 $x \in M$ が {\it 無関係}
\footnote{これは非標準の記法である。} であるとは、
$$
(A_+ x)_{nd} = 0\text{ for all }n \gg 0
$$
が成り立つことをいう。$x \in M$ が斉次かつ無関係であり、$f \in A$ が
斉次ならば、$fx$ も無関係である。従って、無関係な元全体は次数付き部分加群
$M_{irrelevant} \subset M$ を生成する。$M$ は、$M$ のすべての斉次元が
無関係であるとき {\it 無関係} であるという。すなわち、この条件は
$M_{irrelevant} = M$ である。
$M$ が有限生成ならば、これは $M_{nd} = 0$ が $n \gg 0$ に対して
成り立つことと同値である。
$\text{Mod}_A$ で次数付き $A$-加群の圏を、
$\text{Mod}^{fg}_A$ で有限生成なものの充満部分圏を、
$\text{Mod}^{fg}_{A, irrelevant}$ で無関係加群の充満部分圏を表す。

\begin{proposition}
\label{coherent-proposition-coherent-modules-on-proj-general}
$A$ を Noether 次数付き環とする。$X = \text{Proj}(A)$ とおく。関手
$M \mapsto \widetilde M$ は同値
$$
\text{Mod}^{fg}_A/\text{Mod}^{fg}_{A, irrelevant}
\longrightarrow
\textit{Coh}(\mathcal{O}_X)
$$
を誘導し、その擬逆関手は
$\mathcal{F} \longmapsto \bigoplus_{n \geq 0} \Gamma(X, \mathcal{F}(n))$
で与えられる。
\end{proposition}

\begin{proof}
まず $A$ が次数 $1$ で生成される場合の証明を読むことを読者に強く勧める。
命題 \ref{coherent-proposition-coherent-modules-on-proj} を参照。
$f_1, \ldots, f_r \in A$ を正の次数をもつ斉次元で、$A$ を $A_0$ 上
生成するものとする。$d$ を、$d_i$ で表される $f_i$ の次数の最小公倍数とする。
$M$ を有限 $A$-加群とする。
$\widetilde{M}$ が零であることと、$M$ が（命題の主張より上で定義した）
無関係な次数付き $A$-加群であることが同値であると示そう。実際、
$x \in M$ を斉次元とする。
十分小さい $k \in \mathbf{Z}$ を選び、補題
\ref{coherent-lemma-recover-tail-graded-module-general} における
$N \to N'$ および $M \to N'$ をとる。さらに、十分大きい $l$ を選んで
$M_n \to N_n$ が $n \geq l$ に対して同型となるようにしてよい。
$\widetilde{M}$ が零ならば $N = 0$ である。従って、斉次な任意の
$f \in A_+$ で
$\deg(f) + \deg(x) = nd$ かつ $nd > l$ を満たすものに対して、$fx$ は零である。
実際、$N_{nd} \to N'_{nd}$ および $M_{nd} \to N'_{nd}$ は同型だからである。
ゆえに $x$ は無関係である。
逆に、$M$ は無関係であると仮定する。このとき $M_{nd}$ は
$n \gg 0$ に対して零である（命題の前の議論を参照）。
これは明らかに
$M_{(f_i)} = M_{(f_i^{d/\deg(f_i)})} = 0$ を含意し、従って構成により
$\widetilde{M} = 0$ である。

\medskip\noindent
従って部分圏 $\text{Mod}^{fg}_{A, irrelevant}$ は
$\text{Mod}^{fg}_A$ の Serre 部分圏であり、完全関手
$M \mapsto \widetilde M$ の核である。Homology, Lemma
\ref{homology-lemma-kernel-exact-functor} および Constructions, Lemma
\ref{constructions-lemma-proj-sheaves} を参照。従って矢印の左側の商圏は
Homology, Lemma \ref{homology-lemma-serre-subcategory-is-kernel} で
定義されている。命題の関手を定義するには、関手
$M \mapsto \widetilde M$ が無関係加群を $0$ に送ることを示せば十分であり、
これは上で示した。

\medskip\noindent
補題 \ref{coherent-lemma-coherent-on-proj-general} により、提示した擬逆関手は意味をもつ。
実際、この補題は
$\mathcal{F} \longmapsto \bigoplus_{n \geq 0} \Gamma(X, \mathcal{F}(n))$
が関手 $\textit{Coh}(\mathcal{O}_X) \to \text{Mod}^{fg}_A$ であることを示す。
これを商関手
$\text{Mod}^{fg}_A \to \text{Mod}^{fg}_A/\text{Mod}^{fg}_{A, irrelevant}$
と合成できる。

\medskip\noindent
補題 \ref{coherent-lemma-recover-tail-graded-module-general} により、左から右へ進み、
さらに左へ戻る合成は恒等関手と同型である。実際、$M$ を有限次数付き
$A$-加群とし、十分小さい $k \in \mathbf{Z}$ をとり、補題
\ref{coherent-lemma-recover-tail-graded-module-general} における
$N \to N'$ および $M \to N'$ をとる。このとき $M \to N'$ の核と余核は
有限個の次数でしか非零でなく、従って無関係である。さらに、写像
$N \to N'$ の核と余核は、十分大きく $d$ で割り切れるすべての次数で
零であるから、これらも無関係加群である。ゆえに $M \to N'$ と
$N \to N'$ は、ともに商圏で同型であり、望む結果を得る。

\medskip\noindent
最後に、$\mathcal{F}$ を連接 $\mathcal{O}_X$-加群とする。
$M = \bigoplus_{n \in \mathbf{Z}} \Gamma(X, \mathcal{F}(n))$ とおき、
これを次数付き $A$-加群とみなす。このとき本関手は $\mathcal{F}$ を
$M_{\geq 0} = \bigoplus_{n \geq 0} M_n$ に送る。
Properties, Lemma \ref{properties-lemma-proj-quasi-coherent} により、
標準写像 $\widetilde M \to \mathcal{F}$ は同型である。包含写像
$M_{\geq 0} \to M$ は同型
$\widetilde{M_{\geq 0}} \to \widetilde M$ を定めるので、右から左へ進み、
さらに右へ戻る合成も恒等関手と同型であると結論する。
\end{proof}

\section{射影的射に沿う高次順像}
\label{coherent-section-projective-pushforward}

\noindent
まず基底がアフィンである場合の結果を述べて証明し、次に射影的射に関する
いくつかの結果を導く。

\begin{lemma}
\label{coherent-lemma-coherent-proper-ample}
$R$ を Noether 環とする。$X \to \Spec(R)$ を固有射とする。
$\mathcal{L}$ を $X$ 上の豊富な可逆層とする。$\mathcal{F}$ を連接
$\mathcal{O}_X$-加群とする。
\begin{enumerate}
\item 次数付き環
$A = \bigoplus_{d \geq 0} H^0(X, \mathcal{L}^{\otimes d})$
は有限生成 $R$-代数である。
\item ある $r \geq 0$ と
$d_1, \ldots, d_r \in \mathbf{Z}$、および全射
$$
\bigoplus\nolimits_{j = 1, \ldots, r} \mathcal{L}^{\otimes d_j}
\longrightarrow \mathcal{F}.
$$
が存在する。
\item 任意の $p$ に対して、コホモロジー群 $H^p(X, \mathcal{F})$ は有限
$R$-加群である。
\item $p > 0$ ならば、
$H^p(X, \mathcal{F} \otimes_{\mathcal{O}_X} \mathcal{L}^{\otimes d}) = 0$
が十分大きいすべての $d$ に対して成り立つ。
\item 任意の $k \in \mathbf{Z}$ に対して、次数付き $A$-加群
$$
\bigoplus\nolimits_{d \geq k}
H^0(X, \mathcal{F} \otimes_{\mathcal{O}_X} \mathcal{L}^{\otimes d})
$$
は有限 $A$-加群である。
\end{enumerate}
\end{lemma}

\begin{proof}
Morphisms, Lemma
\ref{morphisms-lemma-finite-type-over-affine-ample-very-ample} により、
ある $d > 0$ と浸入 $i : X \to \mathbf{P}^n_R$ が存在して
$\mathcal{L}^{\otimes d} \cong i^*\mathcal{O}_{\mathbf{P}^n_R}(1)$ を満たす。
$X$ は $R$ 上固有なので、射 $i$ は閉埋め込みである
（Morphisms, Lemma \ref{morphisms-lemma-image-proper-scheme-closed}）。
従って $H^i(X, \mathcal{G}) = H^i(\mathbf{P}^n_R, i_*\mathcal{G})$ である。
ここで $\mathcal{G}$ は $X$ 上の任意の準連接層である
（補題 \ref{coherent-lemma-relative-affine-cohomology}、および閉埋め込みが
アフィンであることを参照。Morphisms, Lemma
\ref{morphisms-lemma-closed-immersion-affine}）。さらに
$\mathcal{G}$ が連接ならば $i_*\mathcal{G}$ も連接である
（補題 \ref{coherent-lemma-i-star-equivalence}）。以下ではこれらの事実を断りなく用いる。

\medskip\noindent
(1) の証明。$S = R[T_0, \ldots, T_n]$ とおき、
$\mathbf{P}^n_R = \text{Proj}(S)$ とする。$A$ は $S$-代数であることに
注意する（ただし環写像 $S \to A$ は次数付き環の準同型ではない。
$S_n$ が $A_{dn}$ に写るためである）。射影公式
（Cohomology, Lemma \ref{cohomology-lemma-projection-formula}）により、
$$
i_*(\mathcal{L}^{\otimes nd + q}) =
i_*(\mathcal{L}^{\otimes q})
\otimes_{\mathcal{O}_{\mathbf{P}^n_R}}
\mathcal{O}_{\mathbf{P}^n_R}(n)
$$
がすべての $n \in \mathbf{Z}$ に対して成り立つ。補題
\ref{coherent-lemma-coherent-projective} により、
$\bigoplus_{n \geq 0} A_{nd + q}$ は有限次数付き $S$-加群である。
$A = \bigoplus_{q \in \{0, \ldots, d - 1} \bigoplus_{n \geq 0} A_{nd + q}$
なので、$A$ は $S$-代数として有限である。従って (1) は真である。

\medskip\noindent
(2) の証明。これは Properties, Proposition
\ref{properties-proposition-characterize-ample} から従う。

\medskip\noindent
(3) の証明。補題 \ref{coherent-lemma-coherent-projective} を適用し、
$H^p(X, \mathcal{F}) = H^p(\mathbf{P}^n_R, i_*\mathcal{F})$ を用いる。

\medskip\noindent
(4) の証明。$p > 0$ を固定する。射影公式により
$$
i_*(\mathcal{F} \otimes_{\mathcal{O}_X} \mathcal{L}^{\otimes nd + q}) =
i_*(\mathcal{F} \otimes_{\mathcal{O}_X} \mathcal{L}^{\otimes q})
\otimes_{\mathcal{O}_{\mathbf{P}^n_R}}
\mathcal{O}_{\mathbf{P}^n_R}(n)
$$
がすべての $n \in \mathbf{Z}$ に対して成り立つ。補題
\ref{coherent-lemma-coherent-projective} により、
$H^p(X, \mathcal{F} \otimes \mathcal{L}^{nd + q}) = 0$ が
$n \gg 0$ に対して成り立つ。整数の $d$ を法とする合同類は有限個しか
ないので、(4) が証明された。

\medskip\noindent
(5) の証明。整数 $k$ を固定する。
$M = \bigoplus_{n \geq k} H^0(X, \mathcal{F} \otimes \mathcal{L}^{\otimes n})$
とおく。上と同様に論じると、$\bigoplus_{nd + q \geq k} M_{nd + q}$ は
有限次数付き $S$-加群である。
$M = \bigoplus_{q \in \{0, \ldots, d - 1\}}
\bigoplus_{nd + q \geq k} M_{nd + q}$
なので、$M$ は有限 $S$-加群である。$S$-加群構造は環写像
$S \to A$ を経由するので、$M$ は有限 $A$-加群でもある。
\end{proof}

\begin{lemma}
\label{coherent-lemma-kill-by-twisting}
$f : X \to S$ をスキームの射とする。
$\mathcal{F}$ を準連接 $\mathcal{O}_X$-加群とする。
$\mathcal{L}$ を $X$ 上の可逆層とする。次を仮定する。
\begin{enumerate}
\item $S$ は Noether である、
\item $f$ は固有である、
\item $\mathcal{F}$ は連接である、かつ
\item $\mathcal{L}$ は $X/S$ 上相対的に豊富である。
\end{enumerate}
このとき、ある $n_0$ が存在し、すべての $n \geq n_0$ に対して
$$
R^pf_*\left(\mathcal{F} \otimes_{\mathcal{O}_X} \mathcal{L}^{\otimes n}\right)
=
0
$$
がすべての $p > 0$ について成り立つ。
\end{lemma}

\begin{proof}
有限アフィン開被覆 $S = \bigcup V_j$ を選び、$X_j = f^{-1}(V_j)$ とおく。
有限個の各系 $(X_j \to V_j, \mathcal{L}|_{X_j}, \mathcal{F}|_{V_j})$
について問題を解けば結果が従うことは明らかである。従って $S$ は
アフィンであると仮定してよい。この場合、
$R^pf_*(\mathcal{F} \otimes \mathcal{L}^{\otimes n})$ の消滅は
$H^p(X, \mathcal{F} \otimes \mathcal{L}^{\otimes n})$ の消滅と同値である。
補題 \ref{coherent-lemma-quasi-coherence-higher-direct-images-application} を参照。
従って必要な消滅は補題 \ref{coherent-lemma-coherent-proper-ample} から従う
（Morphisms, Lemma
\ref{morphisms-lemma-finite-type-over-affine-ample-very-ample} により
$\mathcal{L}$ は $X$ 上豊富なので、この補題を適用できる）。
\end{proof}

\begin{lemma}
\label{coherent-lemma-locally-projective-pushforward}
$S$ を局所 Noether スキームとする。
$f : X \to S$ を局所射影的射とする。
$\mathcal{F}$ を連接 $\mathcal{O}_X$-加群とする。
このとき $R^if_*\mathcal{F}$ は連接 $\mathcal{O}_S$-加群であり、
これはすべての $i \geq 0$ に対して成り立つ。
\end{lemma}

\begin{proof}
まず、局所射影的射は固有であり
（Morphisms, Lemma \ref{morphisms-lemma-locally-projective-proper}）、
従って有限型であることに注意する。特に $X$ は局所 Noether であり
（Morphisms, Lemma \ref{morphisms-lemma-finite-type-noetherian}）、
従って主張は意味をもつ。さらに、補題
\ref{coherent-lemma-quasi-coherence-higher-direct-images} により、層
$R^pf_*\mathcal{F}$ は準連接である。

\medskip\noindent
以上を踏まえると、主張は $S$ 上局所的である（例えば Cohomology,
Lemma \ref{cohomology-lemma-localize-higher-direct-images} による）。
従って $S = \Spec(R)$ は Noether 環のスペクトルであり、$X$ は
$\mathbf{P}^n_R$（ある $n$ に対するもの）の閉部分スキームであると仮定してよい。
Morphisms, Lemma \ref{morphisms-lemma-characterize-locally-projective} を参照。
この場合、層 $R^pf_*\mathcal{F}$ は $R$-加群
$H^p(X, \mathcal{F})$ に付随する準連接層である。補題
\ref{coherent-lemma-quasi-coherence-higher-direct-images-application} を参照。
従って、$R$-加群 $H^p(X, \mathcal{F})$ が有限 $R$-加群であることを
示せば十分である（補題 \ref{coherent-lemma-coherent-Noetherian}）。
これは補題 \ref{coherent-lemma-coherent-proper-ample} から従う
（$\mathcal{O}_{\mathbf{P}^n_R}(1)$ の $X$ への制限は $X$ 上豊富だからである）。
\end{proof}

\section{豊富な可逆層とコホモロジー}
\label{coherent-section-ample-cohomology}

\noindent
ここでは、アフィン基底上の固有スキームにおける豊富性を、捻った後の
コホモロジーの消滅によって特徴づける判定法を与える。

\begin{lemma}
\label{coherent-lemma-vanshing-gives-ample}
\begin{reference}
\cite[III Proposition 2.6.1]{EGA}
\end{reference}
$R$ を Noether 環とする。$f : X \to \Spec(R)$ を固有射とする。
$\mathcal{L}$ を可逆 $\mathcal{O}_X$-加群とする。次は同値である。
\begin{enumerate}
\item $\mathcal{L}$ は $X$ 上豊富である（これは他の多くの条件と同値である。
Properties, Proposition \ref{properties-proposition-characterize-ample} および
Morphisms, Lemma
\ref{morphisms-lemma-finite-type-over-affine-ample-very-ample} を参照）、
\item 任意の連接 $\mathcal{O}_X$-加群 $\mathcal{F}$ に対して、
ある $n_0 \geq 0$ が存在し、
$H^p(X, \mathcal{F} \otimes \mathcal{L}^{\otimes n}) = 0$ がすべての
$n \geq n_0$ および $p > 0$ に対して成り立つ、かつ
\item 任意の準連接イデアル層
$\mathcal{I} \subset \mathcal{O}_X$ に対して、ある $n \geq 1$ が存在し、
$H^1(X, \mathcal{I} \otimes \mathcal{L}^{\otimes n}) = 0$ が成り立つ。
\end{enumerate}
\end{lemma}

\begin{proof}
(1) $\Rightarrow$ (2) は補題
\ref{coherent-lemma-coherent-proper-ample} から従う。
(2) $\Rightarrow$ (3) は自明である。
(3) $\Rightarrow$ (1) は補題
\ref{coherent-lemma-quasi-compact-h1-zero-invertible} である。
\end{proof}

\begin{lemma}
\label{coherent-lemma-surjective-finite-morphism-ample}
$R$ を Noether 環とする。$f : Y \to X$ を、$R$ 上固有なスキームの間の
射とする。$\mathcal{L}$ を可逆 $\mathcal{O}_X$-加群とする。
$f$ は有限かつ全射であると仮定する。このとき $\mathcal{L}$ が豊富である
ことと $f^*\mathcal{L}$ が豊富であることは同値である。
\end{lemma}

\begin{proof}
豊富な可逆層の準アフィン射による引き戻しは豊富である。Morphisms, Lemma
\ref{morphisms-lemma-pullback-ample-tensor-relatively-ample} を参照。
有限射は定義によりアフィンなので、これで一方の含意が証明された。

\medskip\noindent
$f^*\mathcal{L}$ が豊富であると仮定する。次の性質を $P$ とする。連接
$\mathcal{O}_X$-加群 $\mathcal{F}$ に対し、
ある $n_0$ が存在し、
$H^p(X, \mathcal{F} \otimes \mathcal{L}^{\otimes n}) = 0$ が
すべての $n \geq n_0$ および $p > 0$ に対して成り立つ。
この性質 $P$ が任意の連接 $\mathcal{O}_X$-加群 $\mathcal{F}$ に対して
成り立つことを
証明する。補題 \ref{coherent-lemma-vanshing-gives-ample} により、これは
$\mathcal{L}$ が豊富であることを意味する。
補題 \ref{coherent-lemma-property-higher-rank-cohomological} を適用する。
従って、その補題の (1)、(2)、(3) を $P$ について確かめなければならない。
性質 (1) は、層の短完全列に付随するコホモロジー長完全列と、可逆層との
テンソル積が完全関手であることから従う。性質 (2) は
$H^p(X, -)$ が加法的関手であることから従う。(3) を示すため、
$Z \subset X$ を、生成点が $\xi$ である整閉部分スキームとする。
$\mathcal{F}$ を $Y$ 上の連接層で、
$f_*\mathcal{F}$ の台が $Z$ に等しく、
$(f_*\mathcal{F})_\xi$ が $\mathfrak m_\xi$ により零化されるものとする。
補題 \ref{coherent-lemma-finite-morphism-Noetherian} を参照。
$\mathcal{G} = f_*\mathcal{F}$ ととればよいと主張する。
補題 \ref{coherent-lemma-property-higher-rank-cohomological} の (3)(c) だけを
確かめればよい。そこで、$\mathcal{J} \subset \mathcal{O}_X$ を
$\mathcal{J}_\xi = \mathcal{O}_{X, \xi}$ を満たす準連接イデアル層とする。
有限射はアフィンなので、補題
\ref{coherent-lemma-affine-morphism-projection-ideal} により
$\mathcal{J}\mathcal{G} = f_*(f^{-1}\mathcal{J}\mathcal{F})$ である。
また、補題 \ref{coherent-lemma-affine-morphism-projection-ideal} の証明で
指摘したように、層 $f^{-1}\mathcal{J}\mathcal{F}$ は連接
$\mathcal{O}_Y$-加群である。$\mathcal{L}$ は豊富なので、補題
\ref{coherent-lemma-vanshing-gives-ample} により、ある $n_0$ が存在して
$$
H^p(Y, f^{-1}\mathcal{J}\mathcal{F}
\otimes_{\mathcal{O}_Y} f^*\mathcal{L}^{\otimes n}) = 0,
$$
が $n \geq n_0$ および $p > 0$ に対して成り立つ。$f$ は有限、従って
アフィンなので、
\begin{align*}
H^p(X, \mathcal{J}\mathcal{G} \otimes_{\mathcal{O}_X}
\mathcal{L}^{\otimes n})
& =
H^p(X, f_*(f^{-1}\mathcal{J}\mathcal{F}) \otimes_{\mathcal{O}_X}
\mathcal{L}^{\otimes n}) \\
& =
H^p(X, f_*(f^{-1}\mathcal{J}\mathcal{F} \otimes_{\mathcal{O}_Y}
f^*\mathcal{L}^{\otimes n})) \\
& =
H^p(Y, f^{-1}\mathcal{J}\mathcal{F} \otimes_{\mathcal{O}_Y}
f^*\mathcal{L}^{\otimes n}) = 0
\end{align*}
となる。ここでは射影公式
（Cohomology, Lemma \ref{cohomology-lemma-projection-formula}）と
補題 \ref{coherent-lemma-relative-affine-cohomology} を用いた。
従って、準連接部分層 $\mathcal{G}' = \mathcal{J}\mathcal{G}$ は $P$ を
満たす。これで、求める補題
\ref{coherent-lemma-property-higher-rank-cohomological} の性質 (3)(c) が確認された。
\end{proof}

\noindent
コホモロジーは関手的である。特に、環付き空間 $X$、可逆
$\mathcal{O}_X$-加群 $\mathcal{L}$、および切断
$s \in \Gamma(X, \mathcal{L})$ が与えられると、写像
$$
H^p(X, \mathcal{F})
\longrightarrow
H^p(X, \mathcal{F} \otimes_{\mathcal{O}_X} \mathcal{L}), \quad
\xi \longmapsto s\xi
$$
を得る。これは写像
$\mathcal{F} \to \mathcal{F} \otimes_{\mathcal{O}_X} \mathcal{L}$、すなわち
$s$ による乗法が誘導する。次数付き環として
$\Gamma_*(X, \mathcal{L}) =
\bigoplus_{n \geq 0} \Gamma(X, \mathcal{L}^{\otimes n})$
とおく。Modules, Definition \ref{modules-definition-gamma-star} を参照。
$\mathcal{O}_X$-加群の層 $\mathcal{F}$ と整数 $p \geq 0$ に対して
$$
H^p_*(X, \mathcal{L}, \mathcal{F}) =
\bigoplus\nolimits_{n \in \mathbf{Z}} H^p(X,
\mathcal{F} \otimes_{\mathcal{O}_X} \mathcal{L}^{\otimes n})
$$
とおく。上で定めた乗法により、これは次数付き
$\Gamma_*(X, \mathcal{L})$-加群である。注意：
$H^p_*(X, \mathcal{L}, \mathcal{F})$ という記法は標準的でない。

\begin{lemma}
\label{coherent-lemma-invert-s-cohomology}
$X$ をスキームとする。$\mathcal{L}$ を $X$ 上の可逆層とする。
$s \in \Gamma(X, \mathcal{L})$ とする。$\mathcal{F}$ を準連接
$\mathcal{O}_X$-加群とする。$X$ が準コンパクトかつ準分離的ならば、
$\xi/s^n$ を $s^{-n}\xi$ に写す標準写像
$$
H^p_*(X, \mathcal{L}, \mathcal{F})_{(s)}
\longrightarrow
H^p(X_s, \mathcal{F})
$$
は同型である。
\end{lemma}

\begin{proof}
$p = 0$ の場合、これは Properties, Lemma
\ref{properties-lemma-invert-s-sections} である。
帰納原理（補題 \ref{coherent-lemma-induction-principle}）を用いて主張を証明する。
準コンパクト開集合 $U \subset X$ に対して、次の性質を $P(U)$ とする：
すべての $p \geq 0$ について写像
$$
H^p_*(U, \mathcal{L}, \mathcal{F})_{(s)}
\longrightarrow
H^p(U_s, \mathcal{F})
$$
が同型である。

\medskip\noindent
$U$ がアフィンならば、補題
\ref{coherent-lemma-quasi-coherent-affine-cohomology-zero} と
Properties, Lemma \ref{properties-lemma-affine-cap-s-open} により、
上の矢印の両辺は $p > 0$ に対して零であり、主張は成り立つ。
$P$ が $U$、$V$、および $U \cap V$ に対して成り立つならば、
Mayer--Vietoris 系列
（Cohomology, Lemma \ref{cohomology-lemma-mayer-vietoris}）を用いて、
長完全系列の写像
$$
\xymatrix{
H^{p - 1}_*(U \cap V, \mathcal{L}, \mathcal{F})_{(s)} \ar[r] \ar[d] &
H^p_*(U \cup V, \mathcal{L}, \mathcal{F})_{(s)} \ar[r] \ar[d] &
H^p_*(U, \mathcal{L}, \mathcal{F})_{(s)}
\oplus
H^p_*(V, \mathcal{L}, \mathcal{F})_{(s)} \ar[d] \\
H^{p - 1}(U_s \cap V_s, \mathcal{F}) \ar[r]&
H^p(U_s \cup V_s, \mathcal{F}) \ar[r] &
H^p(U_s, \mathcal{F})
\oplus
H^p(V_s, \mathcal{F})
}
$$
を得る（図には一部分だけを示した）。
$U_s \cap V_s = (U \cap V)_s$ かつ
$U_s \cup V_s = (U \cup V)_s$ であることに注意する。
従って左右の垂直写像は同型である（図に示していない、さらに右側の一つと
さらに左側の一つの写像も同型である）。五補題
（Homology, Lemma \ref{homology-lemma-five-lemma}）により
$P(U \cup V)$ が成り立つ。これで証明は完了する。
\end{proof}

\begin{lemma}
\label{coherent-lemma-section-affine-open-kills-classes}
$X$ をスキームとする。$\mathcal{L}$ を可逆 $\mathcal{O}_X$-加群とする。
$s \in \Gamma(X, \mathcal{L})$ を切断とする。次を仮定する。
\begin{enumerate}
\item $X$ は準コンパクトかつ準分離的である、かつ
\item $X_s$ はアフィンである。
\end{enumerate}
このとき、任意の準連接 $\mathcal{O}_X$-加群 $\mathcal{F}$、
任意の $p > 0$、およびすべての $\xi \in H^p(X, \mathcal{F})$ に対して、
ある $n \geq 0$ が存在して $s^n\xi = 0$ が
$H^p(X, \mathcal{F} \otimes_{\mathcal{O}_X} \mathcal{L}^{\otimes n})$
において成り立つ。
\end{lemma}

\begin{proof}
$H^p(X_s, \mathcal{G})$ は任意の準連接加群 $\mathcal{G}$ に対して
零であることを思い出そう。補題
\ref{coherent-lemma-quasi-coherent-affine-cohomology-zero} を参照。
従って、この補題は補題 \ref{coherent-lemma-invert-s-cohomology} から従う。
\end{proof}

\noindent
次の補題のより一般的な形については Limits, Lemma
\ref{limits-lemma-ample-on-reduction} を参照。

\begin{lemma}
\label{coherent-lemma-ample-on-reduction}
$i : Z \to X$ を、台位相空間の同相写像を誘導する Noether スキームの
閉埋め込みとする。$\mathcal{L}$ を $X$ 上の可逆層とする。
このとき、$i^*\mathcal{L}$ が $Z$ 上豊富であることと
$\mathcal{L}$ が $X$ 上豊富であることは同値である。
\end{lemma}

\begin{proof}
$\mathcal{L}$ が豊富ならば、例えば Morphisms, Lemma
\ref{morphisms-lemma-pullback-ample-tensor-relatively-ample} により
$i^*\mathcal{L}$ は豊富である。$i^*\mathcal{L}$ が豊富であると仮定する。
$\mathcal{L}$ が $X$ 上豊富であることを示さなければならない。
$\mathcal{I} \subset \mathcal{O}_X$ を、閉部分スキーム $Z$ を切り出す
連接イデアル層とする。集合論的に $i(Z) = X$ なので、補題
\ref{coherent-lemma-power-ideal-kills-sheaf} により、$\mathcal{I}^n = 0$ がある
$n$ に対して成り立つ。
次の閉部分スキームの列を考える。
$$
X = Z_n \supset Z_{n - 1} \supset Z_{n - 2} \supset \ldots \supset Z_1 = Z
$$
これらは
$0 = \mathcal{I}^n \subset \mathcal{I}^{n - 1} \subset \ldots \subset
\mathcal{I}$ により切り出される。このとき、閉埋め込み
$Z_{i - 1} \to Z_{i}$ のそれぞれは、平方零の連接イデアル層により
定義される。このようにして $\mathcal{I}^2 = 0$ の場合に帰着する。

\medskip\noindent
準連接 $\mathcal{O}_X$-加群の短完全列
$$
0 \to \mathcal{I} \to \mathcal{O}_X \to i_*\mathcal{O}_Z \to 0
$$
を考える。$\mathcal{L}^{\otimes n}$ とテンソル積をとると、短完全列
\begin{equation}
\label{coherent-equation-ses}
0 \to \mathcal{I} \otimes_{\mathcal{O}_X} \mathcal{L}^{\otimes n}
\to \mathcal{L}^{\otimes n} \to i_*i^*\mathcal{L}^{\otimes n} \to 0
\end{equation}
を得る。$\mathcal{I}^2 = 0$ なので、Morphisms, Lemma
\ref{morphisms-lemma-i-star-equivalence} を用いて
$\mathcal{I}$ を準連接 $\mathcal{O}_Z$-加群とみなすことができ、明らかな
記法の濫用のもとで
$\mathcal{I} \otimes_{\mathcal{O}_X} \mathcal{L}^{\otimes n} =
\mathcal{I} \otimes_{\mathcal{O}_Z} i^*\mathcal{L}^{\otimes n}$
となる。さらに、この層の $Z$ 上のコホモロジーは、その $X$ 上の
コホモロジーと標準的に同じである（$i$ は同相写像だからである）。

\medskip\noindent
$x \in X$ を点とし、対応する点を $z \in Z$ と書く。
$i^*\mathcal{L}$ は豊富なので、ある $n$ と切断
$s \in \Gamma(Z, i^*\mathcal{L}^{\otimes n})$ が存在して
$z \in Z_s$ かつ $Z_s$ はアフィンとなる。
$s$ を $\mathcal{L}^{\otimes n}$ の $X$ 上の切断へ持ち上げる障害は、
層の短完全列 (\ref{coherent-equation-ses}) から生じる境界元
$$
\xi = \partial s \in
H^1(X, \mathcal{I} \otimes_{\mathcal{O}_X} \mathcal{L}^{\otimes n}) =
H^1(Z, \mathcal{I} \otimes_{\mathcal{O}_Z} i^*\mathcal{L}^{\otimes n})
$$
である。$s$ を $s^{e + 1}$ で置き換えると、$\xi$ は
$\partial(s^{e + 1}) = (e + 1) s^e \xi$ で置き換えられ、この元は
$H^1(Z, \mathcal{I} \otimes_{\mathcal{O}_Z} i^*\mathcal{L}^{\otimes (e + 1)n})$
に属する。実際、
$$
0 \to
\bigoplus\nolimits_{m \geq 0}
\mathcal{I} \otimes_{\mathcal{O}_X} \mathcal{L}^{\otimes m} \to
\bigoplus\nolimits_{m \geq 0}
\mathcal{L}^{\otimes m} \to
\bigoplus\nolimits_{m \geq 0}
i_*i^*\mathcal{L}^{\otimes m} \to 0
$$
の境界写像は、Cohomology, Lemma
\ref{cohomology-lemma-boundary-derivation-over-cup-product} により導分である。
補題 \ref{coherent-lemma-section-affine-open-kills-classes} により、
$s^e \xi$ は十分大きい $e$ に対して零である。
従って $s$ をある冪で置き換えた後、$s$ は切断
$s' \in \Gamma(X, \mathcal{L}^{\otimes n})$ の像であると仮定できる。
このとき $X_{s'}$ は開部分スキームであり、
$Z_s \to X_{s'}$ は、$Z_s$ がアフィンである Noether スキームの
全射閉埋め込みである。従って補題
\ref{coherent-lemma-image-affine-finite-morphism-affine-Noetherian} により
$X_s$ はアフィンであり、$\mathcal{L}$ は豊富であると結論する。
\end{proof}

\noindent
次の補題のより一般的な形については Limits, Lemma
\ref{limits-lemma-thickening-quasi-affine} を参照。

\begin{lemma}
\label{coherent-lemma-thickening-quasi-affine}
$i : Z \to X$ を、台位相空間の同相写像を誘導する Noether スキームの
閉埋め込みとする。このとき $X$ が準アフィンであることと
$Z$ が準アフィンであることは同値である。
\end{lemma}

\begin{proof}
スキームが準アフィンであることと、その構造層が豊富であることは同値である。
Properties, Lemma \ref{properties-lemma-quasi-affine-O-ample} を参照。
従って $Z$ が準アフィンならば $\mathcal{O}_Z$ は豊富であり、補題
\ref{coherent-lemma-ample-on-reduction} により $\mathcal{O}_X$ は豊富であり、
従って $X$ は準アフィンである。逆向きは、いま与えた議論を逆に読むことで
証明できる。このことは初等的にも分かる。
\end{proof}

\begin{lemma}
\label{coherent-lemma-affine-in-presence-ample}
$X$ をスキームとする。$\mathcal{L}$ を豊富な可逆
$\mathcal{O}_X$-加群とする。$n_0$ を整数とする。
$H^p(X, \mathcal{L}^{\otimes -n}) = 0$ が $n \geq n_0$ および
$p > 0$ に対して成り立つならば、$X$ はアフィンである。
\end{lemma}

\begin{proof}
$H^p(X, \mathcal{F}) = 0$ が任意の準連接
$\mathcal{O}_X$-加群および $p > 0$ に対して成り立つと主張する。
Properties, Definition \ref{properties-definition-ample} により
$X$ は準コンパクトなので、この主張と補題
\ref{coherent-lemma-quasi-compact-h1-zero-covering} から証明が完了する。
Properties, Lemma \ref{properties-lemma-ample-separated} により
スキーム $X$ は分離的である。$X$ は $e + 1$ 個のアフィン開集合で
覆われるとする。このとき、補題 \ref{coherent-lemma-vanishing-nr-affines} により
$H^p(X, \mathcal{F}) = 0$ が $p > e$ に対して成り立つ。
従って、$p$ に関する降下帰納法で主張を証明してよい。
$\mathcal{F}$ を有限型準連接加群のフィルター余極限として書き
（Properties, Lemma
\ref{properties-lemma-quasi-coherent-colimit-finite-type}）、
Cohomology, Lemma
\ref{cohomology-lemma-quasi-separated-cohomology-colimit} を用いることで、
$\mathcal{F}$ は有限型であると仮定してよい。Properties, Proposition
\ref{properties-proposition-characterize-ample} により、ある $n > n_0$ を
選んで
$\mathcal{F} \otimes_{\mathcal{O}_X} \mathcal{L}^{\otimes n}$
を大域生成にできる。これは、ある集合 $I$ に対する短完全列
$$
0 \to \mathcal{F}' \to
\bigoplus\nolimits_{i \in I} \mathcal{L}^{\otimes -n}
\to \mathcal{F} \to 0
$$
が存在することを意味する（実際、$I$ は有限に選べる）。
帰納法の仮定により $H^{p + 1}(X, \mathcal{F}') = 0$ であり、
仮定と、すでに用いたコホモロジーと余極限の可換性を合わせると
$H^p(X, \bigoplus_{i \in I} \mathcal{L}^{\otimes -n}) = 0$
である。コホモロジー長完全列から
$H^p(X, \mathcal{F}) = 0$ が従い、主張が証明された。
\end{proof}

\begin{lemma}
\label{coherent-lemma-affine-if-quasi-affine}
$X$ を準アフィンスキームとする。
$H^p(X, \mathcal{O}_X) = 0$ が $p > 0$ に対して成り立つならば、
$X$ はアフィンである。
\end{lemma}

\begin{proof}
Properties, Lemma \ref{properties-lemma-quasi-affine-O-ample} により
$\mathcal{O}_X$ は豊富なので、これは補題
\ref{coherent-lemma-affine-in-presence-ample} から従う。
\end{proof}






\section{Chow の補題}
\label{coherent-section-chows-lemma}

\noindent
本節では Noether の場合の Chow の補題を証明する
（補題 \ref{coherent-lemma-chow-Noetherian}）。
Limits, Section \ref{limits-section-chows-lemma} では、非 Noether の場合の
いくつかの変種を証明する。

\begin{lemma}
\label{coherent-lemma-chow-Noetherian}
\begin{reference}
\cite[II Theorem 5.6.1(a)]{EGA}
\end{reference}
$S$ を Noether スキームとする。
$f : X \to S$ を有限型の分離的射とする。このとき、ある $n \geq 0$ と図式
$$
\xymatrix{
X \ar[rd] & X' \ar[d] \ar[l]^\pi \ar[r] & \mathbf{P}^n_S \ar[dl] \\
& S &
}
$$
が存在し、$X' \to \mathbf{P}^n_S$ は浸入、
$\pi : X' \to X$ は固有かつ全射である。さらに、稠密開部分スキーム
$U \subset X$ が存在して $\pi^{-1}(U) \to U$ が同型となるように
選ぶことができる。
\end{lemma}

\begin{proof}
以下の証明で現れるスキームはすべて Noether スキーム $S$ 上有限型であり、
従って Noether である
（Morphisms, Lemma \ref{morphisms-lemma-finite-type-noetherian} を参照）。
それらの間の射はすべて自動的に準コンパクト、局所有限型、かつ準分離的である。
Morphisms, Lemma \ref{morphisms-lemma-permanence-finite-type} および
Properties, Lemmas
\ref{properties-lemma-locally-Noetherian-quasi-separated}、
\ref{properties-lemma-morphism-Noetherian-schemes-quasi-compact} を参照。

\medskip\noindent
スキーム $X$ は有限個の既約成分しかもたない
（Properties, Lemma \ref{properties-lemma-Noetherian-irreducible-components}）。
$X = X_1 \cup \ldots \cup X_r$ を $X$ の既約成分への分解とする。
$\eta_i \in X_i$ を生成点とする。任意の点 $x \in X$ に対して、
アフィン開集合 $U_x \subset X$ が存在し、これは $x$ とすべての生成点
$\eta_i$ を含む。Properties, Lemma
\ref{properties-lemma-point-and-maximal-points-affine} を参照。
$X$ は準コンパクトなので、有限アフィン開被覆
$X = U_1 \cup \ldots \cup U_m$ で、各 $U_i$ が
$\eta_1, \ldots, \eta_r$ を含むものをとれる。特に、開集合
$U = U_1 \cap \ldots \cap U_m \subset X$ は稠密である。
以下で用いるのは、このことと、$U_i$ が $X$ を覆うアフィン開集合である
ことだけである。

\medskip\noindent
$X^* \subset X$ を $U \to X$ のスキーム論的閉包とする。Morphisms,
Definition \ref{morphisms-definition-scheme-theoretic-image} を参照。
$U_i^* = X^* \cap U_i$ とおく。$U_i^*$ は $U_i$ の閉部分スキームで
あることに注意する。従って $U_i^*$ はアフィンである。$U$ は $X$ で
稠密なので、射 $X^* \to X$ は全射閉埋め込みであり、$U$ 上では
同型である。従って $X$ を $X^*$ で、$U_i$ を $U_i^*$ で置き換え、
$U$ が $X$ でスキーム論的に稠密であると仮定してよい。Morphisms,
Definition \ref{morphisms-definition-scheme-theoretically-dense} を参照。

\medskip\noindent
Morphisms, Lemma \ref{morphisms-lemma-quasi-projective-finite-type-over-S}
により、浸入 $j_i : U_i \to \mathbf{P}_S^{n_i}$ を各 $i$ に対して
とれる。Morphisms, Lemma
\ref{morphisms-lemma-quasi-compact-immersion} により、閉部分スキーム
$Z_i \subset \mathbf{P}_S^{n_i}$ で、
$j_i : U_i \to Z_i$ がスキーム論的に稠密な開埋め込みとなるものを
とれる。$Z_i \to S$ は固有であることに注意する。Morphisms, Lemma
\ref{morphisms-lemma-locally-projective-proper} を参照。射
$$
j = (j_1|_U, \ldots, j_m|_U) : U \longrightarrow
\mathbf{P}_S^{n_1} \times_S \ldots \times_S \mathbf{P}_S^{n_m}.
$$
を考える。上で引用した補題により、閉部分スキーム $Z$ を
$\mathbf{P}_S^{n_1} \times_S \ldots \times_S \mathbf{P}_S^{n_m}$ の中に
とることができ、$j : U \to Z$ は開埋め込みとなり、かつ $U$ は
$Z$ でスキーム論的に稠密となるものをとれる。射 $Z \to S$ は固有で
ある。$i$ 番目の射影
$$
\text{pr}_i|_Z : Z \longrightarrow \mathbf{P}^{n_i}_S.
$$
を考える。この射は $Z_i$ を経由する（Morphisms, Lemma
\ref{morphisms-lemma-factor-factor} を参照）。誘導された射を
$p_i : Z \to Z_i$ と書く。これは固有射である。例えば Morphisms,
Lemma \ref{morphisms-lemma-image-proper-scheme-closed} を参照。この時点で、
$U \subset U_i \subset Z_i$ はスキーム論的に稠密な開埋め込みである。
さらに、$Z$ は「対角」射
$U \to Z_1 \times_S \ldots \times_S Z_m$ のスキーム論的像とみなせる。

\medskip\noindent
$V_i = p_i^{-1}(U_i)$ とおく。$p_i|_{V_i} : V_i \to U_i$ は固有である
ことに注意する。$X' = V_1 \cup \ldots \cup V_m$ とおく。構成により
$X'$ は
$\mathbf{P}^{n_1}_S \times_S \ldots \times_S \mathbf{P}^{n_m}_S$
への浸入をもつ。従って Segre 埋め込み
（Constructions, Lemma \ref{constructions-lemma-segre-embedding} を参照）
により、$X'$ は $S$ 上の射影空間への浸入をもつ。

\medskip\noindent
射 $p_i|_{V_i}: V_i \to U_i$ は貼り合わさって射 $X' \to X$ を与えると
主張する。実際、$p_i|_U$ が $U$ から $U$ への恒等写像であることは
明らかである。$U \subset X'$ は構成によりスキーム論的に稠密なので、
開部分スキーム $V_i \cap V_j$ においてもスキーム論的に稠密である。
従って、等式 $p_i|_{V_i \cap V_j} = p_j|_{V_i \cap V_j}$ が、分離的
$S$-スキーム $X$ への射として成り立つ。Morphisms, Lemma
\ref{morphisms-lemma-equality-of-morphisms} を参照。得られる射を
$\pi : X' \to X$ と書く。

\medskip\noindent
$\pi^{-1}(U_i) = V_i$ であると主張する。
$\pi|_{V_i} = p_i|_{V_i}$ なので、$V_i \subset \pi^{-1}(U_i)$ である。
図式
$$
\xymatrix{
V_i \ar[r] \ar[rd]_{p_i|_{V_i}} & \pi^{-1}(U_i) \ar[d] \\
& U_i
}
$$
を考える。$V_i \to U_i$ は固有なので、水平矢印の像は閉である。
Morphisms, Lemma \ref{morphisms-lemma-image-proper-scheme-closed} を参照。
$V_i \subset \pi^{-1}(U_i)$ は（$U$ を含むので）スキーム論的に稠密で
あるから、主張どおり $V_i = \pi^{-1}(U_i)$ である。

\medskip\noindent
これにより、$\pi^{-1}(U_i) \to U_i$ は固有射
$p_i|_{V_i} : V_i \to U_i$ と同一視される。従って $X$ は有限アフィン被覆
$X = \bigcup U_i$ をもち、この被覆の各要素への $\pi$ の制限は固有で
ある。従って Morphisms, Lemma
\ref{morphisms-lemma-proper-local-on-the-base} により $\pi$ は固有である。

\medskip\noindent
最後に $\pi^{-1}(U) = U$ を示さなければならない。このためには、上と同じ
ように図式
$$
\xymatrix{
U \ar[r] \ar[rd] & \pi^{-1}(U) \ar[d] \\
& U
}
$$
を用い、$\text{id}_U : U \to U$ が固有であること、および $U$ が
$\pi^{-1}(U)$ でスキーム論的に稠密であることを用いればよい。
\end{proof}

\begin{remark}
\label{coherent-remark-chow-Noetherian}
Chow の補題 \ref{coherent-lemma-chow-Noetherian} の状況では、次が成り立つ。
\begin{enumerate}
\item 射 $\pi$ は実際 H-射影的である（従って射影的である。Morphisms,
Lemma \ref{morphisms-lemma-H-projective} を参照）。実際、射
$X' \to \mathbf{P}^n_S \times_S X = \mathbf{P}^n_X$ は閉埋め込みである
（$\pi$ が固有であることを用いる。Morphisms, Lemma
\ref{morphisms-lemma-image-proper-scheme-closed} を参照）。
\item $\pi^{-1}(U)$ が $X'$ でスキーム論的に稠密であると仮定してよい。
実際、$X'$ を $\pi^{-1}(U)$ のスキーム論的閉包で置き換えればよい。
この場合、$U$ を $X'$ のスキーム論的に稠密な開部分スキームとみなせる。
Morphisms, Section \ref{morphisms-section-scheme-theoretic-image} を参照。
\item $X$ が被約ならば、$X'$ も被約に選べる。これは (2) から明らかである。
\end{enumerate}
\end{remark}






\section{連接層の高次順像}
\label{coherent-section-proper-pushforward}

\noindent
本節では、固有射による連接層の高次順像が連接であるという基本的事実を
証明する。

\begin{proposition}
\label{coherent-proposition-proper-pushforward-coherent}
\begin{reference}
\cite[III Theorem 3.2.1]{EGA}
\end{reference}
$S$ を局所 Noether スキームとする。
$f : X \to S$ を固有射とする。
$\mathcal{F}$ を連接 $\mathcal{O}_X$-加群とする。このとき
$R^if_*\mathcal{F}$ は連接 $\mathcal{O}_S$-加群であり、これはすべての
$i \geq 0$ に対して成り立つ。
\end{proposition}

\begin{proof}
問題は $S$ 上局所的なので、$S$ は Noether スキームであると仮定してよい。
固有射は有限型だから、この場合 $X$ も Noether スキームである。
性質 $\mathcal{P}$ を $X$ 上の連接層について、次の規則により定める。
$$
\mathcal{P}(\mathcal{F}) \Leftrightarrow
R^pf_*\mathcal{F}\text{ is coherent for all }p \geq 0
$$
補題 \ref{coherent-lemma-property} の結果を用いて、$\mathcal{P}$ が $X$ 上の
すべての連接層について成り立つことを証明する。

\medskip\noindent
$X$ 上の連接層の短完全列
$$
0 \to \mathcal{F}_1 \to \mathcal{F}_2 \to \mathcal{F}_3 \to 0
$$
をとる。高次順像の長完全列
$$
R^{p - 1}f_*\mathcal{F}_3 \to
R^pf_*\mathcal{F}_1 \to
R^pf_*\mathcal{F}_2 \to
R^pf_*\mathcal{F}_3 \to
R^{p + 1}f_*\mathcal{F}_1
$$
を考える。層 $\mathcal{F}_i$ のうち二つが性質 $\mathcal{P}$ をもつならば、
第三の層の高次順像は、この完全複体において二つの連接層の間に挟まれている。
従って、補題 \ref{coherent-lemma-coherent-abelian-Noetherian} および
\ref{coherent-lemma-coherent-Noetherian-quasi-coherent-sub-quotient} により、これらの
高次順像も連接である。従って性質 $\mathcal{P}$ は第三の層についても成り立つ。

\medskip\noindent
$Z \subset X$ を整閉部分スキームとする。連接層 $\mathcal{F}$ を $X$ 上で、
台が $Z$ に含まれ、生成点 $\xi$（$Z$ の生成点）における茎が
$1$-次元の $\kappa(\xi)$-ベクトル空間であり、性質 $\mathcal{P}$ が
$\mathcal{F}$ について成り立つものとして見つけなければならない。
射 $g = f|_Z : Z \to S$ を考える。これは $f$ の制限である。
連接層 $\mathcal{G}$ を $Z$ 上で、次を満たすものとして見つけられると
仮定する。
(a) $\mathcal{G}_\xi$ は $1$-次元の $\kappa(\xi)$-ベクトル空間である、
(b) $R^pg_*\mathcal{G} = 0$ が $p > 0$ に対して成り立つ、かつ
(c) $g_*\mathcal{G}$ は連接である。このとき
$\mathcal{F} = (Z \to X)_*\mathcal{G}$ を考えられる。$Z \to X$ は閉埋め込み
なので、$(Z \to X)_*\mathcal{G}$ は $X$ 上連接であり、
$R^p(Z \to X)_*\mathcal{G} = 0$ が $p > 0$ に対して成り立つ
（補題 \ref{coherent-lemma-finite-pushforward-coherent}）。
従って、相対 Leray スペクトル系列
（Cohomology, Lemma \ref{cohomology-lemma-relative-Leray}）により、
$R^pf_*\mathcal{F} = R^pg_*\mathcal{G} = 0$ が $p > 0$ に対して成り立ち、
$f_*\mathcal{F} = g_*\mathcal{G}$ は連接となる。最後に
$\mathcal{F}_\xi = ((Z \to X)_*\mathcal{G})_\xi = \mathcal{G}_\xi$ であり、
$\xi$ における茎についての条件が確認される。従ってすべては、性質
(a)、(b)、(c) をもつ連接層 $\mathcal{G}$ を $Z$ 上に見つけることに帰着する。

\medskip\noindent
射 $Z \to S$ に Chow の補題 \ref{coherent-lemma-chow-Noetherian} を適用できる。
従って、Chow の補題の主張にある図式
$$
\xymatrix{
Z \ar[rd]_g & Z' \ar[d]^-{g'} \ar[l]^\pi \ar[r]_i & \mathbf{P}^m_S \ar[dl] \\
& S &
}
$$
を得る。また、$U \subset Z$ を、$\pi^{-1}(U) \to U$ が同型となる稠密
開部分スキームとする。注意 \ref{coherent-remark-chow-Noetherian} の議論により、
$i' = (i, \pi) : Z' \to \mathbf{P}^m_Z$ は閉埋め込みである。従って
$$
\mathcal{L} = i^*\mathcal{O}_{\mathbf{P}^m_S}(1) \cong
(i')^*\mathcal{O}_{\mathbf{P}^m_Z}(1)
$$
は $g'$-相対的に豊富かつ $\pi$-相対的に豊富である（例えば Morphisms,
Lemma \ref{morphisms-lemma-characterize-ample-on-finite-type} による）。
従って補題 \ref{coherent-lemma-kill-by-twisting} により、ある $n \geq 0$ が存在して、
$R^p\pi_*\mathcal{L}^{\otimes n} = 0$ がすべての $p > 0$ に対して成り立ち、
$R^p(g')_*\mathcal{L}^{\otimes n} = 0$ がすべての $p > 0$ に対して成り立つ。
$\mathcal{G} = \pi_*\mathcal{L}^{\otimes n}$ とおく。性質 (a) は、
$\pi_*\mathcal{L}^{\otimes n}|_U$ が可逆層であることから従う
（$\pi^{-1}(U) \to U$ は同型だからである）。性質 (b) と (c) は、相対
Leray スペクトル系列
（Cohomology, Lemma \ref{cohomology-lemma-relative-Leray}）により
$$
E_2^{p, q} = R^pg_* R^q\pi_*\mathcal{L}^{\otimes n}
\Rightarrow
R^{p + q}(g')_*\mathcal{L}^{\otimes n}
$$
となることから従う。実際、$n$ の選び方により、$E_2^{p, q}$ の非零項は
$q = 0$ のものだけであり、$R^{p + q}(g')_*\mathcal{L}^{\otimes n}$ の
非零項は $p = q = 0$ のものだけである。これは、
$R^pg_*\mathcal{G} = 0$ が $p > 0$ に対して成り立ち、かつ
$g_*\mathcal{G} = (g')_*\mathcal{L}^{\otimes n}$ となることを含意する。
最後に、直前の補題 \ref{coherent-lemma-locally-projective-pushforward} を適用すると、
$g_*\mathcal{G} = (g')_*\mathcal{L}^{\otimes n}$ は求めるとおり連接である。
\end{proof}

\begin{lemma}
\label{coherent-lemma-proper-over-affine-cohomology-finite}
$S = \Spec(A)$ とし、$A$ を Noether 環とする。
$f : X \to S$ を固有射とする。
$\mathcal{F}$ を連接 $\mathcal{O}_X$-加群とする。このとき
$H^i(X, \mathcal{F})$ は有限 $A$-加群であり、これはすべての
$i \geq 0$ に対して成り立つ。
\end{lemma}

\begin{proof}
これは命題 \ref{coherent-proposition-proper-pushforward-coherent} のアフィンの場合に
ほかならない。実際、補題 \ref{coherent-lemma-quasi-coherence-higher-direct-images}
および \ref{coherent-lemma-quasi-coherence-higher-direct-images-application} により、
$R^if_*\mathcal{F}$ は $A$-加群 $H^i(X, \mathcal{F})$ に付随する準連接層
であることが分かる。また、補題 \ref{coherent-lemma-coherent-Noetherian} により、
これが連接層であることと $H^i(X, \mathcal{F})$ が有限型 $A$-加群である
ことは同値である。
\end{proof}

\begin{lemma}
\label{coherent-lemma-graded-finiteness}
$A$ を Noether 環とする。
$B$ を有限生成次数付き $A$-代数とする。
$f : X \to \Spec(A)$ を固有射とする。
$\mathcal{B} = f^*\widetilde B$ とおく。
$\mathcal{F}$ を有限型の準連接次数付き $\mathcal{B}$-加群とする。
\begin{enumerate}
\item すべての $p \geq 0$ に対して、次数付き $B$-加群
$H^p(X, \mathcal{F})$ は有限 $B$-加群である。
\item $\mathcal{L}$ が豊富な可逆 $\mathcal{O}_X$-加群ならば、整数
$d_0$ が存在して、$H^p(X, \mathcal{F} \otimes \mathcal{L}^{\otimes d}) = 0$
がすべての $p > 0$ および $d \geq d_0$ に対して成り立つ。
\end{enumerate}
\end{lemma}

\begin{proof}
これを証明するため、ファイバー積の図式
$$
\xymatrix{
X' = \Spec(B) \times_{\Spec(A)} X
\ar[r]_-\pi \ar[d]_{f'} &
X \ar[d]^f \\
\Spec(B) \ar[r] &
\Spec(A)
}
$$
を考える。$f'$ は固有射であることに注意する。Morphisms, Lemma
\ref{morphisms-lemma-base-change-proper} を参照。また、$B$ は有限生成
$A$-代数であり、従って Noether である
（Algebra, Lemma \ref{algebra-lemma-Noetherian-permanence}）。これは
$X'$ が Noether スキームであることを含意する
（Morphisms, Lemma \ref{morphisms-lemma-finite-type-noetherian}）。
$X'$ は、Constructions, Lemma \ref{constructions-lemma-spec-properties}
により、準連接 $\mathcal{O}_X$-代数 $\mathcal{B}$ の相対スペクトルである
ことに注意する。$\mathcal{F}$ は準連接 $\mathcal{B}$-加群なので、一意な
準連接 $\mathcal{O}_{X'}$-加群 $\mathcal{F}'$ で
$\pi_*\mathcal{F}' = \mathcal{F}$ を満たすものが存在する。Morphisms,
Lemma \ref{morphisms-lemma-affine-equivalence-modules} を参照。
$\mathcal{F}$ は有限型 $\mathcal{B}$-加群なので、$\mathcal{F}'$ は有限型
$\mathcal{O}_{X'}$-加群である（詳細は省略する）。言い換えると、
$\mathcal{F}'$ は連接 $\mathcal{O}_{X'}$-加群である
（補題 \ref{coherent-lemma-coherent-Noetherian}）。射 $\pi : X' \to X$ はアフィン
なので、
$$
H^p(X, \mathcal{F}) = H^p(X', \mathcal{F}')
$$
が補題 \ref{coherent-lemma-relative-affine-cohomology} により成り立つ。従って (1) は
補題 \ref{coherent-lemma-proper-over-affine-cohomology-finite} から従う。(2) のような
$\mathcal{L}$ が与えられたとき、$\mathcal{L}' = \pi^*\mathcal{L}$ とおく。
$\mathcal{L}'$ は $X'$ 上豊富であることに注意する。Morphisms, Lemma
\ref{morphisms-lemma-pullback-ample-tensor-relatively-ample} を参照。
射影公式（Cohomology, Lemma \ref{cohomology-lemma-projection-formula}）により
$\pi_*(\mathcal{F}' \otimes \mathcal{L}') = \mathcal{F} \otimes \mathcal{L}$
である。従って (2) も、上と同じ議論により補題
\ref{coherent-lemma-kill-by-twisting} から従う。
\end{proof}

\section{形式関数の定理}
\label{coherent-section-theorem-formal-functions}

\noindent
本節では、$\{I^n\mathcal{F}\}_{n \geq 0}$ または
$\{\mathcal{F}/I^n\mathcal{F}\}_{n \geq 0}$ という形の層の列について、
$n$ が変化するときのコホモロジーの振る舞いを調べる。

\medskip\noindent
以下では次の記法を用いる。スキームの射 $f : X \to Y$、準連接層
$\mathcal{F}$（$X$ 上）、および準連接イデアル層
$\mathcal{I} \subset \mathcal{O}_Y$ が与えられたとき、
$\mathcal{I}^n\mathcal{F}$ により、$f^{-1}(\mathcal{I}^n)$ の局所切断と
$\mathcal{F}$ の局所切断との積で生成される準連接部分層を表す。式で表せば
$$
\mathcal{I}^n\mathcal{F}
=
\Im\left(
f^*(\mathcal{I}^n) \otimes_{\mathcal{O}_X} \mathcal{F}
\longrightarrow
\mathcal{F}
\right).
$$
次の自然な写像があることに注意する：
$$
f^{-1}(\mathcal{I}^n) \otimes_{f^{-1}\mathcal{O}_Y} \mathcal{I}^m\mathcal{F}
\longrightarrow
f^*(\mathcal{I}^n) \otimes_{\mathcal{O}_X} \mathcal{I}^m\mathcal{F}
\longrightarrow
\mathcal{I}^{n + m}\mathcal{F}
$$
従って、高次順像の関手性により、$\mathcal{I}^n$ の切断から写像
$R^pf_*(\mathcal{I}^m\mathcal{F}) \to
R^pf_*(\mathcal{I}^{n + m}\mathcal{F})$ が生じる。局所化してから層化すると、
次の $\mathcal{O}_Y$-加群写像が存在することが分かる：
$$
\mathcal{I}^n \otimes_{\mathcal{O}_Y} R^pf_*(\mathcal{I}^m\mathcal{F})
\longrightarrow
R^pf_*(\mathcal{I}^{n + m}\mathcal{F}).
$$
言い換えると、$\bigoplus_{n \geq 0} R^pf_*(\mathcal{I}^n\mathcal{F})$ は
次数付き $\bigoplus_{n \geq 0} \mathcal{I}^n$-加群である。

\medskip\noindent
$Y = \Spec(A)$ かつ $\mathcal{I} = \widetilde{I}$ のとき、
$\mathcal{I}^n\mathcal{F}$ を単に $I^n\mathcal{F}$ と書く。上で導入した写像は
$M = \bigoplus H^p(X, I^n\mathcal{F})$ に次数付き
$S = \bigoplus I^n$-加群の構造を与える。$f$ が固有で、$A$ が Noether、
かつ $\mathcal{F}$ が連接ならば、これは有限型加群となる。

\begin{lemma}
\label{coherent-lemma-cohomology-powers-ideal-times-F}
\begin{reference}
\cite[III Cor 3.3.2]{EGA}
\end{reference}
$A$ を Noether 環とする。$I \subset A$ をイデアルとする。
$B = \bigoplus_{n \geq 0} I^n$ とおく。
$f : X \to \Spec(A)$ を固有射とし、連接層 $\mathcal{F}$（$X$ 上）をとる。
このとき、すべての $p \geq 0$ に対して、次数付き $B$-加群
$\bigoplus_{n \geq 0} H^p(X, I^n\mathcal{F})$ は有限 $B$-加群である。
\end{lemma}

\begin{proof}
$\mathcal{B} = f^*\widetilde{B}$ とおく。$\mathcal{B}$ は
$\bigoplus I^n\mathcal{O}_X$ へ全射するので、$\bigoplus I^n\mathcal{F}$ が
有限型次数付き $\mathcal{B}$-加群であることは明らかである。従って結果は
補題 \ref{coherent-lemma-graded-finiteness} の (1) から従う。
\end{proof}

\begin{lemma}
\label{coherent-lemma-cohomology-powers-ideal-times-sheaf}
スキームの射 $f : X \to Y$、準連接層 $\mathcal{F}$（$X$ 上）、および準連接
イデアル層 $\mathcal{I} \subset \mathcal{O}_Y$ が与えられているとする。
$Y$ は局所 Noether、$f$ は固有、かつ $\mathcal{F}$ は連接であると仮定する。
このとき
$$
\mathcal{M} =
\bigoplus\nolimits_{n \geq 0} R^pf_*(\mathcal{I}^n\mathcal{F})
$$
は準連接かつ有限型である次数付き
$\mathcal{A} = \bigoplus_{n \geq 0} \mathcal{I}^n$-加群である。
\end{lemma}

\begin{proof}
主張は $Y$ 上局所的なので、$Y$ がアフィンの場合に帰着する。アフィンの場合、
結果は補題 \ref{coherent-lemma-cohomology-powers-ideal-times-F} から従う。詳細は省略する。
\end{proof}

\begin{lemma}
\label{coherent-lemma-cohomology-powers-ideal-application}
$A$ を Noether 環とする。$I \subset A$ をイデアルとする。
$f : X \to \Spec(A)$ を固有射とし、$\mathcal{F}$ を $X$ 上の連接層とする。
このとき、すべての $p \geq 0$ に対して、次を満たす整数 $c \geq 0$ が存在する。
\begin{enumerate}
\item 乗法写像
$I^{n - c} \otimes H^p(X, I^c\mathcal{F}) \to H^p(X, I^n\mathcal{F})$
は、すべての $n \geq c$ に対して全射である。
\item $H^p(X, I^{n + m}\mathcal{F}) \to H^p(X, I^n\mathcal{F})$ の像は部分加群
$I^{m - e} H^p(X, I^n\mathcal{F})$ に含まれる。ただし
$e = \max(0, c - n)$ であり、$n + m \geq c$、$n, m \geq 0$ とする。
\item $n \geq c$ に対して
$$
\Ker(H^p(X, I^n\mathcal{F}) \to H^p(X, \mathcal{F})) =
\Ker(H^p(X, I^n\mathcal{F}) \to H^p(X, I^{n - c}\mathcal{F}))
$$
が成り立つ。
\item 写像 $I^nH^p(X, \mathcal{F}) \to H^p(X, I^{n - c}\mathcal{F})$
が $n \geq c$ に対して存在し、$n \geq 2c$ に対する合成
$$
H^p(X, I^n\mathcal{F}) \to 
I^{n - c}H^p(X, \mathcal{F}) \to
H^p(X, I^{n - 2c}\mathcal{F})
$$
および
$$
I^nH^p(X, \mathcal{F}) \to
H^p(X, I^{n - c}\mathcal{F}) \to
I^{n - 2c}H^p(X, \mathcal{F})
$$
は標準的なものである。
\item 逆系 $(H^p(X, I^n\mathcal{F}))$ と
$(I^nH^p(X, \mathcal{F}))$ はプロ同型である。
\end{enumerate}
\end{lemma}

\begin{proof}
$M_n = H^p(X, I^n\mathcal{F})$ と $n \geq 1$ に対して書き、
$M_0 = H^p(X, \mathcal{F})$ とおくと、写像
$\ldots \to M_3 \to M_2 \to M_1 \to M_0$ がある。
$B = \bigoplus_{n \geq 0} I^n$ とおけば、
$M = \bigoplus_{n \geq 0} M_n$ は有限次数付き $B$-加群である。補題
\ref{coherent-lemma-cohomology-powers-ideal-times-F} を参照。
積 $B_n \otimes M_m \to M_{m + n}$、$a \otimes m \mapsto a \cdot m$ は、
逆系の写像と次の意味で両立することに注意する。すなわち
$$
\xymatrix{
B_n \otimes_A M_m \ar[r] \ar[d] & M_{n + m} \ar[d] \\
B_n \otimes_A M_{m'} \ar[r] & M_{n + m'}
}
$$
は $n, m' \geq 0$ および $m \geq m'$ に対して可換である。

\medskip\noindent
(1) の証明。$d_1, \ldots, d_t \geq 0$ と $x_i \in M_{d_i}$ を、
$M$ が $x_1, \ldots, x_t$ により $B$ 上生成されるように選ぶ。
任意の $c \geq \max\{d_i\}$ に対して
$B_{n - c} \cdot M_c = M_n$ が $n \geq c$ のとき成り立つので、(1) を得る。

\medskip\noindent
(2) の証明。$c$ を (1) の証明におけるものとする。$n + m \geq c$ とする。
$M_{n + m} = B_{n + m - c} \cdot M_c$ である。$c > n$ ならば、
$M_c \to M_n$ と、上で述べた積と遷移写像との両立性を用いることにより、
$M_{n + m} \to M_n$ の像が $I^{n + m - c}M_n$ に含まれることが分かる。
$c \leq n$ ならば、
$M_{n + m} = B_m \cdot B_{n - c} \cdot M_c =
B_m \cdot M_n$ と書くことにより、
像が $I^m M_n$ に含まれることが分かる。これで (2) が示された。

\medskip\noindent
$K_n \subset M_n$ を写像 $M_n \to M_0$ の核とする。上で述べた積と遷移写像との
両立性から、$K = \bigoplus K_n \subset M$ は次数付き $B$-部分加群である。
$B$ は Noether で、$M$ は有限生成次数付き $B$-加群なので、
$K$ は有限生成次数付き $B$-加群である。
$d'_1, \ldots, d'_{t'} \geq 0$ と $y_i \in K_{d'_i}$ を、
$K$ が $y_1, \ldots, y_{t'}$ により $B$ 上生成されるように選ぶ。
$c = \max(d'_i, d'_j)$ とおく。$y_i \in \Ker(M_{d'_i} \to M_0)$ なので、
$B_n \cdot y_i \subset \Ker(M_{n + d'_i} \to M_n)$ である。
こうして、$K_n = \Ker(M_n \to M_{n - c})$ であることが
$n \geq c$ に対して分かる。これで (3) が示された。

\medskip\noindent
次の実線からなる可換図式を考える：
$$
\xymatrix{
I^n \otimes_A M_0 \ar[r] \ar[d] &
I^nM_0 \ar[r] \ar@{..>}[d] &
M_0 \ar[d] \\
M_n \ar[r] &
M_{n - c} \ar[r] &
M_0
}
$$
全射 $I^n \otimes_A M_0 \to I^nM_0$ の核は、上で示したことにより
$K_n$ へ写るので、点線の矢印が得られ、合成
$I^nM_0 \to M_{n - c} \to M_0$ は標準写像である。
従って、合成 $I^nM_0 \to M_{n - c} \to I^{n - 2c}M_0$ が
$n \geq 2c$ に対して標準写像であることは明らかである。
合成 $M_n \to I^{n - c}M_0 \to M_{n - 2c}$ を考える。第1の写像は、
$a \cdot m$ という形の元（$a \in I^{n - c}$、$m \in M_c$）を
$a m'$ へ送る。ただし $m'$ は $m$ の $M_0$ における像である。
次に第2の写像はこれを $a \cdot m'$ へ、$M_{n - 2c}$ において送るので、
(4) が成り立つことが分かる。

\medskip\noindent
(5) は (4) とプロ対象の射の定義から直ちに従う。
\end{proof}

\noindent
補題 \ref{coherent-lemma-cohomology-powers-ideal-times-F} および
\ref{coherent-lemma-cohomology-powers-ideal-application} の状況で、逆系
$$
\mathcal{F}/I\mathcal{F} \leftarrow
\mathcal{F}/I^2\mathcal{F} \leftarrow
\mathcal{F}/I^3\mathcal{F} \leftarrow \ldots
$$
を考える。そのコホモロジー群がどうなるかを知りたい。まず次の結果を示す。

\begin{lemma}
\label{coherent-lemma-ML-cohomology-powers-ideal}
$A$ を Noether 環とする。$I \subset A$ をイデアルとする。
$f : X \to \Spec(A)$ を固有射とし、$\mathcal{F}$ を $X$ 上の連接層とする。
$p \geq 0$ を固定する。このとき、次を満たす $c \geq 0$ が存在する。
\begin{enumerate}
\item すべての $n \geq c$ に対して
$$
\Ker(H^p(X, \mathcal{F}) \to H^p(X, \mathcal{F}/I^n\mathcal{F})) \subset
I^{n - c}H^p(X, \mathcal{F}).
$$
\item 逆系
$$
\left(H^p(X, \mathcal{F}/I^n\mathcal{F})\right)_{n \in \mathbf{N}}
$$
は Mittag-Leffler 条件を満たす（Homology, Definition
\ref{homology-definition-Mittag-Leffler} を参照）。
\item すべての $k \geq n + c$ に対して
$$
\Im(H^p(X, \mathcal{F}/I^k\mathcal{F})
\to H^p(X, \mathcal{F}/I^n\mathcal{F}))
=
\Im(H^p(X, \mathcal{F})
\to H^p(X, \mathcal{F}/I^n\mathcal{F}))
$$
が成り立つ。
\end{enumerate}
\end{lemma}

\begin{proof}
$c = \max\{c_p, c_{p + 1}\}$ とおく。ただし $c_p, c_{p + 1}$ は、$H^p$ および
$H^{p + 1}$ に対して補題 \ref{coherent-lemma-cohomology-powers-ideal-application}
で得られる整数である。

\medskip\noindent
(1) を示そう。短完全列
$$
0 \to I^n\mathcal{F} \to \mathcal{F} \to \mathcal{F}/I^n\mathcal{F} \to 0
$$
を考える。コホモロジー長完全列から
$$
\Ker(
H^p(X, \mathcal{F}) \to H^p(X, \mathcal{F}/I^n\mathcal{F})
)
=
\Im(
H^p(X, I^n\mathcal{F}) \to H^p(X, \mathcal{F})
)
$$
であることが分かる。従って補題
\ref{coherent-lemma-cohomology-powers-ideal-application} の (2) により、これは
$I^{n - c}H^p(X, \mathcal{F})$ に $n \geq c$ に対して含まれる。

\medskip\noindent
(3) は、Mittag-Leffler 系の定義により (2) を含意することに注意する。

\medskip\noindent
(3) を示そう。$n$ を固定する。次の可換図式を考える：
$$
\xymatrix{
0 \ar[r] &
I^n\mathcal{F} \ar[r] &
\mathcal{F} \ar[r] &
\mathcal{F}/I^n\mathcal{F} \ar[r] &
0 \\
0 \ar[r] &
I^{n + m}\mathcal{F} \ar[r] \ar[u] &
\mathcal{F} \ar[r] \ar[u] &
\mathcal{F}/I^{n + m}\mathcal{F} \ar[r] \ar[u] &
0
}
$$
これから、次の可換図式が得られる：
$$
\xymatrix{
H^p(X, \mathcal{F}) \ar[r] &
H^p(X, \mathcal{F}/I^n\mathcal{F}) \ar[r]_\delta &
H^{p + 1}(X, I^n\mathcal{F}) \ar[r] &
H^{p + 1}(X, \mathcal{F}) \\
H^p(X, \mathcal{F}) \ar[r] \ar[u]^1 &
H^p(X, \mathcal{F}/I^{n + m}\mathcal{F}) \ar[r] \ar[u]^\gamma &
H^{p + 1}(X, I^{n + m}\mathcal{F}) \ar[u]^\alpha \ar[r]^-\beta &
H^{p + 1}(X, \mathcal{F}) \ar[u]_1
}
$$
両行は完全である。補題 \ref{coherent-lemma-cohomology-powers-ideal-application}
の (4) により、$\beta$ の核は $\alpha$ の核に $m \geq c$ に対して等しい。
図式追跡により、$\gamma$ の像は $\delta$ の核に含まれる。これは (3) を示す
（$k = n + m$ とおけばよい）。
\end{proof}

\begin{theorem}[形式関数の定理]
\label{coherent-theorem-formal-functions}
$A$ を Noether 環とする。$I \subset A$ をイデアルとする。
$f : X \to \Spec(A)$ を固有射とし、$\mathcal{F}$ を $X$ 上の連接層とする。
$p \geq 0$ を固定する。写像の系
$$
H^p(X, \mathcal{F})/I^nH^p(X, \mathcal{F})
\longrightarrow
H^p(X, \mathcal{F}/I^n\mathcal{F})
$$
は極限の同型
$$
H^p(X, \mathcal{F})^\wedge
\longrightarrow
\lim_n H^p(X, \mathcal{F}/I^n\mathcal{F})
$$
を定める。ただし左辺は、$A$-加群 $H^p(X, \mathcal{F})$ のイデアル
$I$ に関する完備化である。Algebra, Section
\ref{algebra-section-completion} を参照。さらに、これは実際に極限位相に関する
同相写像である。
\end{theorem}

\begin{proof}
これは補題 \ref{coherent-lemma-ML-cohomology-powers-ideal} から次のように従う。
$M = H^p(X, \mathcal{F})$、$M_n = H^p(X, \mathcal{F}/I^n\mathcal{F})$ とおき、
$N_n = \Im(M \to M_n)$ と書く。補題
\ref{coherent-lemma-ML-cohomology-powers-ideal} の (2) および (3) により、
$(M_n)$ は Mittag-Leffler 系であり、$N_n \subset M_n$ は $M_k$ の像に、
すべての $k \gg n$ に対して等しい。従って位相加群として
$\lim M_n = \lim N_n$ である（極限位相を入れる）。一方、$N_n$ は加群 $M$ の
商からなる逆系なので、$\lim N_n$ は $M$ の完備化であり、その位相は核
$K_n = \Ker(M \to N_n)$ により与えられる。補題
\ref{coherent-lemma-ML-cohomology-powers-ideal} の
(1) により $K_n \subset I^{n - c}M$ であり、また $N_n \subset M_n$ は
$I^n$ で零化されるので $I^n M \subset K_n$ である。従って部分加群 $K_n$ を
$0$ の開近傍の基本系として用いて定められる位相は $I$-進位相と同じであり、
誘導される写像
$M^\wedge = \lim M/I^nM \to \lim N_n = \lim M_n$
は位相加群の同型である\footnote{
$M^\wedge$ 上の極限位相がその $I$-進位相と同じであることは、
Algebra, Lemma \ref{algebra-lemma-hathat-finitely-generated} から従う。
More on Algebra, Section \ref{more-algebra-section-topological-ring} も参照。}。
\end{proof}

\begin{lemma}
\label{coherent-lemma-spell-out-theorem-formal-functions}
$A$ を環とし、$I \subset A$ をイデアルとする。$A$ は Noether であり、
$I$ に関して完備であると仮定する。$f : X \to \Spec(A)$ を固有射とし、
$\mathcal{F}$ を $X$ 上の連接層とする。このとき、すべての $p \geq 0$ に対して
$$
H^p(X, \mathcal{F}) = \lim_n H^p(X, \mathcal{F}/I^n\mathcal{F})
$$
が成り立つ。
\end{lemma}

\begin{proof}
これは、完備 Noether 基礎環の場合における形式関数の定理
（定理 \ref{coherent-theorem-formal-functions}）の言い換えである。実際、この場合、
$A$-加群 $H^p(X, \mathcal{F})$ は有限であり
（補題 \ref{coherent-lemma-proper-over-affine-cohomology-finite}）、従って
$I$-進完備である（Algebra, Lemma
\ref{algebra-lemma-completion-tensor}）。ゆえに左辺で完備化する必要はない。
\end{proof}

\begin{lemma}
\label{coherent-lemma-formal-functions-stalk}
スキームの射 $f : X \to Y$ および準連接層 $\mathcal{F}$（$X$ 上）が与えられ、
次を仮定する。
\begin{enumerate}
\item $Y$ は局所 Noether である。
\item $f$ は固有である。
\item $\mathcal{F}$ は連接である。
\end{enumerate}
$y \in Y$ を点とする。ファイバー $X_1 = X_y$ の無限小近傍
$$
\xymatrix{
X_n =
\Spec(\mathcal{O}_{Y, y}/\mathfrak m_y^n) \times_Y X
\ar[r]_-{i_n} \ar[d]_{f_n} &
X \ar[d]^f \\
\Spec(\mathcal{O}_{Y, y}/\mathfrak m_y^n) \ar[r]^-{c_n} & Y
}
$$
を考え、$\mathcal{F}_n = i_n^*\mathcal{F}$ とおく。このとき
$$
\left(R^pf_*\mathcal{F}\right)_y^\wedge
\cong
\lim_n H^p(X_n, \mathcal{F}_n)
$$
が $\mathcal{O}_{Y, y}^\wedge$-加群として成り立つ。
\end{lemma}

\begin{proof}
これは形式関数の定理、定理 \ref{coherent-theorem-formal-functions} の特殊な場合の
言い換えにすぎない。詳しく述べよう。$\mathcal{O}_{Y, y}$ は Noether 局所環
であることに注意する。標準射
$c : \Spec(\mathcal{O}_{Y, y}) \to Y$ を考える。Schemes, Equation
(\ref{schemes-equation-canonical-morphism}) を参照。これは局所環を同一視する
ので平坦射である。ここでは一時的に
$f' : X' \to \Spec(\mathcal{O}_{Y, y})$ により、$f$ のこの局所環への基底変換を
表す。補題 \ref{coherent-lemma-flat-base-change-cohomology} により
$c^*R^pf_*\mathcal{F} = R^pf'_*\mathcal{F}'$ であることが分かる。さらに、
ファイバー $X_y$ と $X'_y$ の無限小近傍は同一視される
（検証は省略する。ヒント：射 $c_n$ は $c$ を経由する）。

\medskip\noindent
従って、$Y = \Spec(A)$ は Noether 局所環 $A$（極大イデアル
$\mathfrak m$）のスペクトルであり、$y \in Y$ は閉点（すなわち $\mathfrak m$）に対応すると
仮定してよい。特に
$$
\left(R^pf_*\mathcal{F}\right)_y =
\Gamma(Y, R^pf_*\mathcal{F}) =
H^p(X, \mathcal{F}).
$$

\medskip\noindent
この場合、射 $c_n$ はいずれも閉埋め込みである。従ってその基底変換 $i_n$ も
閉埋め込みである。$i_{n, *}\mathcal{F}_n = i_{n, *}i_n^*\mathcal{F}
= \mathcal{F}/\mathfrak m^n\mathcal{F}$ であることに注意する。
$i_n$ に対する Leray スペクトル系列と補題
\ref{coherent-lemma-finite-pushforward-coherent} により
$$
H^p(X_n, \mathcal{F}_n) =
H^p(X, i_{n, *}\mathcal{F}_n) =
H^p(X, \mathcal{F}/\mathfrak m^n\mathcal{F})
$$
であることが分かる。従って、補題の主張に現れる極限を計算するために、確かに
形式関数の定理を適用でき、結論を得る。
\end{proof}

\noindent
次の補題は、後で次の補題、すなわち次元 $ > 0$ のファイバーの場合へ一般化する。

\begin{lemma}
\label{coherent-lemma-higher-direct-images-zero-finite-fibre}
$f : X \to Y$ をスキームの射とし、$y \in Y$ とする。次を仮定する。
\begin{enumerate}
\item $Y$ は局所 Noether である。
\item $f$ は固有である。
\item $f^{-1}(\{y\})$ は有限である。
\end{enumerate}
このとき、任意の連接層 $\mathcal{F}$（$X$ 上）に対して、
$(R^pf_*\mathcal{F})_y = 0$ がすべての $p > 0$ について成り立つ。
\end{lemma}

\begin{proof}
ファイバー $X_y$ は有限であり、Morphisms, Lemma
\ref{morphisms-lemma-finite-fibre} により有限離散空間である。さらに、各無限小
近傍 $X_n$ の台位相空間は同じである。従って各スキーム $X_n$ は、Schemes,
Lemma \ref{schemes-lemma-scheme-finite-discrete-affine} によりアフィンである。
ゆえに、$H^p(X_n, \mathcal{F}_n) = 0$ がすべての $p > 0$ に対して成り立つ。
従って補題
\ref{coherent-lemma-formal-functions-stalk} により
$(R^pf_*\mathcal{F})_y^\wedge = 0$ であることが分かる。
$R^pf_*\mathcal{F}$ は命題
\ref{coherent-proposition-proper-pushforward-coherent} により連接であり、従って
$R^pf_*\mathcal{F}_y$ は有限 $\mathcal{O}_{Y, y}$-加群であることに注意する。
Nakayama の補題（Algebra, Lemma \ref{algebra-lemma-NAK}）により、局所環上の
有限加群の完備化が零ならば、その加群は零である。従って
$(R^pf_*\mathcal{F})_y = 0$ である。
\end{proof}

\begin{lemma}
\label{coherent-lemma-higher-direct-images-zero-above-dimension-fibre}
$f : X \to Y$ をスキームの射とし、$y \in Y$ とする。次を仮定する。
\begin{enumerate}
\item $Y$ は局所 Noether である。
\item $f$ は固有である。
\item $\dim(X_y) = d$ である。
\end{enumerate}
このとき、任意の連接層 $\mathcal{F}$（$X$ 上）に対して、
$(R^pf_*\mathcal{F})_y = 0$ がすべての $p > d$ について成り立つ。
\end{lemma}

\begin{proof}
ファイバー $X_y$ は $\Spec(\kappa(y))$ 上有限型である。従って $X_y$ は
Morphisms, Lemma \ref{morphisms-lemma-finite-type-noetherian} により
Noether スキームである。ゆえに、$X_y$ の台位相空間は Noether である。
Properties, Lemma \ref{properties-lemma-Noetherian-topology} を参照。
さらに、各無限小近傍 $X_n$ の台位相空間は $X_y$ のそれと同じである。
従って Cohomology, Proposition
\ref{cohomology-proposition-vanishing-Noetherian} により、
$H^p(X_n, \mathcal{F}_n) = 0$ がすべての $p > d$ に対して成り立つ。
従って、$(R^pf_*\mathcal{F})_y^\wedge = 0$ であることが、補題
\ref{coherent-lemma-formal-functions-stalk} により $p > d$ に対して分かる。
$R^pf_*\mathcal{F}$ は命題
\ref{coherent-proposition-proper-pushforward-coherent} により連接であり、従って
$R^pf_*\mathcal{F}_y$ は有限 $\mathcal{O}_{Y, y}$-加群であることに注意する。
Nakayama の補題（Algebra, Lemma \ref{algebra-lemma-NAK}）により、局所環上の
有限加群の完備化が零ならば、その加群は零である。従って
$(R^pf_*\mathcal{F})_y = 0$ である。
\end{proof}

\section{形式関数の定理の応用}
\label{coherent-section-applications-formal-functions}


\noindent
必要に応じて、ここにはさらに結果を加える。当面必要なのは、Noether の場合における
有限射の次の特徴づけである。

\begin{lemma}
\label{coherent-lemma-characterize-finite}
（より一般的な版については More on Morphisms, Lemma
\ref{more-morphisms-lemma-characterize-finite} を参照。）
$f : X \to S$ をスキームの射とする。$S$ は局所 Noether であると仮定する。
次は同値である。
\begin{enumerate}
\item $f$ は有限である。
\item $f$ は有限ファイバーをもつ固有射である。
\end{enumerate}
\end{lemma}

\begin{proof}
有限射は Morphisms, Lemma \ref{morphisms-lemma-finite-proper} により固有である。
有限射は Morphisms, Lemma \ref{morphisms-lemma-finite-quasi-finite} により
準有限である。準有限射は有限ファイバーをもつ。Morphisms, Lemma
\ref{morphisms-lemma-quasi-finite} を参照。従って有限射は固有かつ有限ファイバーを
もつ。

\medskip\noindent
$f$ は有限ファイバーをもつ固有射であると仮定する。$f$ が有限であることを
示したい。実際、$f$ がアフィンであることを示せば十分である。というのも、
$f$ がアフィンならば、Morphisms, Lemma
\ref{morphisms-lemma-integral-universally-closed} により $f$ は整であり、
従って Morphisms, Lemma \ref{morphisms-lemma-finite-integral} により
$f$ は有限だからである。

\medskip\noindent
$f$ がアフィンであることを示すため、$S$ がアフィンであると仮定してよく、
目標は $X$ もアフィンであると示すことである。$f$ は固有なので、$X$ は分離的かつ
準コンパクトである。従って補題
\ref{coherent-lemma-quasi-separated-h1-zero-covering} の判定法を用いて、$X$ がアフィン
であることを示せばよい。このため、$\mathcal{I} \subset \mathcal{O}_X$ を有限型
イデアル層とする。特に $\mathcal{I}$ は $X$ 上の連接層である。補題
\ref{coherent-lemma-higher-direct-images-zero-finite-fibre} により
$R^1f_*\mathcal{I}_s = 0$ がすべての $s \in S$ に対して成り立つ。
言い換えると、$R^1f_*\mathcal{I} = 0$ である。従って $f$ に対する Leray
スペクトル系列から
$H^1(X , \mathcal{I}) = H^1(S, f_*\mathcal{I})$ が分かる。$S$ はアフィンであり、
$f_*\mathcal{I}$ は準連接なので
（Schemes, Lemma \ref{schemes-lemma-push-forward-quasi-coherent}）、
補題 \ref{coherent-lemma-quasi-coherent-affine-cohomology-zero} から、望んだとおり
$H^1(S, f_*\mathcal{I}) = 0$ を得る。従って、望んだとおり
$H^1(X, \mathcal{I}) = 0$ である。
\end{proof}

\noindent
その帰結として、次の有用な結果を得る。

\begin{lemma}
\label{coherent-lemma-proper-finite-fibre-finite-in-neighbourhood}
\begin{slogan}
固有射は有限ファイバーの近傍で有限である。
\end{slogan}
（より一般的な版については More on Morphisms, Lemma
\ref{more-morphisms-lemma-proper-finite-fibre-finite-in-neighbourhood} を参照。）
$f : X \to S$ をスキームの射とし、$s \in S$ とする。次を仮定する。
\begin{enumerate}
\item $S$ は局所 Noether である。
\item $f$ は固有である。
\item $f^{-1}(\{s\})$ は有限集合である。
\end{enumerate}
このとき、開近傍 $V \subset S$（$s$ の）であって
$f|_{f^{-1}(V)} : f^{-1}(V) \to V$ が有限となるものが存在する。
\end{lemma}

\begin{proof}
射 $f$ は、Morphisms, Lemma \ref{morphisms-lemma-finite-fibre} により、
$f^{-1}(\{s\})$ のすべての点で準有限である。Morphisms, Lemma
\ref{morphisms-lemma-quasi-finite-points-open} により、$f$ が準有限である点の
集合は開集合 $U \subset X$ である。$Z = X \setminus U$ とおく。このとき
$s \not \in f(Z)$ である。$f$ は固有なので、集合 $f(Z) \subset S$ は閉である。
任意の開近傍 $V \subset S$（$s$ の）で $Z \cap V = \emptyset$ を満たすものを選ぶ。
すると $f^{-1}(V) \to V$ は局所準有限かつ固有である。従って準有限であり
（Morphisms, Lemma
\ref{morphisms-lemma-quasi-finite-locally-quasi-compact}）、ゆえに有限ファイバーを
もち（Morphisms, Lemma \ref{morphisms-lemma-quasi-finite}）、従って補題
\ref{coherent-lemma-characterize-finite} により有限である。
\end{proof}

\begin{lemma}
\label{coherent-lemma-ample-on-fibre}
$f : X \to Y$ を、$Y$ が Noether であるようなスキームの固有射とする。
$\mathcal{L}$ を可逆 $\mathcal{O}_X$-加群とし、$\mathcal{F}$ を連接
$\mathcal{O}_X$-加群とする。$y \in Y$ を、$\mathcal{L}_y$ が $X_y$ 上豊富である
ような点とする。このとき、ある $d_0$ が存在し、すべての $d \geq d_0$ に対して
$$
R^pf_*(\mathcal{F} \otimes_{\mathcal{O}_X} \mathcal{L}^{\otimes d})_y = 0
\text{ for }p > 0
$$
かつ、写像
$$
f_*(\mathcal{F} \otimes_{\mathcal{O}_X} \mathcal{L}^{\otimes d})_y
\longrightarrow
H^0(X_y, \mathcal{F}_y \otimes_{\mathcal{O}_{X_y}} \mathcal{L}_y^{\otimes d})
$$
が全射となる。
\end{lemma}

\begin{proof}
$\mathcal{O}_{Y, y}$ は Noether 局所環であることに注意する。標準射
$c : \Spec(\mathcal{O}_{Y, y}) \to Y$ を考える。Schemes, Equation
(\ref{schemes-equation-canonical-morphism}) を参照。これは局所環を同一視する
ので平坦射である。ここでは一時的に
$f' : X' \to \Spec(\mathcal{O}_{Y, y})$ により、$f$ のこの局所環への基底変換を
表す。補題 \ref{coherent-lemma-flat-base-change-cohomology} により
$c^*R^pf_*\mathcal{F} = R^pf'_*\mathcal{F}'$ であることが分かる。さらに、
ファイバー $X_y$ と $X'_y$ は同一視される。従って、
$Y = \Spec(A)$ は Noether 局所環 $(A, \mathfrak m, \kappa)$ のスペクトルであり、
$y \in Y$ は $\mathfrak m$ に対応すると仮定してよい。この場合、
$R^pf_*(\mathcal{F} \otimes_{\mathcal{O}_X} \mathcal{L}^{\otimes d})_y =
H^p(X, \mathcal{F} \otimes_{\mathcal{O}_X} \mathcal{L}^{\otimes d})$
がすべての $p \geq 0$ に対して成り立つ。
$f_y : X_y \to \Spec(\kappa)$ により射影を表す。

\medskip\noindent
$B = \text{Gr}_\mathfrak m(A) =
\bigoplus_{n \geq 0} \mathfrak m^n/\mathfrak m^{n + 1}$ とおく。
層 $\mathcal{B} = f_y^*\widetilde{B}$ を準連接次数付き
$\mathcal{O}_{X_y}$-代数の層として考える。$I$ を $\mathfrak m$ に置き換えて、
節 \ref{coherent-section-theorem-formal-functions} と同じ記法を用いる。
$X_y$ は $X$ の、$\mathfrak m\mathcal{O}_X$ が切り出す閉部分スキームなので、
$\mathfrak m^n\mathcal{F}/\mathfrak m^{n + 1}\mathcal{F}$ を連接
$\mathcal{O}_{X_y}$-加群とみなせる。補題 \ref{coherent-lemma-i-star-equivalence} を参照。
このとき
$\bigoplus_{n \geq 0} \mathfrak m^n\mathcal{F}/\mathfrak m^{n + 1}\mathcal{F}$ は
有限型準連接次数付き $\mathcal{B}$-加群である。実際、これは
$\mathcal{B}$ 上次数零で生成され、その次数零部分
$\mathcal{F}_y = \mathcal{F}/\mathfrak m \mathcal{F}$ は連接
$\mathcal{O}_{X_y}$-加群だからである。従って補題
\ref{coherent-lemma-graded-finiteness} の (2) により
$$
H^p(X_y, \mathfrak m^n \mathcal{F}/ \mathfrak m^{n + 1}\mathcal{F}
\otimes_{\mathcal{O}_{X_y}} \mathcal{L}_y^{\otimes d}) = 0
$$
が、すべての $p > 0$、$d \geq d_0$、および $n \geq 0$ に対して成り立つ。
補題 \ref{coherent-lemma-relative-affine-cohomology} により、これは
$
H^p(X, \mathfrak m^n \mathcal{F}/ \mathfrak m^{n + 1}\mathcal{F}
\otimes_{\mathcal{O}_X} \mathcal{L}^{\otimes d}) = 0
$
が、すべての $p > 0$、$d \geq d_0$、および $n \geq 0$ に対して成り立つという
主張と同じである。

\medskip\noindent
連接 $\mathcal{O}_X$-加群の短完全列
$$
0 \to \mathfrak m^n\mathcal{F}/\mathfrak m^{n + 1} \mathcal{F}
\to \mathcal{F}/\mathfrak m^{n + 1} \mathcal{F}
\to \mathcal{F}/\mathfrak m^n \mathcal{F} \to 0
$$
を考える。$\mathcal{L}^{\otimes d}$ とのテンソル積は完全関手なので、短完全列
$$
0 \to
\mathfrak m^n\mathcal{F}/\mathfrak m^{n + 1} \mathcal{F}
\otimes_{\mathcal{O}_X} \mathcal{L}^{\otimes d}
\to \mathcal{F}/\mathfrak m^{n + 1} \mathcal{F}
\otimes_{\mathcal{O}_X} \mathcal{L}^{\otimes d}
\to \mathcal{F}/\mathfrak m^n \mathcal{F}
\otimes_{\mathcal{O}_X} \mathcal{L}^{\otimes d} \to 0
$$
を得る。コホモロジー長完全列と上の消滅を用いれば、帰納法により次を得る。
\begin{enumerate}
\item $H^p(X, \mathcal{F}/\mathfrak m^n \mathcal{F}
\otimes_{\mathcal{O}_X} \mathcal{L}^{\otimes d}) = 0$
が、すべての $p > 0$、$d \geq d_0$、および $n \geq 0$ に対して成り立つ。
\item $H^0(X, \mathcal{F}/\mathfrak m^n \mathcal{F}
\otimes_{\mathcal{O}_X} \mathcal{L}^{\otimes d}) \to
H^0(X_y, \mathcal{F}_y \otimes_{\mathcal{O}_{X_y}} \mathcal{L}_y^{\otimes d})$
は、すべての $d \geq d_0$ および $n \geq 1$ に対して全射である。
\end{enumerate}
形式関数の定理（定理 \ref{coherent-theorem-formal-functions}）により、
$\mathfrak m$-進完備化した
$H^p(X, \mathcal{F} \otimes_{\mathcal{O}_X} \mathcal{L}^{\otimes d})$ が、
すべての $d \geq d_0$ および $p > 0$ に対して
零であることが分かる。
$H^p(X, \mathcal{F} \otimes_{\mathcal{O}_X} \mathcal{L}^{\otimes d})$ は補題
\ref{coherent-lemma-proper-over-affine-cohomology-finite} により有限 $A$-加群なので、
Nakayama の補題（Algebra, Lemma \ref{algebra-lemma-NAK}）から、
$H^p(X, \mathcal{F} \otimes_{\mathcal{O}_X} \mathcal{L}^{\otimes d})$ は
すべての $d \geq d_0$ および $p > 0$ に対して零である。$p = 0$ に対しては、
補題 \ref{coherent-lemma-ML-cohomology-powers-ideal} の (3) から
$H^0(X, \mathcal{F} \otimes_{\mathcal{O}_X} \mathcal{L}^{\otimes d}) \to
H^0(X_y, \mathcal{F}_y \otimes_{\mathcal{O}_{X_y}} \mathcal{L}_y^{\otimes d})$
が全射であることを得る。これが補題の最後の主張である。
\end{proof}

\begin{lemma}
\label{coherent-lemma-ample-in-neighbourhood}
（より一般的な版については More on Morphisms, Lemma
\ref{more-morphisms-lemma-ample-in-neighbourhood} を参照。）
$f : X \to Y$ を、$Y$ が Noether であるようなスキームの固有射とする。
$\mathcal{L}$ を可逆 $\mathcal{O}_X$-加群とする。$y \in Y$ を、
$\mathcal{L}_y$ が $X_y$ 上豊富であるような点とする。このとき、開近傍
$V \subset Y$（$y$ の）であって、$\mathcal{L}|_{f^{-1}(V)}$ が $f^{-1}(V)/V$ 上豊富と
なるものが存在する。
\end{lemma}

\begin{proof}
$d_0$ を補題 \ref{coherent-lemma-ample-on-fibre} のとおり、
$\mathcal{F} = \mathcal{O}_X$ の場合について選ぶ。閉埋め込み
$$
\varphi_y =
\varphi_{\mathcal{L}_y^{\otimes d}, (s_{y, 0}, \ldots, s_{y, r})} :
X_y \to \mathbf{P}^r_{\kappa(y)}.
$$
を定めるものを見つけられるように、$d \geq d_0$ を選ぶ。すなわち、
$r \geq 0$ および切断
$s_{y, 0}, \ldots, s_{y, r} \in H^0(X_y, \mathcal{L}_y^{\otimes d})$ を選ぶ。
これは Morphisms, Lemma
\ref{morphisms-lemma-finite-type-over-affine-ample-very-ample} により可能だが、
$\varphi_y$ が閉埋め込みであることには Morphisms, Lemma
\ref{morphisms-lemma-image-proper-scheme-closed} も用い、可逆層と切断による
射影空間への射の記述には Constructions, Section
\ref{constructions-section-projective-space} を用いる。$d_0$ の選び方により、
$Y$ を $y$ の開近傍で置き換えた後、
$s_0, \ldots, s_r \in H^0(X, \mathcal{L}^{\otimes d})$ を、
$s_{y, 0}, \ldots, s_{y, r}$ へ写るように選べる。
$X_{s_i} \subset X$ を、$s_i$ が $\mathcal{L}^{\otimes d}$ の生成元となる開部分集合
とする。$s_{y, i}$ は $\mathcal{L}_y^{\otimes d}$ を生成するので、
$X_y \subset U = \bigcup X_{s_i}$ である。$X \to Y$ は閉なので、開近傍
$y \in V \subset Y$ であって $f^{-1}(V) \subset U$ となるものが存在する。
$Y$ を $V$ で置き換えた後、$s_i$ が $\mathcal{L}^{\otimes d}$ を生成すると
仮定してよい。従って射
$$
\varphi = \varphi_{\mathcal{L}^{\otimes d}, (s_0, \ldots, s_r)} :
X \longrightarrow \mathbf{P}^r_Y
$$
で、$\mathcal{L}^{\otimes d} \cong \varphi^*\mathcal{O}_{\mathbf{P}^r_Y}(1)$
を満たし、その $y$ への基底変換が $\varphi_y$ となるものを得る。

\medskip\noindent
小技を用いて証明を終える。「正しい」証明では、$\varphi$ が閉埋め込みであることを、
$y$ の開近傍へ基底変換した後に直接示す。実際、補題
\ref{coherent-lemma-proper-finite-fibre-finite-in-neighbourhood} により、$\varphi$ は
ファイバー $\mathbf{P}^r_{\kappa(y)}$（$\mathbf{P}^r_Y \to Y$ の $y$ 上のもの）
の開近傍上で有限である。
$\mathbf{P}^r_Y \to Y$ が閉であることを用い、$Y$ を縮小した後、
$\varphi$ が有限であると仮定してよい。このとき
$\mathcal{L}^{\otimes d} \cong \varphi^*\mathcal{O}_{\mathbf{P}^r_Y}(1)$ は、
非常に一般的な Morphisms, Lemma
\ref{morphisms-lemma-pullback-ample-tensor-relatively-ample} により豊富である。
\end{proof}












\section{コホモロジーと基底変換 III}
\label{coherent-section-cohomology-and-base-change-perfect}

\noindent
この節では、Derived Categories of Schemes, Section
\ref{perfect-section-cohomology-and-base-change-perfect} で論じる
非常に一般的な現象の最も単純な場合を証明する。次の補題を代数的な言葉で
述べたものについては、Remark \ref{coherent-remark-explain-perfect-direct-image} を参照されたい。

\begin{lemma}
\label{coherent-lemma-perfect-direct-image}
$A$ を Noether 環とし、$S = \Spec(A)$ とおく。$f : X \to S$ を
スキームの固有射とする。$\mathcal{F}$ を連接 $\mathcal{O}_X$-加群であって、
$S$ 上平坦なものとする。このとき
\begin{enumerate}
\item $R\Gamma(X, \mathcal{F})$ は $D(A)$ の完全対象であり、
\item 任意の環準同型 $A \to A'$ に対して、基底変換射
$$
R\Gamma(X, \mathcal{F}) \otimes_A^{\mathbf{L}} A'
\longrightarrow
R\Gamma(X_{A'}, \mathcal{F}_{A'})
$$
は同型である。
\end{enumerate}
\end{lemma}

\begin{proof}
有限アフィン開被覆 $X = \bigcup_{i = 1, \ldots, n} U_i$ を選ぶ。
補題 \ref{coherent-lemma-separated-case-relative-cech} および
\ref{coherent-lemma-base-change-complex} により、{\v C}ech 複体
$K^\bullet = {\check C}^\bullet(\mathcal{U}, \mathcal{F})$ は、
任意の環準同型 $A \to A'$ に対して
$$
K^\bullet \otimes_A A' = R\Gamma(X_{A'}, \mathcal{F}_{A'})
$$
を満たす。$K_{alt}^\bullet = {\check C}_{alt}^\bullet(\mathcal{U}, \mathcal{F})$
を交代 {\v C}ech 複体とする。Cohomology, Lemma
\ref{cohomology-lemma-alternating-usual} により、ホモトピー同値
$K_{alt}^\bullet \to K^\bullet$ が $A$-加群複体の間に存在する。特に、同様に
$$
K_{alt}^\bullet \otimes_A A' = R\Gamma(X_{A'}, \mathcal{F}_{A'})
$$
が成り立つ。$\mathcal{F}$ は $A$ 上平坦なので、各 $K_{alt}^n$ も
$A$ 上平坦である（Morphisms, Lemma
\ref{morphisms-lemma-flat-module-characterize} を参照）。さらに
$K_{alt}^\bullet$ は上に有界なので（このために交代 {\v C}ech 複体へ
切り替えた）、導来テンソル積の定義から
$K_{alt}^\bullet \otimes_A A' = K_{alt}^\bullet \otimes_A^{\mathbf{L}} A'$
である（More on Algebra, Section
\ref{more-algebra-section-derived-tensor-product} を参照）。補題
\ref{coherent-lemma-proper-over-affine-cohomology-finite} により、コホモロジー群
$H^i(K_{alt}^\bullet)$ は有限 $A$-加群である。また
$K_{alt}^\bullet$ は有界なので、More on Algebra, Lemma
\ref{more-algebra-lemma-Noetherian-pseudo-coherent} により
$K_{alt}^\bullet$ は擬連接である。任意の $A$-加群 $M$ に対し、
$A' = A \oplus M$ とおく。ここで $M$ は平方零イデアルである。すなわち、
$(a, m) \cdot (a', m') = (aa', am' + a'm)$ とする。以上から、
$K_{alt}^\bullet \otimes_A^\mathbf{L} A'$ のコホモロジーは次数
$0, \ldots, n$ にのみ存在する。従って
$K_{alt}^\bullet \otimes_A^\mathbf{L} M$ のコホモロジーも次数
$0, \ldots, n$ にのみ存在する。よって $K_{alt}^\bullet$ は有限 Tor 次元をもつ
（More on Algebra, Definition
\ref{more-algebra-definition-tor-amplitude} を参照）。結論は More on Algebra,
Lemma \ref{more-algebra-lemma-perfect} から従う。
\end{proof}

\begin{remark}
\label{coherent-remark-explain-perfect-direct-image}
補題 \ref{coherent-lemma-perfect-direct-image} の帰結として、有限生成射影
$A$-加群からなる有限複体 $M^\bullet$ であって、
$$
H^i(X_{A'}, \mathcal{F}_{A'}) = H^i(M^\bullet \otimes_A A')
$$
が $A'$ に関して関手的に成り立つものが存在する。$\mathcal{F}$ が
$A$ 上平坦であるという条件は本質的である。\cite{Hartshorne} を参照。
\end{remark}







\section{連接形式加群}
\label{coherent-section-coherent-formal}

\noindent
形式スキームの理論をまだ利用できないので、「Noether adic 形式スキーム上の
連接加群」という概念に代わる言葉を少し整備する。

\medskip\noindent
$X$ を Noether スキームとし、$\mathcal{I} \subset \mathcal{O}_X$ を
準連接イデアル層とする。次の条件を満たす逆系 $(\mathcal{F}_n)$ であって、
連接 $\mathcal{O}_X$-加群からなるものを考える。
\begin{enumerate}
\item $\mathcal{F}_n$ は $\mathcal{I}^n$ で零化される。
\item 遷移射は同型
$\mathcal{F}_{n + 1}/\mathcal{I}^n\mathcal{F}_{n + 1} \to \mathcal{F}_n$
を誘導する。
\end{enumerate}
このような逆系の射は通常どおり定義する。これらの逆系の圏を
$\textit{Coh}(X, \mathcal{I})$ と表す。この圏の対象について一連の補題を
証明していく。実際、以下の補題の大部分は、アフィンの場合におけるこの圏の
次の記述から直ちに従う。

\begin{lemma}
\label{coherent-lemma-inverse-systems-affine}
$X = \Spec(A)$ が Noether 環のスペクトルであり、$\mathcal{I}$ がイデアル
$I \subset A$ に付随する準連接イデアル層であるとする。このとき
$\textit{Coh}(X, \mathcal{I})$ は有限 $A^\wedge$-加群の圏と同値である。
ここで $A^\wedge$ は $A$ の $I$ に関する完備化である。
\end{lemma}

\begin{proof}
$\text{Mod}^{fg}_{A, I}$ を、次を満たす逆系 $(M_n)$ であって有限 $A$-加群から
なるものの圏とする：
(1) $M_n$ は $I^n$ で零化され、(2)
$M_{n + 1}/I^nM_{n + 1} = M_n$ である。$X$ 上の連接層と有限 $A$-加群との
対応（補題 \ref{coherent-lemma-coherent-Noetherian}）により、
$\text{Mod}^{fg}_{A, I}$ が有限 $A^\wedge$-加群の圏と同値であることを
示せば十分である。そのためには、任意の $(M_n)$ が
$\text{Mod}^{fg}_{A, I}$ の対象であるとき、加群
$$
M = \lim M_n
$$
が有限 $A^\wedge$-加群であり、$M/I^nM = M_n$ であることを示せばよい。
遷移射は全射なので、$M \to M_1$ は全射である。
$x_1, \ldots, x_t \in M$ を選び、それらの $M_1$ での像が生成元となるようにする。これにより系の射
$(A/I^n)^{\oplus t} \to M_n$ が誘導される。Nakayama の補題（Algebra,
Lemma \ref{algebra-lemma-NAK}）により、これらの射は全射である。
$K_n \subset (A/I^n)^{\oplus t}$ をその核とする。性質 (2) から
$K_{n + 1} \to K_n$ は全射であり、特に系 $(K_n)$ は Mittag-Leffler 条件を
満たす。Homology, Lemma \ref{homology-lemma-Mittag-Leffler} により、完全列
$0 \to K \to (A^\wedge)^{\oplus t} \to M \to 0$
を得る。ここで $K = \lim K_n$ である。従って $M$ は有限 $A^\wedge$-加群である。
$K \to K_n$ は全射なので
$$
M/I^nM = \Coker(K \to (A/I^n)^{\oplus t}) =
(A/I^n)^{\oplus t}/K_n = M_n
$$
となり、主張が従う。
\end{proof}

\begin{lemma}
\label{coherent-lemma-inverse-systems-abelian}
$X$ を Noether スキームとし、$\mathcal{I} \subset \mathcal{O}_X$ を
準連接イデアル層とする。
\begin{enumerate}
\item 圏 $\textit{Coh}(X, \mathcal{I})$ は Abel 圏である。
\item 開集合 $U \subset X$ に対して、制限関手
$\textit{Coh}(X, \mathcal{I}) \to \textit{Coh}(U, \mathcal{I}|_U)$
は完全である。
\item $\textit{Coh}(X, \mathcal{I})$ における完全性は、$X$ の開被覆の
各要素へ制限することによって判定できる。
\end{enumerate}
\end{lemma}

\begin{proof}
$\alpha =(\alpha_n) : (\mathcal{F}_n) \to (\mathcal{G}_n)$ を
$\textit{Coh}(X, \mathcal{I})$ の射とする。$\alpha$ の余核は逆系
$(\Coker(\alpha_n))$ である（詳細は省略する）。核を記述するため、
$l \geq m$ に対して
$$
\mathcal{K}'_{l, m} = \Im(\Ker(\alpha_l) \to \mathcal{F}_m)
$$
とおく。次を主張する。
\begin{enumerate}
\item[(a)] 逆系 $(\mathcal{K}'_{l, m})_{l \geq m}$ は終局的に定常であり、
その値を $\mathcal{K}'_m$ とする。
\item[(b)] 系 $(\mathcal{K}'_m/\mathcal{I}^n\mathcal{K}'_m)_{m \geq n}$ は
終局的に定常であり、その値を $\mathcal{K}_n$ とする。
\item[(c)] 系 $(\mathcal{K}_n)$ は $\textit{Coh}(X, \mathcal{I})$ の対象をなす。
\item[(d)] この対象は $\alpha$ の核である。
\end{enumerate}
(a)、(b)、(c) を示すには、アフィン局所的に作業してよい。そこで
$X = \Spec(A)$ とし、$\mathcal{I}$ はイデアル $I \subset A$ に対応するとする。
補題 \ref{coherent-lemma-inverse-systems-affine} により、$\alpha$ は射
$f : M \to N$ に対応する。これは有限 $A^\wedge$-加群の射である。
$K = \Ker(f)$ と書く。
$A^\wedge$ は Noether 環であることに注意する（Algebra, Lemma
\ref{algebra-lemma-completion-Noetherian-Noetherian}）。
$c \geq 0$ を、$K \cap I^n M \subset I^{n - c}K$ が $n \geq c$ に対して
成り立つような整数として選び
（Algebra, Lemma \ref{algebra-lemma-Artin-Rees}）、さらにこの整数が射 $f$ と
イデアル $I^\wedge = IA^\wedge$ に対して Algebra, Lemma
\ref{algebra-lemma-map-AR} を満たすようにする。このとき
$\mathcal{K}'_{l, m}$ は $A$-加群
$$
K'_{l, m} = \frac{a^{-1}(I^lN) + I^mM}{I^mM} =
\frac{K + I^{l - c}f^{-1}(I^cN) + I^mM}{I^mM} =
\frac{K + I^mM}{I^mM}
$$
に対応する。ここで最後の等号は $l \geq m + c$ のとき成り立つ。従って
$\mathcal{K}'_m$ は $A$-加群 $K/K \cap I^mM$ に対応し、
$\mathcal{K}'_m/\mathcal{I}^n\mathcal{K}'_m$ は
$$
\frac{K}{K \cap I^mM + I^nK} = \frac{K}{I^nK}
$$
に対応する。これは $m \geq n + c$ のとき成り立つ（上で選んだ $c$ による）。
従って $\mathcal{K}_n$ は $K/I^nK$ に対応する。

\medskip\noindent
(d) を証明する。上のアフィン上の記述から、合成
$(\mathcal{K}_n) \to (\mathcal{F}_n) \to (\mathcal{G}_n)$ が零であることは
明らかである。射 $\beta : (\mathcal{H}_n) \to (\mathcal{F}_n)$ であって、
$\alpha \circ \beta = 0$ を満たすものをとる。このとき
$\mathcal{H}_l \to \mathcal{F}_l$ の像は $\Ker(\alpha_l)$ に含まれる。
$\mathcal{H}_m = \mathcal{H}_l/\mathcal{I}^m\mathcal{H}_l$ が $l \geq m$ に対して
成り立つので、射の系
$\mathcal{H}_m \to \mathcal{K}'_{l, m}$ を得る。従って射
$\mathcal{H}_m \to \mathcal{K}_m'$ を得る。さらに
$\mathcal{H}_n = \mathcal{H}_m/\mathcal{I}^n\mathcal{H}_m$ なので、射の系
$\mathcal{H}_n \to \mathcal{K}'_m/\mathcal{I}^n\mathcal{K}'_m$、従って所望の射
$\mathcal{H}_n \to \mathcal{K}_n$ を得る。

\medskip\noindent
(1) の証明を終えるには、さらに $\Coim = \Im$ が
$\textit{Coh}(X, \mathcal{I})$ において成り立つことを示す必要がある。上で、アフィン上では核および余核を
とる操作が $\textit{Coh}(X, \mathcal{I})$ を加群の圏として記述することと
可換であることを見た。加群の圏では $\Im = \Coim$ が成り立つので、
$\Coim = \Im$ を $\textit{Coh}(X, \mathcal{I})$ において得る。
補題の (2) と (3) は、核および余核の構成から直ちに従う。
\end{proof}

\begin{lemma}
\label{coherent-lemma-inverse-systems-surjective}
$X$ を Noether スキームとし、$\mathcal{I} \subset \mathcal{O}_X$ を
準連接イデアル層とする。射
$(\mathcal{F}_n) \to (\mathcal{G}_n)$ が
$\textit{Coh}(X, \mathcal{I})$ において全射であるための必要十分条件は、
$\mathcal{F}_1 \to \mathcal{G}_1$ が全射であることである。
\end{lemma}

\begin{proof}
省略する。ヒント：アフィン開集合上で補題
\ref{coherent-lemma-inverse-systems-affine} と Algebra, Lemma
\ref{algebra-lemma-NAK} を用いる。
\end{proof}

\noindent
$X$ を Noether スキームとし、$\mathcal{I} \subset \mathcal{O}_X$ を
準連接イデアル層とする。関手
\begin{equation}
\label{coherent-equation-completion-functor}
\textit{Coh}(\mathcal{O}_X) \longrightarrow \textit{Coh}(X, \mathcal{I}), \quad
\mathcal{F} \longmapsto \mathcal{F}^\wedge
\end{equation}
があり、連接 $\mathcal{O}_X$-加群 $\mathcal{F}$ に
$\mathcal{F}^\wedge = (\mathcal{F}/\mathcal{I}^n\mathcal{F})$ を対応させる。
これは $\textit{Coh}(X, \mathcal{I})$ の対象である。

\begin{lemma}
\label{coherent-lemma-exact}
関手 (\ref{coherent-equation-completion-functor}) は完全である。
\end{lemma}

\begin{proof}
$X$ 上局所的に判定すれば十分である。従って $X$ はアフィン、すなわち補題
\ref{coherent-lemma-inverse-systems-affine} の状況にあると仮定してよい。この関手は関手
$\text{Mod}^{fg}_A \to \text{Mod}^{fg}_{A^\wedge}$ であり、有限 $A$-加群
$M$ にその完備化 $M^\wedge$ を対応させる。従って主張は
Algebra, Lemma \ref{algebra-lemma-completion-flat} から従う。
\end{proof}

\begin{lemma}
\label{coherent-lemma-completion-internal-hom}
$X$ を Noether スキームとし、$\mathcal{I} \subset \mathcal{O}_X$ を
準連接イデアル層とする。$\mathcal{F}$、$\mathcal{G}$ を連接
$\mathcal{O}_X$-加群とする。
$\mathcal{H} = \SheafHom_{\mathcal{O}_X}(\mathcal{G}, \mathcal{F})$ とおく。このとき
$$
\lim H^0(X, \mathcal{H}/\mathcal{I}^n\mathcal{H}) =
\Mor_{\textit{Coh}(X, \mathcal{I})}
(\mathcal{G}^\wedge, \mathcal{F}^\wedge).
$$
\end{lemma}

\begin{proof}
これを証明するには $X$ 上アフィン局所的に作業してよい。従って
$X = \Spec(A)$ と仮定し、$\mathcal{F}$、$\mathcal{G}$ は有限 $A$-加群
$M$、$N$ によって与えられるとする。このとき $\mathcal{H}$ は有限
$A$-加群 $H = \Hom_A(M, N)$ に対応する。補題の主張は、補題
\ref{coherent-lemma-inverse-systems-affine} の同値を通じて
$$
H^\wedge = \Hom_{A^\wedge}(M^\wedge, N^\wedge)
$$
という主張になる。Algebra, Lemma
\ref{algebra-lemma-completion-flat} を（3 回）用いると
$$
H^\wedge = \Hom_A(M, N) \otimes_A A^\wedge =
\Hom_{A^\wedge}(M \otimes_A A^\wedge, N \otimes_A A^\wedge) =
\Hom_{A^\wedge}(M^\wedge, N^\wedge)
$$
を得る。ここで二つ目の等号には $A^\wedge$ が $A$ 上平坦であることを用いた
（More on Algebra, Lemma
\ref{more-algebra-lemma-pseudo-coherence-and-base-change-ext} を参照）。補題が従う。
\end{proof}

\noindent
$X$ を Noether スキームとする。$\mathcal{I} \subset \mathcal{O}_X$ を
準連接イデアル層とする。対象 $(\mathcal{F}_n)$ が
$\textit{Coh}(X, \mathcal{I})$ において {\it $\mathcal{I}$-冪トーション}、または
{\it $\mathcal{I}$ のある冪で零化される}とは、ある $c \geq 1$ が存在して、
$\mathcal{F}_n = \mathcal{F}_c$ がすべての $n \geq c$ に対して成り立つことをいう。
この場合、$(\mathcal{F}_n)$ は {\it $\mathcal{I}^c$ で零化される}ともいう。
$X = \Spec(A)$ がアフィンなら、補題 \ref{coherent-lemma-inverse-systems-affine} の同値を
通じて、これらの対象は有限 $A$-加群であって、$I$ のある冪、または
$I^c$ で零化されるものにちょうど対応する。

\begin{lemma}
\label{coherent-lemma-existence-easy}
$X$ を Noether スキームとし、$\mathcal{I} \subset \mathcal{O}_X$ を
準連接イデアル層とする。$\mathcal{G}$ を連接 $\mathcal{O}_X$-加群とする。
$(\mathcal{F}_n)$ を $\textit{Coh}(X, \mathcal{I})$ の対象とする。
\begin{enumerate}
\item $\alpha : (\mathcal{F}_n) \to \mathcal{G}^\wedge$ が、その核および余核を
$\mathcal{I}$ のある冪で零化されるような射であるとする。このとき、次を満たす
三つ組 $(\mathcal{F}, a, \beta)$ が一意な同型を除いて一意に存在する。
\begin{enumerate}
\item $\mathcal{F}$ は連接 $\mathcal{O}_X$-加群である。
\item $a : \mathcal{F} \to \mathcal{G}$ は $\mathcal{O}_X$-加群の射であり、
その核および余核は $\mathcal{I}$ のある冪で零化される。
\item $\beta : (\mathcal{F}_n) \to \mathcal{F}^\wedge$ は同型である。
\item $\alpha = a^\wedge \circ \beta$ である。
\end{enumerate}
\item $\alpha : \mathcal{G}^\wedge \to (\mathcal{F}_n)$ が、その核および余核を
$\mathcal{I}$ のある冪で零化されるような射であるとする。このとき、次を満たす
三つ組 $(\mathcal{F}, a, \beta)$ が一意な同型を除いて一意に存在する。
\begin{enumerate}
\item $\mathcal{F}$ は連接 $\mathcal{O}_X$-加群である。
\item $a : \mathcal{G} \to \mathcal{F}$ は $\mathcal{O}_X$-加群の射であり、
その核および余核は $\mathcal{I}$ のある冪で零化される。
\item $\beta : \mathcal{F}^\wedge \to (\mathcal{F}_n)$ は同型である。
\item $\alpha = \beta \circ a^\wedge$ である。
\end{enumerate}
\end{enumerate}
\end{lemma}

\begin{proof}
(1) を証明する。一意性により、$(\mathcal{F}, a, \beta)$ を $X$ 上 Zariski
局所的に構成すれば十分である。従って $X = \Spec(A)$ と仮定してよく、
$\mathcal{I}$ はイデアル $I \subset A$ に対応するとする。この状況では補題
\ref{coherent-lemma-inverse-systems-affine} を適用できる。$M'$ を有限 $A^\wedge$-加群であって
$(\mathcal{F}_n)$ に対応するものとする。$N$ を有限 $A$-加群であって
$\mathcal{G}$ に対応するものとする。このとき $\alpha$ は射
$$
\varphi : M' \longrightarrow N^\wedge
$$
に対応し、その核および余核は $I^t$ で零化される。ただし $t$ はある整数である。
$N^\wedge = N \otimes_A A^\wedge$ であったことを思い出す
（Algebra, Lemma \ref{algebra-lemma-completion-tensor}）。More on Algebra,
Lemma \ref{more-algebra-lemma-application-formal-glueing} により、$A$-加群の射
$\psi : M \to N$ であって核および余核が $I$-冪トーションであるものと、同型
$M \otimes_A A^\wedge = M'$ であって $\varphi$ と整合するものが存在する。
$N$ と $M'$ は有限加群なので、
$M$ は有限 $A$-加群であると結論できる（More on Algebra, Remark
\ref{more-algebra-remark-formal-glueing-algebras} を参照）。従って
$M \otimes_A A^\wedge = M^\wedge$ である。こうして得られた三つ組
$(M, N \to M, M^\wedge \to M')$ が一意な同型を除いて一意であることの確認は
省略する。

\medskip\noindent
(2) の証明もまったく同じなので省略する。
\end{proof}

\begin{lemma}
\label{coherent-lemma-torsion-hom-ext}
$X$ を Noether スキームとし、$\mathcal{I} \subset \mathcal{O}_X$ を
準連接イデアル層とする。$\textit{Coh}(X, \mathcal{I})$ の対象であって、
$\mathcal{I}$ のある冪で零化されるものはすべて
(\ref{coherent-equation-completion-functor}) の本質像に属する。さらに、
$\mathcal{F}$、$\mathcal{G}$ が $\textit{Coh}(\mathcal{O}_X)$ に属し、
$\mathcal{F}$ または $\mathcal{G}$ のいずれかが $\mathcal{I}$ のある冪で
零化されるなら、射
$$
\xymatrix{
\Hom_X(\mathcal{F}, \mathcal{G}) \ar[d] &
\Ext_X(\mathcal{F}, \mathcal{G}) \ar[d] \\
\Hom_{\textit{Coh}(X, \mathcal{I})}(\mathcal{F}^\wedge, \mathcal{G}^\wedge) &
\Ext_{\textit{Coh}(X, \mathcal{I})}(\mathcal{F}^\wedge, \mathcal{G}^\wedge)
}
$$
は同型である。
\end{lemma}

\begin{proof}
$(\mathcal{F}_n)$ を $\textit{Coh}(X, \mathcal{I})$ の対象であって、
$\mathcal{I}^c$ で零化されるものとする。ただし $c \geq 1$ とする。このとき
$\mathcal{F}_n \to \mathcal{F}_c$ は $n \geq c$ に対して同型である。従って
$\mathcal{F} = \mathcal{F}_c$ とおけば、
$\mathcal{F}^\wedge \cong (\mathcal{F}_n)$ であることが分かる。これで最初の
主張が示された。

\medskip\noindent
$\mathcal{F}$、$\mathcal{G}$ を $\textit{Coh}(\mathcal{O}_X)$ の対象とし、
$\mathcal{F}$ または $\mathcal{G}$ のいずれかが $\mathcal{I}^c$ で零化されるとする。
ただし $c \geq 1$ とする。このとき
$\mathcal{H} = \SheafHom_{\mathcal{O}_X}(\mathcal{G}, \mathcal{F})$ は連接
$\mathcal{O}_X$-加群であり、$\mathcal{I}^c$ で零化される。従って
$$
\Hom_X(\mathcal{G}, \mathcal{F}) =
H^0(X, \mathcal{H}) =
\lim H^0(X, \mathcal{H}/\mathcal{I}^n\mathcal{H}) =
\Mor_{\textit{Coh}(X, \mathcal{I})}
(\mathcal{G}^\wedge, \mathcal{F}^\wedge).
$$
補題 \ref{coherent-lemma-completion-internal-hom} を参照。これで準同型についての主張が
示された。

\medskip\noindent
記号 $\Ext$ は Homology, Section \ref{homology-section-extensions} で定義された
拡大を表す。$\Ext$ 上の写像の単射性は、$\Hom$ 上の写像の全単射性から直ちに
従う。全射性について、$\mathcal{F}$ が $I$ のある冪で零化されると仮定する。
このとき補題 \ref{coherent-lemma-existence-easy} の (1) により、
$\textit{Coh}(U, I\mathcal{O}_U)$ における任意の拡大
$$
0 \to \mathcal{G}^\wedge \to (\mathcal{E}_n) \to \mathcal{F}^\wedge \to 0
$$
に対して、射 $\mathcal{G}^\wedge \to (\mathcal{E}_n)$ は
$\mathcal{G} \to \mathcal{E}^\wedge$ と同型である。ここで、ある
$\mathcal{G} \to \mathcal{E}$ が $\textit{Coh}(\mathcal{O}_U)$ に属する。
他方の場合も同様である。
\end{proof}

\begin{lemma}
\label{coherent-lemma-finite-over-rees-algebra}
$X$ を Noether スキームとし、$\mathcal{I} \subset \mathcal{O}_X$ を
準連接イデアル層とする。$(\mathcal{F}_n)$ が
$\textit{Coh}(X, \mathcal{I})$ の対象なら、
$\bigoplus \Ker(\mathcal{F}_{n + 1} \to \mathcal{F}_n)$ は有限型次数付き準連接
$\bigoplus \mathcal{I}^n/\mathcal{I}^{n + 1}$-加群である。
\end{lemma}

\begin{proof}
問題は $X$ 上局所的なので、$X$ はアフィン、すなわち補題
\ref{coherent-lemma-inverse-systems-affine} の状況にあると仮定してよい。この場合、
$(\mathcal{F}_n)$ が有限 $A^\wedge$-加群 $M$ に対応するなら、
$\bigoplus \Ker(\mathcal{F}_{n + 1} \to \mathcal{F}_n)$ は
$\bigoplus I^nM/I^{n + 1}M$ に対応し、これは明らかに
$\bigoplus I^n/I^{n + 1}$ 上の有限加群である。
\end{proof}

\begin{lemma}
\label{coherent-lemma-inverse-systems-pullback}
$f : X \to Y$ を Noether スキームの射とする。
$\mathcal{J} \subset \mathcal{O}_Y$ を準連接イデアル層とし、
$\mathcal{I} = f^{-1}\mathcal{J} \mathcal{O}_X$ とおく。このとき右完全関手
$$
f^* : \textit{Coh}(Y, \mathcal{J}) \longrightarrow \textit{Coh}(X, \mathcal{I})
$$
があり、$(\mathcal{G}_n)$ を $(f^*\mathcal{G}_n)$ へ送る。$f$ が平坦なら、
$f^*$ は完全関手である。
\end{lemma}

\begin{proof}
$f^* : \textit{Coh}(\mathcal{O}_Y) \to \textit{Coh}(\mathcal{O}_X)$ は右完全なので
$$
f^*\mathcal{G}_n =
f^*(\mathcal{G}_{n + 1}/\mathcal{I}^n\mathcal{G}_{n + 1}) =
f^*\mathcal{G}_{n + 1}/f^{-1}\mathcal{I}^nf^*\mathcal{G}_{n + 1} =
f^*\mathcal{G}_{n + 1}/\mathcal{J}^nf^*\mathcal{G}_{n + 1}
$$
であり、従って系の引き戻しは系である。補題
\ref{coherent-lemma-inverse-systems-abelian} の証明における余核の構成から、
$f^* : \textit{Coh}(Y, \mathcal{J}) \to \textit{Coh}(X, \mathcal{I})$ は常に
右完全であることが分かる。$f$ が平坦なら、
$f^* : \textit{Coh}(\mathcal{O}_Y) \to \textit{Coh}(\mathcal{O}_X)$ は完全関手である。
補題 \ref{coherent-lemma-inverse-systems-abelian} の証明における核の構成から、この場合
$f^* : \textit{Coh}(Y, \mathcal{J}) \to \textit{Coh}(X, \mathcal{I})$ は核も核へ
写すことが従う。
\end{proof}

\begin{lemma}
\label{coherent-lemma-inverse-systems-pullback-equivalence}
$f : X' \to X$ を Noether スキームの射とする。$Z \subset X$ を閉部分スキームとし、
$Z' = f^{-1}Z$ をそのスキーム論的逆像と表す。
$\mathcal{I} \subset \mathcal{O}_X$、
$\mathcal{I}' \subset \mathcal{O}_{X'}$ を対応する準連接イデアル層とする。
$f$ が平坦で、誘導される射 $Z' \to Z$ が同型なら、引き戻し関手
$f^* : \textit{Coh}(X, \mathcal{I}) \to \textit{Coh}(X', \mathcal{I}')$
（補題 \ref{coherent-lemma-inverse-systems-pullback}）は同値である。
\end{lemma}

\begin{proof}
$X$ と $X'$ がアフィンなら、これは More on Algebra, Lemma
\ref{more-algebra-lemma-neighbourhood-equivalence} から直ちに従う。一般の場合を
証明するため、$Z_n \subset X$、$Z'_n \subset X'$ を第 $n$ 無限小近傍、すなわち
それぞれ $Z$、$Z'$ のものとする。誘導される射 $Z_n \to Z'_n$ は台となる
位相空間上の同相である。一方、
$z' \in Z'$ が $z \in Z$ へ写るなら、環準同型
$\mathcal{O}_{X, z} \to \mathcal{O}_{X', z'}$ は平坦であり、同型
$\mathcal{O}_{X, z}/\mathcal{I}_z \to \mathcal{O}_{X', z'}/\mathcal{I}'_{z'}$
を誘導する。従って、たとえば More on Algebra, Lemma
\ref{more-algebra-lemma-neighbourhood-isomorphism} により、同型
$\mathcal{O}_{X, z}/\mathcal{I}_z^n \to
\mathcal{O}_{X', z'}/(\mathcal{I}'_{z'})^n$
をすべての $n \geq 1$ に対して誘導する。ゆえに $Z'_n \to Z_n$ はスキームの
同型である。従って $f^*$ は、連接 $\mathcal{O}_X$-加群であって
$\mathcal{I}^n$ で零化されるものの圏と、連接 $\mathcal{O}_{X'}$-加群であって
$(\mathcal{I}')^n$ で零化されるものの圏との同値を誘導する
（補題 \ref{coherent-lemma-i-star-equivalence} を参照）。これは明らかに補題を導く。
\end{proof}

\begin{lemma}
\label{coherent-lemma-inverse-systems-ideals-equivalence}
$X$ を Noether スキームとする。
$\mathcal{I}, \mathcal{J} \subset \mathcal{O}_X$ を準連接イデアル層とする。
$V(\mathcal{I}) = V(\mathcal{J})$ が $X$ の同じ閉部分集合なら、
$\textit{Coh}(X, \mathcal{I})$ と $\textit{Coh}(X, \mathcal{J})$ は同値である。
\end{lemma}

\begin{proof}
まず、$X = \Spec(A)$ がアフィンであると仮定する。
$I, J \subset A$ を $\mathcal{I}, \mathcal{J}$ に対応するイデアルとする。
$V(I) = V(J)$ なら、Zariski 位相の初等的性質により、包含
$I^c \subset J$ および $J^d \subset I$ が、ある $c, d \geq 1$ に対して成り立つ
（Algebra, Section \ref{algebra-section-spectrum-ring} および Lemma
\ref{algebra-lemma-Noetherian-power} を参照）。従って $I$-進および $J$-進という
$A$ の二つの完備化は一致する（Algebra, Lemma
\ref{algebra-lemma-change-ideal-completion} を参照）。ゆえに、この場合の同値は
補題 \ref{coherent-lemma-inverse-systems-affine} から従う。

\medskip\noindent
一般に、上で述べたことと $X$ が準コンパクトであることを用いて、
$c, d \geq 1$ を $\mathcal{I}^c \subset \mathcal{J}$ および
$\mathcal{J}^d \subset \mathcal{I}$ が成り立つように選べる。そこで対象
$(\mathcal{F}_n)$ が $\textit{Coh}(X, \mathcal{I})$ において与えられたとき、逆系
$$
(\mathcal{F}_{cn}/\mathcal{J}^n\mathcal{F}_{cn})
$$
は $\textit{Coh}(X, \mathcal{J})$ に属すると主張する。これはアフィン開被覆の
各要素上で判定できるので、詳細は省略する。同様に、
$\textit{Coh}(X, \mathcal{I})$ の対象を
$\textit{Coh}(X, \mathcal{J})$ の対象から出発して構成できる。これらの構成が互いに擬逆な
関手を定めることの確認は省略する。
\end{proof}





\section{Grothendieck の存在定理 I}
\label{coherent-section-existence}

\noindent
この節では、射影的な場合の Grothendieck の存在定理を論じる。Section
\ref{coherent-section-coherent-formal} で整備した連接形式加群の概念を用いる。
形式スキームに慣れた読者には、\cite{EGA} にある定理の主張と証明を読むことを
勧める。

\begin{lemma}
\label{coherent-lemma-fully-faithful}
$A$ をイデアル $I$ に関して完備な Noether 環とする。
$f : X \to \Spec(A)$ を固有射とする。
$\mathcal{I} = I\mathcal{O}_X$ とおく。このとき関手
(\ref{coherent-equation-completion-functor}) は充満忠実である。
\end{lemma}

\begin{proof}
$\mathcal{F}$、$\mathcal{G}$ を連接 $\mathcal{O}_X$-加群とする。このとき
$\mathcal{H} = \SheafHom_{\mathcal{O}_X}(\mathcal{G}, \mathcal{F})$ は連接
$\mathcal{O}_X$-加群である。Modules, Lemma
\ref{modules-lemma-internal-hom-locally-kernel-direct-sum} を参照。補題
\ref{coherent-lemma-completion-internal-hom} により、写像
$$
\lim_n H^0(X, \mathcal{H}/\mathcal{I}^n\mathcal{H})
\to
\Mor_{\textit{Coh}(X, \mathcal{I})}
(\mathcal{G}^\wedge, \mathcal{F}^\wedge)
$$
は全単射である。従って関手 (\ref{coherent-equation-completion-functor}) の充満忠実性は、
連接層 $\mathcal{H}$ に対する形式関数の定理（補題
\ref{coherent-lemma-spell-out-theorem-formal-functions}）から従う。
\end{proof}

\begin{lemma}
\label{coherent-lemma-vanishing-projective}
$A$ を Noether 環とし、$I \subset A$ をイデアルとする。
$f : X \to \Spec(A)$ を固有射とし、$\mathcal{L}$ を $f$-豊富な可逆層とする。
$\mathcal{I} = I\mathcal{O}_X$ とおく。$(\mathcal{F}_n)$ を
$\textit{Coh}(X, \mathcal{I})$ の対象とする。このとき整数 $d_0$ が存在して
$$
H^1(X, \Ker(\mathcal{F}_{n + 1} \to \mathcal{F}_n)
\otimes \mathcal{L}^{\otimes d} )
= 0
$$
がすべての $n \geq 0$ およびすべての $d \geq d_0$ に対して成り立つ。
\end{lemma}

\begin{proof}
$B = \bigoplus I^n/I^{n + 1}$ および
$\mathcal{B} = \bigoplus \mathcal{I}^n/\mathcal{I}^{n + 1} = f^*\widetilde{B}$ とおく。
補題 \ref{coherent-lemma-finite-over-rees-algebra} により、次数付き準連接
$\mathcal{B}$-加群
$\mathcal{G} = \bigoplus \Ker(\mathcal{F}_{n + 1} \to \mathcal{F}_n)$ は有限型である。
従って補題 \ref{coherent-lemma-graded-finiteness} の (2) から主張が従う。
\end{proof}

\begin{lemma}
\label{coherent-lemma-existence-projective}
$A$ をイデアル $I$ に関して完備な Noether 環とする。
$f : X \to \Spec(A)$ を射影的射とする。
$\mathcal{I} = I\mathcal{O}_X$ とおく。このとき関手
(\ref{coherent-equation-completion-functor}) は同値である。
\end{lemma}

\begin{proof}
関手 (\ref{coherent-equation-completion-functor}) が充満忠実であることは、補題
\ref{coherent-lemma-fully-faithful} ですでに見た。従ってこの関手が本質的全射であることを
示せば十分である。

\medskip\noindent
まず、すべての対象 $(\mathcal{F}_n)$ が $\textit{Coh}(X, \mathcal{I})$ において、関手
(\ref{coherent-equation-completion-functor}) の像に属する対象の商であることを示す。
$\mathcal{L}$ を $f$-豊富な可逆層であって $X$ 上のものとする。補題
\ref{coherent-lemma-vanishing-projective} のとおりに $d_0$ を選ぶ。
$d \geq d_0$ を、$\mathcal{F}_1 \otimes \mathcal{L}^{\otimes d}$ がある切断
$s_{1, 1}, \ldots, s_{t, 1}$ によって大域的に生成されるように選ぶ。系の遷移写像
$$
H^0(X, \mathcal{F}_{n + 1} \otimes \mathcal{L}^{\otimes d})
\longrightarrow
H^0(X, \mathcal{F}_n \otimes \mathcal{L}^{\otimes d})
$$
は $H^1$ の消滅により全射なので、$s_{1, 1}, \ldots, s_{t, 1}$ を整合的な
大域切断系 $s_{1, n}, \ldots, s_{t, n}$ へ持ち上げられる。これらは
$\mathcal{F}_n \otimes \mathcal{L}^{\otimes d}$ の切断である。これらにより整合的な
射の系
$$
(s_{1, n}, \ldots, s_{t, n}) :
(\mathcal{L}^{\otimes -d})^{\oplus t} \longrightarrow \mathcal{F}_n
$$
が定まる。補題 \ref{coherent-lemma-inverse-systems-surjective} を用いると、所望の全射
$$
\left((\mathcal{L}^{\otimes -d})^{\oplus t}\right)^\wedge
\longrightarrow
(\mathcal{F}_n)
$$
を得る。

\medskip\noindent
前段落の結果と $\textit{Coh}(X, \mathcal{I})$ が Abel 圏であること
（補題 \ref{coherent-lemma-inverse-systems-abelian}）により、
$\textit{Coh}(X, \mathcal{I})$ のすべての対象は、
$\textit{Coh}(\mathcal{O}_X)$ に由来する対象の間の射の余核である。関手
(\ref{coherent-equation-completion-functor}) は補題 \ref{coherent-lemma-fully-faithful} および
\ref{coherent-lemma-exact} により充満忠実かつ完全なので、結論を得る。
\end{proof}







\section{Grothendieck の存在定理 II}
\label{coherent-section-existence-proper}

\noindent
この節では、固有な場合の Grothendieck の存在定理を論じる。その主張と証明を
与える前に、Section \ref{coherent-section-coherent-formal} で導入した連接形式加群の圏
$\textit{Coh}(X, \mathcal{I})$ に関する理論をもう少し整備する必要がある。

\begin{remark}
\label{coherent-remark-inverse-systems-kernel-cokernel-annihilated-by}
$X$ を Noether スキームとし、
$\mathcal{I}, \mathcal{K} \subset \mathcal{O}_X$ を準連接イデアル層とする。
$\alpha : (\mathcal{F}_n) \to (\mathcal{G}_n)$ を
$\textit{Coh}(X, \mathcal{I})$ の射とする。上のとおりのアフィン開集合
$\Spec(A) = U \subset X$ が与えられ、$\mathcal{I}|_U, \mathcal{K}|_U$ が
イデアル $I, K \subset A$ に対応するとき、$\alpha_U : M \to N$ を、有限
$A^\wedge$-加群の射であって補題 \ref{coherent-lemma-inverse-systems-affine} を通じて
$\alpha|_U$ に対応するものと表す。次の条件は同値であると主張する。
\begin{enumerate}
\item ある整数 $t \geq 1$ が存在して、$\Ker(\alpha_n)$ および
$\Coker(\alpha_n)$ が $\mathcal{K}^t$ で、すべての $n \geq 1$ に対して零化される。
\item 上のとおりの任意のアフィン開集合 $\Spec(A) = U \subset X$ に対して、
加群 $\Ker(\alpha_U)$ および $\Coker(\alpha_U)$ は $K^t$ で零化される。
ここで $t \geq 1$ はある整数である。
\item 有限アフィン開被覆 $X = \bigcup U_i$ が存在し、(2) の結論が
$\alpha_{U_i}$ に対して成り立つ。
\end{enumerate}
これらの同値な条件が成り立つとき、$\alpha$ を
{\it 核および余核が $\mathcal{K}$ のある冪で零化される射}という。
同値性を示すため、次の可換代数の事実を用いる。$A$-加群の完全列
$$
0 \to T \to M \to N \to Q \to 0
$$
が与えられ、$T$ と $Q$ は $K^t$ で零化されると仮定する。ここで
$K \subset A$ はあるイデアルである。このとき、任意の $f, g \in K^t$ に対して標準的な写像
$"fg": N \to M$ であって、合成 $M \to N \to M$ が $fg$ 倍と等しくなるものが
存在する。実際、$y \in N$ に対し、$x \in M$ を、$fy$ へ $N$ で写るように選び、
$"fg"(y) = gx$ とおけばよい。従って
$\Ker(M/JM \to N/JN)$ および $\Coker(M/JM \to N/JN)$ は $K^{2t}$ で零化される
ことが明らかである。これは任意のイデアル $J \subset A$ に対して成り立つ。

\medskip\noindent
この可換代数の事実を $\alpha_{U_i}$ と $J = I^n$ に適用すると、(3) から (1) が
従う。逆に (1) が成り立ち、$M \to N$ が $\alpha_U$ と等しいと仮定する。
このときある $t \geq 1$ が存在し、
$\Ker(M/I^nM \to N/I^nN)$ および $\Coker(M/I^nM \to N/I^nN)$ は $K^t$ で、
すべての $n$ に対して零化される。写像 $"fg" : N/I^nN \to M/I^nM$ を得て、これらは極限において
$N \to M$ を誘導する。これは $N$ と $M$ が $I$-進完備だからである。合成
$N \to M \to N$ は $fg$ 倍なので、$fg$ は $T$ と $Q$ を零化すると結論できる。
言い換えれば、$T$ と $Q$ は所望のとおり $K^{2t}$ で零化される。
\end{remark}

\begin{lemma}
\label{coherent-lemma-existence-tricky}
$X$ を Noether スキームとする。
$\mathcal{I}, \mathcal{K} \subset \mathcal{O}_X$ を準連接イデアル層とする。
$X_e \subset X$ を $\mathcal{K}^e$ で切り出される閉部分スキームとする。
$\mathcal{I}_e = \mathcal{I}\mathcal{O}_{X_e}$ とおく。
$(\mathcal{F}_n)$ を $\textit{Coh}(X, \mathcal{I})$ の対象とする。次を仮定する。
\begin{enumerate}
\item 関手
$\textit{Coh}(\mathcal{O}_{X_e}) \to \textit{Coh}(X_e, \mathcal{I}_e)$
はすべての $e \geq 1$ に対して同値である。
\item 連接層 $\mathcal{H}$ であって $X$ 上のものと射
$\alpha : (\mathcal{F}_n) \to \mathcal{H}^\wedge$ が存在し、その核および余核は
$\mathcal{K}$ のある冪で零化される。
\end{enumerate}
このとき $(\mathcal{F}_n)$ は関手 (\ref{coherent-equation-completion-functor}) の本質像に
属する。
\end{lemma}

\begin{proof}
この証明では、閉埋め込み $i : Z \to X$ に対して関手 $i_*$ が、$Z$ 上の連接加群の
圏と、$X$ 上の連接加群であって $Z$ のイデアル層で零化されるものの圏との同値を与えることを
断りなく用いる。補題 \ref{coherent-lemma-i-star-equivalence} を参照。特に、
$\textit{Coh}(\mathcal{O}_{X_e})$ を連接 $\mathcal{O}_X$-加群であって
$\mathcal{K}^e$ で零化されるものの圏と同一視し、$\textit{Coh}(X_e, \mathcal{I}_e)$ を
$\textit{Coh}(X, \mathcal{I})$ の対象であって $\mathcal{K}^e$ で零化される
ものからなる充満部分圏と同一視できる。さらに (1) により、これら二つの圏は
完備化関手 (\ref{coherent-equation-completion-functor}) のもとで同値である。

\medskip\noindent
この同値を適用すると、連接 $\mathcal{O}_X$-加群 $\mathcal{G}_e$ であって
$\mathcal{K}^e$ で零化され、系 $(\mathcal{F}_n/\mathcal{K}^e\mathcal{F}_n)$ に
$\textit{Coh}(X, \mathcal{I})$ において対応するものを得る。写像
$\mathcal{F}_n/\mathcal{K}^{e + 1}\mathcal{F}_n \to
\mathcal{F}_n/\mathcal{K}^e\mathcal{F}_n$ は標準的な写像
$\mathcal{G}_{e + 1} \to \mathcal{G}_e$ に対応し、それらは同型
$\mathcal{G}_{e + 1}/\mathcal{K}^e\mathcal{G}_{e + 1} \to \mathcal{G}_e$ を誘導する。
従って $(\mathcal{G}_e)$ は $\textit{Coh}(X, \mathcal{K})$ の対象である。
写像 $\alpha$ は写像の系
$$
\mathcal{F}_n/\mathcal{K}^e\mathcal{F}_n
\longrightarrow
\mathcal{H}/(\mathcal{I}^n + \mathcal{K}^e)\mathcal{H}
$$
を誘導し、従って写像
$\mathcal{G}_e \to \mathcal{H}/\mathcal{K}^e\mathcal{H}$ を誘導する
（再び圏同値による）。仮定 (2) により存在する整数 $t \geq 1$ を、
$\mathcal{K}^t$ が、すべての写像
$\mathcal{F}_n \to \mathcal{H}/\mathcal{I}^n\mathcal{H}$ の核および余核を
零化するようにとる。このとき
$\mathcal{K}^{2t}$ は写像
$\mathcal{F}_n/\mathcal{K}^e\mathcal{F}_n \to
\mathcal{H}/(\mathcal{I}^n + \mathcal{K}^e)\mathcal{H}$ の核および余核を零化する。
Remark \ref{coherent-remark-inverse-systems-kernel-cokernel-annihilated-by} を参照。
従って、写像
$$
\mathcal{G}_e
\longrightarrow
\mathcal{H}/\mathcal{K}^e\mathcal{H},
$$
の核および余核を $\mathcal{K}^{4t}$ が零化すると結論する。Remark
\ref{coherent-remark-inverse-systems-kernel-cokernel-annihilated-by} を参照。補題
\ref{coherent-lemma-existence-easy} を適用して、連接 $\mathcal{O}_X$-加群
$\mathcal{F}$、射 $a : \mathcal{F} \to \mathcal{H}$、および同型
$\beta : (\mathcal{G}_e) \to (\mathcal{F}/\mathcal{K}^e\mathcal{F})$ を得る。
最後の同型は $\textit{Coh}(X, \mathcal{K})$ におけるものである。
逆向きにたどると、与えられた $n$ に対し、三つ組
$(\mathcal{F}/\mathcal{I}^n\mathcal{F}, a \bmod \mathcal{I}^n, \beta
\bmod \mathcal{I}^n)$ は、射 $\alpha_n \bmod \mathcal{K}^e :
(\mathcal{F}_n/\mathcal{K}^e\mathcal{F}_n) \to
(\mathcal{H}/(\mathcal{I}^n + \mathcal{K}^e)\mathcal{H})$ に対する補題中の
三つ組となる。この射は $\textit{Coh}(X, \mathcal{K})$ に属する。従って補題
\ref{coherent-lemma-existence-easy} の一意性から、現れるすべての
射と整合する標準的同型
$\mathcal{F}/\mathcal{I}^n\mathcal{F} \to \mathcal{F}_n$ を得る。これで補題の
証明が終わる。
\end{proof}

\begin{lemma}
\label{coherent-lemma-inverse-systems-push-pull}
$Y$ を Noether スキームとする。
$\mathcal{J}, \mathcal{K} \subset \mathcal{O}_Y$ を準連接イデアル層とする。
$f : X \to Y$ を、$V = Y \setminus V(\mathcal{K})$ 上で同型となる固有射とする。
$\mathcal{I} = f^{-1}\mathcal{J} \mathcal{O}_X$ とおく。
$(\mathcal{G}_n)$ を $\textit{Coh}(Y, \mathcal{J})$ の対象とし、$\mathcal{F}$ を
連接 $\mathcal{O}_X$-加群とし、$\beta : (f^*\mathcal{G}_n)  \to \mathcal{F}^\wedge$
を $\textit{Coh}(X, \mathcal{I})$ における同型とする。このとき射
$$
\alpha :
(\mathcal{G}_n)
\longrightarrow
(f_*\mathcal{F})^\wedge
$$
であって $\textit{Coh}(Y, \mathcal{J})$ に属し、その核および余核が
$\mathcal{K}$ のある冪で零化されるものが存在する。
\end{lemma}

\begin{proof}
$f$ は固有射なので、$f_*\mathcal{F}$ は連接 $\mathcal{O}_Y$-加群である
（Proposition \ref{coherent-proposition-proper-pushforward-coherent}）。従って補題の主張は
意味をもつ。合成
$$
\gamma_n : \mathcal{G}_n \to
f_*f^*\mathcal{G}_n \to
f_*(\mathcal{F}/\mathcal{I}^n\mathcal{F}).
$$
を考える。ここで第一の射は随伴射であり、第二の射は $f_*\beta_n$ である。
補題のとおりの一意な $\alpha$ であって、合成
$$
\mathcal{G}_n \xrightarrow{\alpha_n}
f_*\mathcal{F}/\mathcal{J}^nf_*\mathcal{F} \to
f_*(\mathcal{F}/\mathcal{I}^n\mathcal{F})
$$
が $\gamma_n$ に、すべての $n$ に対して等しくなるものが存在すると主張する。
一意性により、$Y = \Spec(B)$ がアフィンであると仮定してよい。
$J \subset B$ をイデアル $\mathcal{J}$ に対応するものとする。次のようにおく。
$$
M_n = H^0(X, \mathcal{F}/\mathcal{I}^n\mathcal{F})
\quad\text{and}\quad
M = H^0(X, \mathcal{F})
$$
補題 \ref{coherent-lemma-ML-cohomology-powers-ideal} および定理
\ref{coherent-theorem-formal-functions} により、加群 $M_n$ の逆極限は完備化
$M^\wedge = \lim M/J^nM$ に等しい。
$N_n = H^0(Y, \mathcal{G}_n)$ および $N = \lim N_n$ とおく。補題
\ref{coherent-lemma-inverse-systems-affine} の圏同値を通じて、有限 $B^\wedge$-加群
$N$ および $M^\wedge$ は $(\mathcal{G}_n)$ および
$f_*\mathcal{F}^\wedge$ に対応する。従って $\alpha$ は
$\textit{Coh}(Y, \mathcal{J})$ の射であって、有限 $B^\wedge$-加群の準同型
$$
\lim \gamma_n : N = \lim_n N_n \longrightarrow \lim M_n = M^\wedge
$$
に対応するものでなければならない。

\medskip\noindent
さらに $\alpha$ の核および余核が $\mathcal{K}$ のある冪で零化されることを示す
必要がある。$Y' = \Spec(B^\wedge)$ および $X' = Y' \times_Y X$ とおく。
$\mathcal{K}'$、$\mathcal{J}'$、$\mathcal{G}'_n$ および
$\mathcal{I}'$、$\mathcal{F}'$ を、それぞれ $\mathcal{K}$、$\mathcal{J}$、
$\mathcal{G}_n$ および $\mathcal{I}$、$\mathcal{F}$ の $Y'$ および $X'$ への
引き戻しとする。射影射 $f' : X' \to Y'$ は $f$ の、$Y' \to Y$ による
基底変換である。$Y' \to Y$ はスキームの平坦射であることに注意する。これは
$B \to B^\wedge$ が Algebra, Lemma
\ref{algebra-lemma-completion-flat} により平坦だからである。従って補題
\ref{coherent-lemma-flat-base-change-cohomology} により、$f'_*\mathcal{F}'$、それぞれ
$f'_*(f')^*\mathcal{G}_n'$ は、$f_*\mathcal{F}$、それぞれ
$f_*f^*\mathcal{G}_n$ の $Y'$ への引き戻しである。構成の一意性により、
$\alpha$ の $Y'$ への引き戻しは、対応する写像 $\alpha'$、すなわち $Y'$ 上の
状況について構成したものとなる。さらに、$\alpha$ の核および余核が $\mathcal{K}^t$ で
零化されることを判定するには、$\alpha'$ の核および余核が
$(\mathcal{K}')^t$ で零化されることを判定すれば十分である。実際、このためには
写像 $\alpha_n$ および $\alpha'_n$ の核および余核について判定すればよく
（Remark \ref{coherent-remark-inverse-systems-kernel-cokernel-annihilated-by} を参照）、
環準同型 $B \to B^\wedge$ は $J^n$ で零化される加群と $(J')^n$ で零化される
加群との圏同値を誘導する（More on Algebra, Lemma
\ref{more-algebra-lemma-neighbourhood-equivalence} を参照）。従って $B$ は $J$ に
関して完備であると仮定してよい。

\medskip\noindent
$Y = \Spec(B)$ がアフィンであり、$\mathcal{J}$ がイデアル $J \subset B$ に
対応し、$B$ が $J$ に関して完備であると仮定する。この場合
$(\mathcal{G}_n)$ は関手
$\textit{Coh}(\mathcal{O}_Y) \to \textit{Coh}(Y, \mathcal{J})$ の本質像に属する。
そこで $\mathcal{G}$ を連接 $\mathcal{O}_Y$-加群であって
$(\mathcal{G}_n) = \mathcal{G}^\wedge$ を満たすものとする。
$f^*(\mathcal{G}^\wedge) = (f^*\mathcal{G})^\wedge$ であることに注意する。
従って補題 \ref{coherent-lemma-fully-faithful} により、$\beta$ は同型
$b : f^*\mathcal{G} \to \mathcal{F}$ から得られ、$\alpha$ は
$$
\mathcal{G} \to f_*f^*\mathcal{G} \cong f_*\mathcal{F}
$$
に完備化関手を適用して得られる。従って、随伴射
$c : \mathcal{G} \to f_*f^*\mathcal{G}$ の核および余核が $\mathcal{K}$ のある冪で
零化されることを確かめようとしている。しかし制限
$f|_{f^{-1}(V)} : f^{-1}(V) \to V$ は同型なので、$c|_V$ は同型であることが分かる。
従って連接層 $\Ker(c)$ および $\Coker(c)$ は $V(\mathcal{K})$ 上に台をもち、
ゆえに所望のとおり $\mathcal{K}$ のある冪で零化される
（補題 \ref{coherent-lemma-power-ideal-kills-sheaf}）。
\end{proof}

\noindent
次の命題は、実際に最もよく用いられる Grothendieck の存在定理の形である。

\begin{proposition}
\label{coherent-proposition-existence-proper}
$A$ をイデアル $I$ に関して完備な Noether 環とする。
$f : X \to \Spec(A)$ をスキームの固有射とする。
$\mathcal{I} = I\mathcal{O}_X$ とおく。このとき関手
(\ref{coherent-equation-completion-functor}) は同値である。
\end{proposition}

\begin{proof}
関手 (\ref{coherent-equation-completion-functor}) が充満忠実であることは、補題
\ref{coherent-lemma-fully-faithful} ですでに見た。従ってこの関手が本質的全射であることを
示せば十分である。

\medskip\noindent
$\Xi$ を、次の性質をもつ準連接イデアル層 $\mathcal{K} \subset \mathcal{O}_X$ の
集まりとする：すべての対象 $(\mathcal{F}_n)$ であって $\mathcal{K}$ で零化されるものは本質像に
属する。$(0)$ が $\Xi$ に属することを示したい。そうでなければ、$X$ は Noether
なので、極大の準連接イデアル層 $\mathcal{K}$ であって $\Xi$ に属さないものが存在する
（補題 \ref{coherent-lemma-acc-coherent} を参照）。$X$ を $X$ の閉部分スキームであって
$\mathcal{K}$ に対応するものによって置き換えると、零でないすべての $\mathcal{K}$ が $\Xi$ に属すると
仮定してよい。（ここでは、$\mathcal{K}$ で零化される連接加群と、
$\mathcal{K}$ に対応する閉部分スキーム上の連接加群との対応を用いる。補題
\ref{coherent-lemma-i-star-equivalence} を参照。）$(\mathcal{F}_n)$ を
$\textit{Coh}(X, \mathcal{I})$ の対象とする。この対象が関手
(\ref{coherent-equation-completion-functor}) の本質像に属することを示し、それによって
命題の証明を終える。

\medskip\noindent
Chow の補題（補題 \ref{coherent-lemma-chow-Noetherian}）を適用し、固有全射
$f : X' \to X$ であって、稠密開集合 $U \subset X$ 上で同型となり、$X'$ が
$A$ 上射影的となるものを見つける。$\mathcal{K}$ を被約な補集合 $X \setminus U$ を
切り出す準連接イデアル層とする。Grothendieck の存在定理の射影的な場合
（補題 \ref{coherent-lemma-existence-projective}）により、連接加群 $\mathcal{F}'$ であって
$X'$ 上にあり、$(\mathcal{F}')^\wedge \cong (f^*\mathcal{F}_n)$ を
満たすものが存在する。Proposition
\ref{coherent-proposition-proper-pushforward-coherent} により、$\mathcal{O}_X$-加群
$\mathcal{H} = f_*\mathcal{F}'$ は連接であり、補題
\ref{coherent-lemma-inverse-systems-push-pull} により射
$(\mathcal{F}_n) \to \mathcal{H}^\wedge$ であって $\textit{Coh}(X, \mathcal{I})$ に属し、その核および余核が
$\mathcal{K}$ のある冪で零化されるものが存在する。すべての冪
$\mathcal{K}^e$ は $\Xi$ に属するので、関手
(\ref{coherent-equation-completion-functor}) は閉部分スキーム
$X_e = V(\mathcal{K}^e)$ に対して同値である。補題
\ref{coherent-lemma-existence-tricky} により結論を得る。
\end{proof}





\section{基底上固有であること}
\label{coherent-section-proper-over-base}

\noindent
この短い節では、基底上固有な閉集合の有用な性質と、
台が基底上固有である有限型準連接加群の有用な性質をいくつか指摘する。

\begin{lemma}
\label{coherent-lemma-closed-proper-over-base}
$f : X \to S$ を局所有限型なスキームの射とする。
$Z \subset X$ を閉集合とする。次の条件は同値である。
\begin{enumerate}
\item 射 $Z \to S$ は、$Z$ に誘導される被約閉部分スキーム構造を入れると固有である
（Schemes, Definition \ref{schemes-definition-reduced-induced-scheme}）。
\item $Z$ 上のある閉部分スキーム構造について、射 $Z \to S$ は固有である。
\item $Z$ 上の任意の閉部分スキーム構造について、射 $Z \to S$ は固有である。
\end{enumerate}
\end{lemma}

\begin{proof}
(3) $\Rightarrow$ (1) および (1) $\Rightarrow$ (2) は直ちに従う。
従って (2) が (3) を含意することを示せば十分である。
この事実の証明を読者自身で見つけることを勧める。
$Z'$ と $Z''$ を $Z$ 上の閉部分スキーム構造であって、
$Z' \to S$ が固有であるものとする。$Z'' \to S$ が固有であることを示さなければならない。
$Z''' = Z' \cup Z''$ をスキーム論的和とする。Morphisms, Definition
\ref{morphisms-definition-scheme-theoretic-intersection-union} を参照。
このとき $Z'''$ は $Z$ 上の別の閉部分スキーム構造である。
これは、例えば Morphisms, Lemma
\ref{morphisms-lemma-scheme-theoretic-union} におけるスキーム論的和の記述から従う。
$Z'' \to Z'''$ は閉埋入なので、$Z''' \to S$ が固有であることを示せば十分である
（Morphisms, Lemmas \ref{morphisms-lemma-closed-immersion-proper} および
\ref{morphisms-lemma-composition-proper} を参照）。
射 $Z' \to Z'''$ は全単射な閉埋入であり、特に全射かつ普遍閉である。
そこで $Z' \to S$ が分離的であることから $Z''' \to S$ が分離的であることが従う。
Morphisms, Lemma \ref{morphisms-lemma-image-universally-closed-separated} を参照。
さらに、$Z''' \to S$ は、$X \to S$ が局所有限型なので局所有限型である
（Morphisms, Lemmas \ref{morphisms-lemma-immersion-locally-finite-type} および
\ref{morphisms-lemma-composition-finite-type}）。
$Z' \to S$ は準コンパクトであり、$Z' \to Z'''$ は同相写像なので、
$Z''' \to S$ は準コンパクトである。
最後に、$Z' \to S$ は普遍閉なので、Morphisms, Lemma
\ref{morphisms-lemma-image-proper-is-proper} により $Z''' \to S$ も普遍閉である。
これで証明が完了する。
\end{proof}

\begin{definition}
\label{coherent-definition-proper-over-base}
$f : X \to S$ を局所有限型なスキームの射とする。
$Z \subset X$ を閉集合とする。補題
\ref{coherent-lemma-closed-proper-over-base} の同値な条件が満たされるとき、
{\it $Z$ は $S$ 上固有である}という。
\end{definition}

\noindent
上の定義で用いた補題は、射 $f : X \to S$ が局所有限型でなければ偽である。
従って、$f$ が局所有限型でないときにはこの用語を使わないよう読者に強く勧める。

\begin{lemma}
\label{coherent-lemma-closed-closed-proper-over-base}
$f : X \to S$ を局所有限型なスキームの射とする。
$Y \subset Z \subset X$ を閉集合とする。$Z$ が $S$ 上固有ならば、
$Y$ についても同じことが成り立つ。
\end{lemma}

\begin{proof}
省略する。
\end{proof}

\begin{lemma}
\label{coherent-lemma-base-change-closed-proper-over-base}
次のスキームのカルテジアン図式を考える：
$$
\xymatrix{
X' \ar[d]_{f'} \ar[r]_{g'} & X \ar[d]^f \\
S' \ar[r]^g & S
}
$$
ここで $f$ は局所有限型である。$Z$ が $X$ の閉集合で $S$ 上固有ならば、
$(g')^{-1}(Z)$ は $X'$ の閉集合で $S'$ 上固有である。
\end{lemma}

\begin{proof}
Morphisms, Lemma \ref{morphisms-lemma-base-change-finite-type} により
$f'$ は局所有限型なので、この主張には意味がある。
$Z$ に誘導される被約閉部分スキーム構造を入れる。
スキーム論的逆像
$Z' = (g')^{-1}(Z)$ と書く
（Schemes, Definition \ref{schemes-definition-inverse-image-closed-subscheme}）。
このとき
$Z' = X' \times_X Z = (S' \times_S X) \times_X Z = S' \times_S Z$
は $S'$ 上固有である。実際、これは $Z$ の $S$ 上の基底変換である
（Morphisms, Lemma \ref{morphisms-lemma-base-change-proper}）。
\end{proof}

\begin{lemma}
\label{coherent-lemma-functoriality-closed-proper-over-base}
$S$ をスキームとする。$f : X \to Y$ を、$S$ 上局所有限型なスキームの射とする。
\begin{enumerate}
\item $Y$ が $S$ 上分離的であり、$Z \subset X$ が $S$ 上固有な閉集合ならば、
$f(Z)$ は $Y$ の閉集合で $S$ 上固有である。
\item $f$ が普遍閉であり、$Z \subset X$ が $S$ 上固有な閉集合ならば、
$f(Z)$ は $Y$ の閉集合で $S$ 上固有である。
\item $f$ が固有であり、$Z \subset Y$ が $S$ 上固有な閉集合ならば、
$f^{-1}(Z)$ は $X$ の閉集合で $S$ 上固有である。
\end{enumerate}
\end{lemma}

\begin{proof}
(1) の証明。$Y$ が $S$ 上分離的であり、$Z \subset X$ が
$S$ 上固有な閉集合であると仮定する。$Z$ に誘導される被約閉部分スキーム構造を入れ、
射 $Z \to Y$ に $S$ 上で Morphisms, Lemma
\ref{morphisms-lemma-scheme-theoretic-image-is-proper} を適用すればよい。

\medskip\noindent
(2) の証明。$f$ が普遍閉であり、$Z \subset X$ が
$S$ 上固有な閉集合であると仮定する。$Z$ と $Z' = f(Z)$ に、
それぞれ誘導される被約閉部分スキーム構造を入れる。誘導される射
$Z \to Z'$ を得る。
$Z'' = f^{-1}(Z')$ をスキーム論的逆像とする
（Schemes, Definition \ref{schemes-definition-inverse-image-closed-subscheme}）。
このとき $Z'' \to Z'$ は $f$ の基底変換として普遍閉である
（Morphisms, Lemma \ref{morphisms-lemma-base-change-proper}）。
従って $Z \to Z'$ は、閉埋入 $Z \to Z''$ と $Z'' \to Z'$ の合成として普遍閉である
（Morphisms, Lemmas
\ref{morphisms-lemma-closed-immersion-proper} および
\ref{morphisms-lemma-composition-proper}）。
Morphisms, Lemma \ref{morphisms-lemma-image-universally-closed-separated}
により、$Z' \to S$ は分離的である。
$Z \to S$ は準コンパクトであり、$Z \to Z'$ は全射なので、
$Z' \to S$ は準コンパクトである。
$Z' \to S$ は $Z' \to Y$ と $Y \to S$ の合成である。
従って $Z' \to S$ は局所有限型である
（Morphisms, Lemmas \ref{morphisms-lemma-immersion-locally-finite-type} および
\ref{morphisms-lemma-composition-finite-type}）。
最後に、$Z \to S$ は普遍閉なので、Morphisms, Lemma
\ref{morphisms-lemma-image-proper-is-proper} により $Z' \to S$ も普遍閉である。
これで証明が完了する。

\medskip\noindent
(3) の証明。$f$ が固有であり、$Z \subset Y$ が
$S$ 上固有な閉集合であると仮定する。$Z$ に誘導される被約閉部分スキーム構造を入れる。
$Z' = f^{-1}(Z)$ をスキーム論的逆像とする
（Schemes, Definition \ref{schemes-definition-inverse-image-closed-subscheme}）。
このとき $Z' \to Z$ は $f$ の基底変換として固有である
（Morphisms, Lemma \ref{morphisms-lemma-base-change-proper}）。
従って $Z' \to S$ は $Z' \to Z$ と $Z \to S$ の合成として固有である
（Morphisms, Lemma \ref{morphisms-lemma-composition-proper}）。
これで証明が完了する。
\end{proof}

\begin{lemma}
\label{coherent-lemma-union-closed-proper-over-base}
$f : X \to S$ を局所有限型なスキームの射とする。
$Z_i \subset X$, $i = 1, \ldots, n$ を閉集合とする。
$Z_i$, $i = 1, \ldots, n$ が $S$ 上固有ならば、
$Z_1 \cup \ldots \cup Z_n$ についても同じことが成り立つ。
\end{lemma}

\begin{proof}
$Z_i$ に誘導される被約閉部分スキーム構造を入れる。射
$$
Z_1 \amalg \ldots \amalg Z_n \longrightarrow X
$$
は Morphisms, Lemmas
\ref{morphisms-lemma-closed-immersion-finite} および
\ref{morphisms-lemma-finite-union-finite} により有限である。
有限射は普遍閉であり
（Morphisms, Lemma \ref{morphisms-lemma-finite-proper}）、
$Z_1 \amalg \ldots \amalg Z_n$ は $S$ 上固有なので、補題
\ref{coherent-lemma-functoriality-closed-proper-over-base} の (2) により、
像 $Z_1 \cup \ldots \cup Z_n$ は $S$ 上固有である。
\end{proof}





\noindent
$f : X \to S$ を局所有限型なスキームの射とする。
$\mathcal{F}$ を有限型準連接 $\mathcal{O}_X$-加群とする。
このとき、台 $\text{Supp}(\mathcal{F})$（$\mathcal{F}$ のもの）は
$X$ の閉集合である。Morphisms, Lemma
\ref{morphisms-lemma-support-finite-type} を参照。
従って「$\mathcal{F}$ の台は $S$ 上固有である」ということには意味がある。

\begin{lemma}
\label{coherent-lemma-module-support-proper-over-base}
$f : X \to S$ を局所有限型なスキームの射とする。
$\mathcal{F}$ を有限型準連接 $\mathcal{O}_X$-加群とする。
次の条件は同値である。
\begin{enumerate}
\item $\mathcal{F}$ の台は $S$ 上固有である。
\item $\mathcal{F}$ のスキーム論的台
（Morphisms, Definition \ref{morphisms-definition-scheme-theoretic-support}）
は $S$ 上固有である。
\item 閉部分スキーム $Z \subset X$ と有限型準連接
$\mathcal{O}_Z$-加群 $\mathcal{G}$ が存在して、
(a) $Z \to S$ は固有であり、(b)
$(Z \to X)_*\mathcal{G} = \mathcal{F}$ である。
\end{enumerate}
\end{lemma}

\begin{proof}
台 $\text{Supp}(\mathcal{F})$（$\mathcal{F}$ のもの）は $X$ の閉集合である。
Morphisms, Lemma \ref{morphisms-lemma-support-finite-type} を参照。
従って Definition \ref{coherent-definition-proper-over-base} を適用できる。
$\mathcal{F}$ のスキーム論的台は、その台となる閉集合が
$\text{Supp}(\mathcal{F})$ である閉部分スキームなので、
Definition \ref{coherent-definition-proper-over-base} により (1) と (2) は同値である。
(2) が (3) を含意することは明らかである。
逆に (3) が成り立つならば、
$\text{Supp}(\mathcal{F}) \subset Z$
（これは $X$ の閉集合の包含である）であり、従って例えば補題
\ref{coherent-lemma-closed-closed-proper-over-base} により
$\text{Supp}(\mathcal{F})$ は $S$ 上固有である。
\end{proof}

\begin{lemma}
\label{coherent-lemma-base-change-module-support-proper-over-base}
次のスキームのカルテジアン図式を考える：
$$
\xymatrix{
X' \ar[d]_{f'} \ar[r]_{g'} & X \ar[d]^f \\
S' \ar[r]^g & S
}
$$
ここで $f$ は局所有限型である。$\mathcal{F}$ を有限型準連接
$\mathcal{O}_X$-加群とする。$\mathcal{F}$ の台が $S$ 上固有ならば、
$(g')^*\mathcal{F}$ の台は $S'$ 上固有である。
\end{lemma}

\begin{proof}
Modules, Lemma \ref{modules-lemma-pullback-finite-type} により
$(g')*\mathcal{F}$ は有限型なので、この主張には意味がある。
Morphisms, Lemma \ref{morphisms-lemma-support-finite-type} により
$\text{Supp}((g')^*\mathcal{F}) = (g')^{-1}(\text{Supp}(\mathcal{F}))$ である。
従って、補題 \ref{coherent-lemma-base-change-closed-proper-over-base} から結論が従う。
\end{proof}

\begin{lemma}
\label{coherent-lemma-cat-module-support-proper-over-base}
$f : X \to S$ を局所有限型なスキームの射とする。
$\mathcal{F}$、$\mathcal{G}$ を有限型準連接 $\mathcal{O}_X$-加群とする。
\begin{enumerate}
\item $\mathcal{F}$、$\mathcal{G}$ の台が $S$ 上固有ならば、
$\mathcal{F} \oplus \mathcal{G}$、$\mathcal{G}$ の $\mathcal{F}$ による任意の拡大、
$\Im(u)$ と $\Coker(u)$（ここで任意の $\mathcal{O}_X$-加群写像
$u : \mathcal{F} \to \mathcal{G}$ を取る）、および $\mathcal{F}$ または $\mathcal{G}$ の
任意の準連接商についても同じことが成り立つ。
\item $S$ が局所 Noether ならば、連接 $\mathcal{O}_X$-加群で台が
$S$ 上固有なものの圏は、連接 $\mathcal{O}_X$-加群のアーベル圏の
Serre 部分圏である（Homology, Definition
\ref{homology-definition-serre-subcategory}）。
\end{enumerate}
\end{lemma}

\begin{proof}
(1) の証明。$Z$、$Z'$ をそれぞれ $\mathcal{F}$ と $\mathcal{G}$ の台とする。
このとき、(1) で挙げたすべての層の台は $Z \cup Z'$ に含まれる。
従って、これらの層が有限型かつ準連接であることを確認すれば、主張そのものは
補題 \ref{coherent-lemma-closed-closed-proper-over-base} および
\ref{coherent-lemma-union-closed-proper-over-base} から明らかである。
準連接性については Schemes, Section
\ref{schemes-section-quasi-coherent} を参照されたい。
「有限型」については Modules, Section
\ref{modules-section-finite-type} を見ることを勧める。

\medskip\noindent
(2) の証明。証明は (1) の証明と同じである。
Morphisms, Lemma \ref{morphisms-lemma-finite-type-noetherian} により
$X$ は局所 Noether であり、また Section
\ref{coherent-section-coherent-sheaves} に連接加群の圏の記述があるので、
主張には意味があることに注意する。
\end{proof}

\begin{lemma}
\label{coherent-lemma-support-proper-over-base-pushforward}
$S$ を局所 Noether スキームとする。
$f : X \to S$ を局所有限型なスキームの射とする。
$\mathcal{F}$ を連接 $\mathcal{O}_X$-加群で、その台が $S$ 上固有なものとする。
このとき $R^pf_*\mathcal{F}$ は連接 $\mathcal{O}_S$-加群であり、
これはすべての $p \geq 0$ に対して成り立つ。
\end{lemma}

\begin{proof}
補題 \ref{coherent-lemma-module-support-proper-over-base} により、閉埋入
$i : Z \to X$ と有限型準連接 $\mathcal{O}_Z$-加群 $\mathcal{G}$ が存在して、
(a) $g = f \circ i : Z \to S$ は固有であり、(b)
$i_*\mathcal{G} = \mathcal{F}$ である。
Proposition \ref{coherent-proposition-proper-pushforward-coherent} により
$R^pg_*\mathcal{G}$ は $S$ 上連接である。
一方、$R^qi_*\mathcal{G} = 0$ が $q > 0$ に対して成り立つ
（Lemma \ref{coherent-lemma-finite-pushforward-coherent}）。
Cohomology, Lemma \ref{cohomology-lemma-relative-Leray} により
$R^pf_*\mathcal{F} = R^pg_*\mathcal{G}$ を得て、証明が完了する。
\end{proof}

\begin{lemma}
\label{coherent-lemma-systems-with-proper-support}
$S$ を Noether スキームとする。$f : X \to S$ を有限型射とする。
$\mathcal{I} \subset \mathcal{O}_X$ を準連接イデアル層とする。
次のものは $\textit{Coh}(X, \mathcal{I})$ の Serre 部分圏である。
\begin{enumerate}
\item $\textit{Coh}(X, \mathcal{I})$ の対象 $(\mathcal{F}_n)$ であって、
$\mathcal{F}_1$ の台が $S$ 上固有であるものからなる充満部分圏。
\item $\textit{Coh}(X, \mathcal{I})$ の対象 $(\mathcal{F}_n)$ であって、
閉部分スキーム $Z \subset X$ が存在して $S$ 上固有であり、
$\mathcal{I}_Z \mathcal{F}_n = 0$ がすべての $n \geq 1$ に対して
となるものからなる充満部分圏。
\end{enumerate}
\end{lemma}

\begin{proof}
Homology, Lemma \ref{homology-lemma-characterize-serre-subcategory} の
判定法を用いる。さらに、
$0 \to (\mathcal{G}_n) \to (\mathcal{F}_n) \to (\mathcal{H}_n) \to 0$
が $\textit{Coh}(X, \mathcal{I})$ の短完全列ならば、
(a) $\mathcal{G}_n \to \mathcal{F}_n \to \mathcal{H}_n \to 0$ は
すべての $n \geq 1$ に対して完全であり、
(b) $\mathcal{G}_n$ は $\Ker(\mathcal{F}_m \to \mathcal{H}_m)$ の商であり、
これはある $m \geq n$ に対して成り立つことを用いる。
補題 \ref{coherent-lemma-inverse-systems-abelian} の証明を参照。

\medskip\noindent
(1) の証明。$(\mathcal{F}_n)$ を $\textit{Coh}(X, \mathcal{I})$ の対象とする。
このとき $\text{Supp}(\mathcal{F}_n) = \text{Supp}(\mathcal{F}_1)$ が
すべての $n \geq 1$ に対して成り立つ。
従って、上の (a) と (b) により、
任意の短完全列
$0 \to (\mathcal{G}_n) \to (\mathcal{F}_n) \to (\mathcal{H}_n) \to 0$
（$\textit{Coh}(X, \mathcal{I})$ のもの）に対して
$\text{Supp}(\mathcal{G}_1) \cup \text{Supp}(\mathcal{H}_1) =
\text{Supp}(\mathcal{F}_1)$
が成り立つ。これは (1) で定義された圏が
$\textit{Coh}(X, \mathcal{I})$ の Serre 部分圏であることを示す。

\medskip\noindent
(2) の証明。ここでも同様に論じる。
$0 \to (\mathcal{G}_n) \to (\mathcal{F}_n) \to (\mathcal{H}_n) \to 0$
を $\textit{Coh}(X, \mathcal{I})$ の短完全列とする。
$Z \subset X$ が閉部分スキームであり、$\mathcal{I}_Z$ が
$\mathcal{F}_n$ をすべての $n$ に対して零化するならば、上の (a) と (b) により、
$\mathcal{I}_Z$ は $\mathcal{G}_n$ と $\mathcal{H}_n$ も
すべての $n$ に対して零化する。従って、$Z \to S$ が固有ならば、(2) で定義された圏は
$\textit{Coh}(X, \mathcal{I})$ の中で部分対象および商対象を取ることについて閉じている。
最後に、$Z \subset X$ と $Y \subset X$ を $S$ 上固有な閉部分スキームであって、
$\mathcal{I}_Z \mathcal{G}_n = 0$ および
$\mathcal{I}_Y \mathcal{H}_n = 0$ をすべての $n \geq 1$ に対して満たすものとする。
すると上の (a) から、
$\mathcal{I}_{Z \cup Y} = \mathcal{I}_Z \cdot \mathcal{I}_Y$
は $\mathcal{F}_n$ をすべての $n$ に対して零化する。
補題 \ref{coherent-lemma-union-closed-proper-over-base}
（および Definition \ref{coherent-definition-proper-over-base}。これは和に入れる
任意のスキーム構造を選べることを示す）により、
$Z \cup Y \to S$ は固有であり、証明が完了する。
\end{proof}








\section{Grothendieck の存在定理 III}
\label{coherent-section-existence-proper-support}

\noindent
Grothendieck の存在定理の一般形を述べるために、さらに少し記法を導入する。
$A$ をイデアル $I$ に関して完備な Noether 環とする。
$f : X \to \Spec(A)$ を分離的有限型なスキームの射とする。
$\mathcal{I} = I\mathcal{O}_X$ と置く。この状況で
$$
\textit{Coh}_{\text{support proper over } A}(\mathcal{O}_X)
$$
を、$\textit{Coh}(\mathcal{O}_X)$ の充満部分圏であって、連接
$\mathcal{O}_X$-加群のうち台が $\Spec(A)$ 上固有であるものからなるものとする。
これは $\textit{Coh}(\mathcal{O}_X)$ の Serre 部分圏である。
補題 \ref{coherent-lemma-cat-module-support-proper-over-base} を参照。同様に、
$$
\textit{Coh}_{\text{support proper over } A}(X, \mathcal{I})
$$
を、$\textit{Coh}(X, \mathcal{I})$ の充満部分圏であって、対象
$(\mathcal{F}_n)$ のうち $\mathcal{F}_1$ の台が $\Spec(A)$ 上固有であるものからなるものとする。
これは補題 \ref{coherent-lemma-systems-with-proper-support} の (1) により
$\textit{Coh}(X, \mathcal{I})$ の Serre 部分圏である。
商加群の台はもとの加群の台に含まれるので、
(\ref{coherent-equation-completion-functor}) は関手
\begin{equation}
\label{coherent-equation-completion-functor-proper-over-A}
\textit{Coh}_{\text{support proper over }A}(\mathcal{O}_X)
\longrightarrow
\textit{Coh}_{\text{support proper over }A}(X, \mathcal{I})
\end{equation}
を誘導する。これで、この節の主定理を述べる準備が整った。

\begin{theorem}[Grothendieck の存在定理]
\label{coherent-theorem-grothendieck-existence}
\begin{reference}
\cite[III Theorem 5.1.5]{EGA}
\end{reference}
$A$ をイデアル $I$ に関して完備な Noether 環とする。
$X$ を $A$ 上の分離的有限型スキームとする。このとき、関手
(\ref{coherent-equation-completion-functor-proper-over-A})
$$
\textit{Coh}_{\text{support proper over }A}(\mathcal{O}_X)
\longrightarrow
\textit{Coh}_{\text{support proper over }A}(X, \mathcal{I})
$$
は同値である。
\end{theorem}

\begin{proof}
以下、補題 \ref{coherent-lemma-i-star-equivalence} の圏同値を断りなく用いる。
この証明では、閉部分スキーム $Z \subset X$ で $A$ 上固有なものに対して、
$X$ 上の連接加群が「$Z$ に台をもつ」とは、それが $Z$ のイデアル層で
零化されること、または同値に $Z$ 上の連接加群の順像であることをいう。
Proposition \ref{coherent-proposition-existence-proper} により、
$Z$ に台をもつ連接加群と連接加群の系との間の関手について結果が成り立つことは既知である。
従って、
$\textit{Coh}_{\text{support proper over }A}(\mathcal{O}_X)$ のすべての対象と
$\textit{Coh}_{\text{support proper over }A}(X, \mathcal{I})$ のすべての対象が、
ある閉部分スキーム $Z \subset X$ で $A$ 上固有なものに台をもつことを示せば十分である。
$\textit{Coh}_{\text{support proper over }A}(\mathcal{O}_X)$ の対象については、
これは定義から成り立つ。
$\textit{Coh}_{\text{support proper over }A}(X, \mathcal{I})$ の対象について、
Proposition \ref{coherent-proposition-existence-proper} の証明方法を使ってこの主張を示す。
まずその証明を読むことを読者に勧める。

\medskip\noindent
集まり $\Xi$ を次のように定める。その元は準連接イデアル層
$\mathcal{K} \subset \mathcal{O}_X$ であって、すべての対象
$(\mathcal{F}_n)$（$\textit{Coh}_{\text{support proper over }A}(X, \mathcal{I})$
の対象）で $\mathcal{K}$ により零化されるものに対して主張が成り立つものとする。
$(0)$ が $\Xi$ に属することを示したい。そうでないとすると、
$X$ は Noether なので、極大な準連接イデアル層
$\mathcal{K}$ で $\Xi$ に属さないものが存在する。補題
\ref{coherent-lemma-acc-coherent} を参照。
$X$ を、$X$ の閉部分スキームのうち $\mathcal{K}$ に対応するもので置き換えることにより、
任意の非零な $\mathcal{K}$ が $\Xi$ に属すると仮定してよい。
$(\mathcal{F}_n)$ を
$\textit{Coh}_{\text{support proper over }A}(X, \mathcal{I})$ の対象とする。
この対象が閉部分スキーム $Z \subset X$ で $A$ 上固有なものに台をもつことを示す。
これにより定理の証明が完了する。

\medskip\noindent
Chow の補題（Lemma \ref{coherent-lemma-chow-Noetherian}）を適用して、
固有全射 $f : Y \to X$ であって、稠密開集合 $U \subset X$ 上で同型となり、
$Y$ が $A$ 上 H-準射影的であるものを取る。
開埋入 $j : Y \to Y'$ であって $Y'$ が $A$ 上射影的であるものを選ぶ。
Morphisms, Lemma
\ref{morphisms-lemma-H-quasi-projective-open-H-projective} を参照。
次に注意する：
$$
\text{Supp}(f^*\mathcal{F}_n) = f^{-1}\text{Supp}(\mathcal{F}_n) =
f^{-1}\text{Supp}(\mathcal{F}_1)
$$
最初の等式は Morphisms, Lemma
\ref{morphisms-lemma-support-finite-type} による。
仮定と補題 \ref{coherent-lemma-functoriality-closed-proper-over-base} の (3) により、
$f^{-1}\text{Supp}(\mathcal{F}_1)$ は $A$ 上固有である。
従って、$f^{-1}\text{Supp}(\mathcal{F}_1)$ の $j$ による像は
補題 \ref{coherent-lemma-functoriality-closed-proper-over-base} の (1) により
$Y'$ で閉である。従って、補題
\ref{coherent-lemma-pushforward-coherent-on-open} により
$\mathcal{F}'_n = j_*f^*\mathcal{F}_n$ は $Y'$ 上連接である。
ゆえに $(\mathcal{F}_n')$ は
$\textit{Coh}(Y', I\mathcal{O}_{Y'})$ の対象である。
Grothendieck の存在定理の射影的な場合
（Lemma \ref{coherent-lemma-existence-projective}）により、連接
$\mathcal{O}_{Y'}$-加群 $\mathcal{F}'$ と同型
$(\mathcal{F}')^\wedge \cong (\mathcal{F}'_n)$ が
$\textit{Coh}(Y', I\mathcal{O}_{Y'})$ に存在する。




$\mathcal{F}'/I\mathcal{F}' = \mathcal{F}'_1$ なので、
$$
\text{Supp}(\mathcal{F}') \cap V(I\mathcal{O}_{Y'}) =
\text{Supp}(\mathcal{F}'_1) = j(f^{-1}\text{Supp}(\mathcal{F}_1))
$$
である。構造射 $p' : Y' \to \Spec(A)$ は固有なので、
$p'(\text{Supp}(\mathcal{F}') \setminus j(Y))$ は $\Spec(A)$ で閉である。
$\Spec(A)$ の空でない閉部分集合は $V(I)$ の点を含む。
実際、Algebra, Lemma \ref{algebra-lemma-radical-completion} により
$I$ は $A$ の Jacobson 根基に含まれる。
上の等式は
$\text{Supp}(\mathcal{F}') \cap (p')^{-1}V(I) \subset j(Y)$
を示すので、$\text{Supp}(\mathcal{F}') \subset j(Y)$ と結論できる。
従って $\mathcal{F}'|_Y = j^*\mathcal{F}'$ は、閉部分スキーム
$Z'$（$Y$ のもの）で $A$ 上固有なものに台をもち、
$(\mathcal{F}'|_Y)^\wedge = (f^*\mathcal{F}_n)$ である。

\medskip\noindent
$\mathcal{K}$ を、被約な補集合 $X \setminus U$ を切り出す準連接イデアル層とする。
Proposition \ref{coherent-proposition-proper-pushforward-coherent} により
$\mathcal{O}_X$-加群 $\mathcal{H} = f_*(\mathcal{F}'|_Y)$ は連接であり、
補題 \ref{coherent-lemma-inverse-systems-push-pull} により、射
$\alpha : (\mathcal{F}_n) \to \mathcal{H}^\wedge$
（$\textit{Coh}(X, \mathcal{I})$ のもの）であって、
その核と余核がある冪 $\mathcal{K}^t$（$\mathcal{K}$ のもの）で零化されるものが存在する。
完全列
$$
0 \to \Ker(\alpha) \to (\mathcal{F}_n) \to
\mathcal{H}^\wedge \to \Coker(\alpha) \to 0
$$
を $\textit{Coh}(X, \mathcal{I})$ で得る。
$Z_0 \subset X$ を $\mathcal{H}$ のスキーム論的台とすれば、
集合論的に $Z_0 \subset f(Z')$ であることは明らかである。
従って、補題 \ref{coherent-lemma-closed-closed-proper-over-base} および
補題 \ref{coherent-lemma-functoriality-closed-proper-over-base} の (2) により、
$Z_0$ は $A$ 上固有である。
従って $\mathcal{H}^\wedge$ は補題
\ref{coherent-lemma-systems-with-proper-support} の (2) で定義された部分圏に属し、
従って特に
$\textit{Coh}_{\text{support proper over }A}(X, \mathcal{I})$ に属する。
補題 \ref{coherent-lemma-systems-with-proper-support} の (1) により、
$\Ker(\alpha)$ と $\Coker(\alpha)$ は
$\textit{Coh}_{\text{support proper over }A}(X, \mathcal{I})$ に属すると結論する。
帰納法の仮定、より正確には $\mathcal{K}^t$ が $\Xi$ に属することにより、
$\Ker(\alpha)$ と $\Coker(\alpha)$ は補題
\ref{coherent-lemma-systems-with-proper-support} の (2) で定義された部分圏に属する。
これは同補題により Serre 部分圏なので、
$(\mathcal{F}_n)$ についても同じことが成り立つ。これが示すべきことであった。
\end{proof}

\begin{remark}[Grothendieck の存在定理の具体化]
\label{coherent-remark-reformulate-existence-theorem}
$A$ をイデアル $I$ に関して完備な Noether 環とする。
$S = \Spec(A)$ および $S_n = \Spec(A/I^n)$ と書く。
$X \to S$ を分離的有限型射とする。
$n \geq 1$ に対して $X_n = X \times_S S_n$ と置く。
図式で表すと、
$$
\xymatrix{
X_1 \ar[r]_{i_1} \ar[d] & X_2 \ar[r]_{i_2} \ar[d] & X_3 \ar[r] \ar[d] &
\ldots & X \ar[d] \\
S_1 \ar[r] & S_2 \ar[r] & S_3 \ar[r] & \ldots & S
}
$$
この状況で、次を満たす系 $(\mathcal{F}_n, \varphi_n)$ を考える。
\begin{enumerate}
\item $\mathcal{F}_n$ は連接 $\mathcal{O}_{X_n}$-加群である。
\item $\varphi_n : i_n^*\mathcal{F}_{n + 1} \to \mathcal{F}_n$
は同型である。
\item $\text{Supp}(\mathcal{F}_1)$ は $S_1$ 上固有である。
\end{enumerate}
Theorem \ref{coherent-theorem-grothendieck-existence} は、完備化関手
$$
\begin{matrix}
\text{coherent }\mathcal{O}_X\text{-modules }\mathcal{F} \\
\text{with support proper over }A
\end{matrix}
\quad
\longrightarrow
\quad
\begin{matrix}
\text{systems }(\mathcal{F}_n) \\
\text{as above}
\end{matrix}
$$
が圏同値であることを述べる。特に $X$ が $A$ 上固有な場合には、
台に関する条件を省略できる。
\end{remark}



\section{Grothendieck の代数化定理}
\label{coherent-section-algebraization}

\noindent
最初の結果は、Grothendieck の存在定理を閉部分スキームと有限射の言葉で
言い換えたものである。

\begin{lemma}
\label{coherent-lemma-algebraize-formal-closed-subscheme}
$A$ をイデアル $I$ に関して完備な Noether 環とする。
$S = \Spec(A)$ および $S_n = \Spec(A/I^n)$ と書く。
$X \to S$ を分離的有限型射とする。
$n \geq 1$ に対して $X_n = X \times_S S_n$ と置く。
Cartesian な正方形からなる可換図式
$$
\xymatrix{
Z_1 \ar[r] \ar[d] & Z_2 \ar[r] \ar[d] & Z_3 \ar[r] \ar[d] & \ldots \\
X_1 \ar[r]^{i_1} & X_2 \ar[r]^{i_2} & X_3 \ar[r] & \ldots
}
$$
がスキームについて与えられているとする。次を仮定する。
\begin{enumerate}
\item $Z_1 \to X_1$ は閉埋入である。
\item $Z_1 \to S_1$ は固有である。
\end{enumerate}
このとき、スキームの閉埋入 $Z \to X$ であって
$Z_n = Z \times_S S_n$ となるものが存在する。さらに $Z$ は $S$ 上固有である。
\end{lemma}

\begin{proof}
垂直射を $j_n : Z_n \to X_n$ と書く。
主張の正方形は Cartesian なので、$j_n$ の $X_1$ への基底変換は
$j_1$ である。従って Morphisms, Lemma
\ref{morphisms-lemma-check-closed-infinitesimally} により
$j_n$ は閉埋入である。
$\mathcal{F}_n = j_{n, *}\mathcal{O}_{Z_n}$ と置く。このとき
$j_n^\sharp$ は全射 $\mathcal{O}_{X_n} \to \mathcal{F}_n$ である。
再び正方形が Cartesian であることから、$\mathcal{F}_{n + 1}$ の
$X_n$ への引き戻しは $\mathcal{F}_n$ である。
従って、Remark \ref{coherent-remark-reformulate-existence-theorem} で
言い換えた Grothendieck の存在定理により、写像
$\mathcal{O}_X \to \mathcal{F}$（連接 $\mathcal{O}_X$-加群のもの）が存在し、
$X_n$ への制限は $\mathcal{O}_{X_n} \to \mathcal{F}_n$ を復元する。
さらに、$\mathcal{F}$ の台は $S$ 上固有である。
完備化関手は完全なので（Lemma \ref{coherent-lemma-exact}）、
余核 $\mathcal{Q}$（$\mathcal{O}_X \to \mathcal{F}$ のもの）の完備化は零である。
$\mathcal{F}$ の台は $S$ 上固有であり、$\mathcal{Q}$ についても同じなので、
$\mathcal{Q} = 0$ である。例えば、これは Grothendieck の存在定理により
(\ref{coherent-equation-completion-functor-proper-over-A}) の関手が同値であることから従う。
従って $\mathcal{F} = \mathcal{O}_X/\mathcal{J}$ であり、
ここで $\mathcal{J}$ はある準連接イデアル層である。
$Z = V(\mathcal{J})$ と置けば証明が完了する。
\end{proof}

\noindent
次の補題では、実際には $Y_1 \to X_1$ が有限であると仮定すれば十分である。
これは $Y_n \to X_n$ も有限であることを含意するからである
（More on Morphisms, Lemma
\ref{more-morphisms-lemma-thicken-property-morphisms-cartesian} を参照）。

\begin{lemma}
\label{coherent-lemma-algebraize-formal-scheme-finite-over-proper}
$A$ をイデアル $I$ に関して完備な Noether 環とする。
$S = \Spec(A)$ および $S_n = \Spec(A/I^n)$ と書く。
$X \to S$ を分離的有限型射とする。
$n \geq 1$ に対して $X_n = X \times_S S_n$ と置く。
Cartesian な正方形からなる可換図式
$$
\xymatrix{
Y_1 \ar[r] \ar[d] & Y_2 \ar[r] \ar[d] & Y_3 \ar[r] \ar[d] & \ldots \\
X_1 \ar[r]^{i_1} & X_2 \ar[r]^{i_2} & X_3 \ar[r] & \ldots
}
$$
がスキームについて与えられているとする。次を仮定する。
\begin{enumerate}
\item $Y_n \to X_n$ は有限射である。
\item $Y_1 \to S_1$ は固有である。
\end{enumerate}
このとき、スキームの有限射 $Y \to X$ であって
$Y_n = Y \times_S S_n$ となるものが存在する。さらに $Y$ は $S$ 上固有である。
\end{lemma}

\begin{proof}
垂直射を $f_n : Y_n \to X_n$ と書く。
$\mathcal{F}_n = f_{n, *}\mathcal{O}_{Y_n}$ と置く。
これは連接 $\mathcal{O}_{X_n}$-加群である。実際、$f_n$ は有限である
（Lemma \ref{coherent-lemma-finite-pushforward-coherent}）。
正方形が Cartesian であることから、$\mathcal{F}_{n + 1}$ の
$X_n$ への引き戻しは $\mathcal{F}_n$ である。
従って、Remark \ref{coherent-remark-reformulate-existence-theorem} で
言い換えた Grothendieck の存在定理により、連接
$\mathcal{O}_X$-加群 $\mathcal{F}$ であって、その $X_n$ への制限が
$\mathcal{F}_n$ を復元するものが存在する。さらに $\mathcal{F}$ の台は
$S$ 上固有である。完備化関手は充満忠実なので
（Theorem \ref{coherent-theorem-grothendieck-existence}）、
乗法写像
$\mathcal{F}_n \otimes_{\mathcal{O}_{X_n}} \mathcal{F}_n \to
\mathcal{F}_n$
は貼り合わさって $\mathcal{F}$ 上の代数構造を与える。
$Y = \underline{\Spec}_X(\mathcal{F})$ と置けば証明が完了する。
\end{proof}



\begin{lemma}
\label{coherent-lemma-algebraize-morphism}
$A$ をイデアル $I$ に関して完備な Noether 環とする。
$S = \Spec(A)$ および $S_n = \Spec(A/I^n)$ と書く。
$X$、$Y$ を $S$ 上のスキームとする。
$n \geq 1$ に対して $X_n = X \times_S S_n$ および
$Y_n = Y \times_S S_n$ と置く。可換図式の両立する系
$$
\xymatrix{
& & X_{n + 1} \ar[rd] \ar[rr]_{g_{n + 1}} & & Y_{n + 1} \ar[ld] \\
X_n \ar[rru] \ar[rd] \ar[rr]_{g_n} & & Y_n \ar[rru] \ar[ld] & S_{n + 1} \\
& S_n \ar[rru]
}
$$
が与えられているとする。次を仮定する。
\begin{enumerate}
\item $X \to S$ は固有である。
\item $Y \to S$ は分離的有限型である。
\end{enumerate}
このとき、一意なスキームの射 $g : X \to Y$（$S$ 上のもの）であって、
$g_n$ が $g$ の $S_n$ への基底変換となるものが存在する。
\end{lemma}

\begin{proof}
射 $(1, g_n) : X_n \to X_n \times_S Y_n$ は閉埋入である。
実際、$Y_n \to S_n$ は分離的である
（Schemes, Lemma \ref{schemes-lemma-section-immersion}）。
従って、補題 \ref{coherent-lemma-algebraize-formal-closed-subscheme} により、
閉部分スキーム $Z \subset X \times_S Y$ で $S$ 上固有であり、
$S_n$ への基底変換が $X_n \subset X_n \times_S Y_n$ を復元するものが存在する。
第1射影 $p : Z \to X$ は固有射である
（$Z$ は $S$ 上固有であるため。Morphisms, Lemma
\ref{morphisms-lemma-image-proper-scheme-closed} を参照）。
その $S_n$ への基底変換はすべての $n$ に対して同型である。
特に、補題 \ref{coherent-lemma-proper-finite-fibre-finite-in-neighbourhood} により、
$p : Z \to X$ は $X_0$ のある開近傍上で有限である。
$X$ は $S$ 上固有なので、この開近傍は $X$ 全体であり、
$p : Z \to X$ は有限であると結論する。
Proposition \ref{coherent-proposition-existence-proper} の同値を適用すると、
$p_*\mathcal{O}_Z = \mathcal{O}_X$ を得る。実際、これは $I^n$ を法として
すべての $n$ について成り立つ。従って $p$ は同型であり、
$g$ を合成 $X \cong Z \to Y$ として得る。
一意性の証明は省略する。
\end{proof}



\noindent
「抽象的な」代数化定理を証明するには、豊富な可逆層があると仮定する必要がある。
そのような仮定なしには結果は偽だからである。

\begin{theorem}[Grothendieck の代数化定理]
\label{coherent-theorem-algebraization}
$A$ をイデアル $I$ に関して完備な Noether 環とする。
$S = \Spec(A)$ および $S_n = \Spec(A/I^n)$ と置く。
Cartesian な正方形からなるスキームの可換図式
$$
\xymatrix{
X_1 \ar[r]_{i_1} \ar[d] & X_2 \ar[r]_{i_2} \ar[d] & X_3 \ar[r] \ar[d] &
\ldots \\
S_1 \ar[r] & S_2 \ar[r] & S_3 \ar[r] & \ldots
}
$$
を考える。$(\mathcal{L}_n, \varphi_n)$ が与えられており、
各 $\mathcal{L}_n$ は $X_n$ 上の可逆層で、
$\varphi_n : i_n^*\mathcal{L}_{n + 1} \to \mathcal{L}_n$ は同型であるとする。
次を仮定する。
\begin{enumerate}
\item $X_1 \to S_1$ は固有である。
\item $\mathcal{L}_1$ は $X_1$ 上豊富である。
\end{enumerate}
このとき、スキームの固有射 $X \to S$、豊富な可逆
$\mathcal{O}_X$-加群 $\mathcal{L}$、および同型
$X_n \cong X \times_S S_n$ と
$\mathcal{L}_n \cong \mathcal{L}|_{X_n}$ が存在し、これらは射
$i_n$ および $\varphi_n$ と両立する。
\end{theorem}

\begin{proof}
図式の正方形は Cartesian であり、射 $S_n \to S_{n + 1}$ は閉埋入なので、
射 $i_n$ も閉埋入である。特に、$X_m$ を $X_n$ の閉部分スキームと
みなしてよく、これは $m < n$ のときに成り立つ。
実際、$X_m$ は準連接イデアル層 $I^m\mathcal{O}_{X_n}$ で切り出される
閉部分スキームである。さらに、スキーム
$X_1, X_2, X_3, \ldots$ の台位相空間はすべて同一視される。
従って、層 $\mathcal{O}_{X_n}$ は同じ台位相空間上にあるものとみなす
（そして実際そうみなす）。連接 $\mathcal{O}_{X_n}$-加群についても同様である。
次のように置く：
$$
\mathcal{F}_n =
\Ker(\mathcal{O}_{X_{n + 1}} \to \mathcal{O}_{X_n})
$$
すると短完全列
$$
0 \to \mathcal{F}_n \to \mathcal{O}_{X_{n + 1}} \to \mathcal{O}_{X_n} \to 0
$$
を得る。上の議論から
$\mathcal{F}_n = I^n\mathcal{O}_{X_{n + 1}}$ である。
従って $\mathcal{F}_n$ は $X_{n + 1}$ 上の連接層で $I$ により零化されるから、
連接 $\mathcal{O}_{X_1}$-加群とみなす（そして実際そうみなす）。
$m > n$ のとき、層
$$
I^n\mathcal{O}_{X_m}/I^{n + 1}\mathcal{O}_{X_m}
$$
は $\mathcal{F}_n$ に同型に写ることに注意する。この写像は
$\mathcal{O}_{X_m} \to \mathcal{O}_{X_{n + 1}}$ から誘導される。
従って $n_1, n_2 \geq 0$ が与えられたとき、
$m > n_1 + n_2$ を選んで乗法写像
$$
I^{n_1}\mathcal{O}_{X_m} \times I^{n_2}\mathcal{O}_{X_m}
\longrightarrow
I^{n_1 + n_2}\mathcal{O}_{X_m} \to \mathcal{F}_{n_1 + n_2}
$$
を考えられる。これは $\mathcal{O}_{X_1}$-双線形写像
$$
\mathcal{F}_{n_1} \times \mathcal{F}_{n_2} \longrightarrow
\mathcal{F}_{n_1 + n_2}
$$
を誘導し、さらに次数付き $\mathcal{O}_{X_1}$-代数の構造を
$\mathcal{F} = \bigoplus_{n \geq 0} \mathcal{F}_n$ 上に定める。

\medskip\noindent
$B = \bigoplus I^n/I^{n + 1}$ と置く。これは有限生成な次数付き
$A/I$-代数である。$\mathcal{B} = (X_1 \to S_1)^*\widetilde{B}$ と置く。
上の議論から、次数付き $\mathcal{O}_{X_1}$-代数の標準的全射
$$
\mathcal{B} \longrightarrow \mathcal{F}
$$
を得る。特に、$\mathcal{F}$ は有限型準連接次数付き
$\mathcal{B}$-加群である。補題 \ref{coherent-lemma-graded-finiteness} により、整数
$d_0$ であって
$H^1(X_1, \mathcal{F} \otimes \mathcal{L}^{\otimes d}) = 0$ が
すべての $d \geq d_0$ に対して成り立つものを取れる。
$d \geq d_0$ を、切断
$s_{0, 1}, \ldots, s_{N, 1} \in \Gamma(X_1, \mathcal{L}_1^{\otimes d})$
が存在するように選び、それらが埋入
$$
\psi_1 : X_1 \to \mathbf{P}^N_{S_1}
$$
を $S_1$ 上で誘導するようにする。Morphisms, Lemma
\ref{morphisms-lemma-finite-type-over-affine-ample-very-ample} を参照。
$X_1$ は $S_1$ 上固有なので、$\psi_1$ は閉埋入である。
Morphisms, Lemma \ref{morphisms-lemma-image-proper-scheme-closed} および
Schemes, Lemma \ref{schemes-lemma-immersion-when-closed} を参照。
$\psi_1$ を、$X_n$ から $\mathbf{P}^N$ への両立する閉埋入の系へ
「持ち上げる」。

\medskip\noindent
証明の第1段落の短完全列を
$\mathcal{L}_{n + 1}^{\otimes d}$ でテンソルすると、短完全列
$$
0 \to \mathcal{F}_n \otimes \mathcal{L}_{n + 1}^{\otimes d} \to
\mathcal{L}_{n + 1}^{\otimes d} \to \mathcal{L}_{n + 1}^{\otimes d} \to 0
$$
を得る。同型 $\varphi_n$ を用いると、同型
$\mathcal{L}_{n + 1} \otimes \mathcal{O}_{X_l} = \mathcal{L}_l$ を
$l \leq n$ に対して得る。
従って上の列は
$$
0 \to \mathcal{F}_n \otimes \mathcal{L}_1^{\otimes d} \to
\mathcal{L}_{n + 1}^{\otimes d} \to \mathcal{L}_n^{\otimes d} \to 0
$$
となる。$H^1(X, \mathcal{F}_n \otimes \mathcal{L}_1^{\otimes d})$ の消滅により、
$s_{0, 1}, \ldots, s_{N, 1} \in \Gamma(X_1, \mathcal{L}_1^{\otimes d})$
を帰納的に切断
$s_{0, n}, \ldots, s_{N, n} \in \Gamma(X_n, \mathcal{L}_n^{\otimes d})$
へ持ち上げられる。従って可換図式
$$
\xymatrix{
X_1 \ar[r]_{i_1} \ar[d]_{\psi_1} &
X_2 \ar[r]_{i_2} \ar[d]_{\psi_2} &
X_3 \ar[r] \ar[d]_{\psi_3} &
\ldots \\
\mathbf{P}^N_{S_1} \ar[r] &
\mathbf{P}^N_{S_2} \ar[r] &
\mathbf{P}^N_{S_3} \ar[r] & \ldots
}
$$
を得る。ここで Constructions, Section
\ref{constructions-section-projective-space} の記法により
$\psi_n = \varphi_{(\mathcal{L}_n, (s_{0, n}, \ldots, s_{N, n}))}$ である。
定理の主張にある正方形は Cartesian なので、上の図式の正方形も Cartesian である。
補題 \ref{coherent-lemma-algebraize-formal-closed-subscheme} を適用すれば証明が完了する。
\end{proof}
